\documentclass[11pt,a4paper]{amsart}

\usepackage{amsmath,amssymb,amsthm}
\usepackage[mathscr]{euscript}
\usepackage[T1]{fontenc}
\usepackage{lmodern}
\usepackage[hidelinks]{hyperref}
\usepackage{cleveref}
\usepackage{enumitem}
\usepackage{booktabs,array}
\usepackage[protrusion=true,expansion=false]{microtype}
\usepackage{mathtools}
\raggedbottom
\emergencystretch=3em
\allowdisplaybreaks

\theoremstyle{plain}
\newtheorem{theorem}{Theorem}[section]
\newtheorem{proposition}[theorem]{Proposition}
\newtheorem{corollary}[theorem]{Corollary}
\newtheorem{lemma}[theorem]{Lemma}
\newtheorem{conjecture}[theorem]{Conjecture}
\newtheorem*{mainstatement}{Main Theorem}
\newtheorem*{mainproposition}{Main Proposition}

\theoremstyle{definition}
\newtheorem{definition}[theorem]{Definition}

\theoremstyle{remark}
\newtheorem{remark}[theorem]{Remark}
\newtheorem*{roadmapnote}{Roadmap}
\newtheorem*{statusnote}{Status note}

\DeclareMathOperator{\Tr}{Tr}
\DeclareMathOperator{\spec}{spec}
\DeclareMathOperator{\Res}{Res}
\DeclareMathOperator{\Irr}{Irr}
\DeclareMathOperator{\diag}{diag}
\DeclareMathOperator{\End}{End}
\DeclareMathOperator{\lcm}{lcm}
\DeclareMathOperator{\FP}{FP}
\DeclareMathOperator{\Log}{Log}
\newcommand{\CC}{\mathbb{C}}
\newcommand{\RR}{\mathbb{R}}
\newcommand{\ZZ}{\mathbb{Z}}
\newcommand{\NN}{\mathbb{N}}
\newcommand{\QQ}{\mathbb{Q}}
\newcommand{\TT}{\mathbb{T}}
\newcommand{\cP}{\mathcal P}
\newcommand{\dd}{\mathrm{d}}
\newcommand{\abs}[1]{\lvert #1 \rvert}
\newcommand{\norm}[1]{\lVert #1 \rVert}
\newcolumntype{L}[1]{>{\raggedright\arraybackslash}p{#1}}

\title[Adams--Mobius primitive inversion]{Adams--Mobius primitive inversion for finite-group extensions of shifts of finite type}

\author[Haobo Ma]{Haobo Ma}
\address{CHRONOAI}

\date{}
\renewcommand{\shortauthors}{Haobo Ma}

\begin{document}

\begin{abstract}
For a mixing subshift of finite type carrying a finite-group cocycle, we
study how finite-dimensional twisted determinants determine primitive orbit
data by conjugacy class.  Periodic labels record powers of primitive labels,
so passing from periodic determinant logarithms to primitive label series
requires Adams operations \(g\mapsto g^m\) together with M\"obius
inversion.  This yields an explicit Peter--Weyl determinant formula for the
primitive channels.  Under the strict twisted spectral gap
\(\eta<\lambda\), the channels have Perron-point values and give
class-indexed Mertens constants.  Combining the classical Frobenius-class
Euler-product exponent \(|C|/|G|\) with the scalar finite-part
normalisation, we recover the Adams-corrected product constants from the
fixed-label Euler transform rather than from the periodic Artin logarithm.
We prove determinant rigidity for these constants and give a finite
\(S_3\) certificate, with explicit values \(F_\varepsilon(1/2)\) and
\(F_\rho(1/2)\), showing that naive Artin substitution fails at the
product-constant level.  Scalar finite-part and canonical cyclic-lift
calculations are retained only as finite-matrix normalisation and comparison
tools.
\end{abstract}

\subjclass[2020]{37B10, 37D35, 11M41, 20C15}
\keywords{Dynamical $\zeta$-function, subshift of finite type, Adams operation,
primitive orbit, finite-group extension, finite-part constant}

\maketitle

% BEGIN INPUT sec_introduction.tex
% BEGIN INPUT sec_introduction_core.tex
\section{Introduction}\label{sec:introduction}

The problem addressed in this paper is primitive orbit reconstruction
for finite-group extensions of mixing subshifts of finite type. Let
\((X_A^+,\sigma)\) be a mixing subshift with primitive adjacency matrix
\(A\), let \(\lambda=\rho(A)\), and let \(\tau:E_A\to G\) be a cocycle
into a finite group. For every irreducible representation \(\varrho\)
of \(G\), the twisted transition matrix \(A_\varrho\) records periodic
labels through its traces. These traces are periodic data. The
primitive labels are subtler: when a primitive orbit with label \(g\)
is repeated \(m\) times, its label becomes \(g^m\). Thus primitive
extraction is not ordinary M\"obius inversion in the non-abelian
setting; it must be coupled to the Adams operations
\(\psi^m\phi(g)=\phi(g^m)\) on class functions.

The centre of the paper is the Adams--M\"obius passage from twisted
dynamical determinants to primitive Peter--Weyl channels and then to the
constants in Frobenius-class Euler products.  The classical
Frobenius-class Mertens asymptotic supplies the exponent
\(|C|/|G|\):
\[
  P_C(N)=K_C N^{-|C|/|G|}(1+o(1)).
\]
The finite-state problem addressed here is the constant \(K_C\).  Under
the strict twisted gap it is governed by the fixed-label Euler transform
\[
  F_\varrho(\lambda^{-1})
  =\sum_{r\ge1}\frac{\Pi_\varrho(\lambda^{-r})}{r},
\]
not by the periodic Artin logarithm \(\mathcal L_\varrho\).  Replacing
\(F_\varrho\) by \(\mathcal L_\varrho\) is already false in the finite
\(S_3\) model of \Cref{app:s3-model-computations}.  The contribution is
therefore the Adams-corrected fixed-label determinant formula for
\(K_C\), its finite recovery, and the obstruction to Artin substitution;
it is not a new claim about the classical class-Mertens exponent.

The finite-group argument has three steps: Main Theorem~I linearises
twisted determinants and performs primitive Adams--M\"obius inversion;
Main Theorem~II evaluates the resulting non-trivial primitive series at
\(z=\lambda^{-1}\), after imposing the twisted spectral gap
\(\eta<\lambda\), and obtains conjugacy-class Mertens constants; Main
Theorem~III identifies the Euler-product constants, while the Fourier
recovery of the corresponding Perron-point primitive channels is a
structural corollary.  The recovery of these constants is the Hermitian
character-orthogonality formula. Its forward class coefficient is
\(\chi_\varrho^\ast(C):=\overline{\chi_\varrho(C)}\); for unitary
characters this is equivalently \(\chi_\varrho(C^{-1})\).

The scalar finite part of \(\zeta_A\) and the cyclic root-of-unity layer
have narrower roles. The finite part supplies only the trivial-channel
normalisation \(\mathcal M(A)\), while root-of-unity sampling belongs
only to the canonical cyclic lift, where every edge is labelled by the
generator of \(\ZZ/q\ZZ\). It is not the same as taking arbitrary cyclic
quotients of the fixed cocycle. The logical direction of the proofs is
\[
  \begin{array}{c}
    \text{twisted determinants}
    \xrightarrow{\text{trace}}
    \text{periodic channels} \\[2pt]
    \xrightarrow{\text{Adams--M\"obius}}
    \text{primitive channels} \\[2pt]
    \xrightarrow[\eta<\lambda]{z=\lambda^{-1}}
    \text{class Mertens constants}
    \xrightarrow{\text{orthogonality}}
    \text{class profile}.
  \end{array}
\]
Here the Adams--M\"obius arrow is the primitive extraction step; the last
arrow is ordinary character orthogonality after the Perron-point values
have been constructed. The formal input--output statements and the scalar
normalisation are collected in the next subsection.

The formal theorem statements below summarize this argument before the
technical development begins.  Main Theorem~I is the determinant-to-primitive
Adams--M\"obius inversion.  Main Theorem~II adds the strict twisted gap
and evaluates the non-trivial primitive channels at the Perron point.
Main Theorem~III inserts those values into the Frobenius-class Euler
product and recovers the corrected class constants.  The later rigidity,
quotient, cyclic-detection, and certificate statements serve
this chain by testing hypotheses or computing its ingredients in finite
models; in particular, the eta-gap obstruction is used only to certify the
strict Perron-point input of Main Theorems~II--III, not to classify all
finite-group cocycles with a strict gap.  The scalar dynamical \(\zeta\)-function
\[
  \zeta_A(z)
  =\exp\!\left(\sum_{n\ge 1}\frac{\Tr(A^n)}{n}z^n\right)
  =\det(I-zA)^{-1}
\]
enters only through the trivial representation channel. Its rationality
and orbit-counting role are classical, going back to
Bowen--Lanford \cite{BowenLanford1970Zeta} and Manning
\cite{Manning1971Axiom}, with further development through Ruelle's
thermodynamic formalism
\cite{Ruelle1976ZetaExpanding,Ruelle1978Thermodynamic} and the
Parry--Pollicott transfer-operator programme
\cite{ParryPollicott1990Zeta}. The finite-part and cyclic-lift material
therefore supplies scalar normalisation and comparison data, not an
alternative source of non-abelian primitive information.

\subsection{Finite-group main results}

Let \(\tau:E_A\to G\) be a cocycle into a finite group, and let
\(A_\varrho\) be the twisted transition matrices. For a class function
\(\phi:G\to\CC\) write
\[
  T_{\phi,n}:=\sum_{x\in\operatorname{Fix}(\sigma^n)}
    \phi(\tau_n(x)),
  \qquad
  \mathcal L_\phi(z):=\sum_{n\ge 1}\frac{T_{\phi,n}}{n}z^n,
\]
and
\[
  \Pi_\phi(z):=\sum_{\gamma\in\cP}\phi(\tau(\gamma))z^{|\gamma|}.
\]
For a conjugacy class \(C\subset G\), let \(p_{n,C}\) denote the number
of primitive orbits of length \(n\) whose cocycle label lies in \(C\).
The scalar orbit-Mertens constant \(\mathcal M(A)\) is the
trivial-channel constant defined in \Cref{def:abel-finite-part}; the
  finite-matrix formula behind this normalisation is the calculation in
  \Cref{lem:finite-part-primitive-closed-form}.

\subsection{Contributions and scope}\label{sec:intro-scope-of-results}

All main results are finite-state statements for a fixed matrix extension
\((A,\tau,G)\).  Main Theorem~I recovers the primitive Peter--Weyl series
from the twisted determinant family
\(\{\det(I-zA_\varrho):\varrho\in\Irr(G)\}\) on
\(|z|<\lambda^{-1}\).  The recovery is not ordinary M\"obius inversion:
when a primitive orbit is repeated its label is raised to a power, and the
correction is the Adams operation \(\psi^m\phi(g)=\phi(g^m)\) on class
functions.  Equivalently, the determinant logarithm records
\[
  \mathcal L_\phi(z)=\sum_{n\ge1}T_{\phi,n}z^n/n,
\]
whereas the primitive product constants require the fixed-label Euler
transform
\[
  F_\phi(z)=\sum_{r\ge1}\frac{\Pi_\phi(z^r)}{r},
  \qquad
  \Pi_\phi(z)=\sum_{\gamma\in\mathcal P}
       \phi(\tau(\gamma))z^{|\gamma|}.
\]
\Cref{prop:three-logarithm-product-interface} gives the local comparison
between these logarithmic objects, and \Cref{prop:class-product-comparison}
separates the classical Frobenius-density exponent from the uniquely forced
Adams-corrected boundary vector.  Together they explain the obstruction to
replacing \(F_\phi\) by the periodic Artin logarithm.

Main Theorem~II adds \(\lambda>1\) and the strict finite-dimensional
spectral gap
\[
  \eta:=\max_{\varrho\ne\mathbf1}\rho(A_\varrho)<\lambda,
\]
with the empty maximum interpreted as \(0\) when \(G\) is trivial.  Under
this hypothesis the non-trivial primitive channels have Abel-path values at
\(z=\lambda^{-1}\), and the conjugacy-class Mertens constants
\(\mathcal M_C(A,\tau;G)\) are obtained by character orthogonality.
Main Theorem~III then inserts those constants into the Frobenius-class Euler
product and gives the Adams-corrected constant \(K_C\).  The classical
Chebotarev and class-Mertens results supply the exponent \(|C|/|G|\) and the
existence of the leading product asymptotic; the new contribution here is
the finite determinant recovery of the constants after that leading term has
been separated.

\begin{proposition}[Early calibration of the class-product contribution]
  \label{prop:intro-class-product-calibration}
  Under the hypotheses of Main Theorem~III, the class product has the form
  \[
    P_C(N)=K_C N^{-|C|/|G|}(1+O(N^{-1}))
  \]
  for every conjugacy class \(C\subset G\), where the exponent
  \(|C|/|G|\) is the Frobenius density term and the non-scalar part of the
  constant is exactly the Hermitian class transform of the fixed-label Euler
  vector:
  \[
    -\frac{|G|}{|C|}\log K_C-(\gamma+
    \log C(A))
    =
    \sum_{\varrho\ne\mathbf1}
      \chi_\varrho^\ast(C)F_\varrho(\lambda^{-1}).
  \]
  This vector is recovered from the finite twisted determinant family by the
  Adams--M\"obius determinant formula.  If one replaces it by the periodic
  Artin vector
  \((\mathcal L_\varrho(\lambda^{-1}))_{\varrho\ne\mathbf1}\), then all
  class-product constants remain unchanged if and only if
  \(F_\varrho(\lambda^{-1})=
  \mathcal L_\varrho(\lambda^{-1})\) for every non-trivial irreducible
  \(\varrho\).
\end{proposition}
\begin{proof}
  The product asymptotic and the displayed formula for \(K_C\) are precisely
  the conclusion of Main Theorem~III, proved in the body as
  \Cref{thm:class-euler-product-mertens}.  Taking logarithms of its constant
  formula gives the displayed class transform.  The determinant recovery of
  each \(F_\varrho(\lambda^{-1})\) is
  \Cref{prop:fixed-label-determinant-perron}, equivalently the finite
  coordinate form in \Cref{thm:class-product-determinant-rigidity}.  Finally,
  the irreducible characters form an orthonormal basis of class functions;
  applying the inverse formula of
  \Cref{cor:class-product-artin-obstruction-recovery} to the difference
  between the fixed-label and periodic-Artin vectors proves that equality of
  all resulting class constants is equivalent to vanishing of every
  non-trivial coordinate difference.  This is exactly the comparison isolated
  later in \Cref{prop:class-product-comparison,thm:adams-obstruction-class-product}.
\end{proof}

The scalar finite part of \(\zeta_A\), the cyclic-lift comparison
calculations, and the finite certificates in the appendices are support
material only.  The scalar calculation supplies the trivial-channel value
\(\mathcal M(A)\), while the canonical cyclic lift is a fixed-matrix
comparison model.  Conditional secondary-edge refinements are left outside the
compiled article except for the non-interference statement
\Cref{prop:main-chain-secondary-noninterference}.  None of these packages
provides another route to the finite-group product constants or adds a
hypothesis to Main Theorems~I--III.  Thus the proofs of Main
Theorems~I--III use finite-matrix determinant identities, finite-group
character orthogonality, Adams--M\"obius inversion, and Perron--Frobenius
theory for finite non-negative matrices.

For reference, the proof dependencies of the three main statements are as
follows.  Main Theorem~I is the finite determinant linearisation and
Adams--M\"obius primitive inversion of
\Cref{lem:twisted-trace-identity,thm:class-function-linearisation,cor:class-function-adams-mobius,cor:primitive-peter-weyl-determinant,thm:primitive-peter-weyl-trace,thm:class-count-trace-recovery}.  Main
Theorem~II adds the strict Perron-point gap, whose finite certificate,
Perron-disc analytic separation, and explicitly conditional scope are
\Cref{prop:eta-gap-finite-state-certificate,prop:eta-gap-perron-separation,prop:eta-gap-conditional-scope},
and is proved by
\Cref{thm:class-mertens-explicit,prop:class-mertens-deviations,thm:class-mertens-fourier}.  Main Theorem~III is the product-level assembly
from \Cref{thm:class-euler-product-mertens,prop:fixed-label-determinant-perron,prop:fixed-label-perron-finite-algorithm,prop:class-product-comparison,thm:adams-obstruction-class-product}; the concrete \(S_3\) obstruction is
certified by \Cref{prop:s3-naive-artin-obstruction,lem:s3-example-perron-certificates,lem:appendix-s3-f-level-certificates} through the same finite determinant certificate and
does not introduce an additional hypothesis.

\begin{proposition}[Finite-dimensional transfer split]
  \label{prop:intro-finite-state-transfer-split}
  For the finite-group extension of the primitive edge shift fixed in this
  paper, Main Theorems~I--III depend only on the finite matrices
  \[
    A_\varrho=\bigl(\sum_{e:i\to j}\varrho(\tau(e))\bigr)_{i,j\in V},
    \qquad \varrho\in\Irr(G),
  \]
  the finite character algebra of \(G\), the Adams operations
  \(\psi^m\phi(g)=\phi(g^m)\), and the spectral gap stated in Main
  Theorem~II.  More precisely:
  \begin{enumerate}[label=\rm(\alph*)]
  \item On the open Perron disc \(|z|<\lambda^{-1}\), every determinant
    logarithm appearing in Main Theorem~I is the ordinary finite-matrix
    logarithm \(\log\det(I-zA_\varrho)^{-1}\), and its periodic trace
    coefficients are the fixed-point sums of
    \Cref{lem:twisted-trace-identity}.
  \item The passage from periodic determinant data to primitive
    fixed-label series is exactly the finite Peter--Weyl expansion followed
    by Adams--M\"obius inversion of
    \Cref{thm:class-function-linearisation,cor:class-function-adams-mobius,cor:primitive-peter-weyl-determinant,cor:primitive-peter-weyl-trace,thm:class-count-trace-recovery}.
  \item The Perron-point evaluations in Main Theorems~II--III use only the
    finite-dimensional gap
    \(\eta=\max_{\varrho\ne\mathbf1}\rho(A_\varrho)<\lambda\). By
    \Cref{lem:twisted-gap-skew-product-peripheral}, this is equivalent to
    the statement that the skew-product adjacency matrix \(\widetilde A\)
    has the single simple peripheral eigenvalue \(\lambda\), or,
    equivalently, a simple strictly dominant Perron eigenvalue
    \(\lambda\); if \(\widetilde A\) is irreducible, it is equivalent to
    primitivity of \(\widetilde A\).
  \end{enumerate}
  Consequently the hypotheses and constants of Main Theorems~I--III are
  determined inside this finite-dimensional system.  In particular the theorem
  chain contains no hidden premise consisting of a H\"older/Ruelle transfer
  operator theorem, a Banach-space determinant convention, a pressure-boundary
  rigidity statement, a completed-zeta or xi normalisation, a functional
  equation, or a spectral-window continuation theorem.
\end{proposition}
\begin{proof}
  The matrices \(A_\varrho\) are finite matrices indexed by the finite vertex
  set of the base graph, with entries in the finite-dimensional vector spaces
  \(\End(V_\varrho)\).  For \(|z|<\lambda^{-1}\) the identity
  \(\log\det(I-zA_\varrho)^{-1}
  =\sum_{n\ge1}\Tr(A_\varrho^n)z^n/n\) is therefore the ordinary finite-matrix
  determinant expansion.  The coefficient identity relating
  \(\Tr(A_\varrho^n)\) to periodic skew-product weights is exactly
  \Cref{lem:twisted-trace-identity}.  Applying the finite character
  orthogonality of \(G\) gives the class-function linearisation of
  \Cref{thm:class-function-linearisation}; applying the Adams operations
  \(\psi^m\phi(g)=\phi(g^m)\) and scalar M\"obius inversion gives
  \Cref{cor:class-function-adams-mobius}.  Thus the primitive determinant and
  trace recoveries in Main Theorem~I are precisely
  \Cref{cor:primitive-peter-weyl-determinant,cor:primitive-peter-weyl-trace,thm:class-count-trace-recovery}, all of which use only the finite character
  algebra and the finite matrices \(A_\varrho\).

  The Perron-boundary constants in Main Theorems~II and~III are evaluated only
  after imposing the explicit finite-dimensional hypothesis
  \(\eta=\max_{\varrho\ne\mathbf1}\rho(A_\varrho)<\lambda\).  Under this
  hypothesis, \Cref{thm:class-mertens-explicit,prop:class-mertens-deviations,thm:class-mertens-fourier} give the additive class constants from the same
  determinant logarithms, while
  \Cref{thm:class-euler-product-mertens,prop:fixed-label-determinant-perron,prop:class-product-comparison,thm:adams-obstruction-class-product} give the
  product constants and the Adams-corrected obstruction.  The equivalence
  between the displayed gap and the peripheral statement for the finite
  skew-product adjacency matrix is \Cref{lem:twisted-gap-skew-product-peripheral}.
  These cited results exhaust the dependencies listed immediately before the
  proposition.  None invokes a Banach-space transfer operator, a Ruelle
  resonance theorem, a completed zeta function, or a functional equation, so no
  such theorem is a hidden input to Main Theorems~I--III.
\end{proof}

The paper's main contribution is recorded in the following three numbered
headline theorems.  The boundary with the classical literature is as follows.
The inherited symbolic-dynamical layer is the Chebotarev/Mertens existence and
scale layer: Frobenius classes have density 
\(|C|/|G|\), and the associated products have the classical power
\(N^{-|C|/|G|}\) under the usual hypotheses.  The new finite-state layer begins
only after that leading term has been removed: it identifies the residual
constant by a fixed-label Euler transform, proves its Adams--M\"obius
determinant formula, and shows that substituting the ordinary periodic
Artin--Ruelle logarithm gives the wrong coordinate in non-abelian examples.
Thus the theorem chain below is not presented as a new proof of the
Parry--Pollicott, Noorani--Parry, Sharp, or Mesliza--Noorani asymptotic scale;
it is the finite determinant recovery of the primitive-label constants left by
that scale.  Main Theorem~I is the open-disc determinant
inversion theorem: it converts periodic twisted determinants into primitive
Peter--Weyl channels before any Perron-boundary evaluation.  Main
Theorem~II is the Perron-point class-constant theorem: after the strict
twisted gap, it evaluates those primitive channels at
\(z=\lambda^{-1}\) and gives the additive Frobenius-class Mertens constants.
Main Theorem~III is the Adams-corrected symbolic-dynamics novelty of the paper:
it identifies the Adams-corrected fixed-label Euler-product constants and
proves that the ordinary periodic Artin logarithm is not the right primitive
class-product substitute.  Thus the paper is positioned in the symbolic
dynamics and dynamical-zeta literature by source role rather than by a blanket
claim of new prime-orbit asymptotics.  Parry--Pollicott's monograph and
Chebotarev work give the transfer-operator, Artin--Ruelle determinant, and
Galois-cover orbit-counting framework
\cite{ParryPollicott1990Zeta,ParryPollicott1986Chebotarev}.  Noorani--Parry
give the finite homogeneous-extension Chebotarev theorem for shifts
\cite{NooraniParry1992ChebotarevShifts}.  Sharp supplies the closed-orbit
Mertens model and homology-class orbit technology
\cite{Sharp1991Mertens,Sharp1993ClosedOrbitsHomologyAnosov}, while
Mesliza--Noorani give the Frobenius-class Mertens comparison theorem
\cite{MeslizaNoorani1999FrobeniusMertens}.  Those papers account for the
leading density, the exponent \(|C|/|G|\), and the existence layer of
Frobenius-class products.  They do not contain the finite-label
Adams--M\"obius Euler transform used here to pass from periodic determinant
logarithms to primitive fixed-label product constants.  Recent published ETDS papers by Coles, Humbert, O'Hare, and
Crismale--Del Vecchio--Griseta--Rossi are used only to calibrate the present
symbolic-dynamics setting against current work on periodic-orbit invariants,
Ruelle resonances, finite-data rigidity, and non-commutative skew products
\cite{Coles2025VassilievWritheAxiomA,Humbert2025FirstRuelleResonance,OHare2025FiniteDataRigidity,CrismaleDelVecchioGrisetaRossi2025NoncommutativeSkewProduct};
they are not proof inputs.  The new material here is the fixed-label
Adams--M\"obius Euler transform, its finite determinant recovery, the
Hermitian class transform, and the \(S_3\) obstruction showing that the
periodic Artin logarithm gives the wrong product constant.

This positioning leaves a narrow new contribution.  Main Theorem~III keeps the
classical exponent \(|C|/|G|\), but identifies the non-scalar coordinate of the
product constant as the fixed-label Euler value
\(F_\varrho(\lambda^{-1})\).  The determinant formula is
\Cref{prop:fixed-label-determinant-perron,thm:class-product-determinant-rigidity},
and the separation from the periodic Artin logarithm is
\Cref{prop:three-logarithm-product-interface,thm:adams-obstruction-class-product}.
Secondary-edge asymptotics are treated only as a conditional appendix
refinement of the tail after the constants have been identified; the appendix
records the finite-length comparison and is not a fourth main theorem.

\begin{table}[t]
  \centering
  \footnotesize
  \begin{tabular}{L{0.21\textwidth}L{0.34\textwidth}L{0.36\textwidth}}
    \toprule
    Result in this paper & Classical or recent comparison & Calibrated contribution \\
    \midrule
    Main Theorem~I and
    \Cref{cor:primitive-peter-weyl-determinant,cor:primitive-peter-weyl-trace,thm:class-count-trace-recovery}
      & Parry--Pollicott's Artin--Ruelle and Chebotarev framework and
        Noorani--Parry's finite homogeneous-extension theorem
        \cite{ParryPollicott1990Zeta,ParryPollicott1986Chebotarev,NooraniParry1992ChebotarevShifts}
      & Uses only finite matrices \(A_\varrho\) and finite character algebra to
        invert periodic twisted determinant data to primitive Peter--Weyl
        channels on \(|z|<\lambda^{-1}\); the new step is the Adams correction
        forced by primitive repetitions, not a new leading-density theorem. \\
    Main Theorem~II and
    \Cref{thm:class-mertens-explicit,prop:class-mertens-deviations,thm:class-mertens-fourier}
      & Sharp's orbit-Mertens theorem and the Chebotarev/Mertens comparison
        lineage for finite extensions
        \cite{Sharp1991Mertens,ParryPollicott1986Chebotarev,NooraniParry1992ChebotarevShifts,MeslizaNoorani1999FrobeniusMertens}
      & Takes the exponent \(|C|/|G|\) and the Perron gap as inputs, then gives
        the class-indexed constant \(\mathcal M_C(A,\tau;G)\) by explicit
        non-trivial primitive-channel determinant values at
        \(z=\lambda^{-1}\); it is not a new Tauberian proof of Mertens. \\
    Main Theorem~III and
    \Cref{thm:class-euler-product-mertens,prop:fixed-label-determinant-perron,thm:adams-obstruction-class-product}
      & Frobenius-class product asymptotics and Artin-logarithm determinant
        analogues in the same symbolic-dynamical lineage
        \cite{MeslizaNoorani1999FrobeniusMertens,ParryPollicott1990Zeta}
      & Identifies the fixed-label Euler transform
        \(F_\varrho(z)=\sum_{r\ge1}\Pi_\varrho(z^r)/r\), recovers it from
        determinant data, and proves that substituting the ordinary periodic
        Artin logarithm can change non-abelian class-product constants. \\
    Hermitian readout, quotient/shadow boundaries, and \(S_3\) certificates
      & Recent ETDS comparison points on periodic-orbit observables,
        resonances, finite-data rigidity, and non-commutative skew products
        \cite{Coles2025VassilievWritheAxiomA,Humbert2025FirstRuelleResonance,OHare2025FiniteDataRigidity,CrismaleDelVecchioGrisetaRossi2025NoncommutativeSkewProduct}
      & Provides finite-state determinant readout and a concrete non-abelian
        obstruction inside the fixed SFT cocycle model; those comparison works
        motivate the venue calibration but are not invoked as proof inputs. \\
    Secondary-edge split boundary
      & No classical comparison input and no hypothesis for Main Theorems~I--III
      & Conditional finite-length tail refinements belong to a separate
        refinement package; the retained statement
        \Cref{prop:main-chain-secondary-noninterference} records that such
        tails do not change \(\mathcal M_C\),
        \(F_\varrho(\lambda^{-1})\), or \(K_C\). \\
    \bottomrule
  \end{tabular}
  \caption{Theorem-by-theorem calibration of the paper's contribution.  The
  table separates inherited symbolic-dynamical inputs from the finite-label
  Adams--M\"obius and fixed-label product calculations proved here.}
  \label{tab:novelty-calibration}
\end{table}

\begin{theorem}[Main Theorem I: finite determinant recovery of primitive channels]
  \phantomsection\label{stmt:main-theorem-i-adams-mobius-linearisation}
  \label{thm:intro-adams-mobius-linearisation}
  \label{thm:intro-finite-group}
  \label{thm:intro-adams-mobius}
  Let \(A\) be primitive and let \(\tau:E_A\to G\) be a finite-group edge
  cocycle.  For every class function \(\phi:G\to\CC\), the primitive
  label series \(\Pi_\phi(z)\) on \(|z|<\lambda^{-1}\) is determined by
  the finite family of twisted determinant polynomials
  \[
    \{\det(I-zA_\varrho):\varrho\in\Irr(G)\}
  \]
  together with the Adams operations \(\psi^m\phi(g)=\phi(g^m)\).  More
  precisely,
  \[
    \Pi_\phi(z)=
    \sum_{m\ge1}\frac{\mu(m)}{m}\,
      \mathcal L_{\psi^m\phi}(z^m),
    \qquad
    \mathcal L_\phi(z)=
    \sum_{\varrho\in\Irr(G)}\widehat\phi(\varrho)
      \log\det(I-zA_\varrho)^{-1}.
  \]
  Consequently the primitive Frobenius-class counts \(p_{n,C}\) are
  obtained by finite character-table operations and coefficient
  extraction from the same determinant family.
\end{theorem}
\begin{proof}
  This is precisely \Cref{prop:main-theorem-i-assembly-map}.  That proposition
  first invokes the twisted trace identity and the determinant logarithm
  expansion to obtain the periodic series \(\mathcal L_\phi\), then applies
  Adams--M\"obius inversion to pass from periodic labels to primitive labels,
  and finally uses Peter--Weyl inversion to recover the class counts.  These
  are exactly the operations displayed in the statement, so no additional
  finite-group hypothesis or limiting orbit datum is being imported here.
\end{proof}

\begin{theorem}[Main Theorem II: Perron-point class Mertens constants]
  \phantomsection\label{stmt:main-theorem-ii-class-mertens-constants}
  \label{thm:intro-class-mertens-constants}
  \label{thm:intro-class-mertens}
  Assume \(\lambda>1\) and the twisted spectral gap
  \(\eta=\max_{\varrho\neq\mathbf1}\rho(A_\varrho)<\lambda\).  Then for
  every conjugacy class \(C\subset G\),
  \[
    \sum_{n\le N}\frac{p_{n,C}}{\lambda^n}
    =
    \frac{|C|}{|G|}\log N+\mathcal M_C(A,\tau;G)+O(N^{-1}),
  \]
  where
  \[
    \mathcal M_C(A,\tau;G)=
    \frac{|C|}{|G|}\biggl(\mathcal M(A)+
      \sum_{\varrho\neq\mathbf1}\chi_\varrho^\ast(C)
        \Pi_\varrho(\lambda^{-1})\biggr).
  \]
  The non-trivial Perron-point values
  \(\Pi_\varrho(\lambda^{-1})\) are the absolutely convergent boundary
  values, equivalently the Abel regular values of the primitive
  Peter--Weyl series under the strict finite-matrix gap
  \(\eta<\lambda\).  Their Hermitian class Fourier transform recovers the full
  normalised deviation profile and gives the Parseval identity proved in
  \Cref{thm:class-mertens-fourier}.
\end{theorem}
\begin{proof}
  The asymptotic and the displayed constant are
  \Cref{thm:class-mertens-explicit}.  The normalised deviation profile is then
  defined and identified in \Cref{prop:class-mertens-deviations}, and the
  Hermitian Fourier--Parseval readout is \Cref{thm:class-mertens-fourier}.
  Each of these results assumes exactly \(\lambda>1\) and the strict twisted
  gap \(\eta<\lambda\), and the only boundary values used are the regular
  non-trivial Perron-point values furnished by Main Theorem~I.
\end{proof}

\begin{theorem}[Main Theorem III: Adams-corrected Frobenius-class Euler-product constants]
  \phantomsection\label{stmt:main-theorem-iii-adams-frobenius-products}
  \label{thm:intro-adams-frobenius-products}
  Under the hypotheses of Main Theorem~II, define the fixed-label Euler
  transform
  \[
    F_\varrho(z):=\sum_{r\ge1}\frac{\Pi_\varrho(z^r)}{r}.
  \]
  This is distinct from the Artin periodic logarithm
  \(\mathcal L_\varrho(z)=\sum_{r\ge1}\Pi_{\psi^r\varrho}(z^r)/r\) in
  general.  For every conjugacy class \(C\),
  \[
    P_C(N):=
    \prod_{\substack{\gamma\ \mathrm{primitive}\; |\gamma|\le N\\
      \tau(\gamma)\in C}}
      (1-\lambda^{-|\gamma|})
    =K_C\,N^{-|C|/|G|}\bigl(1+O(N^{-1})\bigr),
  \]
  with
  \[
    K_C=\exp\biggl(-\frac{|C|}{|G|}
      \biggl(\gamma+\log C(A)+
        \sum_{\varrho\neq\mathbf1}\chi_\varrho^\ast(C)
          F_\varrho(\lambda^{-1})\biggr)\biggr).
  \]
  The factor \(|C|/|G|\) multiplies the entire Hermitian class-indicator
  bracket, including the non-trivial fixed-label sum.  Each
  \(F_\varrho(\lambda^{-1})\) is computable from the same finite
  twisted determinant data as Main Theorem~I by substituting the
  Adams--M\"obius determinant formula into the displayed series.  Thus
  the leading exponent and the existence of such Frobenius-class products
  are the classical part of the statement; the new assertion is the
  primitive-label fixed-Euler determinant recovery of \(K_C\) and the
  resulting obstruction to replacing \(F_\varrho\) by the periodic Artin
  logarithm.  As a downstream structural refinement, the non-trivial
  Peter--Weyl boundary is a rank-sharp coordinate record only for the
  scalar-normalised product constants in fixed branch-compatible determinant
  coordinates:
  \Cref{cor:class-product-minimal-obstruction-skeleton} proves that every
  real linear record, after the scalar baseline is fixed, determining all
  branch-compatible product constants has rank
  at least the number of non-trivial conjugacy-class degrees of freedom,
  and that the Adams-corrected determinant coordinates attain this bound.
  It is not asserted as a universal minimality theorem for arbitrary
  observations of the extension, nor as a branch-free determinant statement.
  This distinction is part of the statement: the exponent
  \(|C|/|G|\) and the existence of Frobenius-class Mertens products belong
  to the Parry--Pollicott, Noorani--Parry, Sharp, and Mesliza--Noorani
  orbit-counting lineage cited above; the novelty asserted here begins
  only after that classical leading term has been separated, namely at the
  Adams--M\"obius fixed-label determinant formula for the product
  constant.
\end{theorem}
\begin{proof}
  The fixed-label Euler transform is established in
  \Cref{thm:fixed-label-euler-transform}; applying it class by class gives the
  product asymptotic and the displayed expression for \(K_C\) in
  \Cref{thm:class-euler-product-mertens}.  The claim that every value
  \(F_\varrho(\lambda^{-1})\), and hence every \(K_C\), is recovered from the
  finite twisted determinant data is proved in the finite-determinant recovery
  package
  \Cref{prop:fixed-label-determinant-perron,prop:fixed-label-perron-finite-algorithm,thm:class-product-determinant-rigidity,cor:assembled-adams-frobenius-product-formula}.
  The scalar-normalised coordinate and rank statements are precisely
  \Cref{cor:class-product-minimal-obstruction-skeleton}, while the comparison
  with the periodic Artin logarithm is
  \Cref{prop:class-product-comparison,thm:adams-obstruction-class-product}.
  Thus the proof uses only the chain cited in the statement; the finite
  \(S_3\) certificates enter later only as explicit witnesses to the
  obstruction and not as hypotheses for the product theorem.
\end{proof}

The product constants are unaffected by later conditional secondary-edge
refinements: \Cref{prop:main-chain-secondary-noninterference} records that any
finite secondary refinement whose compensated tails vanish leaves
\(\mathcal M_C(A,\tau;G)\), \(K_C\), and the recovered values
\(F_\varrho(\lambda^{-1})\) unchanged.  The split secondary package therefore
supplies at most a conditional refinement of the error term and a detection
criterion, rather than a fourth main theorem or an alternative proof of the
product constants.

The corollary-level refinements
\Cref{cor:class-product-minimal-obstruction-skeleton,cor:adams-obstruction-bound}
then assemble the already proved constants into a fixed-coordinate
Peter--Weyl coordinate vector and a Parseval obstruction bound.  They are used below
only as structural consequences and finite witnesses for the
scalar-normalised product constants, not as observation-system minimality
claims or as additional parts of Main Theorem~III.

\paragraph{Proof guide.}
  The trace identity underlying Main Theorem~I is
  \Cref{lem:twisted-trace-identity}; the determinant linearisation and
  Adams--M\"obius primitive inversion are
  \Cref{thm:class-function-linearisation,cor:class-function-adams-mobius}.
  The determinant-germ formula and coefficient forms are
  \Cref{cor:primitive-peter-weyl-determinant,thm:primitive-peter-weyl-trace,thm:class-count-trace-recovery}.  Main Theorem~II is proved by
  \Cref{thm:class-mertens-explicit,prop:class-mertens-deviations,thm:class-mertens-fourier}.  Main Theorem~III is
  \Cref{thm:class-euler-product-mertens}. The determinant evaluation in
  \Cref{prop:fixed-label-determinant-perron,cor:class-euler-product-determinant-evaluation}
  identifies the constant without adding hypotheses to the main theorems.
  The eta-gap diagnostic \Cref{thm:eta-gap-rigidity} explains when the
  strict-gap hypothesis fails for the fixed finite-state cocycle; its one-dimensional loop-character and
  perfect-loop-group consequences are
  \Cref{cor:eta-gap-character-classification,cor:eta-gap-perfect-group}.
  These results certify the Perron-point hypothesis but are not separate
  mechanisms for product constants or a classification of all strict-gap
  cocycles.

  Secondary spectral-edge refinements are downstream from this proof guide and
  are retained in the appendix only under the additional separation hypotheses
  stated there.
  The core protection against leakage is the finite-tail non-interference
  proposition \Cref{prop:main-chain-secondary-noninterference}: once the
  determinant coordinates and product constants have been fixed by the main
  chain, any auxiliary finite spectral edge whose compensated tail tends to
  zero contributes no new Perron-boundary coordinate, no new Mertens constant,
  and no new Euler-product constant.  Any sharper secondary-tail expansion,
  isolation statement, or dominant-modulus detection result is appendix-level
  conditional material and supplies no hypothesis or constant used by Main
  Theorems~I--III.

As a structural consequence of Main Theorem~II, quotient maps
\(q:G\twoheadrightarrow Q\) intertwine both the periodic logarithms and
primitive series,
\[
  \mathcal L_{q^\ast f}^{G}(z)=\mathcal L_f^{Q}(z),
  \qquad
  \Pi_{q^\ast f}^{G}(z)=\Pi_f^{Q}(z)
  \qquad(|z|<\lambda^{-1}).
\]
For
\[
  \mathcal B_{A,\tau;G}(C):=\frac{|G|}{|C|}\Delta_C(A,\tau;G),
\]
the Peter--Weyl decomposition splits the Perron-point class profile into
its one-dimensional sector and its higher-dimensional sector. The span of
all cyclic quotient observables is exactly the one-dimensional sector;
therefore its projection is the abelian shadow, orthogonal to the genuine
non-abelian defect. Under \(\eta<\lambda\), the full primitive class
profile is generated by cyclic quotient shadows exactly when every
higher-dimensional primitive channel
\(\Pi_\varrho(\lambda^{-1})\) vanishes. These assertions are proved in
\Cref{thm:quotient-functoriality,cor:abelian-cyclic-shadow,thm:abelian-shadow-defect,thm:cyclic-detection-boundary}.

Similarly, a scalar record formed only from the fixed matrix \(A\), such
as \(\FP(A)\), \(\mathcal M(A)\), the reduced polynomial \(F_A\), the
branch-lifted cyclic-lift record \(\{\Psi_q(A)\}_{q\ge1}\), or a torsion
layer \(\{F_A(\omega_q^j)\}_{0\le j<q}\), cannot by itself recover the
non-trivial finite-group profile for a family of cocycles over the same
base matrix. The precise projection boundary is
\Cref{thm:scalar-normalisation-projection-boundary}: scalar data determine
a single normalised class profile across such a family exactly when all
non-trivial Perron-point channels
\(\Pi_\varrho^\tau(\lambda^{-1})\), \(\varrho\neq\mathbf1\), are already
independent of \(\tau\). Thus the scalar finite-part and cyclic-lift
statements below are fixed-matrix inputs and do not enlarge the
finite-group determinant and product results; the two-cocycle separation is
\Cref{cor:fixed-matrix-record-no-descent}.

An explicit \(S_3\)-extension separates the two obstruction types at the
level relevant to Main Theorem~III.  Its fixed-label Euler-transform
certificates
\(F_\varepsilon(1/2)\in(-381/1000,-380/1000)\) and
\(F_\rho(1/2)\in(-923/1000,-922/1000)\) force the product-constant
obstruction to naive Artin substitution.  The accompanying primitive
certificates show that the sign-character trace vanishes but its primitive
value is non-zero by the quadratic Adams mechanism, while the standard
two-dimensional channel gives the certified higher-dimensional cyclic
detection defect. This is a finite-state obstruction model, not a new
periodic-orbit distribution theorem.

The quotient/shadow consequence is proved in
\Cref{thm:quotient-functoriality,thm:scalar-normalisation-projection-boundary,cor:fixed-matrix-record-no-descent,cor:abelian-cyclic-shadow,thm:abelian-shadow-defect,thm:cyclic-detection-boundary,thm:s3-certified-nonabelian-defect}; the canonical cyclic-lift root layer is used only as a scalar comparison boundary, not as part of the quotient-detection theorem.

% END INPUT sec_introduction_core.tex
% BEGIN INPUT sec_introduction_scalar_and_plan.tex
\subsection{Scalar finite-matrix inputs}

We record only the scalar inputs needed for notation and comparison.
They are proved later in \Cref{sec:finite-part}. After removing the
Perron factor one obtains the reduced determinant
\[
  C(A):=\det\nolimits'(I-A/\lambda)^{-1}
\]
and the reduced spectral polynomial
\[
  F_A(t):=\frac{\det(I-tA/\lambda)}{1-t}
  =\prod_{j=2}^d\left(1-\frac{\mu_j}{\lambda}t\right),
\]
where \(\mu_2,\dots,\mu_d\) are the non-Perron eigenvalues. The
Abel-regularised finite-part constant
\[
  \FP(A):=\lim_{r\uparrow 1}\bigl(P_A(r)+\log(1-r)\bigr),
  \qquad
  P_A(r):=\sum_{n\ge 1}p_n(A)\lambda^{-n}r^n,
\]
measures the deviation of the primitive orbit sum from the harmonic
series.

\begin{proposition}[Scalar finite-part computation at the
  Perron pole]
  \label{prop:intro-scalar-finite-part-computation}
  \label{thm:intro-finite-part}
  Let \(A\) be a primitive non-negative real matrix with Perron root
  \(\lambda=\rho(A)>1\), reduced determinant \(C(A)\), and M\"obius
  coefficients \(p_n(A)\) defined by \eqref{eq:primitive-mobius}. Then
  \[
    \FP(A)
    =\log C(A)
      +\sum_{k\ge 2}\frac{\mu(k)}{k}
        \log\det(I-\lambda^{-k}A)^{-1}.
  \]
  Equivalently,
  \[
    \FP(A)
    =\log C(A)
      +\sum_{n\ge 1}p_n(A)
        \bigl(\lambda^{-n}+\log(1-\lambda^{-n})\bigr).
  \]
  In particular, the orbit-Mertens constant
  \(\mathcal M(A)=\gamma+\FP(A)\) satisfies
  \[
    \sum_{n\le N}\frac{p_n(A)}{\lambda^n}
    =\log N+\mathcal M(A)+o(1).
  \]
\end{proposition}

\begin{proof}
  This proposition records the scalar normalisation used in the class-product
  theorem chain.
  The determinant and primitive-orbit closed formulas for \(\FP(A)\) are
  precisely \Cref{lem:finite-part-primitive-closed-form}. The identity
  \(\mathcal M(A)=\gamma+\FP(A)\) and the displayed scalar orbit-Mertens
  asymptotic are \Cref{cor:finite-part-mertens-asymptotic}. Thus the
  proposition records those two proved scalar facts with the normalisations
  needed later in
\Cref{thm:class-mertens-explicit,cor:assembled-adams-frobenius-product-formula,cor:class-product-stability-scalar-independence}; the rank-zero boundary of
  these scalar records for the non-trivial class profile is isolated in
  \Cref{cor:scalar-record-rank-zero-bound}.
\end{proof}


\begin{corollary}[Introductory quotient-shadow consequences]
  \label{thm:intro-quotient-shadow}
  Under the twisted spectral gap, the body proves a quotient-shadow theorem
description of the normalised primitive class profile.  It includes scalar
  projection onto the one-dimensional Peter--Weyl sector, functoriality for
  finite quotients, the orthogonal decomposition into abelian shadow and
  genuine non-abelian defect, and exact detection of the one-dimensional
  sector by cyclic quotients of the fixed cocycle.  In particular, the
  one-dimensional shadow is obtained by pulling back the primitive profile
  of the abelianised cocycle.  The body also records a boundary warning:
  the root-of-unity determinant layers
  \(\log\det(I-\omega_q^j zA)^{-1}\) occur only in the separate canonical
  cyclic-lift comparison model, not as invariants of arbitrary cyclic
  quotients of the fixed cocycle.
\end{corollary}

\begin{proof}
This is the front-matter summary of the body results in
  \Cref{thm:quotient-functoriality,cor:abelian-cyclic-shadow,lem:cyclic-quotient-versus-root-layer,thm:abelian-shadow-defect,lem:cyclic-quotients-one-dimensional-characters}.  The quotient theorem
  gives the functorial finite-quotient formula for primitive profiles.  The
  abelian-shadow corollary proves that the one-dimensional characters of
  \(G\) are precisely the pullbacks of the characters of
  \(G_{\mathrm{ab}}\), that the twisted gap descends to the abelianised
  cocycle, and that substituting the resulting Perron-point primitive
  channels into the normalised class-profile formula gives the pullback
  identity \eqref{eq:abelian-shadow-pullback}.  The defect theorem records
  the orthogonal split between the abelian and higher-dimensional sectors,
  while the cyclic-quotient lemma identifies the span of all cyclic quotient
  coordinates with the one-dimensional character span.  Finally,
  \Cref{lem:cyclic-quotient-versus-root-layer} separates inherited cyclic
  quotient data from the canonical cyclic lift \(\tau_q(e)=1\), for which
  the twisted matrices are \(\omega_q^j A\) and the determinant channels are
  the displayed root-of-unity layers.  This last comparison is used only to
  prevent a false descent from cyclic quotients to scalar root layers.  Thus
  the quotient/shadow statement is a proved structural consequence supporting
  Main Theorems~II--III, not an informal corollary, a cyclic-arithmetic
  refinement program, or an additional descent principle.
\end{proof}


This is the trivial-channel scalar input singled out in
\Cref{sec:intro-scope-of-results};
  the scalar statement above is proved in
\Cref{prop:intro-scalar-finite-part-computation,lem:finite-part-primitive-closed-form,cor:finite-part-mertens-asymptotic}.

The second scalar component is the separate fixed-matrix lift model,
related to inverse spectral questions for SFTs
\cite{LindMarcus1995}. Appendix~\ref{app:cyclic-lift-comparison} records the
branch-lifted reconstruction and one-layer Fourier recovery for this fixed
matrix comparison problem; the stable synopsis labels are
\Cref{thm:intro-cyclic-reconstruction,cor:intro-one-layer-reconstruction};
these are appendix-support labels, not additional main theorem labels.
As already specified in
\Cref{sec:intro-scope-of-results}, those results are not used in the
finite-group determinant recovery or product-constant proofs; they serve
only to prevent the scalar cyclic-lift record from being mistaken for a
non-abelian class-profile invariant.  A broader cyclic-lift inverse-spectral
or identifiability theory would require separate hypotheses and examples;
it is not developed here.  The terminology of cyclic, torsion,
and exponential spectral records is fixed in
\Cref{def:spectral-fingerprints}.

\subsection{Relation to existing work}

The natural comparison class is the symbolic-dynamics and dynamical
\(\zeta\)-function literature on prime cycles, finite extensions, and
twisted determinant factorizations.  In that literature the leading
Chebotarev distribution and Frobenius-class product asymptotics are
classical inputs.  Parry--Pollicott's Chebotarev theorem for Axiom~A
Galois coverings \cite{ParryPollicott1986Chebotarev}, Noorani--Parry's
finite homogeneous-extension theorem for shifts
\cite{NooraniParry1992ChebotarevShifts}, Sharp's Mertens theorem for
Axiom~A flows \cite{Sharp1991Mertens}, and Mesliza--Noorani's
Frobenius-class Mertens product theorem
\cite{MeslizaNoorani1999FrobeniusMertens} supply the background against
which Main Theorems~II--III should be read.  The exponent
\(|C|/|G|\) and the existence of Frobenius-class products are therefore
not claimed as new here.  The theorem-local contribution is the finite
Adams--M\"obius determinant recovery of the class constants after that
classical leading term has been separated.

At the determinant level, the paper is closer to Fried's twisted periodic
point formulae \cite{Fried1983PeriodicPointsTwistedCoefficients,Fried1983HomologicalIdentitiesClosedOrbits}, the twisted
Perron--Frobenius framework of Adachi--Sunada
\cite{AdachiSunada1987TwistedPF}, and graph-Artin determinant theories of
Hashimoto, Sunada, Stark--Terras, and related finite-graph covering work
\cite{Hashimoto1989ZetaFiniteGraphs,Hashimoto1990ZetaLFunctionsFiniteGraphs,Hashimoto1992ArtinDensityPrimeCycles,Sunada1986LFunctionsGeometry,StarkTerras1996FiniteGraphsCoverings,StarkTerras2000FiniteGraphsCoveringsII,StarkTerras2007FiniteGraphsCoveringsIII}.  Those works encode periodic
or prime-cycle data by determinants or representation-theoretic
\(L\)-functions.  The present finite-state calculation fixes one
finite-group SFT cocycle and asks the inverse question needed for the
constants: from the full family of twisted determinants, recover the
primitive Peter--Weyl channels.  The extra operation is the finite-group
Adams correction \(\psi^m\phi(g)=\phi(g^m)\), forced by repetitions of a
primitive orbit.

Finite group actions and extensions of SFTs form the adjacent symbolic
setting.  Adler--Kitchens--Marcus \cite{AdlerKitchensMarcus1985FiniteGroupActions},
Fiebig \cite{Fiebig1993PeriodicFiniteGroupActions}, Silver--Williams
\cite{SilverWilliams2005InvariantFiniteGroupActions}, Boyle--Schmieding
\cite{BoyleSchmieding2017FiniteGroupExtensions}, Boyle--Carlsen--Eilers
\cite{BoyleCarlsenEilers2020FlowEquivalenceGSFT}, and Matsumoto
\cite{Matsumoto2016ExtensionsSubshiftsFiniteGroups} treat finite actions,
quotients, periodic data, and classification or obstruction questions.
Here the extension is fixed; the issue is not existence or
flow-equivalence classification but determinant-level recovery of
\(\Pi_\varrho(\lambda^{-1})\), \(\mathcal M_C(A,\tau;G)\), and the
Adams-corrected Euler-product constants.

The result isolated below is the fixed-label Euler transform
\[
  F_\varrho(z)=\sum_{r\ge1}\frac{\Pi_\varrho(z^r)}{r},
\]
its Adams-corrected determinant expression, and the obstruction to
replacing it by the periodic Artin logarithm
\(\mathcal L_\varrho\) in Frobenius-class product constants.

Thus this paper fixes one finite SFT cocycle and proves the determinant
recovery of primitive Peter--Weyl channels, the Perron-point class Mertens
constants, and the Adams-corrected fixed-label Euler-product constants.  The
quotient, shadow, eta-gap, and \(S_3\) statements later in the paper are
supporting corollaries and finite certificates: they explain functoriality,
scalar no-descent, the boundary of the strict twisted gap, and a concrete
obstruction to replacing \(F_\varrho\) by \(\mathcal L_\varrho\).  They are
not separate classification theorems competing with Main Theorems~I--III.
\subsection{Plan of the paper}

\Cref{sec:preliminaries} fixes notation for primitive orbit counts,
reduced determinants, cyclic lifts, and finite-group extensions.  The core
argument is in \Cref{sec:chebotarev}: the Adams--M\"obius determinant
recovery gives the primitive Peter--Weyl channels, the Perron evaluation
gives the Frobenius-class Mertens constants, and the same finite data give
the Adams-corrected product constants.  \Cref{sec:eta-gap} isolates the
strict eta-gap certificate required for those constants, and
\Cref{sec:shadows} records the quotient and shadow structure that explains
which scalar or cyclic data are insufficient.  \Cref{sec:finite-part}
then places the scalar finite part and fixed-matrix cyclic-lift comparison
inside that scope.

The appendices collect supporting finite-state computations and their role is
described once in the dependency table.  Conditional secondary-edge material is
kept outside the compiled article; the retained non-interference statement
\Cref{prop:main-chain-secondary-noninterference} records why later secondary
tails cannot change the constants and is not used in the proofs of Main
Theorems~I--III.
\Cref{app:cyclic-lift-comparison,app:scalar-boundaries}
support the scalar finite-part and cyclic-lift comparisons used around the
finite-part normalisation.  \Cref{app:s3-model-computations,app:s3-log-certificates} supply the finite \(S_3\) matrices and printed
rational inequalities used only as the concrete witness for the
Artin-substitution obstruction in Main Theorem~III.



% END INPUT sec_introduction_scalar_and_plan.tex
% END INPUT sec_introduction.tex
% BEGIN INPUT sec_preliminaries.tex
\section{Preliminaries}\label{sec:preliminaries}

The notation in this section follows the hierarchy of the paper. We
first fix primitive orbit counts for the base shift, then the
finite-group extension and Adams operations used in the main theorem
chain. The reduced determinant, finite part, and cyclic-lift notation
come afterwards because they serve only as scalar normalisation and
canonical cyclic-lift comparison data.

\medskip
\noindent\textbf{Finite-data guide.}
All main objects are finite.  The base datum is a primitive adjacency
matrix \(A\) of a finite edge shift, with Perron root
\(\lambda=\rho(A)>0\).  A cocycle \(\tau:E_A\to G\) assigns labels in a
finite group.  For \(\varrho\in\Irr(G)\), the twisted matrix
\(A_\varrho\) is the finite block transition matrix whose traces record
periodic labels in the \(\varrho\)-channel, and
\[
  \eta:=\max_{\varrho\ne\mathbf1}\rho(A_\varrho)
\]
is the non-trivial twisted spectral radius, with the empty maximum read
as \(0\).  The strict inequality \(\eta<\lambda\) is an additional
finite-matrix hypothesis in Main Theorems~II and~III; it is not inferred
from Main Theorem~I.  For a class function \(\phi\),
\(\Pi_\phi(z)\) denotes the primitive fixed-label series,
\(\mathcal L_\phi(z)\) the periodic Artin logarithm, and
\[
  F_\phi(z)=\sum_{r\ge1}\frac{\Pi_\phi(z^r)}{r}
\]
the fixed-label Euler transform.  The distinction between
\(F_\phi\) and \(\mathcal L_\phi\) is the Adams distinction: repetitions
change labels by \(g\mapsto g^r\).  The Adams operation is
\(\psi^r\phi(g)=\phi(g^r)\), and the Adams--M\"obius coefficients are the
finite character-table coefficients obtained when each
\(\psi^r\chi_\pi\) is expanded in the irreducible characters.

The Fourier convention is Hermitian.  The forward class coefficient is
\(\chi_\varrho^\ast(C):=\overline{\chi_\varrho(C)}
=\chi_\varrho(C^{-1})\), while inverse recovery uses
\(\chi_\sigma(C)\) without conjugation.  The scalar normalisations are
\(C(A)=\det'(I-A/\lambda)^{-1}\),
\(\FP(A)\), and \(\mathcal M(A)=\gamma+\FP(A)\).  The trivial determinant
\(\det(I-zA)\) is singular at \(z=\lambda^{-1}\), so the trivial channel
is never evaluated there as an ordinary determinant branch; it enters
only through the finite part \(\mathcal M(A)\) and the reduced constant
\(C(A)\).  Non-trivial channels may be evaluated at
\(z=\lambda^{-1}\) only after the strict gap \(\eta<\lambda\) has made
their determinant branches regular.

\subsection{Mixing subshifts and primitive orbit counts}

Our standing symbolic system is a finite directed multigraph
\(\Gamma=(V,E_A)\).  We write \(o(e)\) and \(t(e)\) for the initial and
terminal vertices of an edge and
\[
  A_{ij}:=\#\{e\in E_A:o(e)=i,\ t(e)=j\}
\]
for its non-negative integer adjacency matrix.  The matrix \(A\) is
assumed primitive.  The associated one-sided edge shift is
\[
  X_A^+
  =
  \bigl\{(e_n)_{n\ge 0}\in E_A^{\NN}:
    t(e_n)=o(e_{n+1})\ \text{for all }n\ge 0\bigr\},
\]
equipped with the left shift \(\sigma\).  Thus closed length-\(n\) edge
paths are the same as fixed points of \(\sigma^n\), and we write
\[
  a_n(A):=\Tr(A^n)=\#\operatorname{Fix}(\sigma^n)
\]
for the number of points of period dividing \(n\). The number of
primitive orbits of least period \(n\) is
\begin{equation}\label{eq:primitive-mobius}
  p_n(A)=\frac{1}{n}\sum_{d\mid n}\mu(d)\,a_{n/d}(A),
\end{equation}
and the necklace identity reads
\begin{equation}\label{eq:trace-primitive-identity}
  a_n(A)=\sum_{d\mid n}d\,p_d(A).
\end{equation}

This is the convention in the finite-group sections: \(E_A\) is an edge
set and primitive orbits are orbits in the edge shift.  In the scalar
finite-matrix sections and appendices we also allow \(A\) to be an
arbitrary primitive non-negative real matrix.  There \(p_n(A)\) denotes
the algebraic M\"obius coefficient in \eqref{eq:primitive-mobius}; it is
an actual non-negative orbit count only when \(A\) is the adjacency
matrix of a finite directed multigraph.

We use the next proposition only as standard symbolic-dynamics
background: it is the usual Artin--Mazur/Bowen--Lanford zeta identity
and necklace decomposition for a mixing shift of finite type; see
\cite{ArtinMazur1965PeriodicPoints,BowenLanford1970Zeta,LindMarcus1995,ParryPollicott1990Zeta}
and, for the finite homogeneous extension setting in which these
normalisations are later used,
\cite{NooraniParry1992ChebotarevShifts}.  The proof is included to fix
normalisations for the later finite-group Euler products, and no novelty is
claimed for this supporting identity.

\begin{proposition}[Trace, primitive orbits, and Euler product]
  \label{prop:prelim-trace-primitive}
  For the adjacency matrix \(A\) of a primitive finite directed
  multigraph, the dynamical \(\zeta\)-function
  satisfies
  \begin{equation}\label{eq:zeta-def}
    \zeta_A(z)
    :=\exp\!\left(\sum_{n\ge 1}\frac{\Tr(A^n)}{n}z^n\right)
    =\frac{1}{\det(I-zA)}
    =\prod_{n\ge 1}(1-z^n)^{-p_n(A)}.
  \end{equation}
\end{proposition}

\begin{proof}
  The last equality follows from \eqref{eq:trace-primitive-identity},
  since
  \[
    \sum_{n\ge 1}\frac{a_n(A)}{n}z^n
    =\sum_{m\ge 1}\sum_{k\ge 1}\frac{p_m(A)}{k}z^{mk}
    =-\sum_{m\ge 1}p_m(A)\log(1-z^m).
  \]
  The determinant identity is standard for finite matrices.
\end{proof}

The finite-group Adams--M\"obius arguments only need the following
uniform Perron--Frobenius trace bound.  This is classical
Perron--Frobenius background for primitive non-negative matrices, stated
here in the precise form used below; compare the standard finite-matrix
Perron--Frobenius reference \cite{Seneta2006NonnegativeMatrices} and the
treatment of mixing shifts of finite type in
\cite{LindMarcus1995,ParryPollicott1990Zeta}.
The bound remains valid even when \(\lambda=1\), and it is used only as a
supporting finite-matrix estimate.

\begin{lemma}[Perron--Frobenius trace bound]
  \label{lem:pf-trace-bound}
  Let \(A\) be a primitive non-negative \(d\times d\) matrix with
  Perron root \(\lambda=\rho(A)\). Then
  \[
    \abs{\Tr(A^n)}\le d\,\lambda^n
    \qquad (n\ge 1).
  \]
  In particular, there is a constant \(C_A>0\) such that
  \[
    \abs{\Tr(A^n)}\le C_A\lambda^n
    \qquad (n\ge 1).
  \]
\end{lemma}

\begin{proof}
  Let \(\lambda,\mu_2,\dots,\mu_d\) be the eigenvalues of \(A\), counted
  with algebraic multiplicity. Since \(A\) is primitive,
  the Perron--Frobenius theorem for primitive non-negative matrices
  \cite{Seneta2006NonnegativeMatrices} gives \(\lambda>0\), the eigenvalue
  \(\lambda\) is simple, and \(|\mu_j|<\lambda\) for \(j\ge 2\). Hence
  \[
    \abs{\Tr(A^n)}
    =
    \abs{\lambda^n+\sum_{j=2}^d\mu_j^n}
    \le
    \lambda^n+\sum_{j=2}^d|\mu_j|^n
    \le
    d\,\lambda^n,
  \]
  which proves the claim.
\end{proof}

If \(\lambda=\rho(A)\) is the Perron root and \(\lambda>1\), then
\[
  \Lambda(A):=\max\{|\mu|:\mu\in\spec(A),\ \mu\neq\lambda\},
\]
with the maximum interpreted as \(0\) if there are no non-Perron
eigenvalues,
then the M\"obius coefficients \(p_n(A)\) enjoy the usual square-root
loss from the higher divisors in \eqref{eq:primitive-mobius}.  This is
again a standard SFT/Perron--Frobenius estimate, recorded with constants
because the later finite-part and twisted-determinant arguments need a
uniform error term; see, for example,
\cite{Parry1983PrimeOrbit,LindMarcus1995,ParryPollicott1990Zeta} for the
classical orbit-counting setting and
\cite{NooraniParry1992ChebotarevShifts} for finite homogeneous extensions.
The underlying spectral input is the standard primitive-matrix
Perron--Frobenius theorem \cite{Seneta2006NonnegativeMatrices}.  The
contribution here is not this
background asymptotic, but its later use inside the finite determinant and
class-product normalisations.

\begin{lemma}[Primitive M\"obius asymptotic]
  \label{lem:primitive-orbit-asymptotic}
  Let \(A\) be a primitive non-negative matrix with Perron root
  \(\lambda>1\), and let
  \[
    \Lambda(A):=\max\{|\mu|:\mu\in\spec(A),\ \mu\neq\lambda\}.
  \]
  If the set is empty, put \(\Lambda(A)=0\).
  Then
  \[
    \Tr(A^n)=\lambda^n+O(\Lambda(A)^n)
  \]
  and
  \[
    K_{\mathrm{PF}}(A)
    :=
    \begin{cases}
      \displaystyle
      \sup_{n\ge 1}
        \frac{\abs{\Tr(A^n)-\lambda^n}}{\Lambda(A)^n},
      & \Lambda(A)>0, \\[2ex]
      0,
      & \Lambda(A)=0,
    \end{cases}
    \qquad
    K_{\mathrm{tr}}(A)
    :=
    \sup_{\ell\ge 1}\frac{\abs{\Tr(A^\ell)}}{\lambda^\ell}
  \]
  are finite. Moreover,
  \[
    \abs{p_n(A)-\lambda^n/n}
    \le
    \frac{K_{\mathrm{PF}}(A)\Lambda(A)^n}{n}
    +\frac{K_{\mathrm{tr}}(A)}{n}
      \sum_{\substack{m\mid n\\ m>1}}\lambda^{n/m}
  \]
  and therefore
  \[
    \abs{p_n(A)-\lambda^n/n}
    \le
    \frac{K_{\mathrm{PF}}(A)\Lambda(A)^n}{n}
    +\frac{\lambda K_{\mathrm{tr}}(A)}{\lambda-1}
      \frac{\lambda^{n/2}}{n}.
  \]
  Hence
  \[
    p_n(A)
    =
    \frac{\lambda^n}{n}
    +O\!\left(
      \frac{\max\{\Lambda(A)^n,\lambda^{n/2}\}}{n}
    \right).
  \]
  Consequently, with
  \[
    \vartheta_A:=\max\{\Lambda(A)/\lambda,\lambda^{-1/2}\}<1,
  \]
  one has
  \[
    \frac{p_n(A)}{\lambda^n}
    =
    \frac{1}{n}+O\!\left(\frac{\vartheta_A^n}{n}\right),
  \]
  and the series
  \[
    \sum_{n\ge 1}
      \left(\frac{p_n(A)}{\lambda^n}-\frac{1}{n}\right)
  \]
  converges absolutely.
\end{lemma}

\begin{proof}
  List the eigenvalues of \(A\) with algebraic multiplicity as
  \(\lambda,\mu_2,\dots,\mu_d\). Since \(A\) is primitive,
  Perron--Frobenius theory \cite{Seneta2006NonnegativeMatrices} gives that
  \(\lambda\) is simple and
  \(|\mu_j|<\lambda\) for \(j\ge 2\). For a finite matrix,
  the trace of \(A^n\) is the sum of the \(n\)-th powers of its
  eigenvalues, counted with algebraic multiplicity. Hence
  \[
    \Tr(A^n)
    =
    \lambda^n+\sum_{j=2}^d\mu_j^n
    =
    \lambda^n+O(\Lambda(A)^n).
  \]
  If \(\Lambda(A)=0\), then every non-Perron eigenvalue is \(0\), so
  \(\Tr(A^n)=\lambda^n\) for all \(n\ge 1\), and the convention
  \(K_{\mathrm{PF}}(A)=0\) is exact. If \(\Lambda(A)>0\), then
  \[
    \frac{\abs{\Tr(A^n)-\lambda^n}}{\Lambda(A)^n}
    \le
    \sum_{j=2}^d\left|\frac{\mu_j}{\Lambda(A)}\right|^n
    \le d-1,
  \]
  so \(K_{\mathrm{PF}}(A)<\infty\). The trace bound from
  \Cref{lem:pf-trace-bound} gives \(K_{\mathrm{tr}}(A)\le d\), so
  \(K_{\mathrm{tr}}(A)<\infty\).

  Insert this estimate into the M\"obius formula
  \[
    p_n(A)=\frac{1}{n}\sum_{m\mid n}\mu(m)\Tr(A^{n/m}).
  \]
  The divisor \(m=1\) contributes \(\lambda^n/n\), with error bounded
  by \(K_{\mathrm{PF}}(A)\Lambda(A)^n/n\). For the remaining divisors,
  \(m>1\) implies \(n/m\le n/2\), and the definition of
  \(K_{\mathrm{tr}}(A)\) gives
  \[
    \sum_{\substack{m\mid n\\ m>1}}\abs{\Tr(A^{n/m})}
    \le
    K_{\mathrm{tr}}(A)\sum_{1\le k\le n/2}\lambda^k
    \le
    \frac{\lambda K_{\mathrm{tr}}(A)}{\lambda-1}\lambda^{n/2}.
  \]
  Dividing by \(n\) gives the two explicit displayed bounds and hence
  the stated M\"obius-coefficient asymptotic.
  Since \(\Lambda(A)<\lambda\) and \(\lambda>1\), the number
  \(\vartheta_A\) lies in \((0,1)\), and the final two assertions follow
  by dividing the asymptotic by \(\lambda^n\).
\end{proof}

\begin{definition}[Square-root bound]\label{def:kernel-rh}
We say that \(A\) satisfies the \emph{square-root bound} if
\[
  \Lambda(A)\le \lambda^{1/2}.
\]
\end{definition}

\subsection{Finite-group extensions and Adams operations}

Let \(G\) be a finite group, let \(\tau:E_A\to G\) be an edge cocycle,
and let \(\widetilde\Gamma\) be the skew-product extension on vertex set
\(V\times G\), with one edge
\[
  (i,g)\longrightarrow (j,g\tau(e))
  \qquad \text{for each } e\in E_A \text{ with } o(e)=i,\ t(e)=j.
\]
Its adjacency matrix is denoted by \(\widetilde A\).

For an irreducible unitary representation
\(\varrho:G\to U(V_\varrho)\) with character \(\chi_\varrho\), define
the twisted transition matrix
\begin{equation}\label{eq:twisted-matrix-def}
  (A_\varrho)_{ij}:=\sum_{e:i\to j}\varrho(\tau(e)),
\end{equation}
viewed as an operator on \(\CC^V\otimes V_\varrho\).

If \(\gamma=e_1\cdots e_n\) is a primitive orbit, write
\[
  \tau(\gamma)=\tau(e_1)\cdots\tau(e_n)\in G
\]
for its group label. A different choice of the initial edge conjugates
this element by a prefix product, so the conjugacy class of
\(\tau(\gamma)\) is intrinsic. All orbit-level evaluations below are
therefore made with class functions. For a class function
\(\phi:G\to\CC\), write
\begin{equation}\label{eq:class-function-fourier}
  \phi
  =\sum_{\varrho\in\Irr(G)}\widehat\phi(\varrho)\chi_\varrho,
  \qquad
  \widehat\phi(\varrho)
  :=\frac{1}{|G|}\sum_{g\in G}\phi(g)\overline{\chi_\varrho(g)}.
\end{equation}
With this Hermitian Fourier convention, the coefficient attached to a
conjugacy class \(C\subset G\) is the conjugated character value. We
fix the notation
\[
  \chi^\ast(C):=\overline{\chi(C)}=\chi(C^{-1})
\]
for every character \(\chi\) and every conjugacy class \(C\subset G\).
Thus every forward class expansion below uses \(\chi_\varrho^\ast(C)\),
while inversion and orthogonality pair it against the unconjugated
\(\chi_\pi(C)\). We do not identify \(\chi(C)\) with \(\chi^\ast(C)\)
unless \(\chi\) is real-valued.
For \(m\ge 1\), define the Adams operation on class functions by
\begin{equation}\label{eq:adams-class-function}
  (\psi^m\phi)(g):=\phi(g^m).
\end{equation}
For \(\varrho,\pi\in\Irr(G)\) we write
\begin{equation}\label{eq:prelim-adams-character-coefficients}
  \psi^m\chi_\varrho
  =\sum_{\pi\in\Irr(G)}a_{\varrho,\pi}^{(m)}\chi_\pi,
  \qquad
  a_{\varrho,\pi}^{(m)}
  =\frac{1}{|G|}\sum_{g\in G}
    \chi_\varrho(g^m)\overline{\chi_\pi(g)} .
\end{equation}
These are the Adams coefficients used in the determinant formulas below.

\begin{proposition}[Finite-state determinant data]
  \label{prop:finite-state-data-collection}
  \label{prop:finite-state-data-package}
The following objects form the finite-state input data used in
  Main Theorems~I--III:
  \[
    (A,X_A^+,p_n(A),\zeta_A,\Lambda(A),G,\tau,
      \widetilde A,A_\varrho,\psi^m,a_{\varrho,
      \pi}^{(m)}) .
  \]
  More precisely, once a finite directed multigraph \(\Gamma=(V,E_A)\),
  a primitive adjacency matrix \(A\), a finite group \(G\), and an edge
  cocycle \(\tau:E_A\to G\) are fixed, all the displayed objects are
  determined as follows. The shift \(X_A^+\), the primitive orbit counts
  \(p_n(A)\), and the zeta function \(\zeta_A\) are those of
  \Cref{prop:prelim-trace-primitive}. The secondary spectral radius is
  \[
    \Lambda(A)=\max\{|\mu|:\mu\in\spec(A),\ \mu\neq\lambda\},
  \]
  with value \(0\) if the set is empty. The cocycle \(\tau\) determines
  the skew-product adjacency matrix \(\widetilde A\) and the twisted
  matrices \(A_\varrho\) by \eqref{eq:twisted-matrix-def}. The Adams
  operation \(\psi^m\) and the coefficients
  \(a_{\varrho,\pi}^{(m)}\) are given by
  \eqref{eq:adams-class-function} and
  \eqref{eq:prelim-adams-character-coefficients}. Consequently every
  periodic, primitive, twisted determinant, fixed-label Euler transform,
and class-Mertens constant in Main Theorems~I--III is a function of
these finite data, together with the additional Perron-boundary
  hypotheses explicitly imposed where they are used.
\end{proposition}

\begin{proof}
  The adjacency matrix determines the edge shift \(X_A^+\) by the
  incidence rule defining admissible edge sequences. Its traces count
  fixed points of \(\sigma^n\); M\"obius inversion gives
  \(p_n(A)\) through \eqref{eq:primitive-mobius}, and the determinant
  identity in \Cref{prop:prelim-trace-primitive} gives \(\zeta_A\).
  Since \(A\) is a finite matrix, its Perron root \(\lambda\) and the
  finite set \(\spec(A)\setminus\{\lambda\}\) determine
  \(\Lambda(A)\), including the empty-set convention.

  The cocycle \(\tau:E_A\to G\) assigns one group element to each edge,
  so it determines the skew-product graph on \(V\times G\) and hence its
  adjacency matrix \(\widetilde A\). For a fixed irreducible
  representation \(\varrho\), formula \eqref{eq:twisted-matrix-def}
  is a finite sum over edges from \(i\) to \(j\), so it determines the
  finite matrix \(A_\varrho\). For a closed path, multiplying the finitely
  many labels along the path gives \(\tau(\gamma)\), up to conjugacy
  after changing the initial edge; therefore all class-function orbit
  sums used later are well-defined from \(A\) and \(\tau\).

  The map \(g\mapsto g^m\) is a finite power map on \(G\), hence
  \eqref{eq:adams-class-function} determines \(\psi^m\) on all class
  functions. The irreducible characters form an orthonormal basis for
  class functions, so expanding \(\psi^m\chi_\varrho\) in that basis gives
  the unique coefficients in
  \eqref{eq:prelim-adams-character-coefficients}. The periodic
  logarithms, primitive series, fixed-label Euler transforms, and
  determinant expressions in \S\ref{sec:chebotarev} are assembled from
  the resulting finite traces, labels, characters, and determinant
  polynomials. Evaluations at \(z=\lambda^{-1}\) require exactly the
  additional hypotheses stated in the corresponding Perron-boundary
  results, such as \(\lambda>1\) and \(\eta<\lambda\), and are therefore
  not extra data.
\end{proof}

\subsection{Reduced determinants and finite parts}

Let \(z_\ast=\lambda^{-1}\). Since \(\lambda\) is a simple eigenvalue,
\(\zeta_A\) has a simple pole at \(z_\ast\), and we define the residue
constant
\begin{equation}\label{eq:residue-constant-def}
  C(A):=\lim_{z\to z_\ast}(1-\lambda z)\zeta_A(z).
\end{equation}
If \(\mu_2,\dots,\mu_d\) are the non-Perron eigenvalues, then
\begin{equation}\label{eq:reduced-det-poly}
  F_A(t):=\frac{\det(I-tA/\lambda)}{1-t}
  =\prod_{j=2}^d\left(1-\frac{\mu_j}{\lambda}t\right).
\end{equation}
We call \(F_A\) the \emph{reduced spectral polynomial}. Moreover,
\begin{equation}\label{eq:reduced-det-constant}
  C(A)=\det\nolimits'(I-A/\lambda)^{-1}=F_A(1)^{-1}.
\end{equation}
This reduced determinant calculation is only the elementary simple-pole
finite-matrix form of the Perron normalisation used in the zeta literature;
see \cite{ParryPollicott1990Zeta} for the dynamical normalisation and
\cite{Seneta2006NonnegativeMatrices} for the primitive-matrix spectral input.

\begin{lemma}[Reduced constant at a simple eigenvalue]
  \label{lem:simple-eigenvalue-reduced-constant}
  Let \(M\) be a complex square matrix and let
  \(\alpha\neq 0\) be a simple eigenvalue of \(M\). Define
  \[
    C_\alpha(M)
    :=\lim_{z\to \alpha^{-1}}
      (1-\alpha z)\det(I-zM)^{-1}.
  \]
  Then the limit exists and
  \[
    C_\alpha(M)
    =
    \prod_{\substack{\nu\in\spec(M)\\ \nu\neq\alpha}}
      \left(1-\frac{\nu}{\alpha}\right)^{-1},
  \]
  with algebraic multiplicities in the product. In particular, if \(A\)
  is primitive with Perron root \(\lambda\), then
  \(C(A)=C_\lambda(A)\).
\end{lemma}

\begin{proof}
  Since \(\alpha\) has algebraic multiplicity one,
  \[
    \det(I-zM)
    =(1-\alpha z)
      \prod_{\substack{\nu\in\spec(M)\\ \nu\neq\alpha}}(1-\nu z),
  \]
  where the remaining eigenvalues are listed with algebraic
  multiplicity. Multiplying by \(1-\alpha z\) after inverting and then
  taking \(z\to \alpha^{-1}\) gives the displayed formula. For a
  primitive matrix, Perron--Frobenius theory
  \cite{Seneta2006NonnegativeMatrices} makes the Perron root simple, so the
  definition agrees with \eqref{eq:residue-constant-def}.
\end{proof}

When \(\lambda>1\), define the primitive Abel generating series
\[
  P_A(r):=\sum_{n\ge 1}p_n(A)\lambda^{-n}r^n,
  \qquad 0\le r<1.
\]
By \Cref{lem:primitive-orbit-asymptotic},
\(p_n(A)\lambda^{-n}=1/n+O(\vartheta^n/n)\) for some
\(\vartheta\in(0,1)\), the renormalised expression
\[
  P_A(r)+\log(1-r)
\]
has a finite limit as \(r\uparrow 1\).

\begin{definition}[Abel finite-part constant]\label{def:abel-finite-part}
For a primitive matrix \(A\) with Perron root \(\lambda>1\), we write
\[
  \FP(A):=\lim_{r\uparrow 1}\bigl(P_A(r)+\log(1-r)\bigr).
\]
The associated orbit-Mertens constant is
\[
  \mathcal M(A):=\gamma+\FP(A).
\]
\end{definition}

\subsection{Cyclic lifts}

For \(q\ge 1\) let \(P_q\) be the \(q\)-cycle permutation matrix and
define the cyclic lift
\begin{equation}\label{eq:cyclic-lift-def}
  A^{\langle q\rangle}:=A\otimes P_q.
\end{equation}
The basic spectral identity
\[
  \spec(A^{\langle q\rangle})
  =\{\mu\,\omega_q^j:\mu\in\spec(A),\ 0\le j\le q-1\}
\]
will be used repeatedly. In particular,
\(\rho(A^{\langle q\rangle})=\rho(A)\). For \(q>1\) this cyclic lift is
typically periodic rather than primitive, but \(\lambda=\rho(A)\) is
still a simple eigenvalue of \(A^{\langle q\rangle}\). We therefore use
the branch-lifted logarithmic record
\[
  \Psi_q(A):=
  -\sum_{j=2}^d
    \Log\!\left(1-\left(\frac{\mu_j}{\lambda}\right)^q\right)
\]
to encode the non-Perron spectrum. Since \(|\mu_j|<\lambda\), every
factor \(1-(\mu_j/\lambda)^q\) lies in the open right half-plane, so the
principal branch \(\Log\) is well defined here. By
\Cref{lem:simple-eigenvalue-reduced-constant},
\[
  \exp(\Psi_q(A))=q\,C_\lambda(A^{\langle q\rangle}),
\]
so the reduced constants determine only the exponentiated values unless
a branch choice is included as part of the data.
When interpreted as a \(\ZZ/q\ZZ\)-extension, this is the special
all-edges-label-\(1\) cocycle; it is not the same object as an arbitrary
cyclic quotient of a later finite-group cocycle.

\begin{definition}[Fixed-matrix cyclic-lift records]
  \label{def:spectral-fingerprints}
  In the fixed-matrix cyclic-lift comparison, the
  \emph{cyclic spectral record} of \(A\) is the sequence
  \[
    \{\Psi_q(A)\}_{q\ge 1}.
  \]
  For each \(q\ge 1\), the \emph{\(q\)-torsion cyclic record} is the
  root-of-unity layer
  \[
    \{F_A(\omega_q^j)\}_{0\le j\le q-1}.
  \]
  The associated \emph{fixed-matrix exponential record} is the growth rate
  \[
    \limsup_{q\to\infty}\abs{\Psi_q(A)}^{1/q}.
  \]
  These are auxiliary scalar records of this one matrix \(A\). They are
  not invariants of a varying family and not class-profile data for a
  finite-group cocycle unless additional descent information is supplied. In
  particular, the label records fixed-matrix data only; it does not by
  itself produce a cocycle class profile, a cyclic-lift invariant, or a
  stable invariant under changes of matrix presentation.
\end{definition}

\subsection{Standing notation}

Throughout the paper \(\omega_q=e^{2\pi i/q}\) denotes a primitive
\(q\)th root of unity, and \(\Phi_q\) the \(q\)th cyclotomic
polynomial. We reserve the notation \(\varrho\) for irreducible
representations of finite groups and keep \(\rho(\cdot)\) for spectral
radii. Sections~\ref{sec:chebotarev} and~\ref{sec:shadows} use the
graph and cocycle picture for finite-group extensions. The later scalar
section, Section~\ref{sec:finite-part}, uses only the primitive matrix
formalism and cyclic lifts.


% END INPUT sec_preliminaries.tex
% BEGIN INPUT sec_chebotarev.tex
% BEGIN INPUT sec_chebotarev_main_inversion.tex
\section{Adams--M\"obius primitive inversion for finite-group extensions}
\label{sec:chebotarev}

This section proves the finite-group primitive inversion results. Twisted matrices
\(A_\varrho\) first linearise periodic data in Peter--Weyl coordinates;
the Adams--M\"obius identity then removes repetitions of primitive
orbits, where labels change from \(g\) to \(g^m\). We use standard
character-ring and Adams-operation formalism as in
\cite{AtiyahTall1969GroupRepresentationsLambdaRings,Serre1977LinearRepresentations,Knutson1973LambdaRings}, together with
the classical Witt-vector and necklace formalism for passing from
repeated to primitive data
\cite{DressSiebeneicher1988BurnsideRing,MetropolisRota1983Necklaces,Wehrhahn1990AperiodicRings}.

The proof is organised in two labelled blocks and several bridge
statements.  Only the first block is part of Main Theorem~I. The block
\Cref{lem:twisted-trace-identity,thm:class-function-linearisation,cor:class-function-adams-mobius,cor:primitive-peter-weyl-determinant,cor:primitive-peter-weyl-trace},
is the Adams--M\"obius determinant and coefficient
mechanism and does not use a twisted spectral gap. The second block,
\Cref{thm:class-mertens-explicit,prop:class-mertens-deviations,thm:class-mertens-fourier}, belongs to Main Theorem~II and starts only
after the twisted spectral gap
\(\eta=\max_{\varrho\neq\mathbf 1}\rho(A_\varrho)<\lambda\) is imposed
and converts the primitive channels into the class-indexed constants
\(\mathcal M_C(A,\tau;G)\), from which the channels are read back by
non-abelian character orthogonality. The
intervening class-decomposition formulas are bridge statements between
these two blocks, not extra headline theorems.

Throughout the section, the underlying finite-state data are the data
fixed in \Cref{prop:finite-state-data-collection}; the additional boundary
assumptions \(\lambda>1\) and \(\eta<\lambda\) are invoked only in the
statements where Perron-point values or Mertens constants are evaluated.

\subsection{Main Theorem I: Adams--M\"obius primitive determinant
linearisation}

The first main block has five steps. The twisted trace identity converts
matrices into periodic Peter--Weyl sums. The determinant linearisation
rewrites those sums as logarithms of twisted determinants. The
Adams--M\"obius inversion then removes repetitions of primitive orbits,
and the irreducible-channel formula rewrites the result entirely in
terms of the determinant data. Finally, coefficient extraction gives the
primitive Peter--Weyl trace formula. The preliminary class-indicator and
block-diagonalisation propositions below set up this mechanism; they
are not separate main results.

For reference, the logical dependency map for this block is
\[
\begin{array}{c}
  \text{twisted traces}\to\text{determinant linearisation}\\[2pt]
  \to\text{Adams--M\"obius inversion}\to\text{irreducible channels}\\[2pt]
  \to\text{class-count recovery}.
\end{array}
\]
The next proposition records the assembled finite construction with its
hypotheses and outputs, so that the introduction, the Peter--Weyl bridge, and
the later Mertens section all refer to a single proved interface rather than to
an informal chain of aliases.

\begin{proposition}[Main Theorem I finite determinant interface]
  \label{prop:main-theorem-i-assembly-map}
  The following implications assemble Main Theorem~I without any use of
  the Perron-boundary hypothesis \(\eta<\lambda\).
  \begin{enumerate}[label=\textup{(\alph*)},leftmargin=*]
    \item \Cref{lem:twisted-trace-identity,thm:class-function-linearisation}
    identify, for every class function \(\phi\), the periodic logarithm
    \(\mathcal L_\phi(z)\) on \(|z|<\lambda^{-1}\) as a finite
    character-table linear combination of the determinant logarithms
    \(\log\det(I-zA_\pi)^{-1}\), \(\pi\in\Irr(G)\).
    \item \Cref{cor:class-function-adams-mobius} applies the Adams power
    maps in the character ring to recover the primitive class series
    \(\Pi_\phi(z)\) from those periodic determinant logarithms.
    \item \Cref{cor:primitive-peter-weyl-determinant,thm:primitive-peter-weyl-trace} specialize the preceding identity to
    irreducible characters and give, respectively, the determinant-germ
    formula and the coefficient formula for every primitive Peter--Weyl
    channel.
    \item \Cref{thm:class-count-trace-recovery} applies finite character
    orthogonality to those channels and gives the primitive Frobenius-class
    counts from the same finite determinant and Adams data.
  \end{enumerate}
  Consequently Main Theorem~I is exactly the finite-state determinant
  construction consisting of the twisted matrices \(A_\pi\), the finite
  character table of \(G\), and the power maps on conjugacy classes.
\end{proposition}

\begin{proof}
  Part~(a) is the content of
  \Cref{lem:twisted-trace-identity,thm:class-function-linearisation}: the
  trace identity expresses periodic labelled traces in Peter--Weyl
  coordinates, and the usual identity
  \(\log\det(I-zA_\pi)^{-1}=\sum_{n\ge1}\Tr(A_\pi^n)z^n/n\) converts those
  traces into determinant logarithms in the open Perron disc. Part~(b) is
  exactly \Cref{cor:class-function-adams-mobius}; its proof is M\"obius
  inversion of the primitive expansion with the label changed by
  \(g\mapsto g^m\), hence by the Adams operation \(\psi^m\). Part~(c)
  follows by substituting \(\phi=\chi_\varrho\) and then either leaving the
  identity as a determinant germ or extracting coefficients. The Adams
  coefficients are the finite inner products
  \(a_{\varrho,\pi}^{(m)}\) of
  \eqref{eq:adams-coefficient-inner-product}, so no analytic input beyond
  convergence in \(|z|<\lambda^{-1}\) is introduced. Part~(d) is the
  orthogonality inversion for the basis of irreducible characters on class
  functions, applied coefficient by coefficient. These four steps are
  precisely the determinant, primitive-channel, and class-recovery assertions
  claimed as Main Theorem~I.
\end{proof}

For a conjugacy class \(C\subset G\), we write \(\chi_\varrho(C)\) for
the common value of \(\chi_\varrho\) on \(C\), and \(\mathbf 1_C\) for
the indicator function of \(C\). If \(x\in\operatorname{Fix}(\sigma^n)\)
corresponds to the closed edge path \(e_0\cdots e_{n-1}\), define
\[
  \tau_n(x):=\tau(e_0)\cdots\tau(e_{n-1})\in G.
\]
For a class function \(\phi:G\to\CC\), set
\begin{equation}\label{eq:class-logarithm-def}
  T_{\phi,n}:=\sum_{x\in\operatorname{Fix}(\sigma^n)}\phi(\tau_n(x)),
  \qquad
  \mathcal L_\phi(z):=\sum_{n\ge 1}\frac{T_{\phi,n}}{n}z^n,
\end{equation}
and
\begin{equation}\label{eq:primitive-class-series-def}
  \Pi_\phi(z):=\sum_{\gamma\in\cP}\phi(\tau(\gamma))z^{|\gamma|}.
\end{equation}
The fixed-label Euler transform of the primitive series is
\begin{equation}\label{eq:fixed-label-euler-transform-def}
  F_\phi(z):=\sum_{r\ge1}\frac{\Pi_\phi(z^r)}{r}.
\end{equation}
Since \(\phi\) is bounded and \(p_n(A)\ll \lambda^n/n\), the series
\(\Pi_\phi(z)\) converges absolutely for \(|z|<\lambda^{-1}\).
The transform \(F_\phi\) also converges absolutely in this disc. It is
important that \(F_\phi\) is not the Artin logarithm
\(\mathcal L_\phi\): by \Cref{lem:class-logarithm-primitive-expansion},
\(\mathcal L_\phi(z)=\sum_{r\ge1}\Pi_{\psi^r\phi}(z^r)/r\), where the
label is changed by the Adams operation under repetition, whereas
\(F_\phi\) keeps the primitive label fixed in every Euler factor.
These symbols are generating functions for deterministic orbit data.
They are not transition kernels or probabilities attached to repeated
measurements.

\begin{center}
\renewcommand{\arraystretch}{1.18}
\begin{tabular}{L{0.25\textwidth}L{0.32\textwidth}L{0.34\textwidth}}
\toprule
Symbol & Orbit datum & Role in the main chain \\
\midrule
\(\mathcal L_\phi(z)\)
& Periodic fixed points; a primitive label \(g\) appears as \(g^r\) on
the \(r\)-fold repeat.
& Determinant logarithm recovered directly from the twisted matrices. \\
\(\Pi_\phi(z)\)
& Primitive orbits with the primitive label kept once.
& Obtained from \(\mathcal L_{\psi^m\phi}\) by Adams--M\"obius
inversion. \\
\(F_\phi(z)\)
& Euler transform \(\sum_{r\ge1}\Pi_\phi(z^r)/r\) with the primitive label
fixed in every Euler factor.
& Supplies Frobenius-class product constants in Main Theorem~III. \\
\bottomrule
\end{tabular}
\end{center}

\medskip
\noindent\textbf{Fourier convention.}
\textbf{All Fourier coefficients in this section use the Hermitian
character convention}
\[
  \widehat\phi(\varrho)
  =\frac{1}{|G|}\sum_{g\in G}\phi(g)\overline{\chi_\varrho(g)}.
\]
\textbf{Thus every class-indicator expansion and every later Fourier
recovery uses the fixed forward class coefficient
\(\chi_\varrho^\ast(C):=\overline{\chi_\varrho(C)}
=\chi_\varrho(C^{-1})\).
We keep this \(\ast\)-notation throughout, and the inverse recovery is
always written separately with \(\chi_\pi(C)\).}

\medskip
\noindent\textbf{Logarithm convention for \S\ref{sec:chebotarev}.}
For each \(\varrho\in\Irr(G)\) and each \(|z|<\lambda^{-1}\), we set
\[
  \log\det(I-zA_\varrho)^{-1}
  :=
  \sum_{n\ge 1}\frac{\Tr(A_\varrho^n)}{n}z^n,
\]
the analytic branch normalized by
\(\log\det(I-0\cdot A_\varrho)^{-1}=0\). This is well defined on
\(|z|<\lambda^{-1}\) because \(\rho(A_\varrho)\le \lambda\) by
\Cref{prop:twisted-spectral-radius-bound}. Every determinant logarithm
on the disc \(|z|<\lambda^{-1}\) is taken in this branch. If later the
twisted spectral gap
\(\eta=\max_{\varrho\neq\mathbf 1}\rho(A_\varrho)<\lambda\) is assumed,
then for every non-trivial \(\varrho\) the matrix \(I-zA_\varrho\) is
invertible on the closed disc \(|z|\le \lambda^{-1}\); the same analytic
branch therefore extends to a neighbourhood of that closed disc, and
\(\log\det(I-\lambda^{-1}A_\varrho)^{-1}\) means this continuation value.
Likewise, for every \(m\ge 2\) and every \(\pi\in\Irr(G)\), the value
\(\log\det(I-\lambda^{-m}A_\pi)^{-1}\) is the same branch evaluated at
\(z=\lambda^{-m}\). We never evaluate
\(\log\det(I-\lambda^{-1}A_{\mathbf 1})^{-1}\), since
\(I-\lambda^{-1}A\) is singular.

The following class-indicator formula is the standard finite-group character
orthogonality expansion, in the normalization of
\cite{Serre1977LinearRepresentations}.  We record it only to fix the
Peter--Weyl coordinates used by the determinant block.

\begin{proposition}[Conjugacy-class expansion]
  \label{prop:class-indicator-expansion}
  For every conjugacy class \(C\subset G\) and every \(g\in G\),
  \[
    \mathbf 1_C(g)
    =\frac{|C|}{|G|}
      \sum_{\varrho\in\Irr(G)}
        \chi_\varrho^\ast(C)\,\chi_\varrho(g)
  \]
\end{proposition}

\begin{proof}
  The irreducible characters form an orthonormal basis of the class
  functions on \(G\) for the Hermitian inner product
  \[
    \langle f,h\rangle
    :=\frac{1}{|G|}\sum_{g\in G}f(g)\overline{h(g)}.
  \]
  Hence
  \[
    \mathbf 1_C
    =\sum_{\varrho\in\Irr(G)}
      \langle \mathbf 1_C,\chi_\varrho\rangle \chi_\varrho.
  \]
  Since \(\chi_\varrho\) is constant on \(C\),
  \[
    \langle \mathbf 1_C,\chi_\varrho\rangle
    =\frac{1}{|G|}\sum_{g\in C}\overline{\chi_\varrho(g)}
    =\frac{|C|}{|G|}\chi_\varrho^\ast(C).
  \]
  This gives the claim.
\end{proof}

The block diagonalisation below is the corresponding standard Peter--Weyl
decomposition of the regular representation, again in the normalization of
\cite{Serre1977LinearRepresentations}.  The point here is only its
application to the labelled adjacency operator, so this proposition is
supporting finite-group infrastructure rather than a new representation-
theoretic result.

\begin{proposition}[Peter--Weyl determinant factorisation]
  \label{prop:kernel-peter-weyl-block-diagonalisation}
  Let \(\widetilde A\) be the adjacency matrix of the skew-product
  extension attached to \(\tau\). Let matrices act on row vectors. With
  the right regular action \(R_h\delta_g=\delta_{gh}\) on \(\CC[G]\), one
  has
  \[
    \widetilde A
    =
    \sum_{e\in E_A} E_{o(e),t(e)}\otimes R_{\tau(e)},
  \]
  and under the Fourier--Wedderburn decomposition
  \[
    \CC[G]\cong \bigoplus_{\varrho\in\Irr(G)}M_{d_\varrho}(\CC),
    \qquad d_\varrho:=\dim\varrho,
  \]
  the operator \(R_h\) becomes, on the \(\varrho\)-summand, right
  multiplication \(X\mapsto X\varrho(h)\). Equivalently, \(\widetilde A\)
  is similar to the direct sum, over \(\varrho\in\Irr(G)\), of
  \(d_\varrho\) copies of \(A_\varrho\). In particular,
  \[
    \det(I-z\widetilde A)
    =
    \prod_{\varrho\in\Irr(G)}
      \det(I-zA_\varrho)^{\dim\varrho}.
  \]
\end{proposition}

\begin{proof}
  Let \(\{\delta_g\}_{g\in G}\) be the standard row basis of \(\CC[G]\),
  and let \(E_{ij}\) denote the matrix units acting on row vectors in
  \(\CC^V\), so that \(e_iE_{ij}=e_j\). By definition of the skew
  product, an edge \(e:i\to j\) sends \((i,g)\) to \((j,g\tau(e))\).
  Hence the corresponding adjacency operator on row vectors in
  \(\CC^V\otimes\CC[G]\) is
  \[
    \widetilde A
    =
    \sum_{e\in E_A} E_{o(e),t(e)}\otimes R_{\tau(e)},
  \]
  where \(R_h\delta_g=\delta_{gh}\) is the right regular action.

  For each \(\varrho\in\Irr(G)\), define the Fourier coefficient map
  \[
    \mathcal F_\varrho:\CC[G]\to M_{d_\varrho}(\CC),
    \qquad
    \mathcal F_\varrho\!\left(\sum_{g\in G}a_g\delta_g\right)
    :=\sum_{g\in G}a_g\,\varrho(g).
  \]
  Then for every \(h\in G\),
  \[
    \mathcal F_\varrho\bigl((\sum_g a_g\delta_g)R_h\bigr)
    =\sum_g a_g\,\varrho(gh)
    =\mathcal F_\varrho\!\left(\sum_g a_g\delta_g\right)\varrho(h).
  \]
  Thus right multiplication by \(h\) becomes right multiplication by
  \(\varrho(h)\) on the \(\varrho\)-summand. The standard
  Fourier--Wedderburn theorem gives a vector-space isomorphism
  \[
    \mathcal F:\CC[G]\xrightarrow{\sim}
    \bigoplus_{\varrho\in\Irr(G)}M_{d_\varrho}(\CC)
  \]
  intertwining the right regular action with the direct sum of the maps
  \(X\mapsto X\varrho(h)\).

  Now write a matrix \(X\in M_{d_\varrho}(\CC)\) by its rows. Right
  multiplication \(X\mapsto X\varrho(h)\) acts independently on each row,
  so it is the direct sum of \(d_\varrho\) copies of \(\varrho(h)\). After
  tensoring with \(\CC^V\), the operator induced by \(\widetilde A\) on
  the \(\varrho\)-summand is therefore the direct sum of \(d_\varrho\)
  copies of
  \[
    \sum_{e\in E_A} E_{o(e),t(e)}\otimes \varrho(\tau(e))
    =
    A_\varrho.
  \]
  Hence \(\widetilde A\) is similar to
  \(\bigoplus_{\varrho\in\Irr(G)}A_\varrho^{\oplus d_\varrho}\). Taking
  determinants gives
  \[
    \det(I-z\widetilde A)
    =
    \prod_{\varrho\in\Irr(G)}
      \det(I-zA_\varrho)^{\dim\varrho},
  \]
  as claimed.
\end{proof}

The next lemma is a finite-matrix/Peter--Weyl translation of the twisted gap
into the skew-product peripheral spectrum.  It is included as a background
checkable criterion for the hypotheses used later, not as a new dynamical
extension criterion; compare the finite homogeneous extension framework of
\cite{Noorani1997HomogeneousExtensionsMarkov,NooraniParry1992ChebotarevShifts}.

\begin{lemma}[Twisted gap as a skew-product peripheral-spectrum condition]
  \label{lem:twisted-gap-skew-product-peripheral}
  Let \(\widetilde A\) be the skew-product adjacency matrix attached to
  \(\tau\), and set
  \[
    \eta:=\max_{\varrho\neq\mathbf 1}\rho(A_\varrho).
  \]
  If there is no non-trivial irreducible representation, in particular
  if \(G\) is trivial, this maximum is empty and by convention
  \(\eta=0\).
  Then
  \[
    \spec(\widetilde A)
    =
    \bigsqcup_{\varrho\in\Irr(G)}
      \spec(A_\varrho)^{\sqcup \dim\varrho},
  \]
  where the union is taken with algebraic multiplicity. Consequently,
  the following are equivalent:
  \begin{enumerate}[label=\rm(\roman*)]
  \item \(\eta<\lambda\).
  \item Every eigenvalue of every non-trivial Peter--Weyl block
    \(A_\varrho\) (\(\varrho\neq\mathbf 1\)) has modulus strictly less
    than \(\lambda\).
  \item The peripheral spectral multiset of \(\widetilde A\) is exactly
    the single eigenvalue \(\lambda\) with algebraic multiplicity \(1\).
  \item \(\widetilde A\) has a simple strictly dominant Perron
    eigenvalue \(\lambda\).
  \end{enumerate}
  If, in addition, \(\widetilde A\) is irreducible, then these
  conditions are also equivalent to primitivity of \(\widetilde A\).
\end{lemma}

\begin{proof}
  By
  \Cref{prop:kernel-peter-weyl-block-diagonalisation},
  \(\widetilde A\) is similar to
  \(\bigoplus_{\varrho\in\Irr(G)}A_\varrho^{\oplus\dim\varrho}\). Hence
  the multiset of eigenvalues of \(\widetilde A\) is exactly
  \[
    \bigsqcup_{\varrho\in\Irr(G)}
      \spec(A_\varrho)^{\sqcup \dim\varrho},
  \]
  which is the displayed spectral decomposition. The trivial
  representation contributes the block \(A_{\mathbf 1}=A\). Since \(A\)
  is primitive, Perron--Frobenius theory gives
  \[
    \spec(A)\cap\{z\in\CC:\abs{z}=\lambda\}=\{\lambda\},
  \]
  and the eigenvalue \(\lambda\) is simple.

  Therefore \(\eta<\lambda\) holds exactly when every non-trivial block
  \(A_\varrho\) has spectral radius strictly less than \(\lambda\).
  Since every eigenvalue of \(A_\varrho\) has modulus at most
  \(\rho(A_\varrho)\), this is equivalent to saying that every
  eigenvalue of every non-trivial block has modulus strictly less than
  \(\lambda\), which is statement \rm(ii). Because the trivial block
  contributes exactly one peripheral eigenvalue, namely the simple
  eigenvalue \(\lambda\), statement \rm(ii) is equivalent to saying
  that the peripheral spectral multiset of \(\widetilde A\) consists of
  that single copy of \(\lambda\). This is statement \rm(iii). A mere
  point-set identity
  \(\spec(\widetilde A)\cap\{z\in\CC:\abs{z}=\lambda\}=\{\lambda\}\)
  would not exclude extra copies of \(\lambda\) coming from
  non-trivial blocks, so multiplicity is essential. Finally,
  statement \rm(iii) says precisely that \(\widetilde A\) has spectral
  radius \(\lambda\), that \(\lambda\) is algebraically simple, and
  that every other eigenvalue has modulus strictly less than
  \(\lambda\), which is statement \rm(iv).

  Finally, if \(\widetilde A\) is irreducible, then Perron--Frobenius
  theory for irreducible non-negative matrices says that irreducibility
  together with the absence of any peripheral eigenvalue other than the
  simple Perron eigenvalue is equivalent to primitivity. Thus, under
  irreducibility, statements \rm(i)--\rm(iv) are also equivalent to
  primitivity of \(\widetilde A\).
\end{proof}

The following estimate is the standard domination of a unitary-twisted
adjacency block by the untwisted Perron--Frobenius block.  It is recorded only
to support the convergence and branch choices in the determinant formulae.

\begin{proposition}[Twisted spectral-radius bound]
  \label{prop:twisted-spectral-radius-bound}
  For every irreducible unitary representation
  \(\pi\in\Irr(G)\),
  \[
    \rho(A_\pi)\le \lambda.
  \]
\end{proposition}

\begin{proof}
  By \Cref{prop:kernel-peter-weyl-block-diagonalisation}, matrices act
  on row vectors throughout this subsection. Let \(\alpha\) be an
  eigenvalue of \(A_\pi\), and choose a non-zero left eigenvector
  \(v=(v_i)_{i\in V}\) with
  \(vA_\pi=\alpha v\). With the block convention
  \[
    (A_\pi)_{ij}=\sum_{e:i\to j}\pi(\tau(e)),
  \]
  this means
  \[
    \alpha v_j
    =
    \sum_i\sum_{e:i\to j}v_i\pi(\tau(e))
    \qquad (j\in V).
  \]
  Now collect the block norms into the column vector
  \(w=(w_j)_{j\in V}^{\mathsf T}\), where \(w_j:=\|v_j\|\). Since
  \(\pi(\tau(e))\) is unitary for every edge \(e\),
  \[
    |\alpha|\,w_j
    =
    \left\|\sum_i\sum_{e:i\to j}v_i\pi(\tau(e))\right\|
    \le \sum_i\sum_{e:i\to j}\|v_i\|
    =\sum_i A_{ij}w_i
    =(A^{\mathsf T}w)_j.
  \]
  Hence \(w\ge 0\), \(w\neq 0\), and
  \[
    |\alpha|\,w\le A^{\mathsf T}w
  \]
  coordinatewise. Let \(r>0\) be a right Perron vector of \(A\), so
  \(Ar=\lambda r\), equivalently
  \(r^{\mathsf T}A^{\mathsf T}=\lambda r^{\mathsf T}\). Left-multiplying
  the last inequality by \(r^{\mathsf T}\) gives
  \[
    |\alpha|\,r^{\mathsf T}w
    \le
    r^{\mathsf T}A^{\mathsf T}w
    =
    \lambda\,r^{\mathsf T}w.
  \]
  Because \(r>0\) and \(w\neq 0\), one has \(r^{\mathsf T}w>0\), and
  therefore \(|\alpha|\le \lambda\). This holds for every eigenvalue
  \(\alpha\) of \(A_\pi\), so \(\rho(A_\pi)\le \lambda\).
\end{proof}

\noindent
Up to this point, the material is representation-theoretic
infrastructure. The specific content of the present paper begins when
the determinant data are pushed to the primitive level by
Adams--M\"obius inversion.

The next lemma fixes the trace convention used throughout the rest of
the section. It is the only place where the twisted matrices are
identified with fixed-point sums; afterwards all primitive formulas are
obtained by class-function expansion and Adams--M\"obius inversion.
Together with the determinant linearisation and its primitive
corollaries below, it forms Main Theorem~I from the introduction.
The point of spelling it out separately is that the passage from
determinants to primitive data has two genuinely distinct stages:
matrix traces first give periodic Peter--Weyl sums, and only then does
Adams--M\"obius inversion remove repetitions of primitive orbits.
Thus the identity in \Cref{lem:twisted-trace-identity} is not a
primitive-orbit statement; it is the periodic input on which the
primitive inversion operates.

\begin{lemma}[Twisted trace identity]
  \label{lem:twisted-trace-identity}
  For every irreducible representation \(\varrho\in\Irr(G)\) and every
  \(n\ge 1\),
  \[
    \Tr(A_\varrho^n)
    =
    \sum_{x\in\operatorname{Fix}(\sigma^n)}
      \chi_\varrho(\tau_n(x)).
  \]
\end{lemma}

\begin{proof}
  View \(A_\varrho\) as a block matrix indexed by vertices, with
  \((i,j)\)-block
  \[
    (A_\varrho)_{ij}=\sum_{e:i\to j}\varrho(\tau(e))
    \in\End(V_\varrho).
  \]
  The \((i_0,i_0)\)-block of \(A_\varrho^n\) is obtained by multiplying
  blocks along all length-\(n\) paths which start and end at \(i_0\):
  \[
    (A_\varrho^n)_{i_0i_0}
    =
    \sum_{\substack{e_0\cdots e_{n-1}\\
          e_k:i_k\to i_{k+1},\ i_n=i_0}}
      \varrho(\tau(e_0))\cdots\varrho(\tau(e_{n-1})).
  \]
  Since \(\varrho\) is a representation, the product in each summand is
  \[
    \varrho(\tau(e_0)\cdots\tau(e_{n-1})).
  \]
  Taking the matrix trace means summing the ordinary trace of the
  diagonal vertex blocks:
  \[
    \Tr(A_\varrho^n)
    =
    \sum_{i_0}
    \sum_{\substack{e_0\cdots e_{n-1}\\
          e_k:i_k\to i_{k+1},\ i_n=i_0}}
      \operatorname{tr}_{V_\varrho}
      \varrho(\tau(e_0)\cdots\tau(e_{n-1})).
  \]
  The ordinary trace on \(V_\varrho\) is the character
  \(\chi_\varrho\). Finally, because the base is presented by the
  directed multigraph underlying \(A\), a choice of \(i_0\) and of a closed
  length-\(n\) admissible edge path based at \(i_0\) is exactly a
  fixed point of \(\sigma^n\) with its time-zero coordinate fixed. No
  quotient by cyclic rotation is taken here; fixed points are
  basepointed closed paths, and this is why there is no primitive-orbit
  factor \(1/n\) in the trace identity. Equivalently, in the edge-shift
  presentation this is the usual basepointed closed path. The product
  \(\tau(e_0)\cdots\tau(e_{n-1})\) is exactly \(\tau_n(x)\). Hence the
  last display is
  \[
    \sum_{x\in\operatorname{Fix}(\sigma^n)}
      \chi_\varrho(\tau_n(x)),
  \]
  as claimed.
\end{proof}

\begin{theorem}[Main Theorem I(a): class-function determinant
  linearisation]
  \label{thm:class-function-linearisation}
  Let \(\phi:G\to\CC\) be a class function and write
  \[
    \phi=\sum_{\varrho\in\Irr(G)}\widehat\phi(\varrho)\chi_\varrho
  \]
  as in \eqref{eq:class-function-fourier}. Then, for \(|z|<\lambda^{-1}\),
  \begin{equation}\label{eq:class-logarithm-determinant-linearisation}
    \mathcal L_\phi(z)
    =
    \sum_{\varrho\in\Irr(G)}
      \widehat\phi(\varrho)\,
      \log\det(I-zA_\varrho)^{-1}.
  \end{equation}
  Here each determinant logarithm is taken in the branch fixed by the
  logarithm convention above.
\end{theorem}

\begin{proof}
  For \(n\ge 1\), expanding \(\phi\) in the irreducible character basis
  gives
  \[
    T_{\phi,n}
    =\sum_{x\in\operatorname{Fix}(\sigma^n)}
      \sum_{\varrho\in\Irr(G)}
        \widehat\phi(\varrho)\chi_\varrho(\tau_n(x))
    =\sum_{\varrho\in\Irr(G)}
      \widehat\phi(\varrho)\Tr(A_\varrho^n),
  \]
  by \Cref{lem:twisted-trace-identity}. Therefore,
  \[
    \mathcal L_\phi(z)
    =\sum_{n\ge 1}\frac{z^n}{n}
      \sum_{\varrho\in\Irr(G)}
        \widehat\phi(\varrho)\Tr(A_\varrho^n).
  \]
  Moreover,
  \[
    \sum_{n\ge 1}\frac{|z|^n}{n}\abs{T_{\phi,n}}
    \le
    \norm{\phi}_\infty
    \sum_{n\ge 1}\frac{\Tr(A^n)}{n}|z|^n
    =
    \norm{\phi}_\infty \log\det(I-|z|A)^{-1}
    <\infty
  \]
  for \(|z|<\lambda^{-1}\). Since
  \[
    \abs{T_{\phi,n}}
    \le \norm{\phi}_\infty \#\operatorname{Fix}(\sigma^n)
    =\norm{\phi}_\infty \Tr(A^n),
  \]
  the series is absolutely convergent for \(|z|<\lambda^{-1}\), and we
  may exchange the order of summation to obtain
  \[
    \mathcal L_\phi(z)
    =
    \sum_{\varrho\in\Irr(G)}
      \widehat\phi(\varrho)
      \sum_{n\ge 1}\frac{\Tr(A_\varrho^n)}{n}z^n
    =
    \sum_{\varrho\in\Irr(G)}
      \widehat\phi(\varrho)\,
      \log\det(I-zA_\varrho)^{-1}.
  \]
  The last equality uses the power-series definition of the determinant
  logarithm. It is valid in this disc because
  \(\rho(A_\varrho)\le \lambda\) by
  \Cref{prop:twisted-spectral-radius-bound}, so
  \(I-zA_\varrho\) is invertible for \(|z|<\lambda^{-1}\).
  This proves \eqref{eq:class-logarithm-determinant-linearisation}.
\end{proof}

% END INPUT sec_chebotarev_main_inversion.tex
% BEGIN INPUT sec_chebotarev_fixed_label_and_boundary.tex
\begin{lemma}[Primitive expansion of class-function logarithms]
  \label{lem:class-logarithm-primitive-expansion}
  For every class function \(\phi:G\to\CC\) and every
  \(|z|<\lambda^{-1}\),
  \begin{equation}\label{eq:class-logarithm-primitive-expansion}
    \mathcal L_\phi(z)
    =
    \sum_{\gamma\in\cP}\sum_{r\ge 1}
      \frac{\phi(\tau(\gamma)^r)}{r}z^{r|\gamma|}.
  \end{equation}
\end{lemma}

\begin{proof}
  Every fixed point of \(\sigma^n\) lies on a unique primitive orbit
  \(\gamma\) whose length \(\ell=|\gamma|\) divides \(n\). Writing
  \(n=r\ell\), the \(r\)-fold repetition of \(\gamma\) contributes
  exactly \(\ell\) fixed points to \(\operatorname{Fix}(\sigma^n)\), one
  for each starting edge of \(\gamma\). If
  \(g=\tau(e_0)\cdots\tau(e_{\ell-1})\) is the label from one starting
  point and \(h=\tau(e_0)\cdots\tau(e_{a-1})\) is the prefix leading to
  another starting point, then the new primitive label is conjugate to
  \(g\), namely \(h^{-1}gh\). Its \(r\)-fold repeated label is therefore
  conjugate to \(g^r\). Since \(\phi\) is a class function, all these
  \(\ell\) fixed points have the same value, namely
  \(\phi(\tau(\gamma)^r)\). Their total contribution to
  \(T_{\phi,n}z^n/n\) is therefore
  \[
    \frac{\ell}{r\ell}\phi(\tau(\gamma)^r)z^{r\ell}
    =
    \frac{\phi(\tau(\gamma)^r)}{r}z^{r|\gamma|}.
  \]
  Summing over primitive orbits and repetitions gives
  \eqref{eq:class-logarithm-primitive-expansion}.

  The rearrangement is justified by absolute convergence. Indeed,
  \[
    \sum_{\gamma\in\cP}\sum_{r\ge 1}
      \frac{|z|^{r|\gamma|}}{r}
    =
    \sum_{n\ge 1}\frac{\#\operatorname{Fix}(\sigma^n)}{n}|z|^n
    \le
    \sum_{n\ge 1}\frac{\Tr(A^n)}{n}|z|^n
    <\infty
  \]
  for \(|z|<\lambda^{-1}\).
\end{proof}

\begin{remark}[Where ordinary M\"obius inversion stops]
  \label{rem:ordinary-mobius-stops}
  If \(\phi=\mathbf 1\), the factor
  \(\phi(\tau(\gamma)^r)\) in
  \eqref{eq:class-logarithm-primitive-expansion} is independent of
  \(r\), and \Cref{cor:class-function-adams-mobius} reduces to the
  usual scalar M\"obius inversion for primitive orbit counts. For a
  non-trivial class function this independence fails: the \(r\)-fold
  repetition changes the label from \(g\) to \(g^r\). The Adams operator
  \(\psi^r\) is exactly the bookkeeping needed for this change. Thus the
  primitive finite-group step below is not obtained by applying the
  scalar formula representation by representation; it uses the Adams
  action on the character ring before M\"obius inversion can remove
  repetitions. For the general ``primitive from repetitions'' viewpoint,
  compare the necklace/Witt formalism in
  \cite{MetropolisRota1983Necklaces,Wehrhahn1990AperiodicRings}.
\end{remark}

\begin{corollary}[Main Theorem I(b): Adams--M\"obius primitive
  inversion]
  \label{cor:class-function-adams-mobius}
  For every class function \(\phi:G\to\CC\) and every \(|z|<\lambda^{-1}\),
  \begin{equation}\label{eq:class-adams-mobius-inversion}
    \Pi_\phi(z)
    =
    \sum_{m\ge 1}\frac{\mu(m)}{m}\,
      \mathcal L_{\psi^m\phi}(z^m).
  \end{equation}
\end{corollary}

\begin{proof}
  For each \(m\ge 1\), apply
  \Cref{lem:class-logarithm-primitive-expansion} to \(\psi^m\phi\) and
  then substitute \(z^m\). Since
  \((\psi^m\phi)(\tau(\gamma)^r)=\phi(\tau(\gamma)^{mr})\), one gets
  \[
    \mathcal L_{\psi^m\phi}(z^m)
    =
    \sum_{\gamma\in\cP}\sum_{r\ge 1}
      \frac{\phi(\tau(\gamma)^{mr})}{r}z^{mr|\gamma|}.
  \]
  Multiply by \(\mu(m)/m\) and sum over \(m\ge 1\). The resulting
  \(m,r,\gamma\) series may be rearranged absolutely: since \(\phi\) is
  bounded,
  \[
    \sum_{m\ge 1}\frac{1}{m}
      \sum_{\gamma\in\cP}\sum_{r\ge 1}
        \frac{\abs{\phi(\tau(\gamma)^{mr})}}{r}|z|^{mr|\gamma|}
    \le
    \norm{\phi}_\infty
    \sum_{m\ge 1}\frac{1}{m}
      \sum_{n\ge 1}\frac{\Tr(A^n)}{n}|z|^{mn},
  \]
  where the inner identity
  \[
    \sum_{\gamma\in\cP}\sum_{r\ge 1}\frac{|z|^{mr|\gamma|}}{r}
    =
    \sum_{n\ge 1}\frac{\Tr(A^n)}{n}|z|^{mn}
  \]
  is \eqref{eq:class-logarithm-primitive-expansion} with
  \(\phi=\mathbf 1\). Tonelli's theorem then gives
  \[
    \sum_{m\ge 1}\frac{1}{m}
      \sum_{n\ge 1}\frac{\Tr(A^n)}{n}|z|^{mn}
    =
    \sum_{n\ge 1}\frac{\Tr(A^n)}{n}
      \bigl(-\log(1-|z|^n)\bigr).
  \]
  By \Cref{lem:pf-trace-bound},
  \(\Tr(A^n)\le C_A\lambda^n\) for some \(C_A>0\), and
  \[
    -\log(1-|z|^n)\le \frac{|z|^n}{1-|z|^n}\le \frac{|z|^n}{1-|z|}.
  \]
  Hence the majorant is bounded by
  \[
    \frac{C_A}{1-|z|}
    \sum_{n\ge 1}\frac{(\lambda|z|)^n}{n},
  \]
  which converges for \(|z|<\lambda^{-1}\). Therefore
  \[
    \sum_{m\ge 1}\frac{\mu(m)}{m}\mathcal L_{\psi^m\phi}(z^m)
    =
    \sum_{\gamma\in\cP}\sum_{m,r\ge 1}
      \frac{\mu(m)}{mr}\phi(\tau(\gamma)^{mr})z^{mr|\gamma|}.
  \]
  Now set \(s:=mr\). For fixed \(\gamma\) and \(s\), the coefficient of
  \(\phi(\tau(\gamma)^s)z^{s|\gamma|}\) is
  \[
    \sum_{m\mid s}\frac{\mu(m)}{m(s/m)}
    =
    \frac{1}{s}\sum_{m\mid s}\mu(m).
  \]
  Hence
  \[
    \sum_{m\ge 1}\frac{\mu(m)}{m}\mathcal L_{\psi^m\phi}(z^m)
    =
    \sum_{\gamma\in\cP}\sum_{s\ge 1}
      \frac{\phi(\tau(\gamma)^s)}{s}z^{s|\gamma|}
      \sum_{m\mid s}\mu(m).
  \]
  The divisor sum is \(1\) when \(s=1\) and \(0\) when \(s>1\), so only
  the primitive term remains:
  \[
    \sum_{m\ge 1}\frac{\mu(m)}{m}\mathcal L_{\psi^m\phi}(z^m)
    =
    \sum_{\gamma\in\cP}\phi(\tau(\gamma))z^{|\gamma|}
    =
    \Pi_\phi(z).
  \]
  This is \eqref{eq:class-adams-mobius-inversion}.
\end{proof}

\subsection{Fixed-label Euler transform versus periodic Artin logarithm}
\label{subsec:fixed-label-vs-artin-logarithm}

The next theorem is the notational hinge between Main Theorem~I and the
class-product constants.  The periodic determinant logarithm
\(\mathcal L_\phi\) sums repeated primitive labels as
\(\phi(\tau(\gamma)^r)\), while the fixed-label Euler transform
\(F_\phi\) keeps the primitive label fixed and then forms the Euler-product
tail \(\sum_{r\ge1}\Pi_\phi(z^r)/r\).  Thus \(\Pi_\phi\) is the primitive
Peter--Weyl channel, \(\mathcal L_\phi\) is the Artin/periodic logarithm,
and \(F_\phi\) is the quantity entering the Frobenius-class product
constant.  For non-trivial characters its Perron-point value is always read
with the Abel-path branch convention of the determinant formula below: the
scalar Perron determinant is deleted only in the singular triple
\((r,m,\pi)=(1,1,\mathbf1)\), whereas scalar terms with \(rm\ge2\) remain
regular and are retained.

\begin{theorem}[Fixed-label Euler transform]
  \label{thm:fixed-label-euler-transform}
  Let \(\phi:G\to\CC\) be a class function. For every
  \(|z|<\lambda^{-1}\), the fixed-label Euler transform satisfies
  \begin{equation}\label{eq:fixed-label-euler-transform-adams-mobius}
    F_\phi(z)
    =
    \sum_{r\ge1}\sum_{m\ge1}
      \frac{\mu(m)}{rm}\,
      \mathcal L_{\psi^m\phi}(z^{rm}).
  \end{equation}
  Equivalently,
  \begin{equation}\label{eq:fixed-label-euler-transform-vs-artin}
    \mathcal L_\phi(z)
    =
    \sum_{r\ge1}\frac{\Pi_{\psi^r\phi}(z^r)}{r},
    \qquad
    F_\phi(z)
    =
    \sum_{r\ge1}\frac{\Pi_\phi(z^r)}{r}.
  \end{equation}
  Thus the Artin logarithm \(\mathcal L_\phi\) and the primitive-label
  Euler transform \(F_\phi\) agree only in channels where the Adams powers
  do not change the relevant primitive labels. If \(\lambda>1\),
  \(\eta<\lambda\), and \(\phi\) is orthogonal to the trivial character,
  then \(F_\phi(\lambda^{-1})\) is well defined by absolute convergence.
  Its determinant continuation at the Perron point is the local
  determinant evaluation of \Cref{prop:fixed-label-determinant-perron}:
  the single singular scalar summand
  \((r,m,\pi)=(1,1,\mathbf1)\) is omitted before any logarithm is
  evaluated, while scalar components of \(\psi^m\phi\) with \(rm\ge2\)
  remain in the formula and are regular because
  \(\lambda^{-rm}A_{\mathbf1}\) lies strictly inside the Perron disc.
  Thus non-trivial characters may acquire trivial components under Adams
  powers, but only away from the scalar Perron boundary.  The later
  product-constant role of this fixed-label value, and the obstruction to
  substituting \(\mathcal L_\phi\), are isolated in
  \Cref{prop:three-logarithm-product-interface}.
\end{theorem}

\begin{proof}
  For \(|z|<\lambda^{-1}\), absolute convergence of \(F_\phi\) follows
  from
  \[
    \sum_{r\ge1}\frac{1}{r}
      \sum_{\gamma\in\cP}\abs{\phi(\tau(\gamma))}|z|^{r|\gamma|}
    \le
    \norm{\phi}_\infty
    \sum_{n\ge1}p_n(A)\sum_{r\ge1}\frac{|z|^{rn}}{r}
    <\infty,
  \]
  since \(p_n(A)\ll\lambda^n/n\) and \(\lambda|z|<1\). Substituting
  \Cref{cor:class-function-adams-mobius} into
  \eqref{eq:fixed-label-euler-transform-def} gives
  \[
    F_\phi(z)
    =
    \sum_{r\ge1}\frac{1}{r}
      \sum_{m\ge1}\frac{\mu(m)}{m}
        \mathcal L_{\psi^m\phi}(z^{rm}).
  \]
  The rearrangement is justified by the same majorant used in the proof
  of \Cref{cor:class-function-adams-mobius}, with \(|z|\) replaced by
  \(|z|^r\), giving \eqref{eq:fixed-label-euler-transform-adams-mobius}.
  The first identity in
  \eqref{eq:fixed-label-euler-transform-vs-artin} is exactly
  \Cref{lem:class-logarithm-primitive-expansion}, because
  \((\psi^r\phi)(\tau(\gamma))=\phi(\tau(\gamma)^r)\); the second is the
  definition of \(F_\phi\).

  Now assume \(\lambda>1\), \(\eta<\lambda\), and
  \(\langle\phi,\mathbf1\rangle=0\). Expand
  \(\phi=\sum_{\varrho\neq\mathbf1}\widehat\phi(\varrho)\chi_\varrho\).
  By \Cref{lem:nontrivial-perron-point-convergence}, the term
  \(\Pi_\varrho(\lambda^{-1})\) is absolutely convergent for every
  non-trivial \(\varrho\).  The remaining fixed-label tail is controlled
  directly from primitive-orbit counting: for any bounded class function
  \(\psi\) and every \(r\ge2\),
  \[
    \abs{\Pi_\psi(\lambda^{-r})}
    \le
    \norm{\psi}_\infty\sum_{n\ge1}p_n(A)\lambda^{-rn}
    \ll_\psi
    \sum_{n\ge1}\frac{\lambda^{(1-r)n}}{n},
  \]
  and the last quantity is summable after division by \(r\) over
  \(r\ge2\). Hence \(\sum_{r\ge1}\Pi_\phi(\lambda^{-r})/r\) converges
  absolutely. For the determinant expansion of the same boundary value, apply
  \Cref{lem:fixed-label-perron-double-majorant}. It supplies an explicit
  summable majorant for the whole \((r,m)\)-determinant tail and removes
  the only scalar Perron logarithm by character orthogonality: the term
  \((r,m,\pi)=(1,1,\mathbf1)\) is excluded before the limit is taken, and
  its coefficient is the trivial Fourier coefficient of \(\phi\), hence is
  zero under \(\langle\phi,\mathbf1\rangle=0\). Dominated
  convergence from \(z=t\lambda^{-1}\), \(0<t<1\), therefore gives the
  stated Perron-point continuation and branch convention.
\end{proof}

\begin{proposition}[Periodic, primitive, and fixed-label logarithms]
  \label{prop:logarithmic-channel-separation}
  Let \(\phi:G\to\CC\) be a class function. On \(|z|<\lambda^{-1}\) the
  following three logarithmic channels are related, and only related, by
  the displayed Adams operations:
  \begin{align}
    \mathcal L_\phi(z)
    &=\sum_{r\ge1}\frac{\Pi_{\psi^r\phi}(z^r)}{r},
      \label{eq:channel-separation-artin}\\
    \Pi_\phi(z)
    &=\sum_{m\ge1}\frac{\mu(m)}{m}
      \mathcal L_{\psi^m\phi}(z^m),
      \label{eq:channel-separation-primitive}\\
    F_\phi(z)
    &=\sum_{r,m\ge1}\frac{\mu(m)}{rm}
      \mathcal L_{\psi^m\phi}(z^{rm}).
      \label{eq:channel-separation-fixed}
  \end{align}
  Consequently the Artin logarithm \(\mathcal L_\phi\) and the fixed-label
  Euler transform \(F_\phi\) have the same coefficients for the given
  finite mixing SFT and finite-group cocycle if and only if
  \begin{equation}\label{eq:channel-separation-obstruction}
    \sum_{r\mid n}\frac{1}{r}
      P_{n/r}(\psi^r\phi-\phi)=0
    \qquad(n\ge1),
  \end{equation}
  where
  \(P_k(\psi):=\sum_{|\gamma|=k}\psi(\tau(\gamma))\) is the primitive
  length-\(k\) \(\psi\)-weighted label sum. In particular, no theorem in
  this paper may replace \(F_\phi(\lambda^{-1})\) by
  \(\mathcal L_\phi(\lambda^{-1})\) unless the primitive
  Adams-difference coefficients in \eqref{eq:channel-separation-obstruction}
  vanish in the channel under discussion.
\end{proposition}

\begin{proof}
  The first identity is \Cref{lem:class-logarithm-primitive-expansion}; the
  second is \Cref{cor:class-function-adams-mobius}; the third is
  \Cref{thm:fixed-label-euler-transform}. Absolute convergence on
  \(|z|<\lambda^{-1}\) is proved in those statements, so coefficient
  comparison is legitimate.

  Write
  \[
    \Pi_\psi(z)=\sum_{k\ge1}P_k(\psi)z^k.
  \]
  From \eqref{eq:channel-separation-artin} and the definition of
  \(F_\phi\), the coefficient of \(z^n\) in
  \(\mathcal L_\phi(z)-F_\phi(z)\) is
  \[
    \sum_{r\mid n}\frac{1}{r}
      \bigl(P_{n/r}(\psi^r\phi)-P_{n/r}(\phi)\bigr)
    =
    \sum_{r\mid n}\frac{1}{r}P_{n/r}(\psi^r\phi-\phi).
  \]
  Hence the two analytic germs agree exactly when
  \eqref{eq:channel-separation-obstruction} holds for every \(n\). Since
  the Perron-point constants of \Cref{thm:class-euler-product-mertens} are
  obtained by radial continuation of the fixed-label germ under the strict
  gap, the same obstruction is the only valid passage from the fixed-label
  constant to the periodic Artin logarithm. This proves the proposition.
\end{proof}

\begin{lemma}[Trace-power majorant for Perron-boundary determinant logarithms]
  \label{lem:perron-boundary-trace-power-majorant}
  Assume \(\lambda>1\) and \(\eta<\lambda\). Let
  \(\mathscr A:=\{A_\pi:\\pi\in\Irr(G)\}\). There is a constant
  \(C_{\mathscr A}\) with the following properties.
  \begin{enumerate}
  \item For every \(B\in\mathscr A\), every \(\ell\ge1\), and every
    \(k\ge2\),
    \begin{equation}\label{eq:trace-power-perron-bound}
      \abs{\Tr(B^\ell)}\le C_{\mathscr A}\lambda^\ell,
      \qquad
      \abs{\log\det(I-t^k\lambda^{-k}B)^{-1}}
      \le C_{\mathscr A}\lambda^{1-k}
    \end{equation}
    uniformly for \(0\le t\le1\), where the logarithm is the power-series
    branch normalised to vanish at \(t=0\).
  \item For every non-trivial \(\pi\in\Irr(G)\), the path
    \(0\le t\le1\mapsto I-t\lambda^{-1}A_\pi\) is invertible. Hence the
    normalized determinant logarithm
    \(\log\det(I-t\lambda^{-1}A_\pi)^{-1}\) continues uniquely and remains
    bounded on this Abel path.
  \item If \(c_{m,\pi}\) is any coefficient family with
    \(\sup_{m,\pi}\abs{c_{m,
    \pi}}<\infty\), then every determinant-log series of the form
    \[
      \sum_{m\ge1}\frac{\mu(m)}{m}
      \sum_{\substack{\pi\in\Irr(G)\\(m,\pi)\neq(1,\mathbf1)}}
        c_{m,\pi}\log\det(I-t^m\lambda^{-m}A_\pi)^{-1}
    \]
    has an Abel majorant independent of \(t\in[0,1]\): the finitely many
    \(m=1\), \(\pi\neq\mathbf1\) terms are bounded by item~2, and the
    \(m\ge2\) tail is dominated by a constant multiple of
    \(\sum_{m\ge2}\lambda^{1-m}/m\).
  \item If \(c_{r,m,
    \pi}\) is any coefficient family with
    \(\sup_{r,m,
    \pi}\abs{c_{r,m,
    \pi}}<\infty\), then every determinant-log series of the form
    \[
      \sum_{r,m\ge1}\frac{\mu(m)}{rm}
      \sum_{\substack{\pi\in\Irr(G)\\(r,m,\pi)\neq(1,1,\mathbf1)}}
        c_{r,m,
        \pi}
        \log\det(I-t^{rm}\lambda^{-rm}A_\pi)^{-1}
    \]
    has an Abel majorant independent of \(t\in[0,1]\): the finitely many
    \(rm=1\), \(\pi\neq\mathbf1\) terms are bounded by item~2, and the
    \(rm\ge2\) tail is dominated by a constant multiple of
    \begin{equation}\label{eq:double-trace-power-majorant}
      \sum_{\substack{r,m\ge1\\rm\ge2}}\frac{\lambda^{1-rm}}{rm}<\infty .
    \end{equation}
  \end{enumerate}
  Consequently the only possible scalar Perron singularity in such
  normalized determinant expansions is the single term with
  \((m,\pi)=(1,\mathbf1)\), respectively
  \((r,m,\pi)=(1,1,\mathbf1)\). If its coefficient is zero, dominated
  convergence along any of the displayed Abel paths is legitimate at
  \(t=1\).
\end{lemma}

\begin{proof}
  The set \(\mathscr A\) is finite. If the eigenvalues of
  \(B\in\mathscr A\) are \(\alpha_1,\ldots,\alpha_d\), counted with
  algebraic multiplicity, then
  \(\rho(B)\le\lambda\) because the skew-product adjacency matrix is
  non-negative with Perron root \(\lambda\), and
  \Cref{prop:kernel-peter-weyl-block-diagonalisation} places each
  \(A_\pi\) as a block of it. Thus
  \[
    \abs{\Tr(B^\ell)}
    =\abs{\sum_{j=1}^d\alpha_j^\ell}
    \le d\lambda^\ell
    \le C_0\lambda^\ell
  \]
  for a uniform constant \(C_0\). This is the trace-power bound; no
  diagonalizability is being assumed, since traces of powers are sums of
  eigenvalue powers with algebraic multiplicity, equivalently after Schur
  triangularisation.

  For \(k\ge2\),
  \(\rho(t^k\lambda^{-k}B)\le\lambda^{1-k}<1\). Hence
  \[
    \log\det(I-t^k\lambda^{-k}B)^{-1}
    =\sum_{\ell\ge1}
      \frac{t^{k\ell}\lambda^{-k\ell}}{\ell}\Tr(B^\ell)
  \]
  in the normalized power-series branch. The preceding trace bound gives
  \[
    \abs{\log\det(I-t^k\lambda^{-k}B)^{-1}}
    \le C_0\sum_{\ell\ge1}\frac{\lambda^{(1-k)\ell}}{\ell}.
  \]
  Since \(0<\lambda^{1-k}\le\lambda^{-1}<1\) for all \(k\ge2\), the last
  sum is bounded by a uniform constant times \(\lambda^{1-k}\), after
  enlarging the constant. This proves item~1.

  If \(\pi\neq\mathbf1\), the twisted gap gives
  \(\rho(A_\pi)<\lambda\). Therefore
  \(1/t\lambda^{-1}\) is not an eigenvalue of \(A_\pi\) for
  \(0\le t\le1\), or equivalently \(I-t\lambda^{-1}A_\pi\) is invertible
  on the entire compact Abel path. The determinant has no zero there, so
  the logarithm normalized at \(t=0\) continues uniquely and is bounded.
  This proves item~2.

  For item~3, the \(m=1\), \(\pi\neq\mathbf1\) terms form a finite bounded
  family by item~2. The tail with \(m\ge2\) is bounded by item~1 and the
  assumed coefficient bound by a constant multiple of
  \(\sum_{m\ge2}\lambda^{1-m}/m\), which converges. Item~4 is identical
  with \(k=rm\): the tail is bounded by
  \(\sum_{rm\ge2}\lambda^{1-rm}/(rm)\), and
  \[
    \sum_{\substack{r,m\ge1\\rm\ge2}}\frac{\lambda^{1-rm}}{rm}
    =\sum_{k\ge2}\frac{d(k)}{k}\lambda^{1-k}<\infty
  \]
  because the divisor function has subexponential growth while
  \(\lambda^{1-k}\) decays exponentially. The final statement is then
  exactly the dominated-convergence consequence of the displayed
  majorants, after the scalar Perron term with coefficient zero has been
  omitted before evaluating any logarithm.
\end{proof}

\begin{lemma}[Perron-boundary majorant for the fixed-label double series]
  \label{lem:fixed-label-perron-double-majorant}
  Assume \(\lambda>1\) and \(\eta<\lambda\). Let
  \(\phi:G\to\CC\) be a class function with
  \(\langle\phi,\mathbf1\rangle=0\), and write
  \(\psi^m\phi=\sum_{\pi\in\Irr(G)}b_{\phi,\pi}^{(m)}\chi_\pi\).
  Then the Abel-boundary determinant series
  \begin{equation}\label{eq:fixed-label-perron-double-majorant-series}
    \sum_{q,m\ge1}\frac{\mu(m)}{qm}
      \sum_{\substack{\pi\in\Irr(G)\\(q,m,\pi)\neq(1,1,\mathbf1)}}
        b_{\phi,\pi}^{(m)}
        \log\det(I-t^{qm}\lambda^{-qm}A_\pi)^{-1}
  \end{equation}
  has, for \(0\le t\le1\), a summable majorant independent of \(t\).
  Here \(t\) is the radial Abel parameter and \(q\) is the summation index
  of the fixed-label Euler transform. More precisely, the finitely many
  \(qm=1\), \(\pi\neq\mathbf1\) terms are bounded on the Abel path, while
  the \(qm\ge2\) terms are dominated by
  \begin{equation}\label{eq:fixed-label-perron-double-majorant-bound}
    C_\phi
    \sum_{\substack{q,m\ge1\\qm\ge2}}\frac{\lambda^{1-qm}}{qm}<\infty .
  \end{equation}
  The omitted term \((q,m,\pi)=(1,1,\mathbf1)\) has coefficient
  \(b_{\phi,\mathbf1}^{(1)}=\langle\phi,\mathbf1\rangle=0\); hence the
  singular scalar logarithm
  \(\log\det(I-t\lambda^{-1}A_{\mathbf1})^{-1}\) never enters the
  non-trivial fixed-label boundary value.
\end{lemma}

\begin{proof}
  The Adams coefficients are uniformly bounded in \(m\). Indeed,
  \[
    \abs{b_{\phi,\pi}^{(m)}}
    =\abs{\langle\psi^m\phi,\chi_\pi\rangle}
    \le \norm{\phi}_\infty\chi_\pi(1)
    \le B_\phi,
  \]
  because \(G\) has only finitely many irreducible representations. Since
  \(\psi^1\phi=\phi\), orthogonality gives
  \(b_{\phi,\mathbf1}^{(1)}=0\). Thus the scalar Perron determinant with
  \(qm=1\) is excluded before any logarithm is evaluated. Applying
  \Cref{lem:perron-boundary-trace-power-majorant} to the bounded
  coefficient family \(c_{q,m,
  \pi}=b_{\phi,\pi}^{(m)}\) gives the bounded regular \(qm=1\) terms, the
  summable \(qm\ge2\) majorant in
  \eqref{eq:fixed-label-perron-double-majorant-bound}, and the stated
  dominated-convergence branch exclusion.
\end{proof}

\begin{remark}[Relation to necklace/Witt inversion and twisted-zeta
  precedents]
  \label{rem:adams-mobius-literature-position}
  The cancellation in
  \Cref{cor:class-function-adams-mobius} is the finite-group analogue
  of the classical ``primitive from repetitions'' mechanism in the
  necklace and Witt-vector literature
  \cite{MetropolisRota1983Necklaces,Wehrhahn1990AperiodicRings}. What
  is specific here is the finite-state cocycle input:
  repetitions replace a primitive label \(g\) by \(g^m\), so the
  relevant correction is not scalar M\"obius inversion alone but the
  Adams operation \(\psi^m\phi(g)=\phi(g^m)\) inside the character ring
  of one fixed finite-group extension. Twisted periodic-point formulas
  and determinant identities also appear in Fried's work
  \cite{Fried1983PeriodicPointsTwistedCoefficients,Fried1983HomologicalIdentitiesClosedOrbits} and in the dynamical-zeta
  treatment of Parry--Pollicott \cite{ParryPollicott1990Zeta}. Those
  works provide determinant and orbit-counting frameworks, but they do
  not carry out the finite-state Adams--M\"obius primitive inversion and
  classwise Perron-point constant extraction developed here.
\end{remark}
% END INPUT sec_chebotarev_fixed_label_and_boundary.tex
% END INPUT sec_chebotarev.tex
% BEGIN INPUT sec_chebotarev_peter_weyl_bridge.tex
\subsection{Primitive Peter--Weyl series}

Applying \Cref{cor:class-function-adams-mobius} to irreducible
characters produces the primitive Peter--Weyl coefficients directly
from the twisted determinants.

For \(\varrho\in\Irr(G)\), define
\begin{equation}\label{eq:primitive-peter-weyl-series}
  \Pi_\varrho(z):=\Pi_{\chi_\varrho}(z)
  =\sum_{n\ge 1}\pi_{n,\varrho}(A,\tau)z^n,
  \qquad
  \pi_{n,\varrho}(A,\tau)
  :=\sum_{\substack{\gamma\in\cP\\ |\gamma|=n}}
    \chi_\varrho(\tau(\gamma)).
\end{equation}
For \(m\ge 1\), write
\begin{equation}\label{eq:adams-character-decomposition}
  \psi^m\chi_\varrho
  =\sum_{\pi\in\Irr(G)}a_{\varrho,\pi}^{(m)}\chi_\pi
\end{equation}
for the Adams decomposition in the character ring. Equivalently, with
\(p_m(g)=g^m\),
\begin{equation}\label{eq:adams-coefficient-inner-product}
  a_{\varrho,\pi}^{(m)}
  =
  \left\langle \chi_\varrho\circ p_m,\chi_\pi\right\rangle
  =
  \frac{1}{|G|}\sum_{g\in G}
    \chi_\varrho(g^m)\overline{\chi_\pi(g)}.
\end{equation}
Thus the finite character-ring input is not merely the ordinary
character table: it is the character table together with the power maps
on conjugacy classes, or equivalently the Adams operations on \(R(G)\).

\begin{lemma}[Uniform bound for Adams decomposition coefficients]
  \label{lem:adams-coefficients-bounded}
  For every \(\varrho,\pi\in\Irr(G)\) and every \(m\ge 1\),
  \[
    \abs{a_{\varrho,\pi}^{(m)}}\le \dim\varrho.
  \]
\end{lemma}

\begin{proof}
  Since the irreducible characters are orthonormal,
  \[
    a_{\varrho,\pi}^{(m)}
    =\langle \psi^m\chi_\varrho,\chi_\pi\rangle.
  \]
  Hence
  \[
    \abs{a_{\varrho,\pi}^{(m)}}
    \le \norm{\psi^m\chi_\varrho}_2\norm{\chi_\pi}_2
    =\norm{\psi^m\chi_\varrho}_2.
  \]
  Moreover,
  \[
    \norm{\psi^m\chi_\varrho}_2^2
    =\frac{1}{|G|}\sum_{g\in G}\abs{\chi_\varrho(g^m)}^2
    \le (\dim\varrho)^2,
  \]
  because \(\abs{\chi_\varrho(h)}\le \dim\varrho\) for all \(h\in G\).
  Taking the square root gives the claim.
\end{proof}

\begin{remark}[Periodicity of Adams coefficients]
  \label{rem:adams-coefficients-periodic}
  Let \(\exp(G)\) denote the exponent of \(G\). Since
  \(g^{m+\exp(G)}=g^m\) for every \(g\in G\), one has
  \[
    \psi^{m+\exp(G)}\chi_\varrho=\psi^m\chi_\varrho.
  \]
  Consequently
  \(a_{\varrho,\pi}^{(m+\exp(G))}=a_{\varrho,\pi}^{(m)}\) for all
  \(m\), so only finitely many Adams coefficient vectors occur. In the
  estimates below we use this finiteness only through the uniform bound
  of \Cref{lem:adams-coefficients-bounded}.
\end{remark}

\begin{corollary}[Main Theorem I(c): irreducible-channel determinant
  formula]
  \label{cor:primitive-peter-weyl-determinant}
  For each \(\pi\in\Irr(G)\), write
  \[
    P_\pi(z):=\det(I-zA_\pi),\qquad P_\pi(0)=1,
  \]
  and let \(D_\pi\) be the unique analytic germ on
  \(|z|<\lambda^{-1}\) with \(D_\pi(0)=0\) and
  \(\exp D_\pi=P_\pi^{-1}\). The finite determinant data are the
  polynomials \(P_\pi\), or equivalently the matrices \(A_\pi\); the
  logarithmic branches are these normalized germs derived from them.
  For every \(\varrho\in\Irr(G)\) and every \(|z|<\lambda^{-1}\),
  \begin{equation}\label{eq:primitive-peter-weyl-determinant}
    \Pi_\varrho(z)
    =
    \sum_{m\ge 1}\frac{\mu(m)}{m}
      \sum_{\pi\in\Irr(G)}
        a_{\varrho,\pi}^{(m)}
        D_\pi(z^m).
  \end{equation}
  Equivalently \(D_\pi(z^m)=\log\det(I-z^mA_\pi)^{-1}\) in the branch
  fixed by the logarithm convention for this section.
\end{corollary}

\begin{proof}
  Apply \Cref{cor:class-function-adams-mobius} with
  \(\phi=\chi_\varrho\). Then
  \[
    \Pi_\varrho(z)
    =
    \sum_{m\ge 1}\frac{\mu(m)}{m}
      \mathcal L_{\psi^m\chi_\varrho}(z^m).
  \]
  By \eqref{eq:adams-character-decomposition} and the determinant
  linearisation \eqref{eq:class-logarithm-determinant-linearisation},
  \[
    \mathcal L_{\psi^m\chi_\varrho}(z^m)
    =
    \sum_{\pi\in\Irr(G)}
      a_{\varrho,\pi}^{(m)}
      D_\pi(z^m),
  \]
  because \(D_\pi\) is the normalized determinant logarithm on the
  Perron disc.
  Substituting this identity proves
  \eqref{eq:primitive-peter-weyl-determinant}.
\end{proof}

\noindent
The determinant identity above is the generating-function part of Main
Theorem~I.  The next theorem is the coefficient recovery from the same
identity, while the following corollary is only the data-interface version:
the theorem gives the formula, and the corollary records that the finite
twisted trace, or determinant, data is exactly the input needed to use it.
Neither the twisted spectral gap nor cyclic root-of-unity reconstruction
enters this first main theorem.

\begin{remark}[Abelian toy model: the cyclic lift]
  \label{rem:abelian-toy-model-cyclic}
  Let \(G=\ZZ/q\ZZ\), and let \(\chi_j(h)=\omega_q^{jh}\) be its
  irreducible characters. The canonical cyclic-lift cocycle is the
  special cocycle \(\tau_q:E_A\to\ZZ/q\ZZ\) with
  \(\tau_q(e)=1\) on every edge; its skew-product adjacency matrix is
  \(A\otimes P_q\). For this cocycle, and not for a general
  \(\ZZ/q\ZZ\)-valued edge labelling, one has
  \(A_{\chi_j}=\omega_q^jA\), and therefore
  \[
    \mathcal L_{\chi_j}(z)=\log\det(I-\omega_q^j zA)^{-1}.
  \]
  The Adams--M\"obius inversion
  \eqref{eq:class-adams-mobius-inversion} then reduces to the ordinary
  Fourier decomposition on the \(q\)-th roots of unity. This is the
  concrete bridge between the abelian formulas of
  \Cref{sec:finite-part} and the present finite-group formalism, but
  only after the original cocycle has been replaced by this canonical
  cyclic-lift cocycle.
  This model is also not an example of the twisted spectral gap used
  later when \(q>1\): for every non-trivial character \(\chi_j\),
  \(\rho(A_{\chi_j})=\rho(A)=\lambda\), so the non-trivial twisted
  radius is \(\eta=\lambda\), not \(<\lambda\). The cyclic-lift layer is
  therefore a finite-matrix comparison case, not a source of
  Perron-point non-trivial-channel constants.
\end{remark}

\begin{theorem}[Main Theorem I(d): primitive Peter--Weyl trace formula]
  \label{thm:primitive-peter-weyl-trace}
  For every \(\varrho\in\Irr(G)\) and every \(n\ge 1\),
  \begin{equation}\label{eq:primitive-peter-weyl-trace}
    \pi_{n,\varrho}(A,\tau)
    =
    \frac{1}{n}\sum_{m\mid n}\mu(m)
      \sum_{\pi\in\Irr(G)}
        a_{\varrho,\pi}^{(m)}
        \Tr(A_\pi^{n/m}).
  \end{equation}
\end{theorem}

\begin{proof}
  Extract the coefficient of \(z^n\) in
  \eqref{eq:primitive-peter-weyl-determinant}. Since
  \[
    [z^n]\log\det(I-z^mA_\pi)^{-1}
    =
    \begin{cases}
      \Tr(A_\pi^{n/m})/(n/m), & m\mid n, \\
      0, & m\nmid n,
    \end{cases}
  \]
  the stated formula follows immediately.
\end{proof}

\begin{corollary}[Finite trace formula for primitive coefficients]
  \label{cor:primitive-peter-weyl-trace}
  The coefficient-level primitive data are finite-state data.  More
  precisely, the finite character table of \(G\), the Adams coefficients
  \(a_{\varrho,\pi}^{(m)}\), and the twisted trace sequences
  \(\{\Tr(A_\pi^k):k\ge1,\ \pi\in\Irr(G)\}\) determine every primitive
  Peter--Weyl coefficient \(\pi_{n,\varrho}(A,\tau)\). Equivalently, for
  every \(\varrho\in\Irr(G)\), the series \(\Pi_\varrho(z)
  =\sum_{n\ge1}\pi_{n,\varrho}(A,\tau)z^n\) is determined by the finite
  twisted determinant family \(\{\det(I-zA_\pi):\pi\in\Irr(G)\}\) and the
  power-map action on conjugacy classes.
\end{corollary}

\begin{proof}
  The explicit coefficient formula is exactly
  \Cref{thm:primitive-peter-weyl-trace}.  The coefficients
  \(a_{\varrho,\pi}^{(m)}\) are finite character-table inner products
  defined in \eqref{eq:adams-coefficient-inner-product}, hence are known
  from the Adams power maps on conjugacy classes.  The traces
  \(\Tr(A_\pi^k)\) are determined by the logarithmic derivative, or
  equivalently by the coefficients of the finite polynomial
  \(\det(I-zA_\pi)\). Substitution in
  \eqref{eq:primitive-peter-weyl-trace} determines each coefficient
  \(\pi_{n,\varrho}(A,\tau)\), and therefore determines the whole series
  \(\Pi_\varrho(z)\). The determinant-series formulation is the same
  statement before coefficient extraction, namely
  \Cref{cor:primitive-peter-weyl-determinant}.
\end{proof}

\noindent
This is the coefficient-level form of Main Theorem~I, not a preliminary
estimate: it is the point at which the finite twisted determinant family
becomes primitive orbit data. The trace lemma supplies the periodic
input, the linearisation theorem rewrites it as twisted determinant data,
the Adams--M\"obius corollary removes repetitions, and the determinant
corollary and trace theorem give the primitive Peter--Weyl series and
coefficients.

\subsection{Bridge to classes and Main Theorem II}

The primitive Peter--Weyl coefficients determine the conjugacy-class
counts by Fourier inversion on \(G\). The next proposition is a bridge: it
translates the primitive channels of Main Theorem~I into class counts,
but it does not add a new independent main theorem.

\begin{proposition}[Conjugacy-class primitive decomposition]
  \label{prop:class-primitive-decomposition}
  Let \(C\subset G\) be a conjugacy class. If \(p_{n,C}\) denotes the
  number of primitive orbits of length \(n\) whose \(G\)-label lies in
  \(C\), then for \(|z|<\lambda^{-1}\),
  \begin{equation}\label{eq:class-primitive-generating-function}
    \sum_{n\ge 1}p_{n,C}z^n
    =
    \frac{|C|}{|G|}
      \sum_{\varrho\in\Irr(G)}
        \chi_\varrho^\ast(C)\,\Pi_\varrho(z),
  \end{equation}
  and therefore
  \begin{equation}\label{eq:class-primitive-coefficients}
    p_{n,C}
    =
    \frac{|C|}{|G|}
      \sum_{\varrho\in\Irr(G)}
        \chi_\varrho^\ast(C)\,\pi_{n,\varrho}(A,\tau).
  \end{equation}
\end{proposition}

\begin{proof}
  By definition,
  \[
    \sum_{n\ge 1}p_{n,C}z^n
    =
    \sum_{\gamma\in\cP}\mathbf 1_C(\tau(\gamma))z^{|\gamma|}.
  \]
  Applying \Cref{prop:class-indicator-expansion} to the class function
  \(\mathbf 1_C\) and using the linearity of \(\Pi_\phi\) yields
  \eqref{eq:class-primitive-generating-function}. Comparing coefficients
  gives \eqref{eq:class-primitive-coefficients}.
\end{proof}

\begin{theorem}[Main Theorem I(e): class-count trace recovery]
  \label{thm:class-count-trace-recovery}
  \label{thm:class-count-trace-readout}
  \label{cor:class-artin-mobius-trace}
  For every conjugacy class \(C\subset G\) and every \(n\ge 1\),
  \begin{equation}\label{eq:class-artin-mobius-trace}
    p_{n,C}
    =
    \frac{|C|}{|G|\,n}
      \sum_{\varrho\in\Irr(G)}
        \chi_\varrho^\ast(C)
        \sum_{m\mid n}\mu(m)
          \sum_{\pi\in\Irr(G)}
            a_{\varrho,\pi}^{(m)}
            \Tr(A_\pi^{n/m}).
  \end{equation}
  Equivalently, after the finite character table, the conjugacy-class
  power maps, and the determinant polynomials \(\det(I-zA_\pi)\) are
  known, no additional orbit-level information is needed to determine
  any primitive class count \(p_{n,C}\).
\end{theorem}

\begin{proof}
  By \Cref{prop:class-primitive-decomposition},
  \[
    p_{n,C}
    =
    \frac{|C|}{|G|}
      \sum_{\varrho\in\Irr(G)}
        \chi_\varrho^\ast(C)\,\pi_{n,\varrho}(A,\tau).
  \]
  Substituting the primitive Peter--Weyl trace formula
  \eqref{eq:primitive-peter-weyl-trace} for each channel
  \(\pi_{n,\varrho}(A,\tau)\) gives
  \eqref{eq:class-artin-mobius-trace}.

  The Adams coefficients are determined by the finite character table and
  the conjugacy-class power maps through
  \eqref{eq:adams-coefficient-inner-product}. The trace sequence
  \(\Tr(A_\pi^k)\) is determined by the finite matrix \(A_\pi\). If one
  records only the determinant polynomial \(\det(I-zA_\pi)\), the same
  traces are still determined by the identity
  \[
    \log\det(I-zA_\pi)^{-1}
    =
    \sum_{k\ge1}\frac{\Tr(A_\pi^k)}{k}z^k
  \]
  near \(z=0\), equivalently by Newton identities. Hence the displayed
  formula uses exactly the finite determinant and finite character-ring
  data stated in the theorem.
\end{proof}

\noindent
This completes Main Theorem~I in the sense used in the introduction:
the primitive Peter--Weyl trace theorem gives the irreducible channels,
and \Cref{thm:class-count-trace-recovery} is the direct class-count
recovery that feeds the class-Mertens constants of Main Theorem~II.

From this point onward in the Perron-point part of Main Theorem~II we
assume \(\lambda>1\) and impose the gap hypothesis
\[
  \eta:=\max_{\varrho\neq\mathbf 1}\rho(A_\varrho)<\lambda.
\]
If \(G\) is trivial, this maximum is empty and \(\eta=0\).
Thus the formal determinant identities
\Cref{thm:class-function-linearisation,cor:class-function-adams-mobius,cor:primitive-peter-weyl-determinant,thm:primitive-peter-weyl-trace,thm:class-count-trace-recovery}
are unconditional in the open Perron disc, while the estimates and
constants from \Cref{lem:nontrivial-peter-weyl-channel-bound} through
\Cref{thm:class-mertens-fourier} are conditional on \(\lambda>1\) and
the displayed gap.
The only scalar input used below is the trivial-channel
orbit-Mertens constant \(\mathcal M(A)\); its finite-matrix formula is
proved independently in \Cref{sec:finite-part}.

The class Mertens theorem below is the analytic part of Main
Theorem~II from the introduction. The leading density term is included
to identify the constant term and normalisation; the new ingredient in
this section is the passage from the twisted determinant channels to the
primitive Perron-point class profile. The later Fourier formula is a
recovery of the constants by character orthogonality. The proof is
assembled from the next five lemmas:
class decomposition, the non-trivial channel estimate, absolute
convergence at the Perron point, a geometric-harmonic tail bound, and
the trivial Perron-channel contribution.

\begin{lemma}[Trivial and non-trivial class decomposition]
  \label{lem:class-primitive-trivial-split}
  For every conjugacy class \(C\subset G\) and every \(n\ge 1\),
  \[
    p_{n,C}
    =
    \frac{|C|}{|G|}p_n(A)
    +\frac{|C|}{|G|}
      \sum_{\varrho\neq\mathbf 1}
        \chi_\varrho^\ast(C)\,\pi_{n,\varrho}(A,\tau).
  \]
\end{lemma}

\begin{proof}
  Separate the trivial representation from
  \eqref{eq:class-primitive-coefficients}. Since
  \(\chi_{\mathbf 1}^\ast(C)=1\) and
  \(\pi_{n,\mathbf 1}(A,\tau)=p_n(A)\), the displayed identity follows.
\end{proof}

\begin{lemma}[Non-trivial Peter--Weyl channel bound]
  \label{lem:nontrivial-peter-weyl-channel-bound}
  For every non-trivial \(\varrho\in\Irr(G)\) and every \(n\ge 1\),
  \[
    \abs{\pi_{n,\varrho}(A,\tau)}
    \le
    \frac{|V|\,\dim\varrho\,\eta^n}{n}
    +\frac{|V|\,\dim\varrho}{n}
      \left(\sum_{\kappa\in\Irr(G)}\dim\kappa\right)
      \sum_{\substack{m\mid n\\ m>1}}\lambda^{n/m},
  \]
  and therefore
  \[
    \abs{\pi_{n,\varrho}(A,\tau)}
    \le
    \frac{|V|\,\dim\varrho\,\eta^n}{n}
    +\frac{|V|\,\dim\varrho}{n}
      \left(\sum_{\kappa\in\Irr(G)}\dim\kappa\right)
      \frac{\lambda^{n/2+1}}{\lambda-1}.
  \]
  In particular,
  \[
    \pi_{n,\varrho}(A,\tau)
    =O\!\left(\frac{\max\{\eta^n,\lambda^{n/2}\}}{n}\right).
  \]
\end{lemma}

\begin{proof}
  Fix \(\varrho\neq\mathbf 1\). For each
  \(\kappa\in\Irr(G)\), let
  \(\mu_{\kappa,1},\dots,\mu_{\kappa,|V|\dim\kappa}\) be the eigenvalues
  of \(A_\kappa\), repeated with algebraic multiplicity. By
  \Cref{prop:twisted-spectral-radius-bound},
  \[
    \abs{\mu_{\kappa,j}}\le \rho(A_\kappa)\le \lambda
    \qquad
    (\kappa\in\Irr(G)),
  \]
  and for the present non-trivial \(\varrho\) the twisted gap gives
  \[
    \abs{\mu_{\varrho,j}}\le \rho(A_\varrho)\le \eta.
  \]
  Therefore
  \[
    \abs{\Tr(A_\varrho^n)}
    =
    \left|\sum_{j=1}^{|V|\dim\varrho}\mu_{\varrho,j}^n\right|
    \le |V|\,\dim\varrho\,\eta^n,
  \]
  and, for every \(\kappa\in\Irr(G)\) and every \(\ell\ge 1\),
  \[
    \abs{\Tr(A_\kappa^\ell)}
    =
    \left|\sum_{j=1}^{|V|\dim\kappa}\mu_{\kappa,j}^\ell\right|
    \le |V|\,\dim\kappa\,\lambda^\ell.
  \]

  By \Cref{thm:primitive-peter-weyl-trace},
  \[
    \pi_{n,\varrho}(A,\tau)
    =
    \frac{1}{n}\Tr(A_\varrho^n)
    +\frac{1}{n}
      \sum_{\substack{m\mid n\\ m>1}}\mu(m)
        \sum_{\kappa\in\Irr(G)}
          a_{\varrho,\kappa}^{(m)}\Tr(A_\kappa^{n/m}).
  \]
  Taking absolute values, the first term is bounded by
  \(|V|\,\dim\varrho\,\eta^n/n\). For the divisor sum, use
  \Cref{lem:adams-coefficients-bounded} together with the trace bound
  above:
  \begin{align*}
    \abs{
      \sum_{\kappa\in\Irr(G)}
        a_{\varrho,\kappa}^{(m)}\Tr(A_\kappa^{n/m})
    }
    &\le
    \sum_{\kappa\in\Irr(G)}
      \dim\varrho\,|V|\,\dim\kappa\,\lambda^{n/m} \\
    &=
    |V|\,\dim\varrho
      \left(\sum_{\kappa\in\Irr(G)}\dim\kappa\right)\lambda^{n/m}.
  \end{align*}
  This proves the first displayed estimate. Since \(n/m\le n/2\) for
  \(m>1\),
  \[
    \sum_{\substack{m\mid n\\ m>1}}\lambda^{n/m}
    \le \sum_{k=1}^{\lfloor n/2\rfloor}\lambda^k
    \le \frac{\lambda^{\lfloor n/2\rfloor+1}}{\lambda-1}
    \le \frac{\lambda^{n/2+1}}{\lambda-1}.
  \]
  Substituting this into the first estimate gives the second displayed
  bound, and the \(O\)-estimate follows.
\end{proof}

\begin{lemma}[Absolute convergence at the Perron point]
  \label{lem:nontrivial-perron-point-convergence}
  Assume \(\lambda>1\) and \(\eta<\lambda\). Then, for every non-trivial
  \(\varrho\in\Irr(G)\), the series
  \[
    \Pi_\varrho(\lambda^{-1})
    =\sum_{n\ge 1}\pi_{n,\varrho}(A,\tau)\lambda^{-n}
  \]
  converges absolutely.
\end{lemma}

\begin{proof}
  Let
  \[
    \vartheta:=\max\{\eta/\lambda,\lambda^{-1/2}\}\in(0,1).
  \]
  By \Cref{lem:nontrivial-peter-weyl-channel-bound},
  \[
    \abs{\pi_{n,\varrho}(A,\tau)}\lambda^{-n}
    =O\!\left(\frac{\vartheta^n}{n}\right),
  \]
  and the series \(\sum_{n\ge 1}\vartheta^n/n\) converges.
\end{proof}

\begin{lemma}[Abel limit equals the ordinary correction sum]
  \label{lem:abel-correction-sum}
  Let \(A\) be a primitive non-negative matrix with Perron root
  \(\lambda>1\). If
  \[
    \sum_{n\ge 1}\left|\frac{p_n(A)}{\lambda^n}-\frac{1}{n}\right|
    <\infty,
  \]
  then
  \[
    \FP(A)
    =
    \sum_{n\ge 1}\left(\frac{p_n(A)}{\lambda^n}-\frac{1}{n}\right).
  \]
  Equivalently, for every \(0\le r<1\),
  \[
    P_A(r)+\log(1-r)
    =
    \sum_{n\ge 1}\left(\frac{p_n(A)}{\lambda^n}-\frac{1}{n}\right)r^n,
  \]
  and the limit \(r\uparrow 1\) may be passed termwise.
\end{lemma}

\begin{proof}
  For \(0\le r<1\), the harmonic-series identity
  \[
    -\log(1-r)=\sum_{n\ge 1}\frac{r^n}{n}
  \]
  gives
  \[
    P_A(r)+\log(1-r)
    =
    \sum_{n\ge 1}\left(\frac{p_n(A)}{\lambda^n}-\frac{1}{n}\right)r^n.
  \]
  The right-hand side is dominated by the absolutely summable majorant
  \[
    \sum_{n\ge 1}\left|\frac{p_n(A)}{\lambda^n}-\frac{1}{n}\right|,
  \]
  so dominated convergence allows \(r\uparrow 1\) termwise. By the
  definition of \(\FP(A)\), the limit equals \(\FP(A)\), proving the
  stated ordinary-sum formula.
\end{proof}

\begin{lemma}[Geometric-harmonic tails]
  \label{lem:geometric-harmonic-tail}
  If \(0<\vartheta<1\), then for every \(N\ge 1\),
  \[
    \sum_{n>N}\frac{\vartheta^n}{n}
    \le
    \frac{\vartheta^{N+1}}{(N+1)(1-\vartheta)}.
  \]
  In particular, if \(b_n=O(\vartheta^n/n)\), then
  \[
    \sum_{n>N}b_n=O\!\left(\frac{\vartheta^N}{N}\right).
  \]
\end{lemma}

\begin{proof}
  Since \(n\ge N+1\) throughout the tail,
  \[
    \sum_{n>N}\frac{\vartheta^n}{n}
    \le
    \frac{1}{N+1}\sum_{n>N}\vartheta^n
    =
    \frac{\vartheta^{N+1}}{(N+1)(1-\vartheta)}.
  \]
  The second statement is immediate.
\end{proof}

\begin{lemma}[Harmonic-log remainder bound]
  \label{lem:harmonic-log-remainder}
  For every \(N\ge 1\),
  \[
    0
    <
    H_N-\log N-\gamma
    \le
    \frac{1}{N+1}.
  \]
\end{lemma}

\begin{proof}
  Since
  \[
    \log(N+1)-\log N
    =
    \log\!\left(1+\frac{1}{N}\right)
  \]
  and \(\frac{x}{1+x}\le \log(1+x)\le x\) for \(x>-1\), one has
  \[
    \frac{1}{N+1}
    \le
    \log\!\left(1+\frac{1}{N}\right)
    \le
    \frac{1}{N}.
  \]
  Therefore
  \[
    0
    \le
    \frac{1}{N}-\log\!\left(1+\frac{1}{N}\right)
    \le
    \frac{1}{N(N+1)}.
  \]
  Using the defining limit
  \(\gamma=\lim_{M\to\infty}(H_M-\log M)\), we obtain
  \begin{align*}
    H_N-\log N-\gamma
    &=
    \sum_{m>N}
      \left(
        \frac{1}{m}-\log\!\left(1+\frac{1}{m}\right)
      \right) \\
    &\le
    \sum_{m>N}\frac{1}{m(m+1)}
    =
    \frac{1}{N+1}.
  \end{align*}
  Each summand is strictly positive, so the left-hand inequality follows.
\end{proof}

\begin{lemma}[Perron-channel contribution]
  \label{lem:perron-channel-contribution}
  Assume \(\lambda>1\), \(\eta<\lambda\), and set
  \[
    \Lambda_*:=\max\{\Lambda(A),\eta,\lambda^{1/2}\}.
  \]
  Then \(\vartheta_*:=\Lambda_*/\lambda\in(0,1)\), and
  \[
    \sum_{n\le N}\frac{p_n(A)}{\lambda^n}
    =
    H_N+\mathcal M(A)-\gamma
    +O\!\left(\frac{\vartheta_*^N}{N}\right).
  \]
\end{lemma}

\begin{proof}
  Since \(\Lambda(A)<\lambda\), \(\eta<\lambda\), and
  \(\lambda^{1/2}<\lambda\), one has \(\vartheta_*\in(0,1)\). By
  \Cref{lem:primitive-orbit-asymptotic},
  \[
    \frac{p_n(A)}{\lambda^n}
    =
    \frac{1}{n}
    +O\!\left(\frac{\vartheta_*^n}{n}\right).
  \]
  Hence the correction series
  \[
    \sum_{n\ge 1}\left(\frac{p_n(A)}{\lambda^n}-\frac{1}{n}\right)
  \]
  converges absolutely, and \Cref{lem:abel-correction-sum} gives
  \[
    \sum_{n\ge 1}\left(\frac{p_n(A)}{\lambda^n}-\frac{1}{n}\right)
    =
    \FP(A)
    =
    \mathcal M(A)-\gamma.
  \]
  Therefore
  \[
    \sum_{n\le N}\frac{p_n(A)}{\lambda^n}
    =
    H_N+\mathcal M(A)-\gamma
    -\sum_{n>N}\left(\frac{p_n(A)}{\lambda^n}-\frac{1}{n}\right).
  \]
  The tail is \(O(\vartheta_*^N/N)\) by
  \Cref{lem:geometric-harmonic-tail}.
\end{proof}
% END INPUT sec_chebotarev_peter_weyl_bridge.tex
% BEGIN INPUT sec_chebotarev_class_mertens.tex
% BEGIN INPUT sec_chebotarev_class_mertens_overview.tex
\subsection{Class Mertens constants and Adams-corrected Euler products}
\label{subsec:class-mertens-adams-products}

This subsection is the finite-group Mertens core of the paper.  It first
proves Main Theorem~II, the additive Frobenius-class constant formula, and
then Main Theorem~III, the Adams-corrected Euler-product formula.  The later
determinant, stability, and obstruction results all refer to the same
scalar-normalised product vector.  The proof route is:
\[
\begin{array}{c}
  \text{class Mertens constants}\\
  \downarrow\;\text{Euler product}\\
  \text{class product constants}\\
  \downarrow\;\text{fixed-label transform}\\
  \text{finite determinant evaluation}\\
  \downarrow\;\text{coordinate comparison}\\
  \text{stability and Artin obstruction}.
\end{array}
\]
Here \Cref{thm:class-mertens-explicit} supplies the additive constants, while
\Cref{thm:fixed-label-euler-transform,thm:class-euler-product-mertens} prove
Main Theorem~III.  The remaining results in the subsection have three roles:
finite determinant evaluation
\Cref{prop:fixed-label-determinant-perron,prop:fixed-label-perron-finite-algorithm,thm:class-product-determinant-rigidity},
coordinate comparison and branch conventions
\Cref{prop:class-product-fixed-coordinate-equivalence,thm:class-product-comparison,cor:class-product-hyperplane-stability},
and Artin-substitution obstruction
\Cref{thm:adams-obstruction-class-product,cor:adams-obstruction-bound}.  Thus
the subsequent stability, obstruction, and quotient-shadow statements use the
same recovered product vector without introducing further product constants.

The notation and dependencies for Main Theorem~III are as follows.
\begin{center}
\footnotesize
\begin{tabular}{@{}L{0.18\textwidth}L{0.40\textwidth}L{0.34\textwidth}@{}}
\toprule
Symbol & Meaning in this subsection & Role in the theorem chain \\
\midrule
\(\Pi_\varrho(z)\) & Primitive Peter--Weyl channel obtained from Adams--M\"obius determinant inversion & Perron value gives the non-trivial additive class constant coordinate. \\
\(\mathcal L_\varrho(z)\) & Periodic Artin logarithm \(\sum_{r\ge1}\Pi_{\psi^r\varrho}(z^r)/r=\log\det(I-zA_\varrho)^{-1}\) & Periodic comparison object; it is not the product-coordinate input. \\
\(F_\varrho(z)\) & Fixed-label Euler transform \(\sum_{r\ge1}\Pi_\varrho(z^r)/r\) & New product-coordinate input in Main Theorem~III. \\
\(\mathcal M_C(A,\tau;G)\) & Additive Frobenius-class Mertens constant & Output of Main Theorem~II and scalar finite-part normalisation input for products. \\
\(K_C\) & Positive Frobenius-class Euler-product constant & Output of Main Theorem~III, recovered from \(C(A)\) and \((F_\varrho(\lambda^{-1}))_{\varrho\ne\mathbf1}\). \\
\(\chi_\varrho^\ast(C)\) & Hermitian convention \(\overline{\chi_\varrho(C)}\) & Forward class transform; inverse readout uses \(\chi_\varrho(C)\). \\
Adams coefficients & Multiplicities \(a_{\varrho,\pi}^{(m)}\) in \(\psi^m\chi_\varrho=\sum_\pi a_{\varrho,\pi}^{(m)}\chi_\pi\) & Convert periodic determinant logarithms into fixed-label Perron values. \\
Branch and deletion & Principal determinant branches in the open disc; Abel-path boundary values at \(\lambda^{-1}\); omitted singular triple \((r,m,\pi)=(1,1,\mathbf1)\); positive real branch for \(K_C\) & Keeps the scalar Perron determinant out of non-trivial fixed-label formulas while retaining all regular scalar terms with \(rm\ge2\). \\
\bottomrule
\end{tabular}
\end{center}

With this notation, the proof blocks are additive constants
\Cref{thm:class-mertens-explicit}, product tails
\Cref{subsubsec:class-product-tail-reduction}, fixed-label product constants
\Cref{subsubsec:main-iii-fixed-label-products}, determinant-only computation
\Cref{subsubsec:determinant-only-fixed-label-values}, local observation and
rigidity
\Cref{subsubsec:class-product-vector-rigidity,subsubsec:class-product-local-observation},
classical comparison
\Cref{subsubsec:finite-determinant-closure-classical-comparison}, and Artin
obstruction certificates \Cref{thm:adams-obstruction-class-product}.

The constant normalisation is therefore visible in one line:
\[
  \{\mathcal M_C\}_C
  \Longleftrightarrow
  \{\Pi_\varrho(\lambda^{-1})\}_{\varrho\ne\mathbf1},
  \qquad
  \{K_C\}_C+C(A)
  \Longleftrightarrow
  \{F_\varrho(\lambda^{-1})\}_{\varrho\ne\mathbf1}.
\]
The first equivalence is the additive class-Mertens formula together with the
Hermitian class transform.  The second is the scalar finite-part identity
\(\mathcal M(A)=\gamma+\FP(A)=\gamma+\log C(A)\) combined with the fixed-label
Euler transform and the positive-branch product formula for \(K_C\).

For reference, the finite-data closure statement below is placed at the
front because the section intentionally proves several consequences of the
same scalar-normalised vector rather than scattering them into independent
black boxes.

For dependency checking, the proof spine should be read in the following
order.  First, \Cref{thm:class-mertens-explicit} proves the additive
Chebotarev--Mertens asymptotic under the two headline hypotheses
\(\lambda>1\) and \(\eta<\lambda\).  Second,
\Cref{thm:fixed-label-euler-transform,thm:class-euler-product-mertens} converts
those additive constants into the fixed-label Euler-product constants; this is
the only step where the distinction between the fixed-label transform
\(F_\varrho\) and the periodic Artin logarithm \(\mathcal L_\varrho\) enters.
Third, the determinant and rigidity results evaluate the non-trivial Perron
coordinates from finite determinant data and record the effective truncation
convention.  Finally, the comparison, observation-budget, and obstruction
results are consequences of the already recovered product vector.  Thus the
later rigidity, stability, and \(S_3\) certificate statements do not feed back
into the proofs of Main Theorems~II--III; they only test or separate the
product constants produced by this ordered chain.

\begin{proposition}[Finite-data closure for Main Theorems II--III]
  \label{prop:class-mertens-product-dependency-chain}
  Assume the standing hypotheses for the class-Mertens theorem, namely
  \(\lambda>1\) and the strict twisted spectral gap
  \(\eta<\lambda\).  Let
  \[
  \begin{aligned}
    \mathfrak D(A,\tau;G)=
    \Bigl(&G,\{C\subset G\},(C,m)\mapsto C^m,\\
      &C(A),\{\det(I-zA_\pi):\pi\in\Irr(G)\}\Bigr)
  \end{aligned}
  \]
  denote the finite determinant data together with the class-power table
  and Perron normalisation.  Then the results in this subsection form the
  following ordered chain, and every later stability, comparison, and
  obstruction statement factors through this chain:
  \[
  \begin{array}{ccl}
    \mathfrak D(A,\tau;G)&\Longrightarrow&
      \{p_{n,C}\}_{n,C}\quad\text{and}\quad
      \{\Pi_\varrho(\lambda^{-1})\}_{\varrho\ne\mathbf1}\\[0.4ex]
    &\Longrightarrow&
      \{\mathcal M_C(A,\tau;G)\}_C\\[0.4ex]
    &\Longrightarrow&
      \{F_\varrho(\lambda^{-1})\}_{\varrho\ne\mathbf1}\\[0.4ex]
    &\Longleftrightarrow&
      \{K_C\}_C\quad\text{after the scalar normalisation }C(A).
  \end{array}
  \]
  More precisely:
  \begin{enumerate}[label=\textup{(\roman*)},leftmargin=*]
    \item the first arrow is exactly Main Theorem~I together with
    \Cref{thm:class-mertens-explicit}; it uses coefficient extraction,
    Adams coefficients, finite character orthogonality, and only the regular
    Abel-path Perron boundary values of the non-trivial twisted determinant
    series;
    \item the second arrow is the class-Mertens identity
    \eqref{eq:class-mertens-constant};
    \item the third arrow is the fixed-label Euler transform of
    \Cref{thm:fixed-label-euler-transform};
    \item the final equivalence is the scalar-normalised Fourier inversion
    \eqref{eq:comparison-fixed-label-inverse} of
    \Cref{thm:class-product-comparison}, equivalently the rigidity formula
    of \Cref{thm:class-product-determinant-rigidity}.
  \end{enumerate}
  Consequently the finite determinant data, together with the stated
  branch and scalar-normalisation conventions, determines the whole output of
  Main Theorems~II--III.  Conversely, once \(C(A)\) is fixed, equality of two
  product-constant vectors \((K_C)_C\) is equivalent to equality of all
  non-trivial fixed-label Perron coordinates
  \((F_\varrho(\lambda^{-1}))_{\varrho\ne\mathbf1}\).  Thus any attempted
  substitute for the Adams-corrected product constants is correct for every
  conjugacy class if and only if it matches the same non-trivial coordinate
  vector.

  The source closure is literal.  Let \((A',\tau',G)\) be another finite-group
  extension satisfying the same hypotheses with the same Perron value
  \(\lambda\), the same scalar normalisation \(C(A')=C(A)\), and the same
  group \(G\).  Suppose that the branch-compatible determinant values which
  occur in
  \Cref{prop:fixed-label-determinant-perron} agree for the two systems,
  namely
  \[
    \Log\det(I-\lambda^{-k}A_\pi)^{-1}
    =
    \Log\det(I-\lambda^{-k}A'_\pi)^{-1}
  \]
  for every \(k=rm\ge2\) and every \(\pi\in\Irr(G)\), and also for the
  Abel-path non-trivial boundary values with \(k=1\) and
  \(\pi\ne\mathbf1\).  Then the two extensions have identical values
  \[
    F_\varrho(\lambda^{-1}),\qquad
    \mathcal M_C(A,\tau;G),
    \qquad K_C
  \]
  for all non-trivial \(\varrho\) and all conjugacy classes \(C\).  Conversely,
  if \(C(A')=C(A)\) and the positive product vectors \((K_C)_C\) agree, then
  all non-trivial fixed-label Perron coordinates
  \(F_\varrho(\lambda^{-1})\) agree.  Hence any later finite auxiliary
  calculation, secondary-edge refinement, cyclic-lift scalar comparison, or
  notational substitute can affect the main constants only by changing either the scalar
  normalisation or one of these non-trivial determinant boundary coordinates.
\end{proposition}

\begin{proof}
  Main Theorem~I expresses the primitive Peter--Weyl coefficient
  \(\pi_{n,\varrho}(A,\tau)\) by Adams--M\"obius inversion from the finite
  twisted determinant logarithms.  Finite character orthogonality then
  recovers every primitive class count \(p_{n,C}\).  Under the strict twisted
  spectral gap, \Cref{thm:class-mertens-explicit} evaluates the Perron-point
  primitive values \(\Pi_\varrho(\lambda^{-1})\) for
  \(\varrho\ne\mathbf1\) along the specified Abel path and gives the
  class-Mertens constants by \eqref{eq:class-mertens-constant}.  This proves
  the first two arrows.

  The fixed-label Euler transform is a convergent Adams-weighted transform of
  the primitive channels at the Perron point; this is precisely
  \Cref{thm:fixed-label-euler-transform}.  Substituting those values into
  \Cref{thm:class-euler-product-mertens} gives the product constants.  After
  the scalar base \(C(A)\) has been separated, the normalised logarithmic
  vector
  \[
    \mathfrak K(C)=-\frac{|G|}{|C|}\log K_C-(\gamma+\log C(A))
  \]
  has Fourier expansion
  \(\mathfrak K(C)=\sum_{\varrho\ne\mathbf1}
  \chi_\varrho^\ast(C)F_\varrho(\lambda^{-1})\).  Orthogonality of the
  irreducible characters gives the inverse formula
  \eqref{eq:comparison-fixed-label-inverse}.  Hence the product vector and
  the non-trivial fixed-label coordinate vector determine one another once
  the scalar normalisation is fixed.

  The final assertion follows immediately from the same inverse Fourier
  formula.  If a substitute has all the same class constants, its normalised
  logarithmic vector is the same and therefore all non-trivial Fourier
  coordinates agree.  If the coordinates agree, the displayed Fourier
  expansion gives the same normalised logarithmic vector and hence the same
  class constants.  The stability, comparison, and obstruction results later
  in the subsection use only these coordinates, so they factor through the
  stated chain.

  It remains to prove the two-extension source-closure assertion.  The Adams
  coefficients \(a_{\varrho,\pi}^{(m)}\), the M\"obius factors, and the class
  characters depend only on the common finite group \(G\).  Therefore the
  determinant formula of \Cref{prop:fixed-label-determinant-perron} is the
  same absolutely convergent linear combination of the displayed
  branch-compatible determinant values for the two systems.  Equality of
  those values gives
  \(F_\varrho(\lambda^{-1})=F'_{\varrho}(\lambda^{-1})\) for every
  \(\varrho\ne\mathbf1\).  The same determinant-value hypothesis applies to
  the boundary formula of \Cref{lem:perron-boundary-determinant-evaluation},
  so
  \(\Pi_\varrho(\lambda^{-1})=\Pi'_{\varrho}(\lambda^{-1})\) for every
  non-trivial \(\varrho\).  The class-Mertens formula
  \eqref{eq:class-mertens-constant} then gives equality of the additive
  constants \(\mathcal M_C\), and the product formula
  \eqref{eq:class-euler-product-k-f}, together with
  \(C(A')=C(A)\), gives equality of all \(K_C\).

  Conversely, if the scalar normalisations and positive product vectors agree,
  then the normalised logarithmic vectors \(\mathfrak K\) agree.  Applying the
  inverse Fourier formula \eqref{eq:comparison-fixed-label-inverse}, or
  equivalently \eqref{eq:class-product-rigidity-inverse}, gives equality of
  every non-trivial fixed-label coordinate.  Thus an auxiliary record which
  leaves the scalar normalisation and these determinant boundary coordinates
  fixed is invisible to the main constants, while changing a main constant
  necessarily changes one of those two pieces of data.  This proves the stated
  source closure.
\end{proof}

\begin{proof}[Assembly of the three main statements]
  Under the hypotheses stated in the introduction, the three displayed
  introduction theorems are exactly the following body theorem chain.
  \begin{enumerate}[label=\textup{(\roman*)},leftmargin=*]
    \item Main Theorem~I is the finite determinant interface of
    \Cref{prop:main-theorem-i-assembly-map}: twisted traces give the
    class-function determinant logarithms, Adams--M\"obius inversion gives
    the primitive channels, and finite character orthogonality gives the
    primitive Frobenius-class counts.
    \item Main Theorem~II is the Perron-boundary continuation of those same
    primitive channels under \(\lambda>1\) and \(\eta<\lambda\), namely
    \Cref{thm:class-mertens-explicit,prop:class-mertens-deviations,thm:class-mertens-fourier}.
    It produces the class constants \(\mathcal M_C(A,\tau;G)\) and their
    Hermitian Fourier--Parseval readout.
    \item Main Theorem~III is the fixed-label Euler-product evaluation and
    rigidity package
    \Cref{thm:class-euler-product-mertens,prop:fixed-label-determinant-perron,prop:fixed-label-perron-finite-algorithm,thm:class-product-determinant-rigidity,prop:class-product-comparison,cor:class-product-hyperplane-stability,cor:class-product-artin-obstruction-recovery,cor:assembled-adams-frobenius-product-formula,thm:adams-obstruction-class-product},
    with the finite \(S_3\) certificates used only as explicit witnesses to
    the Artin-substitution obstruction.
  \end{enumerate}
  Consequently the introduction theorem environments are assembly statements:
  no additional hypothesis, boundary value, branch convention, scalar
  normalisation, or appendix input is hidden in their restatement.  In
  particular, the authoritative proof obligations are the body statements
  listed above, and every output in Main Theorems~I--III factors through the
  finite-data chain of \Cref{prop:class-mertens-product-dependency-chain}.

  For \textup{(i)}, \Cref{prop:main-theorem-i-assembly-map} is proved before
  any Perron-boundary assumption is imposed.  Its proof cites precisely
  \Cref{lem:twisted-trace-identity,thm:class-function-linearisation,cor:class-function-adams-mobius,cor:primitive-peter-weyl-determinant,thm:primitive-peter-weyl-trace,thm:class-count-trace-recovery}.
  These statements establish, in order, the fixed-point trace identity, the
  finite determinant linearisation of \(\mathcal L_\phi\), the Adams--M\"obius
  identity for \(\Pi_\phi\), the irreducible determinant-germ formula, the
  primitive Peter--Weyl coefficient formula, and class-count recovery by finite
  character orthogonality.  This is exactly the content of the first
  introduction theorem.

  For \textup{(ii)}, imposing \(\lambda>1\) and \(\eta<\lambda\) gives the
  non-trivial Perron-point values of the primitive channels by
  \Cref{thm:class-mertens-explicit}.  The same theorem proves the displayed
  class-Mertens asymptotic and the formula
  \eqref{eq:class-mertens-constant};
  \Cref{prop:class-mertens-deviations,thm:class-mertens-fourier} then identify
  the normalised deviation profile and its Hermitian Fourier--Parseval
  transform.  These are exactly the additional assertions made in the second
  introduction theorem, and their only new analytic hypothesis is the strict
  finite-dimensional twisted gap already stated there.

  For \textup{(iii)}, \Cref{thm:class-euler-product-mertens} proves the
  Frobenius-class Euler-product asymptotic with the fixed-label transform
  \(F_\varrho(\lambda^{-1})\) and the product constant \(K_C\).  The determinant
  evaluation and finite certificate for those fixed-label Perron values are
  \Cref{prop:fixed-label-determinant-perron,prop:fixed-label-perron-finite-algorithm};
  substituting them into the class product gives
  \Cref{thm:class-product-determinant-rigidity,cor:assembled-adams-frobenius-product-formula}.
  The comparison, stability, scalar-baseline readout, and obstruction clauses
  are precisely
  \Cref{prop:class-product-comparison,cor:class-product-hyperplane-stability,cor:class-product-artin-obstruction-recovery,thm:adams-obstruction-class-product}.
  The cited \(S_3\) statements supply concrete finite witnesses to that last
  obstruction and do not add hypotheses to the product theorem.  Finally,
  \Cref{prop:class-mertens-product-dependency-chain} proves that these three
  tiers are connected by one finite-data dependency chain, so the introduction
  restatements introduce no independent theorem beyond the assembled body
  chain.
\end{proof}

\begin{statusnote}[Product-constant refinements]
  Under the hypotheses of \Cref{thm:class-euler-product-mertens}, fix the
  scalar normalisation \(\alpha_A=\gamma+\log C(A)\) and the non-trivial
  Perron-boundary vector
  \[
    \mathsf F(A,\tau;G):=
    \bigl(F_\varrho(\lambda^{-1})\bigr)_{\varrho\ne\mathbf1}.
  \]
  Then the hyperplane stability of
  \Cref{cor:class-product-hyperplane-stability}, the assembled finite
  determinant formula and truncation of
  \Cref{cor:assembled-adams-frobenius-product-formula}, the fixed-coordinate
  rank obstruction of \Cref{cor:class-product-minimal-obstruction-skeleton},
  and the Parseval
  obstruction budget of \Cref{cor:adams-obstruction-bound} are downstream
  refinements of the Main Theorem~III coordinate vector
  \((\alpha_A,\mathsf F(A,\tau;G))\).  More precisely, after the periodic
  Artin comparison vector
  \((\mathcal L_\varrho(\lambda^{-1}))_{\varrho\ne\mathbf1}\) is fixed for
  the obstruction statements, each listed refinement is functorially
  determined by these already recovered coordinates.  Conversely, equality of
  the scalar-normalised fixed-coordinate product vector, or equality of the
  positive product
  constants with the same scalar normalisation, forces equality of
  \(\mathsf F(A,\tau;G)\).  Thus none of the listed refinements supplies an
  independent main contribution, an additional hypothesis for the product
  asymptotic, or a new source of class-product constants.
\end{statusnote}

Indeed, Main Theorem~III is \Cref{thm:class-euler-product-mertens}.  It gives
  \[
    K_C=
    \exp\!\left(-\frac{|C|}{|G|}
      \left(\alpha_A+
        \sum_{\varrho\ne\mathbf1}\chi_\varrho^\ast(C)
          F_\varrho(\lambda^{-1})\right)
      \right),
  \]
  so \((\alpha_A,\mathsf F(A,\tau;G))\) determines the positive product
  vector.  The hyperplane statement is the same formula rewritten in the
  zero-mean class-function chart, and its stability estimate is Parseval's
  identity in that chart.  The assembled finite determinant formula is obtained
  by substituting the determinant expression of
  \Cref{prop:fixed-label-determinant-perron} for each coordinate of
  \(\mathsf F(A,\tau;G)\); its truncation bound is the tail estimate supplied
  there and in \Cref{lem:fixed-label-perron-double-majorant}.  Hence both are
  functions of the already fixed Main Theorem~III coordinates and the finite
  determinant data used to recover them.

  For the fixed-coordinate rank statement, set
  \(H_K(C)=-|G|\log K_C/|C|-\alpha_A\).  The displayed formula identifies the
  non-trivial Fourier coefficients of \(H_K\) with
  \(\mathsf F(A,\tau;G)\).  Therefore
  \Cref{cor:class-product-minimal-obstruction-skeleton} is precisely the
  finite-dimensional injectivity and rank statement for this zero-mean
  coordinate chart; in particular its inverse Fourier formula gives the stated
  converse.  If the periodic Artin comparison vector is also fixed, then the
  error vector
  \(D_\varrho=F_\varrho(\lambda^{-1})-
  \mathcal L_\varrho(\lambda^{-1})\) is fixed, and
  \Cref{cor:adams-obstruction-bound} is exactly Parseval's identity applied to
  that already determined vector.These verifications use no determinant channel, asymptotic input, or product
constant beyond those in Main Theorem~III.

% END INPUT sec_chebotarev_class_mertens_overview.tex
% BEGIN INPUT sec_chebotarev_class_mertens_additive.tex
\subsubsection{Additive class constants}

Classical Chebotarev orbit-counting and Frobenius-class Mertens product
theorems for hyperbolic and symbolic extensions are due in this setting to
Parry--Pollicott, Noorani--Parry, Sharp, and Mesliza--Noorani
\cite{ParryPollicott1986Chebotarev,NooraniParry1992ChebotarevShifts,Sharp1991Mertens,MeslizaNoorani1999FrobeniusMertens}.  The next theorem is
not presented as a new Tauberian input: under the explicit finite-dimensional
spectral gap \(\eta<\lambda\), it records the elementary Perron--Frobenius and
Peter--Weyl version needed here, with constants normalized for the later
Adams--M\"obius determinant calculation.

\begin{theorem}[Main Theorem II(a): class Mertens theorem under the
  stated twisted spectral gap]
  \label{thm:class-mertens-explicit}
  Assume that \(\lambda>1\) and
  \[
    \eta:=\max_{\varrho\neq\mathbf 1}\rho(A_\varrho)<\lambda,
    \qquad
    \Lambda_*:=\max\{\Lambda(A),\eta,\lambda^{1/2}\}.
  \]
  If \(G\) is trivial, the maximum defining \(\eta\) is empty and
  \(\eta=0\). For non-trivial \(G\), this is the finite-state hypothesis
  certified equivalently by
  \Cref{prop:eta-gap-finite-state-certificate}; it is not asserted for
  all cocycles.
  Write \(\theta:=\Lambda_*/\lambda\in(0,1)\).
  Then, for every conjugacy class \(C\subset G\),
  \begin{equation}\label{eq:class-primitive-asymptotic}
    p_{n,C}
    =
    \frac{|C|}{|G|}\frac{\lambda^n}{n}
    +O\!\left(\frac{\Lambda_*^n}{n}\right).
  \end{equation}
  In particular, for all sufficiently large \(n\),
  \begin{equation}\label{eq:class-chebotarev-density}
    \frac{p_{n,C}}{p_n(A)}
    =
    \frac{|C|}{|G|}
    +O\!\left(\left(\frac{\Lambda_*}{\lambda}\right)^n\right).
  \end{equation}
  Moreover, the series \(\Pi_\varrho(\lambda^{-1})\) converges
  absolutely for every non-trivial \(\varrho\), and
  \begin{equation}\label{eq:class-mertens-harmonic}
    \sum_{n\le N}\frac{p_{n,C}}{\lambda^n}
    =
    \frac{|C|}{|G|}H_N
    +\mathcal M_C(A,\tau;G)-\frac{|C|}{|G|}\gamma
    +O\!\left(\frac{(\Lambda_*/\lambda)^N}{N}\right),
  \end{equation}
  where \(H_N=\sum_{n\le N}1/n\) and the corresponding class-indexed
  constant is
  \begin{equation}\label{eq:class-mertens-constant}
    \mathcal M_C(A,\tau;G)
    =
    \frac{|C|}{|G|}\mathcal M(A)
    +\frac{|C|}{|G|}
      \sum_{\varrho\neq\mathbf 1}
        \chi_\varrho^\ast(C)\,\Pi_\varrho(\lambda^{-1}).
  \end{equation}
  Let \(K_{\mathrm{PF}}(A)\) and \(K_{\mathrm{tr}}(A)\) be the base
  constants from \Cref{lem:primitive-orbit-asymptotic}. Then
  \begin{equation}\label{eq:class-mertens-base-constant}
    \abs{p_n(A)-\lambda^n/n}
    \le
    \frac{K_{\mathrm{PF}}(A)\Lambda(A)^n}{n}
    +\frac{\lambda K_{\mathrm{tr}}(A)}{\lambda-1}
      \frac{\lambda^{n/2}}{n}
    \qquad (n\ge 1),
  \end{equation}
  Define
  \[
    K_{\mathrm{base}}(A)
    :=
    K_{\mathrm{PF}}(A)
    +\frac{\lambda K_{\mathrm{tr}}(A)}{\lambda-1}.
  \]
  Then
  \[
    \abs{p_n(A)-\lambda^n/n}
    \le
    \frac{K_{\mathrm{base}}(A)}{n}
      \max\{\Lambda(A)^n,\lambda^{n/2}\}.
  \]
  Set
  \[
    D_G:=\sum_{\pi\in\Irr(G)}\dim\pi,
    \qquad
    S_G:=\sum_{\varrho\neq\mathbf 1}(\dim\varrho)^2.
  \]
  Define
  \begin{equation}\label{eq:class-mertens-explicit-constant}
    C_{\mathrm{cls}}(A,\tau;G)
    :=
    K_{\mathrm{base}}(A)
    +|V|\,S_G
    +\frac{\lambda |V|\,D_G}{\lambda-1}\,S_G.
  \end{equation}
  Then, for every conjugacy class \(C\subset G\) and every \(n\ge 1\),
  \begin{equation}\label{eq:class-primitive-explicit-bound}
    \abs{
      p_{n,C}
      -\frac{|C|}{|G|}\frac{\lambda^n}{n}
    }
    \le
    \frac{|C|}{|G|}
    \frac{C_{\mathrm{cls}}(A,\tau;G)\Lambda_*^n}{n}.
  \end{equation}
  The displayed constant in
  \eqref{eq:class-primitive-explicit-bound} is uniform in the conjugacy
  class \(C\) and in \(n\).  The summatory conclusion of Main
  Theorem~II(a) is the harmonic form
  \eqref{eq:class-mertens-harmonic}; the logarithmic \(O(N^{-1})\)
  form used in the Euler-product passage is separated in
  \Cref{cor:class-mertens-log-form}.
\end{theorem}

\begin{proof}
  Set \(\vartheta_*:=\theta=\Lambda_*/\lambda\in(0,1)\). By
  \Cref{lem:class-primitive-trivial-split},
  \[
    p_{n,C}
    =
    \frac{|C|}{|G|}p_n(A)
    +\frac{|C|}{|G|}
      \sum_{\varrho\neq\mathbf 1}
        \chi_\varrho^\ast(C)\,\pi_{n,\varrho}(A,\tau).
  \]
  The two inputs are
  \[
    \abs{p_n(A)-\lambda^n/n}
    \le
    \frac{K_{\mathrm{base}}(A)}{n}\max\{\Lambda(A)^n,\lambda^{n/2}\}
  \]
  from \eqref{eq:class-mertens-base-constant}, and, for each
  non-trivial \(\varrho\),
  \[
    \abs{\pi_{n,\varrho}(A,\tau)}
    \le
    \frac{|V|\,\dim\varrho\,\eta^n}{n}
    +\frac{|V|\,\dim\varrho\,D_G}{n}
      \frac{\lambda^{n/2+1}}{\lambda-1}
  \]
  from \Cref{lem:nontrivial-peter-weyl-channel-bound}. Since
  \(\abs{\chi_\varrho(C)}\le \dim\varrho\) and
  \(\Lambda_*=\max\{\Lambda(A),\eta,\lambda^{1/2}\}\), one has
  \[
    \max\{\Lambda(A)^n,\lambda^{n/2}\}\le \Lambda_*^n,
    \qquad
    \eta^n\le \Lambda_*^n,
    \qquad
    \lambda^{n/2+1}\le \lambda\Lambda_*^n.
  \]
  Therefore
  \[
    \abs{
      p_{n,C}
      -\frac{|C|}{|G|}\frac{\lambda^n}{n}
    }
    \le
    \frac{|C|}{|G|}
    \frac{\Lambda_*^n}{n}
    \left(
      K_{\mathrm{base}}(A)
      +|V|\,S_G
      +\frac{\lambda |V|\,D_G}{\lambda-1}\,S_G
    \right),
  \]
  which is \eqref{eq:class-primitive-explicit-bound}. In particular,
  \[
    p_{n,C}
    =
    \frac{|C|}{|G|}\frac{\lambda^n}{n}
    +O\!\left(\frac{\Lambda_*^n}{n}\right),
  \]
  which is \eqref{eq:class-primitive-asymptotic}.
  Since
  \[
    p_n(A)=\frac{\lambda^n}{n}+O\!\left(\frac{\Lambda_*^n}{n}\right),
  \]
  equivalently
  \(p_n(A)=\frac{\lambda^n}{n}\bigl(1+O(\vartheta_*^n)\bigr)\) with
  \(0<\vartheta_*<1\), one has \(p_n(A)>0\) (hence \(p_n(A)\neq 0\))
  for all sufficiently large \(n\). Therefore
  dividing \eqref{eq:class-primitive-asymptotic} by \(p_n(A)\) yields
  \eqref{eq:class-chebotarev-density} for all sufficiently large \(n\).

  By \Cref{lem:nontrivial-perron-point-convergence}, the series
  \(\Pi_\varrho(\lambda^{-1})\) is absolutely convergent for every
  non-trivial \(\varrho\). Since
  \eqref{eq:class-primitive-explicit-bound} gives
  \[
    \abs{
      \frac{p_{n,C}}{\lambda^n}
      -\frac{|C|}{|G|}\frac{1}{n}
    }
    \le
    \frac{|C|}{|G|}
    \frac{C_{\mathrm{cls}}(A,\tau;G)\vartheta_*^n}{n},
  \]
  the correction series
  \[
    \sum_{n\ge 1}
      \left(
        \frac{p_{n,C}}{\lambda^n}
        -\frac{|C|}{|G|}\frac{1}{n}
      \right)
  \]
  converges absolutely. Using the Peter--Weyl decomposition termwise,
  together with \Cref{lem:abel-correction-sum} and the definition of
  \(\Pi_\varrho(\lambda^{-1})\), we obtain
  \begin{align*}
    \sum_{n\ge 1}
      \left(
        \frac{p_{n,C}}{\lambda^n}
        -\frac{|C|}{|G|}\frac{1}{n}
      \right)
    &=
    \frac{|C|}{|G|}
      \sum_{n\ge 1}
        \left(
          \frac{p_n(A)}{\lambda^n}-\frac{1}{n}
        \right) \\
    &\quad
    +
    \frac{|C|}{|G|}
      \sum_{\varrho\neq\mathbf 1}
        \chi_\varrho^\ast(C)
        \sum_{n\ge 1}\pi_{n,\varrho}(A,\tau)\lambda^{-n} \\
    &=
    \frac{|C|}{|G|}\bigl(\mathcal M(A)-\gamma\bigr)
    +\frac{|C|}{|G|}
      \sum_{\varrho\neq\mathbf 1}
        \chi_\varrho^\ast(C)\,\Pi_\varrho(\lambda^{-1}),
  \end{align*}
  which is exactly
  \eqref{eq:class-mertens-constant}. Therefore, for every \(N\ge 1\),
  \[
    \sum_{n\le N}\frac{p_{n,C}}{\lambda^n}
    =
    \frac{|C|}{|G|}H_N
    +\mathcal M_C(A,\tau;G)-\frac{|C|}{|G|}\gamma
    -\sum_{n>N}
      \left(
        \frac{p_{n,C}}{\lambda^n}
        -\frac{|C|}{|G|}\frac{1}{n}
      \right).
  \]
  Applying \Cref{lem:geometric-harmonic-tail} to the explicit
  correction bound above yields
  \[
    \abs{
      \sum_{n\le N}\frac{p_{n,C}}{\lambda^n}
      -\frac{|C|}{|G|}H_N
      -\mathcal M_C(A,\tau;G)+\frac{|C|}{|G|}\gamma
    }
    \le
    \frac{|C|}{|G|}
    \frac{
      C_{\mathrm{cls}}(A,\tau;G)\vartheta_*^{N+1}
    }{
      (N+1)(1-\vartheta_*)
    },
  \]
  which implies \eqref{eq:class-mertens-harmonic}.
\end{proof}

\begin{corollary}[Logarithmic form of Main Theorem II]
  \label{cor:class-mertens-log-form}
  Assume the hypotheses of \Cref{thm:class-mertens-explicit}.  Then, for
  every conjugacy class \(C\subset G\),
  \begin{equation}\label{eq:class-mertens-log}
    \sum_{n\le N}\frac{p_{n,C}}{\lambda^n}
    =
    \frac{|C|}{|G|}\log N+\mathcal M_C(A,\tau;G)+O(N^{-1}).
  \end{equation}
  More explicitly, with the same \(\theta=\Lambda_*/\lambda\) and
  \(C_{\mathrm{cls}}(A,\tau;G)\) as in
  \Cref{thm:class-mertens-explicit}, every \(N\ge1\) satisfies
  \begin{equation}\label{eq:class-mertens-log-explicit}
    \abs{
      \sum_{n\le N}\frac{p_{n,C}}{\lambda^n}
      -\frac{|C|}{|G|}\log N
      -\mathcal M_C(A,\tau;G)
    }
    \le
    \frac{|C|}{|G|}
    \left(
      \frac{1}{N+1}
      +
      \frac{
        C_{\mathrm{cls}}(A,\tau;G)\theta^{N+1}
      }{
        (N+1)(1-\theta)
      }
    \right).
  \end{equation}
  Thus \Cref{thm:class-mertens-explicit} supplies the harmonic form with
  exponentially small finite-state tail, while the present corollary is the
  \(\log N\) form and class-uniform \(O(N^{-1})\) remainder used in the
  proof of \Cref{thm:class-euler-product-mertens}.
\end{corollary}

\begin{proof}
  Subtract \((|C|/|G|)\gamma\) from both sides of
  \eqref{eq:class-mertens-harmonic} and add
  \((|C|/|G|)(H_N-\log N-\gamma)\).  The explicit harmonic remainder from
  \Cref{thm:class-mertens-explicit} gives the second term in the bound, and
  \Cref{lem:harmonic-log-remainder} gives
  \[
    \abs{H_N-\log N-\gamma}
    \le
    \frac{1}{N+1},
  \]
  which gives the first term.  This proves
  \eqref{eq:class-mertens-log-explicit}, and
  \eqref{eq:class-mertens-log} follows immediately.
\end{proof}

% END INPUT sec_chebotarev_class_mertens_additive.tex
% BEGIN INPUT sec_chebotarev_class_mertens_products.tex
\subsubsection{From Main Theorem II constants to product tails}
\label{subsubsec:class-product-tail-reduction}

This part ends the additive Main Theorem~II input and isolates the Euler-product
tail estimates needed before the fixed-label constants enter.  The constants
\(\mathcal M_C(A,\tau;G)\) are already fixed by
\Cref{thm:class-mertens-explicit}; the lemmas below only convert the
logarithmic class-Mertens asymptotic into a product asymptotic with a separate
absolutely convergent tail.

\begin{lemma}[Explicit product-tail bound]
  \label{lem:class-euler-product-tail-bound}
  Assume the hypotheses of \Cref{thm:class-mertens-explicit}.  For every
  conjugacy class \(C\subset G\), put
  \[
    T_C(N):=
      \sum_{r\ge2}\frac1r\sum_{n\le N}p_{n,C}\lambda^{-rn}.
  \]
  Then
  \begin{equation}\label{eq:class-product-tail-bound}
    T_C(N)=
      \sum_{r\ge2}\frac{\Pi_{\mathbf1_C}(\lambda^{-r})}{r}
      +O(\lambda^{-N}/N),
  \end{equation}
  uniformly in \(C\). Consequently, if
  \begin{equation}\label{eq:class-product-log-bound-input}
    \sum_{n\le N}\frac{p_{n,C}}{\lambda^n}
    =
    \frac{|C|}{|G|}\log N+\mathcal M_C(A,\tau;G)+O(N^{-1}),
  \end{equation}
  then
  \begin{equation}\label{eq:class-product-log-bound-output}
    \log P_C(N)=
    -\frac{|C|}{|G|}\log N
    -\mathcal M_C(A,\tau;G)
    -\sum_{r\ge2}\frac{\Pi_{\mathbf1_C}(\lambda^{-r})}{r}
    +O(N^{-1}).
  \end{equation}
\end{lemma}

\begin{proof}
  The primitive counting functions satisfy \(p_{n,C}=O(\lambda^n/n)\),
  uniformly in the finite set of conjugacy classes. Hence the full tail
  \(T_C(\infty)-T_C(N)\) is bounded by a constant multiple of
  \[
    \sum_{r\ge2}\frac1r\sum_{n>N}\frac{\lambda^{(1-r)n}}{n}.
  \]
  Since \(n\ge N+1\) on the inner tail,
  \[
    \sum_{r\ge2}\frac1r\sum_{n>N}\frac{\lambda^{(1-r)n}}{n}
    \le
    \frac{1}{N+1}
    \sum_{r\ge2}\frac{\lambda^{(1-r)(N+1)}}{r(1-\lambda^{1-r})}.
  \]
  Write \(q=\lambda^{-1}<1\).  The last sum equals
  \[
    \sum_{r\ge2}\frac{q^{(r-1)(N+1)}}{r(1-q^{r-1})}
    \le
    \frac{1}{2(1-q)}\sum_{s\ge1}q^{s(N+1)}
    =O(q^{N+1}),
  \]
  with a constant independent of \(C\).  This proves
  \eqref{eq:class-product-tail-bound}.  Taking logarithms in the finite
  Euler product gives
  \[
    \log P_C(N)=
    -\sum_{n\le N}\frac{p_{n,C}}{\lambda^n}-T_C(N).
  \]
  Substituting \eqref{eq:class-product-log-bound-input} and
  \eqref{eq:class-product-tail-bound} gives
  \eqref{eq:class-product-log-bound-output}, because
  \(O(\lambda^{-N}/N)\) is absorbed by \(O(N^{-1})\).
\end{proof}

\begin{lemma}[Exponentiating the class product logarithm]
  \label{lem:class-product-log-to-product-error}
  Let \(a\in\RR\), \(L\in\RR\), and let \(R_N=O(N^{-1})\) be a real sequence.
  If a positive sequence \(Q_N\) satisfies
  \[
    \log Q_N=-a\log N+L+R_N,
  \]
  then
  \[
    Q_N=e^L N^{-a}\bigl(1+O(N^{-1})\bigr).
  \]
  If the implied constant in \(R_N=O(N^{-1})\) is uniform over a finite set of
  indices, then the resulting multiplicative error is uniform over that set.
\end{lemma}

\begin{proof}
  Choose \(B>0\) with \(|R_N|\le B/N\) for all large \(N\).  Then
  \[
    Q_N=e^L N^{-a}e^{R_N}.
  \]
  For all large \(N\), \(|R_N|\le1\), and the elementary estimate
  \(|e^x-1|\le e|x|\) for \(|x|\le1\) gives
  \(e^{R_N}=1+O(N^{-1})\).  The same bound with the maximum of the finitely
  many constants gives the uniform assertion.
\end{proof}

\begin{proposition}[Isolated reduction to Main Theorem II and the product tail]
  \label{prop:class-euler-product-main-chain-reduction}
  Assume the hypotheses of \Cref{thm:class-mertens-explicit}.  For a fixed
  conjugacy class \(C\), put
  \[
    P_C(N):=
    \prod_{\substack{\gamma\in\cP\,:\\ |\gamma|\le N\\
      \tau(\gamma)\in C}}
      \left(1-\lambda^{-|\gamma|}\right).
  \]
  Then the logarithmic Main Theorem~II estimate
  \eqref{eq:class-mertens-log}, the \(r\ge2\) product-tail identity
  \eqref{eq:class-product-tail-bound}, and
  \Cref{lem:class-product-log-to-product-error} imply
  \begin{equation}\label{eq:class-euler-product-reduction-asymptotic}
    P_C(N)=
    \exp\!\left(
      -\mathcal M_C(A,\tau;G)
      -\sum_{r\ge2}\frac{\Pi_{\mathbf1_C}(\lambda^{-r})}{r}
    \right)
    N^{-|C|/|G|}\bigl(1+O(N^{-1})\bigr).
  \end{equation}
  Moreover the two constant formulas in Main Theorem~III are equivalent
  exactly through the fixed-label normalisation identity
  \begin{equation}\label{eq:class-euler-product-reduction-normalisation}
    \mathcal M_C(A,\tau;G)
    +\sum_{r\ge2}\frac{\Pi_{\mathbf1_C}(\lambda^{-r})}{r}
    =
    \frac{|C|}{|G|}
    \left(
      \gamma+
      \log C(A)+
      \sum_{\varrho\ne\mathbf1}
        \chi_\varrho^\ast(C)F_\varrho(\lambda^{-1})
    \right).
  \end{equation}
  Thus the proof of \Cref{thm:class-euler-product-mertens} has a short
  verifiable spine: Main Theorem~II controls the \(r=1\) logarithm, the
  product-tail lemma controls all \(r\ge2\) repetitions, and the fixed-label
  transform supplies the determinant-ready coordinates for the same constant.
\end{proposition}

\begin{proof}
  The logarithm of the finite Euler product is
  \[
    \log P_C(N)
    =
    -\sum_{n\le N}\frac{p_{n,C}}{\lambda^n}
    -
    \sum_{r\ge2}\frac1r\sum_{n\le N}p_{n,C}\lambda^{-rn}.
  \]
  Substituting \eqref{eq:class-mertens-log} into the first term and
  \eqref{eq:class-product-tail-bound} into the second term gives
  \[
    \log P_C(N)
    =
    -\frac{|C|}{|G|}\log N
    -\mathcal M_C(A,\tau;G)
    -\sum_{r\ge2}\frac{\Pi_{\mathbf1_C}(\lambda^{-r})}{r}
    +O(N^{-1}).
  \]
  Since \(P_C(N)>0\), \Cref{lem:class-product-log-to-product-error}
  exponentiates this estimate and proves
  \eqref{eq:class-euler-product-reduction-asymptotic}.

  The constants in \eqref{eq:class-euler-product-k-mc} and
  \eqref{eq:class-euler-product-k-f} are both written on the positive real
  branch.  Taking real logarithms of those two displayed expressions shows
  that their equality is precisely
  \eqref{eq:class-euler-product-reduction-normalisation}.  Conversely,
  substituting \eqref{eq:class-euler-product-reduction-normalisation} into
  the exponent of \eqref{eq:class-euler-product-k-mc} gives exactly the
  exponent of \eqref{eq:class-euler-product-k-f}.  Hence the two formulas are
  equivalent through the stated normalisation and no additional analytic input
  is hidden in the comparison.
\end{proof}

\begin{proposition}[Three logarithmic interfaces at the product constant]
  \label{prop:three-logarithm-product-interface}
  Assume the hypotheses of \Cref{thm:class-mertens-explicit}, and let
  \(\phi:G\to\CC\) be a class function orthogonal to the trivial
  character.  The three logarithmic objects used in Main Theorem~III have
  the following distinct roles:
  \begin{center}
  \begin{tabular}{@{}L{0.25\textwidth}L{0.36\textwidth}L{0.29\textwidth}@{}}
  \toprule
  Object & Formula & Product-constant role \\
  \midrule
  Periodic Artin logarithm
  & \(\mathcal L_\phi(z)=\sum_{r\ge1}\Pi_{\psi^r\phi}(z^r)/r\)
  & Encodes periodic determinant traces, where repetitions replace
    \(\phi\) by \(\psi^r\phi\). \\
  Primitive channel
  & \(\Pi_\phi(z)=\sum_{\gamma\in\cP}\phi(\tau(\gamma))z^{|\gamma|}\)
  & Is recovered from determinant logarithms by Adams--M\"obius inversion. \\
  Fixed-label Euler transform
  & \(F_\phi(z)=\sum_{r\ge1}\Pi_\phi(z^r)/r\)
  & Is the quantity evaluated at \(z=\lambda^{-1}\) in the non-trivial
    Fourier coordinates of \(K_C\). \\
  \bottomrule
  \end{tabular}
  \end{center}
  More precisely, \(F_\phi(\lambda^{-1})\) is absolutely convergent and is
  determined by the determinant data through
  \begin{equation}\label{eq:three-logarithm-product-interface}
    F_\phi(\lambda^{-1})
    =
    \sum_{r,m\ge1}\frac{\mu(m)}{rm}\,
      \mathcal L_{\psi^m\phi}(\lambda^{-rm}),
  \end{equation}
  using the Perron-boundary branches of \Cref{thm:fixed-label-euler-transform}.
  In the determinant expansion of the right-hand side the scalar Perron
  logarithm is not evaluated: the term \((r,m,\pi)=(1,1,\mathbf1)\) has zero
  coefficient because \(\phi\perp\mathbf1\), and it is deleted before the
  boundary value \(z=\lambda^{-1}\) is taken.  All remaining terms are either
  non-trivial Abel-boundary values made regular by \(\eta<\lambda\), or
  interior values with \(rm\ge2\).
  Consequently the non-trivial term in the class-product constant is a
  fixed-label Euler transform, and it may be replaced by the periodic
  Artin logarithm only in the channels satisfying the coefficient
  obstruction \eqref{eq:channel-separation-obstruction}.
  The fixed-coordinate rank obstruction
  \Cref{thm:class-product-minimal-obstruction-skeleton} below makes this
  last sentence sharp inside the chosen branch-compatible chart: after the
  scalar coordinate is fixed, every bad replacement of the determinant
  boundary data is exactly a non-zero vector in the \((q-1)\)-dimensional
  zero-mean class-function coordinate space, where \(q\) is the number of
  conjugacy classes of \(G\), and no scalar, exponent-level, or
  notation-level record can hide it there.
\end{proposition}

\begin{proof}
  The displayed formulas are the definitions of the primitive channel and
  fixed-label Euler transform together with
  \Cref{lem:class-logarithm-primitive-expansion}, which gives the Artin
  logarithm formula.  Since \(\phi\) is orthogonal to the trivial character
  and \(\eta<\lambda\), \Cref{thm:fixed-label-euler-transform} gives
  absolute convergence of \(F_\phi(\lambda^{-1})\) and the determinant
  expression \eqref{eq:three-logarithm-product-interface}.  In the expanded
  determinant expression the coefficient of the single scalar boundary term is
  \(\langle\phi,\mathbf1\rangle_G=0\); hence that singular logarithm is omitted
  before the Abel limit is formed.  The remaining boundary terms have
  non-trivial twist and are regular by the strict twisted gap, while every
  term with \(rm\ge2\) lies strictly inside the Perron disc.  This is exactly
  the branch-compatible continuation used here.
  Finally, \Cref{prop:logarithmic-channel-separation} proves that equality
  between \(F_\phi\) and \(\mathcal L_\phi\) is equivalent to
  \eqref{eq:channel-separation-obstruction}.  Main Theorem~III uses
  precisely the values \(F_\varrho(\lambda^{-1})\) in its non-trivial
  Fourier coordinates, so the stated product-constant roles follow.  The
  final roadmap sentence is the specialisation of
  \Cref{thm:class-product-minimal-obstruction-skeleton} to the boundary-error
  vector
  \(B_\varrho=\mathcal L_\varrho(\lambda^{-1})\), or more generally to any
  branch-compatible substitute \((B_\varrho)_{\varrho\ne\mathbf1}\).
\end{proof}

Before stating the product theorem, we separate the classical and finite-label
inputs.  The exponent \(|C|/|G|\), Chebotarev equidistribution, and the
existence/scale of Frobenius-class Mertens products belong to the dynamical
orbit-counting lineage of Parry--Pollicott, Noorani--Parry, Sharp, and
Mesliza--Noorani
\cite{ParryPollicott1986Chebotarev,NooraniParry1992ChebotarevShifts,Sharp1991Mertens,MeslizaNoorani1999FrobeniusMertens}.
Those results calibrate the first row of the theorem below; they are not the
novelty claimed here.  The new finite-state contribution begins only after
that leading term and the scalar normalisation \(C(A)\) have been separated:
the class vector \((K_C)_C\) is recovered from the fixed-label coordinates
\(F_\varrho(\lambda^{-1})\), with their Adams--M\"obius determinant formula,
and not from the periodic Artin logarithms
\(\mathcal L_\varrho(\lambda^{-1})\).  Recent finite-data, resonance, and
skew-product comparison papers are likewise only calibration points for the
finite SFT cocycle setting, not proof inputs to this product theorem.

\subsubsection{Main Theorem III: fixed-label Euler-product constants}
\label{subsubsec:main-iii-fixed-label-products}

The notation in the theorem is the fixed-label notation of
\Cref{subsec:fixed-label-vs-artin-logarithm}.  Here \(K_C\) is the positive
real product constant, \(C(A)\) is the scalar finite-part normalisation,
\(\chi_\varrho^\ast(C)=\overline{\chi_\varrho(C)}\) is the Hermitian class
coefficient, and the non-trivial coordinates are
\(F_\varrho(\lambda^{-1})\).  Determinant recovery of these same coordinates
is a later step, not an additional definition of the class product.

\begin{theorem}[Main Theorem III: Adams-corrected Frobenius-class
  Euler-product constants]
  \label{thm:class-euler-product-mertens}
  Assume the hypotheses of \Cref{thm:class-mertens-explicit}. For a
  conjugacy class \(C\subset G\), set
  \begin{equation}\label{eq:class-euler-product-def}
    P_C(N):=
    \prod_{\substack{\gamma\in\cP\,:\\ |\gamma|\le N\\
      \tau(\gamma)\in C}}
      \left(1-\lambda^{-|\gamma|}\right).
  \end{equation}
  Then
  \begin{equation}\label{eq:class-euler-product-asymptotic}
    P_C(N)=K_C\,N^{-|C|/|G|}\bigl(1+O(N^{-1})\bigr),
  \end{equation}
  The exponent \(|C|/|G|\) and the existence of a Frobenius-class product
  asymptotic are the inherited classical layer; the local contribution of this
  theorem is the Adams-corrected determinant expression for the constant
  \(K_C\) in the fixed finite-state cocycle model.
  Here the positive branch of the constant is
  \begin{equation}\label{eq:class-euler-product-k-mc}
    K_C=
    \exp\!\left(
      -\mathcal M_C(A,\tau;G)
      -\sum_{r\ge2}\frac{\Pi_{\mathbf1_C}(\lambda^{-r})}{r}
    \right).
  \end{equation}
  Equivalently, the same positive constant is given by the branch-compatible
  character formula
  \begin{equation}\label{eq:class-euler-product-k-f}
    K_C=
    \exp\!\left(
      -\frac{|C|}{|G|}
      \left(
        \gamma+\log C(A)
        +\sum_{\varrho\neq\mathbf1}
          \chi_\varrho^\ast(C)F_\varrho(\lambda^{-1})
      \right)
    \right).
  \end{equation}
  The summation in \eqref{eq:class-euler-product-k-mc} is absolutely
  convergent, and \(F_\varrho(\lambda^{-1})\) is the fixed-label Euler
  transform of \eqref{eq:fixed-label-euler-transform-def}, not the Artin
  logarithm \(\mathcal L_\varrho(\lambda^{-1})\); the local comparison,
  classical/new separation, and obstruction are recorded in
  \Cref{prop:three-logarithm-product-interface,prop:class-product-comparison}.
  The individual summands
  \(\chi_\varrho^\ast(C)F_\varrho(\lambda^{-1})\) need not be real when
  \(\varrho\) is not self-conjugate.  The bracket is real only after the
  conjugate irreducible pairs are combined; equivalently it is the value of a
  real class function under the Hermitian character expansion.  The
  exponential in \eqref{eq:class-euler-product-k-f} is therefore taken in the
  ordinary positive real branch.  The factor \(|C|/|G|\) in
  \eqref{eq:class-euler-product-k-f} multiplies the entire Hermitian
  class-indicator bracket, including the non-trivial fixed-label Fourier sum.
  The determinant version of the non-trivial values
  \(F_\varrho(\lambda^{-1})\) is exactly the branch-compatible one in
  \Cref{prop:fixed-label-determinant-perron}: the scalar Perron determinant
  \((r,m,\pi)=(1,1,\mathbf1)\) is deleted before evaluation, non-trivial
  \(rm=1\) determinants are taken by the Abel branch, and the \(rm\ge2\)
  determinant tail is absolutely convergent at \(z=\lambda^{-1}\).
\end{theorem}

\begin{proof}
  We spell out the short proof chain isolated in
  \Cref{prop:class-euler-product-main-chain-reduction}.  That proposition
  reduces the product asymptotic to the logarithmic form of Main Theorem~II,
  the explicit \(r\ge2\) product tail, and the fixed-label normalisation.
  It remains here only to verify those inputs with the branch and positivity
  conventions stated in the theorem.

  Taking logarithms in \eqref{eq:class-euler-product-def} gives
  \[
    \log P_C(N)
    =
    \sum_{\substack{n\le N}}
      p_{n,C}\log(1-\lambda^{-n})
    =
    -\sum_{n\le N}\frac{p_{n,C}}{\lambda^n}
      -
      \sum_{r\ge2}\frac{1}{r}
        \sum_{n\le N}p_{n,C}\lambda^{-rn}.
  \]
  The full \(r\ge2\) tail, its truncation error, and the passage from the
  additive class-Mertens estimate to the product logarithm are exactly
  \Cref{lem:class-euler-product-tail-bound}.  Substituting the logarithmic
  form \eqref{eq:class-mertens-log} from \Cref{cor:class-mertens-log-form}
  into \eqref{eq:class-product-log-bound-output} gives
  \[
    \log P_C(N)
    =
    -\frac{|C|}{|G|}\log N
    -\mathcal M_C(A,\tau;G)
    -\sum_{r\ge2}\frac{\Pi_{\mathbf1_C}(\lambda^{-r})}{r}
    +O(N^{-1}).
  \]
  Exponentiating gives \eqref{eq:class-euler-product-asymptotic} and
  \eqref{eq:class-euler-product-k-mc} by
  \Cref{lem:class-product-log-to-product-error}: take
  \(a=|C|/|G|\) and
  \[
    L_C=-\mathcal M_C(A,\tau;G)
    -\sum_{r\ge2}\frac{\Pi_{\mathbf1_C}(\lambda^{-r})}{r}.
  \]
  The product \(P_C(N)\) is positive, the logarithmic remainder in
  \Cref{lem:class-euler-product-tail-bound} is uniform over the finite set of
  conjugacy classes, and the cited lemma converts that additive
  \(O(N^{-1})\) remainder into the stated multiplicative
  \(1+O(N^{-1})\) remainder.  This also fixes the positive branch uniformly for
  all classes.

  We now record the normalization calculation that produces the single
  global factor \(|C|/|G|\) in the character formula.  The Hermitian
  class-indicator expansion is
  \[
    \mathbf1_C
    =
    \frac{|C|}{|G|}\mathbf1+
    \frac{|C|}{|G|}
      \sum_{\varrho\neq\mathbf1}\chi_\varrho^\ast(C)\chi_\varrho .
  \]
  Applying this identity first to the additive class-Mertens constant and
  then to the absolutely convergent Euler-product tail gives
  \[
  \begin{aligned}
    &\mathcal M_C(A,\tau;G)+
      \sum_{r\ge2}\frac{\Pi_{\mathbf1_C}(\lambda^{-r})}{r} \\
    &\quad=
      \frac{|C|}{|G|}
      \left(
        \mathcal M(A)+\sum_{r\ge2}\frac{\Pi_{\mathbf1}(\lambda^{-r})}{r}
      \right) \\
    &\qquad+
      \frac{|C|}{|G|}
      \sum_{\varrho\neq\mathbf1}\chi_\varrho^\ast(C)
      \left(
        \Pi_\varrho(\lambda^{-1})+
        \sum_{r\ge2}\frac{\Pi_\varrho(\lambda^{-r})}{r}
      \right).
  \end{aligned}
  \]
  The non-trivial parenthesis is exactly
  \(F_\varrho(\lambda^{-1})\), while the scalar parenthesis is
  \(\gamma+\log C(A)\) by the scalar finite-part product normalization.
  Thus the factor \(|C|/|G|\) comes from the class-indicator Fourier
  coefficient and multiplies the scalar term and every fixed-label Fourier
  contribution before the outer minus sign is applied.

  \Cref{lem:class-product-constant-equivalence} rewrites the exponent in
  \eqref{eq:class-euler-product-k-mc} as the exponent in
  \eqref{eq:class-euler-product-k-f}. This proves the equivalent
  character formula.  The same computation is isolated in
   \Cref{lem:class-product-hermitian-normalisation} to make the
   \(|C|/|G|\) normalisation explicit.  Finally,
   \Cref{prop:fixed-label-determinant-perron} supplies the determinant
   evaluation of each non-trivial \(F_\varrho(\lambda^{-1})\) with the local
   singular-term deletion, branch convention, and absolute tail convergence
   stated in the theorem.
\end{proof}

\begin{proposition}[Closed fixed-label coordinate equivalence]
  \label{prop:main-iii-fixed-label-proof-spine}
  Assume the hypotheses of \Cref{thm:class-mertens-explicit}.  The product
  constants in \Cref{thm:class-euler-product-mertens} are determined by the
  following three finite coordinate maps, and each arrow is reversible once the scalar
  normalisation \(\gamma+\log C(A)\) is fixed:
  \begin{enumerate}[label=\textup{(\roman*)},leftmargin=*]
    \item the primitive-label logarithm for a non-trivial character is the
      fixed-label Euler transform
      \[
        F_\varrho(\lambda^{-1})
        =\sum_{r\ge1}\frac{\Pi_\varrho(\lambda^{-r})}{r};
      \]
    \item the same number is the branch-normalised determinant boundary value
      \[
        \sum_{r,m\ge1}\frac{\mu(m)}{rm}
        \sum_{\substack{\pi\in\Irr(G)\\(r,m,\pi)\ne(1,1,\mathbf1)}}
          a_{\varrho,\pi}^{(m)}
          \Log\det(I-\lambda^{-rm}A_\pi)^{-1};
      \]
    \item the class-product constants are the positive real exponentials
      \[
        K_C=
        \exp\!\left(
        -\frac{|C|}{|G|}
        \left(
          \gamma+\log C(A)+
          \sum_{\varrho\ne\mathbf1}\chi_\varrho^\ast(C)
            F_\varrho(\lambda^{-1})
        \right)
        \right).
      \]
  \end{enumerate}
  Consequently the finite twisted determinant polynomials
  \(\det(I-zA_\pi)\), the finite character table, the Adams coefficients, and
  the scalar determinant constant \(C(A)\) form a complete fixed-coordinate
  record of \((K_C)_{C\subset G}\).  Conversely, the positive constants
  \((K_C)_C\), together with \(C(A)\), recover all non-trivial values
  \(F_\varrho(\lambda^{-1})\).  In this sense the fixed-label transform,
  determinant Perron formula, and \(K_C\)-formula form a single closed
  finite-dimensional coordinate system rather than three independent
  computational layers.
\end{proposition}

\begin{proof}
  Step \textup{(i)} is the second identity in
  \eqref{eq:fixed-label-euler-transform-vs-artin}, evaluated at
  \(z=\lambda^{-1}\); the value exists for non-trivial \(\varrho\) by the
  Perron-boundary assertion of \Cref{thm:fixed-label-euler-transform}.  Step
  \textup{(ii)} is exactly \Cref{prop:fixed-label-determinant-perron}: the
  scalar Perron triple \((1,1,\mathbf1)\) is excluded before evaluation,
  non-trivial \(rm=1\) terms use the Abel branch, and all \(rm\ge2\) terms are
  ordinary determinant-germ values inside the Perron disc.  Step
  \textup{(iii)} is formula \eqref{eq:class-euler-product-k-f} in
  \Cref{thm:class-euler-product-mertens}.  Substituting \textup{(ii)} into
  \textup{(iii)} gives the advertised finite determinant record, since
  \(G\) has finitely many irreducible characters and each twisted matrix
  \(A_\pi\) is finite dimensional.

  It remains only to prove the reversibility assertion.  With
  \(\alpha_A:=\gamma+\log C(A)\), define
  \[
    H_K(C):=-\frac{|G|}{|C|}\log K_C-\alpha_A .
  \]
  Formula \textup{(iii)} gives
  \[
    H_K(C)=
    \sum_{\varrho\ne\mathbf1}\chi_\varrho^\ast(C)
      F_\varrho(\lambda^{-1}).
  \]
  Character orthogonality on class functions therefore recovers, for every
  non-trivial \(\pi\),
  \[
    F_\pi(\lambda^{-1})
    =\sum_{C\subset G}\frac{|C|}{|G|}H_K(C)\chi_\pi(C).
  \]
  Thus the constants and the fixed scalar normalisation determine the
  non-trivial fixed-label vector, while the preceding substitutions determine
  the constants from that vector.  This proves the claimed closed coordinate equivalence.
\end{proof}

The exponent \(|C|/|G|\) and the existence of a Frobenius-class product
constant are the classical part of \Cref{thm:class-euler-product-mertens},
parallel to the dynamical Mertens and Chebotarev product theorems of Sharp,
Parry--Pollicott, Noorani--Parry, and especially Mesliza--Noorani's
Frobenius-class Mertens product theorem
\cite{Sharp1991Mertens,ParryPollicott1986Chebotarev,NooraniParry1992ChebotarevShifts,MeslizaNoorani1999FrobeniusMertens}.  The new assertion
used in the finite-state main chain is the branch-compatible
Adams--M\"obius determinant formula for that constant, and hence the
obstruction to replacing the fixed-label Euler transform by an ordinary
Artin logarithm.




\begin{corollary}[Adjacent finite-determinant evaluation of the product constants]
  \label{cor:class-euler-product-determinant-evaluation}
  Assume the hypotheses of \Cref{thm:class-mertens-explicit}.  For each
  non-trivial \(\varrho\in\Irr(G)\), put
  \begin{equation}\label{eq:class-euler-product-determinant-F}
    \mathfrak F_\varrho:=
    \sum_{r,m\ge1}\frac{\mu(m)}{rm}
      \sum_{\substack{\pi\in\Irr(G)\\(r,m,\pi)\neq(1,1,\mathbf1)}}
        a_{\varrho,\pi}^{(m)}
        \log\det(I-\lambda^{-rm}A_\pi)^{-1},
  \end{equation}
  with the same Abel-path branch convention as
  \Cref{prop:fixed-label-determinant-perron}: the only boundary
  evaluations are the non-trivial blocks at \(z=\lambda^{-1}\), and the
  scalar Perron logarithm \(\log\det(I-\lambda^{-1}A_{\mathbf1})^{-1}\)
  is excluded.  Then the product constants in
  \Cref{thm:class-euler-product-mertens} are recovered directly from the
  finite determinant family by
  \begin{equation}\label{eq:class-euler-product-k-det-adjacent}
    K_C=
    \exp\!\left(
      -\frac{|C|}{|G|}
      \left(
        \gamma+\log C(A)
        +\sum_{\varrho\neq\mathbf1}
          \chi_\varrho^\ast(C)\mathfrak F_\varrho
      \right)
    \right).
  \end{equation}
  Thus no cocycle-level orbit enumeration beyond the finite matrices
  \(A_\pi\), the character table, and the Adams coefficients is needed to
  compute \((K_C)_{C\subset G}\).
\end{corollary}

\begin{proof}
  By \Cref{prop:fixed-label-determinant-perron}, the determinant expression
  \eqref{eq:class-euler-product-determinant-F} is exactly the Perron-point
  fixed-label value \(F_\varrho(\lambda^{-1})\) for every non-trivial
  irreducible \(\varrho\).  Substituting these equalities into the
  branch-compatible character formula
  \eqref{eq:class-euler-product-k-f} of
  \Cref{thm:class-euler-product-mertens} gives
  \eqref{eq:class-euler-product-k-det-adjacent}.  The branch convention and
  the exclusion of the scalar Perron determinant are precisely those of
  \Cref{prop:fixed-label-determinant-perron}; hence every determinant in
  the displayed expression is either a non-trivial Abel-boundary value made
  regular by \(\eta<\lambda\), or an ordinary interior value with
  \(rm\ge2\).  The last assertion follows because the right-hand side uses
  only these determinants, the finite character table, and the finite Adams
  coefficient tensors.
\end{proof}


% END INPUT sec_chebotarev_class_mertens_products.tex
% BEGIN INPUT sec_chebotarev_class_mertens_rigidity_core.tex
\subsubsection{Finite-determinant rigidity}

The theorem in this subsubsection is the determinant-coordinate form of Main
Theorem~III.  The supporting corollaries that follow are finite-dimensional
linear algebra in this fixed coordinate chart; they are recorded for
normalisation and reproducibility, not as additional main theorems or as new
hypotheses for the class-product asymptotic.

\begin{theorem}[Determinant-coordinate rigidity for Adams-corrected constants]
  \label{thm:class-product-determinant-rigidity}
  Assume the hypotheses of \Cref{thm:class-mertens-explicit}.  The input
  data for this rigidity statement are exactly the scalar Perron
  normalisation \(C(A)\), the finite character table of \(G\), the finite
  Adams coefficient tensors \(a_{\varrho,\pi}^{(m)}\), and the twisted
  determinant polynomials \(\det(I-zA_\pi)\) evaluated with the Abel branch
  convention of \Cref{prop:fixed-label-determinant-perron}.  For every
  non-trivial \(\varrho\in\Irr(G)\), define the determinant coordinate
  \begin{equation}\label{eq:class-product-rigidity-determinant-coordinate}
    \mathfrak D_\varrho:=
    \sum_{r,m\ge1}\frac{\mu(m)}{rm}
      \sum_{\substack{\pi\in\Irr(G)\\(r,m,\pi)\neq(1,1,\mathbf1)}}
        a_{\varrho,\pi}^{(m)}
        \log\det(I-\lambda^{-rm}A_\pi)^{-1}.
  \end{equation}
  The logarithm is the normalized germ at the origin; non-trivial
  \(rm=1\) terms are Abel-boundary values, all \(rm\ge2\) terms are ordinary
  interior-disc values, and the scalar Perron triple
  \((1,1,\mathbf1)\) is deleted before evaluation.  Then the determinant
  coordinates are exactly the fixed-label Perron values,
  \begin{equation}\label{eq:class-product-rigidity-d-equals-f}
    \mathfrak D_\varrho=F_\varrho(\lambda^{-1})
    \qquad(\varrho\ne\mathbf1),
  \end{equation}
  and the series in
  \eqref{eq:class-product-rigidity-determinant-coordinate} is absolutely
  convergent after the finitely many regular \(rm=1\) non-trivial boundary
  terms have been separated.  Put
  \(\alpha_A:=\gamma+\log C(A)\), and let \((K_C)_{C\subset G}\) be the
  Frobenius-class Euler-product constants of
  \Cref{thm:class-euler-product-mertens}.  Define the normalised logarithmic
  product vector
  \begin{equation}\label{eq:class-product-log-vector}
    H_K(C):=-\frac{|G|}{|C|}\log K_C-\alpha_A .
  \end{equation}
  Then
  \begin{equation}\label{eq:class-product-normalisation-constraint}
    \prod_{C\subset G}K_C=\exp(-\alpha_A),
    \qquad
    \sum_{C\subset G}\frac{|C|}{|G|}H_K(C)=0,
  \end{equation}
  and the non-trivial Peter--Weyl coordinates of this zero-mean vector are
  exactly the determinant coordinates:
  \begin{equation}\label{eq:class-product-fourier-coordinates}
    H_K(C)=
    \sum_{\varrho\ne\mathbf1}
      \chi_\varrho^\ast(C)\mathfrak D_\varrho,
    \qquad
    \sum_{C\subset G}\frac{|C|}{|G|}H_K(C)\chi_\pi(C)
    =\mathfrak D_\pi.
  \end{equation}
  Equivalently, the constants themselves are recovered by the finite
  determinant-coordinate formula
  \begin{equation}\label{eq:class-product-rigidity-product-formula}
    K_C=
    \exp\!\left(
      -\frac{|C|}{|G|}
      \left(
        \alpha_A+
        \sum_{\varrho\ne\mathbf1}
          \chi_\varrho^\ast(C)\mathfrak D_\varrho
      \right)
    \right),
  \end{equation}
  In particular,
  \begin{equation}\label{eq:class-product-parseval}
    \sum_{C\subset G}\frac{|C|}{|G|}|H_K(C)|^2
    =
    \sum_{\varrho\ne\mathbf1}|\mathfrak D_\varrho|^2.
  \end{equation}
  Thus \((K_C)_{C\subset G}\) is determined by precisely the determinant
  coordinates \((\mathfrak D_\varrho)_{\varrho\ne\mathbf1}\) and the scalar
  normalisation \(C(A)\).  Conversely, the positive constant vector
  \((K_C)_{C\subset G}\), together with \(C(A)\), recovers every
  non-trivial determinant coordinate by the inverse formula in
  \eqref{eq:class-product-fourier-coordinates}.  Consequently two finite
  extensions with the same scalar normalisation, character table, Adams
  tensors, branch tags, and twisted determinant polynomials have the same
  Adams-corrected Frobenius-class product constants.  This is a rigidity
  statement only in those fixed coordinates and branch/scalar normalisation:
  it does not assert invariance under arbitrary presentations, covers, folds,
  gauge changes, or conjugacies of the cocycle \(\tau\) or of the skew product.
\end{theorem}

\begin{proof}
  The determinant identity
  \eqref{eq:class-product-rigidity-d-equals-f}, including the local deletion
  of \((1,1,\mathbf1)\), the Abel-boundary convention for non-trivial
  \(rm=1\) terms, the ordinary interior evaluation for \(rm\ge2\), and the
  absolute convergence of the remaining tail, is exactly
  \Cref{prop:fixed-label-determinant-perron}.  Hence the numbers
  \(\mathfrak D_\varrho\) are fixed by the stated finite determinant data.

  The leading product exponent and the existence of the positive constants
  are exactly \Cref{thm:class-euler-product-mertens}; the historical
  calibration is the one stated in \Cref{prop:class-product-comparison}.
  The character formula \eqref{eq:class-euler-product-k-f} expresses
  \((K_C)_C\) in terms of \(C(A)\) and the non-trivial values
  \(F_\varrho(\lambda^{-1})\).  Replacing those values by
  \(\mathfrak D_\varrho\) using
  \eqref{eq:class-product-rigidity-d-equals-f}, and then taking logarithms in
  the positive real branch fixed by \Cref{thm:class-euler-product-mertens},
  gives the first identity in \eqref{eq:class-product-fourier-coordinates}.
  Averaging this identity over conjugacy classes with weights \(|C|/|G|\)
  gives the second identity in
  \eqref{eq:class-product-normalisation-constraint}, because every
  non-trivial irreducible character is orthogonal to the trivial character.
  Multiplying the definition of \(H_K\) by \(|C|/|G|\) and summing over
  classes then gives the product constraint in
  \eqref{eq:class-product-normalisation-constraint}.  The inverse Fourier
  identity in \eqref{eq:class-product-fourier-coordinates} and the Parseval
  identity \eqref{eq:class-product-parseval} are the standard Hermitian
  character orthogonality relations for the zero-mean class-function space.

  Solving the definition \eqref{eq:class-product-log-vector} for \(K_C\) and
  substituting the first identity in
  \eqref{eq:class-product-fourier-coordinates} gives
  \eqref{eq:class-product-rigidity-product-formula}.  This proves that the
  product vector is determined by the listed scalar normalisation, character
  table, Adams tensors, branch tags, and twisted determinant polynomials.
  Conversely, if the positive constants and \(C(A)\) are known, then
  \eqref{eq:class-product-log-vector} constructs \(H_K\), and the inverse
  Hermitian Fourier formula in \eqref{eq:class-product-fourier-coordinates}
  recovers each \(\mathfrak D_\pi\).  The final equal-data assertion follows
  by applying the same displayed formulae to the two systems.  No step uses
  any presentation, cover, fold, gauge, or cocycle-conjugacy invariant beyond
  these fixed coordinates, which proves the asserted scope of the rigidity
  statement.
\end{proof}

\begin{corollary}[Branch-compatible hyperplane classification and fixed-coordinate stability]
  \label{cor:class-product-hyperplane-stability}
  \label{cor:class-product-stability-scalar-independence}
  Assume the hypotheses of \Cref{thm:class-mertens-explicit}.  All
  determinant coordinates below are the branch-compatible finite coordinates
  of \Cref{prop:fixed-label-determinant-perron}: the scalar Perron
  coordinate is fixed by \(C(A)\), non-trivial \(rm=1\) determinant
  logarithms are evaluated by the stated Abel continuation, and every
  \(rm\ge2\) determinant logarithm is taken from the same normalized germ
  inside the Perron disc.  After the scalar value \(\alpha_A\) is fixed, the
  identities of \Cref{thm:class-product-determinant-rigidity} classify all
  positive class-product vectors satisfying the forced scalar product
  constraint in this coordinate chart.  Namely let
  \[
    \mathscr K_{\alpha_A}:=
    \left\{(L_C)_C:\ L_C>0,
      \prod_{C\subset G}L_C=\exp(-\alpha_A)\right\}
  \]
  and let
  \[
    \mathcal H_G:=
    \left\{H\text{ real class function}\;:\;
      \sum_C\frac{|C|}{|G|}H(C)=0\right\}.
  \]
  Then
  \begin{equation}\label{eq:class-product-coordinate-map}
    H\longmapsto
    \left(
      \exp\!\left(-\frac{|C|}{|G|}
        \bigl(\alpha_A+H(C)\bigr)\right)
    \right)_{C\subset G}
  \end{equation}
  is a bijection from \(\mathcal H_G\) onto
  \(\mathscr K_{\alpha_A}\), with inverse
  \begin{equation}\label{eq:class-product-coordinate-map-inverse}
    (L_C)_C\longmapsto H_L,
    \qquad
    H_L(C):=-\frac{|G|}{|C|}\log L_C-\alpha_A .
  \end{equation}
  The non-trivial irreducible characters give an isometric coordinate system
  on this hyperplane:
  \begin{equation}\label{eq:class-product-hyperplane-expansion}
    H(C)=
    \sum_{\varrho\ne\mathbf1}
      \chi_\varrho^\ast(C)
      \sum_D\frac{|D|}{|G|}H(D)\chi_\varrho(D),
  \end{equation}
  and
  \begin{equation}\label{eq:class-product-stability-isometry}
    \sum_{C\subset G}\frac{|C|}{|G|}|H(C)-H'(C)|^2
    =
    \sum_{\varrho\ne\mathbf1}
      \left|
        \sum_C\frac{|C|}{|G|}(H(C)-H'(C))\chi_\varrho(C)
      \right|^2 .
  \end{equation}
  Hence two admissible branch-compatible finite determinant-coordinate data
  sets with the same scalar normalisation and the same Abel/ordinary branch
  convention give the same Frobenius-class product constants if and only if
  their non-trivial fixed-label Perron vectors agree.
  For cocycles \(\tau,\tau'\) over the same base matrix, this equivalence is
  the exact stability identity
  \begin{equation}\label{eq:class-product-stability-forward}
    -\frac{|G|}{|C|}\log K_C^\tau
    +\frac{|G|}{|C|}\log K_C^{\tau'}
    =
    \sum_{\varrho\neq\mathbf1}
      \chi_\varrho^\ast(C)
      \bigl(F_\varrho^\tau(\lambda^{-1})
            -F_\varrho^{\tau'}(\lambda^{-1})\bigr),
  \end{equation}
  together with Parseval for the two sides.  In particular, the zero-mean
  class-product factor has real dimension
  \(\#\{\text{conjugacy classes of }G\}-1\).  Any scalar fixed-matrix
  record, such as \(C(A)\), \(\mathcal M(A)\), \(\FP(A)\), or a fixed-matrix
  cyclic-lift record, has rank \(0\) on this factor and cannot replace the
  non-trivial Peter--Weyl coordinates within the branch-compatible finite
  determinant model.
\end{corollary}

\begin{proof}
  If \(H\in\mathcal H_G\), then the product of the entries in
  \eqref{eq:class-product-coordinate-map} is \(\exp(-\alpha_A)\), because the
  weighted mean of \(H\) is zero.  Conversely, if
  \((L_C)_C\in\mathscr K_{\alpha_A}\), the positive logarithm in
  \eqref{eq:class-product-coordinate-map-inverse} has weighted mean zero
  exactly because \(\prod_C L_C=\exp(-\alpha_A)\).  These two assignments are
  inverse by direct substitution, proving the bijection.

  Fourier inversion on the orthogonal complement of the trivial character gives
  \eqref{eq:class-product-hyperplane-expansion}; applying Parseval to
  \(H-H'\) gives \eqref{eq:class-product-stability-isometry}.  Substituting
  the dynamical vectors \(H_{K^\tau}\) and \(H_{K^{\tau'}}\), and then using
  \eqref{eq:class-product-fourier-coordinates} for each cocycle, gives
  \eqref{eq:class-product-stability-forward} and its Parseval identity; both
  sides use the same branch-compatible determinant-coordinate convention
  fixed in \Cref{prop:fixed-label-determinant-perron}.
  Hence equality of the product constants is equivalent to equality of all
  non-trivial fixed-label Perron coordinates.  The dimension and rank
  assertions follow because \(\mathcal H_G\) is the codimension-one subspace
  orthogonal to the trivial character, while every listed scalar fixed-matrix
  record is constant as the non-trivial Peter--Weyl coordinates vary with the
  scalar normalisation fixed.
\end{proof}

\begin{corollary}[Assembled finite-determinant product formula and truncation]
  \label{cor:assembled-adams-frobenius-product-formula}
  \label{cor:assembled-adams-frobenius-product-package}
  Assume the hypotheses of \Cref{thm:class-mertens-explicit}.  The determinant
  recovery formula \eqref{eq:class-product-rigidity-product-formula} and the
  finite determinant evaluation of
  \Cref{prop:fixed-label-determinant-perron} assemble Main Theorem~III into a
  finite twisted-determinant computation of \((K_C)_{C\subset G}\).  More
  concretely, for
  \(M\ge1\) set
  \begin{equation}\label{eq:class-product-finite-truncation-coordinate}
    F^{(M)}_\varrho
    :=
    \sum_{\substack{r,m\ge1\\rm\le M}}
      \frac{\mu(m)}{rm}
      \sum_{\substack{\pi\in\Irr(G)\\(r,m,\pi)\neq(1,1,\mathbf1)}}
        a_{\varrho,\pi}^{(m)}
        \log\det(I-\lambda^{-rm}A_\pi)^{-1}
  \end{equation}
  and
  \begin{equation}\label{eq:class-product-finite-truncation-constant}
    K_C^{(M)}
    :=
    \exp\!\left(
      -\frac{|C|}{|G|}
      \left(
        \alpha_A+
        \sum_{\varrho\ne\mathbf1}\chi_\varrho^\ast(C)F^{(M)}_\varrho
      \right)
    \right).
  \end{equation}
  Then there are constants \(B_C\), depending only on the finite data
  \((A,\tau,G,C)\), such that
  \begin{equation}\label{eq:class-product-finite-truncation-error}
    \abs{\log K_C-\log K_C^{(M)}}
    \le
    B_C
    \sum_{k>M}\frac{\mathrm d(k)}{k}\lambda^{1-k},
  \end{equation}
  where \(\mathrm d(k)\) is the divisor function.  Thus the determinant
  estimates supply an error-controlled exhaustion of the Adams-corrected
  constants by finitely many evaluations of the finitely many twisted
  determinant polynomials.
\end{corollary}

\begin{proof}
  Fix a non-trivial irreducible representation \(\varrho\).  The finite sum
  in \eqref{eq:class-product-finite-truncation-coordinate} is the truncation
  at \(rm\le M\) of the determinant series in
  \eqref{eq:fixed-label-determinant-perron}, with the same deletion of
  \((r,m,\pi)=(1,1,\mathbf1)\) and with the same branch convention for every
  retained logarithm.  Hence \Cref{prop:fixed-label-determinant-perron}
  gives
  \[
    F_\varrho(\lambda^{-1})
    =\lim_{M\to\infty}F^{(M)}_\varrho .
  \]
  More precisely, the double-series majorant used in
  \Cref{prop:fixed-label-perron-finite-algorithm} gives a constant
  \(C_\varrho\), depending only on \((A,\tau,G,\varrho)\), such that
  \begin{equation}\label{eq:assembled-truncation-coordinate-tail}
    \abs{F_\varrho(\lambda^{-1})-F^{(M)}_\varrho}
    \le
    C_\varrho
    \sum_{k>M}\frac{\mathrm d(k)}{k}\lambda^{1-k}.
  \end{equation}
  Indeed, grouping the omitted pairs \((r,m)\) by \(k=rm\) after the
  finitely many regular \(k=1\) boundary terms leaves at most
  \(\mathrm d(k)\) pairs at level \(k\), and the fixed-label Perron majorant
  bounds the whole inner character sum and determinant logarithm by a
  constant multiple of \(\lambda^{1-k}/k\).  The same estimate includes
  scalar terms with \(k\ge2\), since those are ordinary interior-disc
  determinant values.

  By \Cref{thm:class-product-determinant-rigidity}, the determinant
  coordinate \(\mathfrak D_\varrho\) equals
  \(F_\varrho(\lambda^{-1})\), and the exact constant is
  \[
    \log K_C
    =-
    \frac{|C|}{|G|}
      \left(
        \alpha_A+
        \sum_{\varrho\ne\mathbf1}
          \chi_\varrho^\ast(C)F_\varrho(\lambda^{-1})
      \right).
  \]
  The definition \eqref{eq:class-product-finite-truncation-constant} is the
  identical finite Fourier reconstruction with
  \(F_\varrho(\lambda^{-1})\) replaced by \(F^{(M)}_\varrho\).  Subtracting
  the two logarithms therefore gives the exact identity
  \[
    \log K_C-\log K_C^{(M)}
    =-
    \frac{|C|}{|G|}
      \sum_{\varrho\ne\mathbf1}
        \chi_\varrho^\ast(C)
        \bigl(F_\varrho(\lambda^{-1})-F^{(M)}_\varrho\bigr).
  \]
  Since \(G\) has finitely many irreducible representations, combining this
  identity with \eqref{eq:assembled-truncation-coordinate-tail} yields
  \eqref{eq:class-product-finite-truncation-error} with, for example,
  \[
    B_C:=
    \frac{|C|}{|G|}
    \sum_{\varrho\ne\mathbf1}
      \abs{\chi_\varrho(C)}C_\varrho .
  \]
  The right-hand side tends to zero because \(\lambda>1\) and the divisor
  function has subexponential growth.  Thus the finite sums exhaust the
  determinant coordinates and, through the finite Hermitian class transform,
  the Adams-corrected Frobenius-class product constants.
\end{proof}
% END INPUT sec_chebotarev_class_mertens_rigidity_core.tex
% BEGIN INPUT sec_chebotarev_class_mertens_rigidity_consequences.tex

\begin{proof}
  \Cref{thm:class-product-determinant-rigidity} expresses each \(K_C\) by
  \eqref{eq:class-product-rigidity-product-formula}.  By
  \Cref{prop:fixed-label-determinant-perron}, every coordinate
  \(F_\varrho(\lambda^{-1})\) appearing there is the determinant series over
  the finitely many twisted matrices \(A_\pi\), with the branch convention and
  scalar exclusion already fixed in that proposition.  Substituting those
  determinant values gives the assembled finite-determinant product formula.

  The same determinant formula, truncated to the indices \(rm\le M\), gives
  \eqref{eq:class-product-finite-truncation-coordinate}.  By
  \Cref{lem:fixed-label-perron-double-majorant}, for each non-trivial
  \(\varrho\) there is a constant \(C_\varrho\) such that the omitted part of
  the determinant series is bounded by
  \[
    C_\varrho
    \sum_{\substack{r,m\ge1\\rm>M}}\frac{\lambda^{1-rm}}{rm}
    =
    C_\varrho
    \sum_{k>M}\frac{\mathrm d(k)}{k}\lambda^{1-k}.
  \]
  Multiplying these bounds by the finite coefficients
  \((|C|/|G|)\chi_\varrho^\ast(C)\) in the logarithm of
  \eqref{eq:class-euler-product-k-f}, and summing over the finitely many
  non-trivial irreducible representations, proves
  \eqref{eq:class-product-finite-truncation-error} with
  \[
    B_C=\frac{|C|}{|G|}
    \sum_{\varrho\ne\mathbf1}
      \abs{\chi_\varrho^\ast(C)}C_\varrho .
  \]
  Hence each finite truncation uses only finitely many evaluations of the
  fixed finite family of determinant polynomials, and the displayed
  exponentially summable tail gives the claimed finite-determinant
  computability.
\end{proof}

\subsubsection{Local equivalence and observation of product coordinates}
\label{subsubsec:class-product-local-observation}

The next consequences are local observations on the product-coordinate vector
already recovered in Main Theorem~III.  They identify exactly which
branch-compatible fixed-label coordinates are determined by the positive class
constants and record the finite-rank obstruction when too few class coordinates
are observed.  These are finite Fourier statements on conjugacy classes, not
classification statements for extensions or cocycles.

\begin{proposition}[Branch-compatible fixed-coordinate equivalence for product constants]
  \label{prop:class-product-fixed-coordinate-equivalence}
  Under the hypotheses and notation of
  \Cref{thm:class-product-determinant-rigidity}, consider two finite
  determinant-coordinate data sets with the same finite group \(G\), the same
  character table and Adams tensors, the same scalar normalisation \(C(A)\),
  and the same normalized Abel/ordinary branch convention of
  \Cref{prop:fixed-label-determinant-perron}.  Thus equality of a coordinate
  always means equality of the corresponding branch-compatible logarithmic
  determinant value, not merely equality of an exponentiated determinant.  If
  their non-trivial fixed-label Perron vectors agree,
  \[
    F_\varrho(\lambda^{-1})=F'_\varrho(\lambda^{-1})
    \qquad(\varrho\ne\mathbf1),
  \]
  then they give the same Frobenius-class product constants \(K_C=K'_C\) for
  every conjugacy class \(C\).  Conversely, if the constants agree for every
  \(C\) and the same scalar normalisation \(C(A)\) is used, then the
  non-trivial fixed-label Perron vectors agree.  Thus the recoverable invariant
  at Main Theorem~III is exactly the fixed branch-normalised determinant
  coordinate vector modulo this finite Fourier equivalence, not a
  presentation-, cover-, gauge-, branch-free, or fold-invariant classification
  of the underlying cocycle.
\end{proposition}

\begin{proof}
  With \(C(A)\) fixed, equality of the fixed-label Perron vectors makes the
  right-hand side of \eqref{eq:class-product-fourier-coordinates} identical for
  every class.  Substitution in the exponential formula for \(K_C\) in
  \Cref{thm:class-product-determinant-rigidity} therefore gives
  \(K_C=K'_C\) for all \(C\).  Conversely, equality of the positive constants
  and equality of \(C(A)\) imply equality of the normalised logarithmic class
  functions \(H_K\) defined in \eqref{eq:class-product-log-vector}.  Applying
  the inverse character formula in
  \eqref{eq:class-product-fourier-coordinates} gives
  \[
    F_\pi(\lambda^{-1})
    =\sum_{C\subset G}\frac{|C|}{|G|}H_K(C)\chi_\pi(C)
    =\sum_{C\subset G}\frac{|C|}{|G|}H_{K'}(C)\chi_\pi(C)
    =F'_\pi(\lambda^{-1})
  \]
  for every non-trivial \(\pi\).  Every step uses only the listed finite
  character table and the fixed scalar/branch normalisation.  In particular,
  the implication is false as a branch-free statement: changing logarithmic
  lifts can change the Perron vector even when exponentiated determinant
  values are unchanged, exactly as in the fixed-matrix branch ambiguity
  recorded in \Cref{lem:cyclic-lift-branch-ambiguity}.  Thus no conclusion
  about alternative presentations, gauges, folds, branches, or extensions is
  introduced.
\end{proof}




\begin{corollary}[Fixed-coordinate rank obstruction subordinate to Main Theorem III]
  \label{thm:class-product-minimal-obstruction-skeleton}
  \label{cor:class-product-minimal-obstruction-skeleton}
  Fix a finite group \(G\) and a real scalar normalisation \(\alpha\).  This
  is a rank obstruction inside the scalar-normalised, branch-compatible
  determinant-coordinate chart used in Main Theorem~III and in the
  Adams-obstruction theorem.  It compares only already-proved
  scalar-normalised product constants after the finite determinant
  coordinates and their branches have been fixed.  It is not a minimality
  statement for arbitrary observation systems, for all cocycles, for
  unnormalised determinant data, for branch-free determinant values, or for
  any theorem chain independent of Main Theorem~III.  Let
  \(q\) be the number of conjugacy classes of \(G\), and let
  \[
    \mathcal H_G:=
    \left\{H:G\to\RR\;\middle|\;
      H\text{ is a class function and }
      \sum_{C\subset G}\frac{|C|}{|G|}H(C)=0
    \right\}.
  \]
  For \(H\in\mathcal H_G\), put
  \begin{equation}\label{eq:minimal-skeleton-product-map}
    \mathcal K_\alpha(H)_C:=
    \exp\!\left(-\frac{|C|}{|G|}\bigl(\alpha+H(C)\bigr)\right),
    \qquad C\subset G .
  \end{equation}
  Then \(\mathcal K_\alpha\) is an injective parametrisation of the positive
  scalar-normalised product hyperplane
  \[
    \left\{(L_C)_C:\ L_C>0,
      \prod_{C\subset G}L_C=e^{-\alpha}\right\}.
  \]
  Consequently, if \(R:\mathcal H_G\to W\) is a real linear record, formed in
  these fixed determinant coordinates after \(\alpha\) has been fixed, with
  the property that
  \[
    R(H)=R(H')\quad\Longrightarrow\quad
    \mathcal K_\alpha(H)=\mathcal K_\alpha(H')
    \qquad(H,H'\in\mathcal H_G),
  \]
  then \(R\) is injective and
  \begin{equation}\label{eq:minimal-skeleton-rank-bound}
    \operatorname{rank}R\ge \dim\mathcal H_G=q-1 .
  \end{equation}
  Equivalently, every non-injective such fixed-coordinate record has a
  concrete bad pair: there are \(H,H'\in\mathcal H_G\) with
  \(R(H)=R(H')\) but
  \(\mathcal K_\alpha(H)\ne\mathcal K_\alpha(H')\).  The non-trivial
  Peter--Weyl coordinate map
  \begin{equation}\label{eq:minimal-skeleton-peter-weyl-record}
    H\longmapsto
    \left(
      \sum_{C\subset G}\frac{|C|}{|G|}H(C)\chi_\varrho(C)
    \right)_{\varrho\ne\mathbf1}
  \end{equation}
  is an isometric record of exactly this fixed-coordinate rank, after
  imposing the usual contragredient conjugation relations for real class
  functions.

  Applied with \(\alpha=\gamma+\log C(A)\) and
  \(H(C)=H_K(C)\) from \eqref{eq:class-product-log-vector}, the actual
  Adams-corrected Frobenius-class product constants therefore have the
  following extremal interpretation and no broader one: the non-trivial
  fixed-label determinant boundary coordinates in
  \Cref{thm:class-product-determinant-rigidity} supply a rank-sharp complete
  record of the scalar-normalised product constants in this fixed coordinate
  chart.  In particular, for
  any branch-compatible substitute boundary
  \((B_\varrho)_{\varrho\ne\mathbf1}\), the normalised classwise error
  \[
    E_B(C):=
    \sum_{\varrho\ne\mathbf1}\chi_\varrho^\ast(C)
      \bigl(F_\varrho(\lambda^{-1})-B_\varrho\bigr)
  \]
  satisfies the exact Parseval identity
  \begin{equation}\label{eq:minimal-skeleton-error-bound}
    \sum_{C\subset G}\frac{|C|}{|G|}|E_B(C)|^2
    =
    \sum_{\varrho\ne\mathbf1}
      |F_\varrho(\lambda^{-1})-B_\varrho|^2.
  \end{equation}
  Hence a non-zero non-trivial boundary error cannot be hidden, inside these
  scalar-normalised product constants, by changing notation, by keeping the
  classical exponent \(|C|/|G|\), or by using only scalar fixed-matrix data.
  This statement is a corollary-level structural refinement of Main
  Theorem~III: it assumes the scalar-normalised product constants and their
  Adams-corrected Peter--Weyl determinant boundary have already been
  identified, and it adds no new product asymptotic, determinant evaluation,
  observation-system minimality claim, or Perron hypothesis to Main
  Theorem~III.
\end{corollary}

\begin{proof}
  First observe that \(\dim\mathcal H_G=q-1\), because the real vector space
  of class functions has dimension \(q\) and the displayed weighted-mean
  condition is one non-zero real linear equation.  The image of
  \(\mathcal K_\alpha\) lies in the stated product hyperplane: multiplying
  \eqref{eq:minimal-skeleton-product-map} over all conjugacy classes gives
  \[
    \prod_C\mathcal K_\alpha(H)_C
    =
    \exp\!\left(-\sum_C\frac{|C|}{|G|}(\alpha+H(C))\right)
    =e^{-\alpha}.
  \]
  Conversely, if \((L_C)_C\) is a positive vector in that hyperplane, then
  \[
    H_L(C):=-\frac{|G|}{|C|}\log L_C-\alpha
  \]
  has weighted mean zero, exactly because \(\prod_C L_C=e^{-\alpha}\), and
  \(\mathcal K_\alpha(H_L)=L\).  Thus \(\mathcal K_\alpha\) parametrises the
  whole hyperplane.  It is injective because equality of the positive
  vectors \(\mathcal K_\alpha(H)=\mathcal K_\alpha(H')\), followed by the
  ordinary real logarithm in each coordinate, gives
  \((|C|/|G|)(H(C)-H'(C))=0\) for every class \(C\), hence \(H=H'\).

  Now let \(R:\mathcal H_G\to W\) have the stated determining property.  If
  \(V\in\ker R\), then \(R(V)=R(0)\), so the determining property gives
  \(\mathcal K_\alpha(V)=\mathcal K_\alpha(0)\).  By injectivity of
  \(\mathcal K_\alpha\), \(V=0\).  Thus \(R\) is injective, and its rank is at
  least \(\dim\mathcal H_G=q-1\), proving
  \eqref{eq:minimal-skeleton-rank-bound}.

  Conversely, if a fixed-coordinate linear record is not injective, choose a
  non-zero \(V\in\ker R\) and put \(H=0\), \(H'=V\).  Then
  \(R(H)=R(H')\), but the injectivity of \(\mathcal K_\alpha\), already proved
  above, gives
  \(\mathcal K_\alpha(H)\ne\mathcal K_\alpha(H')\).  This is the advertised bad
  pair and shows that the rank lower bound is exactly the obstruction to
  determining all scalar-normalised product constants in the fixed
  determinant-coordinate chart.

  The finite character orthogonality relations identify the non-trivial
  irreducible characters with an orthonormal basis of the zero-mean
  class-function space, with the standard conjugation symmetry when the
  original class function is real.  Therefore the Peter--Weyl record
  \eqref{eq:minimal-skeleton-peter-weyl-record} is an isometric injection
  from \(\mathcal H_G\) onto a real coordinate space of dimension \(q-1\), so
  it attains the lower bound.

  For the dynamical constants, \Cref{thm:class-product-determinant-rigidity}
  identifies the coordinate tuple of \(H_K\) with
  \((F_\varrho(\lambda^{-1}))_{\varrho\ne\mathbf1}\), and
  \Cref{prop:fixed-label-determinant-perron} evaluates those coordinates
  from the finite twisted determinant family.  This gives the claimed
  rank-sharp fixed-coordinate interpretation of the determinant data.  Finally,
  \(E_B\) is the zero-mean class function whose non-trivial Fourier
  coordinates are
  \(F_\varrho(\lambda^{-1})-B_\varrho\).  Applying Parseval in the same
  zero-mean Peter--Weyl space proves
  \eqref{eq:minimal-skeleton-error-bound}.  If the right-hand side is
  non-zero, some classwise error is non-zero, so the substitute boundary
  cannot give all scalar-normalised product constants; scalar data live in
  the trivial one-dimensional summand and have already been fixed by
  \(\alpha\).
\end{proof}

\begin{corollary}[Scalar baseline and Artin-substitution obstruction]
  \label{cor:class-product-artin-obstruction-recovery}
  \label{cor:class-product-artin-obstruction-readout}
  Assume the hypotheses of \Cref{thm:class-mertens-explicit}.  After the
  same scalar normalisation, any positive class-product vector satisfying
  \eqref{eq:class-product-normalisation-constraint} determines a unique
  zero-mean Peter--Weyl coordinate vector by
  \[
    \widehat H(\pi)
    =
    \sum_{C\subset G}\frac{|C|}{|G|}
      \left(-\frac{|G|}{|C|}\log K_C-\bigl(\gamma+\log C(A)\bigr)
      \right)\chi_\pi(C),
    \qquad \pi\ne\mathbf1 .
  \]
  For the dynamical vector of \Cref{thm:class-euler-product-mertens}, these
  coordinates are exactly
  \(\widehat H(\pi)=F_\pi(\lambda^{-1})\).  The scalar baseline
  \(K_C=\exp(-|C|\alpha_A/|G|)\) for all \(C\) occurs if and only if
  \(F_\varrho(\lambda^{-1})=0\) for every non-trivial \(\varrho\).  If the
  displayed fixed-label values are replaced by regular periodic Artin
  logarithms \(\mathcal L_\varrho(\lambda^{-1})\), then the resulting class
  constants agree with the true constants for every conjugacy class if and
  only if
  \[
    F_\varrho(\lambda^{-1})=
    \mathcal L_\varrho(\lambda^{-1})
    \qquad(\varrho\ne\mathbf1).
  \]
\end{corollary}

\begin{proof}
  Define
  \[
    \mathfrak K(C):=
    -\frac{|G|}{|C|}\log K_C-\bigl(\gamma+\log C(A)\bigr).
  \]
  Taking logarithms in \eqref{eq:class-euler-product-k-f} for the dynamical
  product vector gives
  \[
    \mathfrak K(C)=
    \sum_{\varrho\ne\mathbf1}
      \chi_\varrho^\ast(C)F_\varrho(\lambda^{-1}).
  \]
  Multiplying by \(|C|\chi_\pi(C)/|G|\), summing over conjugacy classes,
  and using Hermitian character orthogonality gives the displayed inverse
  formula for the coordinate \(\widehat H(\pi)\).  For an arbitrary positive
  vector satisfying \eqref{eq:class-product-normalisation-constraint}, the
  same orthogonality applies to the zero-mean profile \(\mathfrak K\), so
  the coordinates are unique.  For the actual dynamical vector the preceding
  Fourier identity identifies this coordinate with
  \(F_\pi(\lambda^{-1})\), \(\pi\ne\mathbf1\).

  The scalar-baseline assertion follows from
  \eqref{eq:class-product-parseval}: the displayed norm vanishes if and
  only if all non-trivial fixed-label coordinates vanish, and then
  \eqref{eq:class-product-log-vector} gives
  \(K_C=\exp(-|C|\alpha_A/|G|)\) for every class \(C\).  Finally,
  \Cref{thm:adams-obstruction-class-product} compares the true fixed-label
  constants with the naive constants obtained by substituting
  \(\mathcal L_\varrho(\lambda^{-1})\).  Its ratio formula shows that all
  class constants agree if and only if the non-trivial Peter--Weyl expansion
  of the discrepancy is zero.  Since the irreducible characters form an
  orthonormal basis of class functions, this is equivalent to
  \(F_\varrho(\lambda^{-1})=\mathcal L_\varrho(\lambda^{-1})\) for every
  non-trivial \(\varrho\).
\end{proof}
\subsubsection{Finite determinant closure and classical comparison}
\label{subsubsec:finite-determinant-closure-classical-comparison}

This final comparison block keeps the novelty calibration explicit.  The
Frobenius-class exponent and the existence of a product constant are classical
inputs; the finite-state contribution proved here is that the constant is the
Adams-corrected fixed-label determinant coordinate and is computable from the
finite twisted determinant family with the stated branch convention.

\begin{proposition}[Finite determinant dependence of the main constants]
  \label{prop:finite-data-main-chain-closure}
  \label{prop:finite-data-main-spine-closure}
  Assume \(\lambda>1\) and the twisted spectral gap
  \(\eta<\lambda\) of \Cref{thm:class-mertens-explicit}.  Encode the finite
  group by its multiplication table and let the finite determinant datum be
  \begin{equation*}
    \begin{aligned}
      \mathfrak D(A,\tau;G):=
      \Bigl(&G,\ \{C\subset G\,\textup{conjugacy class}\},
      (C,m)\mapsto C^m, \\
      &C(A),\ \{\det(I-zA_\pi):\pi\in\Irr(G)\}\Bigr).
    \end{aligned}
  \end{equation*}
  Together with the Abel-path branch convention fixed separately in
  \Cref{thm:class-mertens-explicit,prop:fixed-label-determinant-perron}, this
  finite datum determines the following three nested output tiers.  Each tier
  is asserted only under the hypotheses and branch conventions of the theorem
  package cited for that tier.
  \begin{enumerate}[label=\textup{(\roman*)},leftmargin=*]
    \item \emph{Algebraic determinant recovery:} every primitive Peter--Weyl coefficient
    \(\pi_{n,\varrho}(A,\tau)\) and hence every primitive class count
    \(p_{n,C}\);
    \item \emph{Perron-boundary recovery under the twisted gap:} every non-trivial Perron-point primitive value
    \(\Pi_\varrho(\lambda^{-1})\) and every class-Mertens constant
    \(\mathcal M_C(A,\tau;G)\);
    \item \emph{Fixed-label product recovery:} every fixed-label value \(F_\varrho(\lambda^{-1})\) and every
    Frobenius-class Euler-product constant \(K_C\).
  \end{enumerate}
  The recovery in item~\textup{(i)} is purely algebraic: coefficient
  extraction from the finitely many determinant polynomials, finite
  character-table linear algebra, Adams coefficients, M\"obius inversion, and
  finite Fourier inversion.  The recovery in items~\textup{(ii)} and
  \textup{(iii)} is the branch-compatible Perron-boundary evaluation of those
  same finite determinant polynomials under \(\eta<\lambda\), not a claim that
  the boundary constants are finite algebraic expressions in the polynomial
  coefficients alone.  The singular scalar determinant logarithm is never
  evaluated; the scalar normalisation uses the positive Perron constant
  \(C(A)\).  Conversely, after the scalar normalisation
  \(C(A)\), the product-constant vector \((K_C)_C\) determines exactly the
  non-trivial fixed-label Perron vector
  \((F_\varrho(\lambda^{-1}))_{\varrho\ne\mathbf1}\) by the inverse
  Fourier formula in \eqref{eq:comparison-fixed-label-inverse}.
  Hence two extensions with identical data \(\mathfrak D(A,\tau;G)\), the
  same branch choices, and the same scalar normalisation \(C(A)\) have
  identical outputs in Main Theorems~I--III.  This proposition is only the
  closure of the assembled chain: the raw determinant recovery is
  \Cref{cor:primitive-peter-weyl-determinant,thm:class-count-trace-recovery},
  the Perron-point convergence and additive constants are
  \Cref{thm:class-mertens-explicit,thm:class-mertens-fourier}, and the
  product-level fixed-label constants are
  \Cref{thm:class-euler-product-mertens,prop:fixed-label-determinant-perron,thm:class-product-determinant-rigidity}.
\end{proposition}
\begin{proof}
  For tier~\textup{(i)}, the finite group component of
  \(\mathfrak D(A,\tau;G)\) gives the complete
  character table, the irreducible characters, conjugacy classes, and the Adams
  class maps \(C\mapsto C^m\).  The determinant polynomials determine all
  traces \(\Tr(A_\pi^n)\) by the finite identity
  \(\log\det(I-zA_\pi)^{-1}=\sum_{n\ge1}\Tr(A_\pi^n)z^n/n\) on its formal
  power-series side.  \Cref{lem:twisted-trace-identity} identifies these traces
  with the periodic twisted fixed-point sums.  Finite character orthogonality
  then gives \Cref{thm:class-function-linearisation}, and the Adams--M\"obius
  formula in \Cref{cor:class-function-adams-mobius} converts the periodic
  class-function coefficients to primitive coefficients.  Therefore
  \Cref{cor:primitive-peter-weyl-determinant,cor:primitive-peter-weyl-trace,thm:class-count-trace-recovery} recover all primitive Peter--Weyl
  coefficients and primitive class counts from \(\mathfrak D(A,\tau;G)\).

  For tier~\textup{(ii)}, under \(\eta<\lambda\), the non-trivial determinant series are regular at the
  Perron point along the branch specified in
  \Cref{thm:class-mertens-explicit}.  That theorem, together with
  \Cref{prop:class-mertens-deviations,thm:class-mertens-fourier}, evaluates
  the non-trivial primitive values \(\Pi_\varrho(\lambda^{-1})\) and hence the
  class-Mertens constants.  For tier~\textup{(iii)}, the fixed-label
  Euler-transform theorem \Cref{thm:class-euler-product-mertens} and the finite determinant formula
  \Cref{prop:fixed-label-determinant-perron} evaluate
  \(F_\varrho(\lambda^{-1})\) for every non-trivial \(\varrho\); substituting
  these values into \Cref{thm:class-product-determinant-rigidity} gives the
  Frobenius-class product constants \(K_C\).  The scalar singular determinant
  is not evaluated in these steps: the scalar contribution is the supplied
  Perron constant \(C(A)\), exactly as in the displayed product formulas of
  \Cref{thm:class-euler-product-mertens,thm:class-product-determinant-rigidity}.
  Thus the proof uses exactly the branch conventions and convergence
  hypotheses attached to each tier's theorem package.

  Finally, the inverse formula \eqref{eq:comparison-fixed-label-inverse} is a
  finite Fourier inversion on the class algebra after the scalar baseline has
  been fixed.  It shows that the vector \((K_C)_C\) determines precisely the
  non-trivial vector \((F_\varrho(\lambda^{-1}))_{\varrho\ne\mathbf1}\), and
  the direct product formula gives the converse dependence.  Every operation
  used above is therefore either finite coefficient extraction, finite
  character-table algebra, Adams--M\"obius inversion, or the already proved
  branch-compatible Perron-point evaluation under \(\eta<\lambda\).  The
  proposition adds no new leading asymptotic claim beyond the cited packages:
  the Frobenius-density exponent in the product theorem is classical, while the
  new finite-data content used here is the Adams--M\"obius fixed-label
  determinant recovery of the constants.  Identical
  finite collections with identical branch choices and scalar normalisation pass
  through the same sequence of operations, so they have identical outputs in
  the three main theorems.
\end{proof}

\begin{remark}[Ordered finite recovery]
  \label{rem:assembled-recovery-dependency-order}
  Under the hypotheses of \Cref{prop:finite-data-main-chain-closure}, the
  finite-data recovery factors through the following ordered chain:
  \begin{equation}\label{eq:assembled-recovery-dependency-order}
    \mathfrak D(A,\tau;G)
    \longrightarrow
    \mathcal O_I
    \longrightarrow
    \mathcal O_{II}
    \longrightarrow
    \mathcal O_{III},
  \end{equation}
  where \(\mathcal O_I\) consists of the primitive Peter--Weyl coefficients
  and primitive class counts, \(\mathcal O_{II}\) consists of the
  non-trivial Perron-point primitive values and class-Mertens constants, and
  \(\mathcal O_{III}\) consists of the fixed-label Perron values and
  Frobenius-class product constants.  The first arrow uses only coefficient
  extraction, Adams coefficients, and finite Fourier inversion.  The second
  arrow is the Abel-branch Perron-boundary evaluation supplied by the
  twisted spectral gap.  The third arrow is the fixed-label Euler-transform
  and class-product formula of \Cref{thm:class-euler-product-mertens} and
  \Cref{thm:class-product-determinant-rigidity}.  Thus
  \Cref{prop:finite-data-main-chain-closure,cor:assembled-adams-frobenius-product-formula}
  are finite recoveries of Main Theorems~I--III, while
  \Cref{thm:class-product-determinant-rigidity} is the branch-compatible
  fixed-coordinate rigidity theorem for the Adams-corrected product vector
  after scalar normalisation.
  In particular, equality of two finite determinant data implies equality
  of all three ordered outputs, and after \(C(A)\) is fixed equality of the
  product-constant vectors is equivalent to equality of the non-trivial
  fixed-label Perron vectors.
\end{remark}

The factor \(N^{-|C|/|G|}\) and the existence of a Frobenius-class Mertens
product asymptotic are the classical part of
\eqref{eq:class-euler-product-asymptotic}; compare the Frobenius-class product
theorem of Mesliza--Noorani \cite{MeslizaNoorani1999FrobeniusMertens},
Sharp's dynamical Mertens theorem \cite{Sharp1991Mertens}, and the Chebotarev
orbit-distribution antecedents
\cite{ParryPollicott1986Chebotarev,NooraniParry1992ChebotarevShifts}.
Nearby finite-data and skew-product works such as
\cite{CrismaleDelVecchioGrisetaRossi2025NoncommutativeSkewProduct,OHare2025FiniteDataRigidity}
address different rigidity or extension questions.  The finite-state
primitive-label contribution needed here is the following named comparison
theorem: it is the Adams correction to the Frobenius-class Mertens product,
not merely a restatement of the classical density theorem.

\begin{theorem}[Adams correction to Frobenius-class Mertens products]
  \label{thm:class-product-comparison}
  \label{prop:class-product-comparison}
  Under the hypotheses of \Cref{thm:class-euler-product-mertens}, the
  following assertions hold after the classical exponent and the scalar base
  constant have been separated.
  Equivalently, Main Theorem~III has the following dependency and comparison
  profile:
  \begin{center}
  \small
  \begin{tabular}{>{\raggedright\arraybackslash}p{0.25\textwidth}
                  >{\raggedright\arraybackslash}p{0.34\textwidth}
                  >{\raggedright\arraybackslash}p{0.29\textwidth}}
    \toprule
    layer & input used & output in this paper\\
    \midrule
    classical Frobenius product layer
      & class density and Tauberian existence
      & \(|C|/|G|\) and existence of \(K_C\)\\
    fixed-label determinant layer
      & \(\{\det(I-zA_\pi)\}_{\pi\in\Irr(G)}\) with Adams powers
      & \(F_\varrho(\lambda^{-1})\) and the constants \(K_C\)\\
    Artin-replacement test
      & \(\mathcal L_\varrho(\lambda^{-1})\) substituted for
        \(F_\varrho(\lambda^{-1})\)
      & zero-mean obstruction vector and Parseval bound\\
    \bottomrule
  \end{tabular}
  \end{center}
  The first row is the inherited classical exponent/existence layer.  The
  second row is the determinant-valued contribution of the theorem.  The third
  row is downstream: it tests a proposed replacement of the already recovered
  constants and is not a hypothesis for Main Theorems~I--III.
  \begin{enumerate}[label=\textup{(\Alph*)}]
    \item The three descriptions of the logarithmic constant are equivalent:
  \begin{enumerate}[label=\textup{(\roman*)}]
    \item the primitive product tail expression
    \[
      -\log K_C=
      \mathcal M_C(A,\tau;G)
      +\sum_{r\ge2}\frac{\Pi_{\mathbf1_C}(\lambda^{-r})}{r};
    \]
    \item the fixed-label class-character expression
    \[
      -\log K_C=
      \frac{|C|}{|G|}
      \left(
        \gamma+\log C(A)
        +\sum_{\varrho\neq\mathbf1}
          \chi_\varrho^\ast(C)F_\varrho(\lambda^{-1})
      \right);
    \]
    \item the determinant expression obtained from
    \eqref{eq:fixed-label-determinant-perron} after substituting
    \[
      F_\varrho(\lambda^{-1})
      =
      \sum_{r,m\ge1}\frac{\mu(m)}{rm}
        \sum_{\substack{\pi\in\Irr(G)\\(r,m,\pi)\ne(1,1,\mathbf1)}}
          a_{\varrho,\pi}^{(m)}
          \log\det(I-\lambda^{-rm}A_\pi)^{-1}.
      \]
    \end{enumerate}
    \item Define the normalised logarithmic boundary vector
    \[
      \mathfrak K(C)
      :=-
      \frac{|G|}{|C|}\log K_C-(\gamma+\log C(A)).
    \]
    Then
    \begin{equation}\label{eq:comparison-fixed-label-inverse}
      \mathfrak K(C)=
      \sum_{\varrho\ne\mathbf1}
        \chi_\varrho^\ast(C)F_\varrho(\lambda^{-1}),
      \qquad
      F_\pi(\lambda^{-1})=
      \sum_{C\subset G}\frac{|C|}{|G|}\mathfrak K(C)\chi_\pi(C)
    \end{equation}
    for every non-trivial \(\pi\).  Hence, after scalar normalisation, the
    full class-product constant vector is equivalent to the non-trivial
    fixed-label Perron vector.
    \item If the fixed-label values are replaced by the regular periodic
    Artin logarithms \(\mathcal L_\varrho(\lambda^{-1})\), define
    \[
      K_C^{\mathrm{Art}}
      :=
      \exp\!\left(
        -\frac{|C|}{|G|}
        \left(
          \gamma+\log C(A)
          +\sum_{\varrho\neq\mathbf1}
            \chi_\varrho^\ast(C)\mathcal L_\varrho(\lambda^{-1})
        \right)
      \right).
    \]
    Then the naive constants satisfy
    \begin{equation}\label{eq:comparison-artin-error}
      \log\frac{K_C}{K_C^{\mathrm{Art}}}
      =
      -\frac{|C|}{|G|}
        \sum_{\varrho\neq\mathbf1}
          \chi_\varrho^\ast(C)
          \left(F_\varrho(\lambda^{-1})-
          \mathcal L_\varrho(\lambda^{-1})\right),
    \end{equation}
    and all class constants are unchanged by this substitution if and only if
    \(F_\varrho(\lambda^{-1})=\mathcal L_\varrho(\lambda^{-1})\) for every
    non-trivial \(\varrho\).  Thus the novelty calibrated by this proposition
    is exactly the fixed-label Adams--M\"obius determinant constant and the
    Peter--Weyl obstruction to replacing it by an ordinary Artin logarithm.
    \item The preceding obstruction is quantitatively sharp in the following
    sense.  Let \(B=(B_\varrho)_{\varrho\ne\mathbf1}\) be any proposed
    non-trivial Peter--Weyl boundary vector and set
    \[
      K_C(B):=
      \exp\!\left(
        -\frac{|C|}{|G|}
        \left(
          \gamma+\log C(A)
          +\sum_{\varrho\neq\mathbf1}
            \chi_\varrho^\ast(C)B_\varrho
        \right)
      \right).
    \]
    With the same logarithm branch as above, define the normalised classwise
    error
    \[
      \mathscr E_B(C):=-\frac{|G|}{|C|}
        \log\frac{K_C}{K_C(B)}.
    \]
    Then
    \begin{equation}\label{eq:comparison-parseval-bound}
      \begin{aligned}
      \mathscr E_B(C)
      &=
      \sum_{\varrho\ne\mathbf1}
        \chi_\varrho^\ast(C)
        \left(F_\varrho(\lambda^{-1})-B_\varrho\right),\\
      \sum_{C\subset G}\frac{|C|}{|G|}\abs{\mathscr E_B(C)}^2
      &=
      \sum_{\varrho\ne\mathbf1}
        \abs{F_\varrho(\lambda^{-1})-B_\varrho}^2 .
      \end{aligned}
    \end{equation}
    Consequently a branch-compatible formula with the classical exponent
    \(|C|/|G|\) and scalar base \(\gamma+\log C(A)\) gives the true constants
    for every conjugacy class if and only if
    \(B_\varrho=F_\varrho(\lambda^{-1})\) for every non-trivial
    \(\varrho\).  No non-zero non-trivial Peter--Weyl error can be absorbed
    into the classical exponent, the existence statement, or the scalar base
    constant.
  \end{enumerate}
\end{theorem}

\begin{proof}
  The displayed comparison table is a restatement of the ordered assertions
  proved below: the classical exponent and existence statement are supplied by
  \Cref{thm:class-euler-product-mertens}, the determinant-valued fixed-label
  coordinates are \textup{(A)} and \textup{(B)}, and the Artin-replacement row
  is precisely the error criterion and Parseval bound in \textup{(C)} and
  \textup{(D)}.  Since \textup{(C)} and \textup{(D)} are obtained only after
  subtracting a proposed replacement vector from the fixed-label vector, they
  are downstream tests rather than inputs to the proof of
  \Cref{thm:class-euler-product-mertens}.

  The equivalence between \textup{(A)(i)} and \textup{(A)(ii)} is exactly
  \Cref{lem:class-product-constant-equivalence}. For \textup{(A)(iii)}, substitute the
  determinant formula of \Cref{prop:fixed-label-determinant-perron} into
  the non-trivial summands in (ii). The exceptional scalar Perron term is
  absent there because, for \(\varrho\neq\mathbf1\),
  \(a_{\varrho,\mathbf1}^{(1)}=0\) by character orthogonality; all
  remaining logarithms use the regular Abel branch specified in
  \Cref{prop:fixed-label-determinant-perron}. Hence the new datum is a
  finite twisted-determinant evaluation of fixed-primitive-label Euler
  transforms.

  This proves \textup{(A)}.  Taking logarithms in
  \eqref{eq:class-euler-product-k-f} gives the first formula in
  \eqref{eq:comparison-fixed-label-inverse}. Multiplying that identity by
  \((|C|/|G|)\chi_\pi(C)\), summing over conjugacy classes, and applying
  character orthogonality gives the inverse formula for every
  \(\pi\ne\mathbf1\).  Hence \textup{(B)} follows.

  Finally \Cref{thm:fixed-label-euler-transform} identifies
  \[
    F_\varrho(z)=\sum_{r\ge1}\frac{\Pi_\varrho(z^r)}{r}
  \]
  as the fixed-label Euler transform.  The ordinary Artin logarithm is
  \[
    \mathcal L_\varrho(z)=
    \sum_{r\ge1}\frac{\Pi_{\psi^r\varrho}(z^r)}{r}.
  \]
  The distinction comes from
  repetitions of primitive labels, which apply the power map
  \(g\mapsto g^r\) in the periodic logarithm but not in the fixed-label
  primitive Euler product. Subtracting the logarithm of the Artin-replaced
  formula from the corrected logarithmic formula gives
  \eqref{eq:comparison-artin-error}.  If the right-hand side vanishes for
  every conjugacy class, then the non-trivial character expansion has all
  coefficients zero by orthogonality; conversely, if every coefficient
  \(F_\varrho(\lambda^{-1})-
  \mathcal L_\varrho(\lambda^{-1})\) is zero, the displayed error vanishes
  classwise.  This proves \textup{(C)}.

  For \textup{(D)}, subtract the branch-compatible logarithm defining
  \(K_C(B)\) from the true logarithm in \eqref{eq:class-euler-product-k-f}.
  After multiplying by \(-|G|/|C|\), the scalar term
  \(\gamma+\log C(A)\) cancels and gives the first identity in
  \eqref{eq:comparison-parseval-bound}.  The irreducible characters form an
  orthonormal basis of class functions with inner product
  \[
    \langle f,h\rangle_G=
    \sum_{C\subset G}\frac{|C|}{|G|}f(C)\overline{h(C)}.
  \]
  Applying Parseval to the class function \(\mathscr E_B\) gives the second
  identity.  The right-hand side is zero exactly when
  \(B_\varrho=F_\varrho(\lambda^{-1})\) for every non-trivial
  \(\varrho\), and otherwise the normalised error is non-zero for at least
  one conjugacy class.  This proves the uniqueness and non-absorption claims
  in \textup{(D)}.

\end{proof}

Thus Main Theorems~I--III are finite-state determinant-to-primitive,
Perron-point, and Adams-corrected product-constant theorems.  They do not
assert a new geometric periodic-orbit invariant, a Ruelle-resonance theorem, a
finite-data rigidity theorem, a general certificate format, or an
operator-algebraic skew-product classification.  The shadow results
\Cref{thm:quotient-functoriality,thm:abelian-shadow-defect,thm:cyclic-detection-boundary}, the eta-gap tests
\Cref{thm:eta-gap-rigidity,cor:eta-gap-boundary-classifier,prop:eta-gap-finite-state-certificate}, and the \(S_3\) certificates
\Cref{thm:s3-certified-nonabelian-defect,prop:s3-two-coordinate-product-obstruction-vector,cor:s3-product-constant-obstruction-finite-state,prop:s3-naive-artin-obstruction}
are internal finite-state tools for the same product-coordinate picture,
not comparisons with the cited geometric, resonance, finite-data-rigidity,
computational-certificate, or operator-algebraic papers.

\begin{proposition}[Theorem-by-theorem comparison calibration]
  \label{prop:theorem-by-theorem-comparison-calibration}
  Under the hypotheses of \Cref{thm:class-mertens-explicit}, the logical role of
  the cited comparison literature in Main Theorems~I--III is exactly the
  following.
  \begin{enumerate}[label=\textup{(\alph*)},leftmargin=*]
    \item Main Theorem~I,
    \Cref{thm:class-function-linearisation,cor:class-function-adams-mobius,cor:primitive-peter-weyl-determinant,thm:primitive-peter-weyl-trace},
    is calibrated against Parry--Pollicott's Chebotarev theorem for finite
    Galois coverings \cite{ParryPollicott1986Chebotarev}: the comparison layer
    is orbit distribution and Artin--Ruelle motivation, while the contribution
    proved here is finite character-algebra inversion from periodic twisted
    determinants to primitive Peter--Weyl channels.
    \item The same Main Theorem~I package is calibrated separately against
    Noorani--Parry's Chebotarev theorem for shifts of finite type
    \cite{NooraniParry1992ChebotarevShifts}: the comparison layer is SFT
    Frobenius-class equidistribution, while the contribution proved here is
    Adams--M\"obius recovery of primitive Peter--Weyl channels without invoking
    an additional Tauberian orbit-counting theorem.
    \item Main Theorem~II,
    \Cref{thm:class-mertens-explicit,prop:class-mertens-deviations,thm:class-mertens-fourier},
    is calibrated against Sharp's orbit-Mertens theorem
    \cite{Sharp1991Mertens}: the comparison layer is the classical logarithmic
    scale, while the contribution proved here is the finite-state determinant
    formula for the class-indexed constant term under the stated Perron gap.
    \item Main Theorem~II and the product normalisation theorem
    \Cref{thm:class-mertens-explicit,thm:class-euler-product-mertens} are
    calibrated against Mesliza--Noorani's Frobenius-class Mertens theorem
    \cite{MeslizaNoorani1999FrobeniusMertens}: the comparison layer is product
    existence and scale, while the contribution proved here is explicit finite
    Peter--Weyl boundary evaluation of the classwise constants.
    \item Main Theorem~III,
    \Cref{thm:class-euler-product-mertens,prop:fixed-label-determinant-perron,thm:class-product-determinant-rigidity,prop:class-product-comparison,thm:adams-obstruction-class-product},
    is calibrated against ordinary Artin--Ruelle determinant analogues in the
    symbolic-dynamical lineage \cite{ParryPollicott1990Zeta}: the contribution
    proved here is the Adams-corrected fixed-label Euler value
    \(F_\varrho(\lambda^{-1})\), its determinant recovery, and the obstruction
    to replacing it by \(\mathcal L_\varrho(\lambda^{-1})\).
    \item The quotient/shadow readouts
    \Cref{thm:quotient-functoriality,thm:abelian-shadow-defect} are calibrated
    against Coles's Vassiliev-writhe/Axiom~A invariant setting
    \cite{Coles2025VassilievWritheAxiomA}: no geometric periodic-orbit
    invariant is asserted here; the retained statements are finite SFT cocycle
    determinant-coordinate readouts.
    \item The eta-gap package
    \Cref{thm:eta-gap-rigidity,cor:eta-gap-boundary-classifier,prop:eta-gap-finite-state-certificate}
    is calibrated against Humbert's first-Ruelle-resonance comparison setting
    \cite{Humbert2025FirstRuelleResonance}: no Banach-space or resonance
    theorem is asserted here; the retained assertions are finite-matrix
    spectral-gap tests at the Perron point.
    \item The cyclic-detection and finite-state certificate boundary
    \Cref{thm:cyclic-detection-boundary,prop:eta-gap-finite-state-certificate}
    is calibrated against O'Hare's finite-data rigidity theorem
    \cite{OHare2025FiniteDataRigidity}: no finite-data rigidity theorem is
    asserted here; the paper gives explicit finite-state readouts and detection
    boundaries instead.
    \item The non-abelian obstruction package
    \Cref{thm:s3-certified-nonabelian-defect,prop:s3-two-coordinate-product-obstruction-vector,cor:s3-product-constant-obstruction-finite-state,prop:s3-naive-artin-obstruction}
    is calibrated against Crismale--Del Vecchio--Griseta--Rossi's
    non-commutative skew-product framework
    \cite{CrismaleDelVecchioGrisetaRossi2025NoncommutativeSkewProduct}: no
    operator-algebraic skew-product classification is asserted here; the
    retained results are explicit finite-group character obstructions for the
    fixed SFT cocycle model.
  \end{enumerate}
  Consequently Main Theorems~I--III depend on the classical comparison sources
  only through the explicitly cited symbolic-dynamical orbit-counting,
  Mertens, and determinant framework, while the recent 2025 comparison papers
  are not proof inputs for any implication in the main theorem chain.
\end{proposition}

\begin{proof}
  For Main Theorem~I, the asserted dependency is exactly the finite determinant
  interface in \Cref{prop:main-theorem-i-assembly-map} together
  with the class-function linearisation and Adams--M\"obius inversion in
  \Cref{thm:class-function-linearisation,cor:class-function-adams-mobius}.
  Those results start from the finite matrices \(A_\varrho\) and the character
  identities of the finite group; their proofs use the determinant expansions
  and Adams operations displayed there, not an additional orbit-counting
  theorem.  Thus Parry--Pollicott's finite-covering Chebotarev theorem supplies
  the classical Artin--Ruelle orbit-distribution model, Noorani--Parry supplies
  the corresponding shift-of-finite-type Chebotarev comparison, and the
  primitive-channel inversion is the internal finite-label step.

  For Main Theorem~II, \Cref{thm:class-mertens-explicit} assumes the strict
  non-trivial spectral gap and isolates the leading term
  \((|C|/|G|)\log N\).  The constant term is then evaluated by the finite
  Peter--Weyl formula and the Perron-boundary values in
  \Cref{thm:class-mertens-fourier}.  This proves the Sharp row: Sharp supplies
  the relevant closed-orbit Mertens scale, but the determinant expression for
  \(\mathcal M_C(A,\tau;G)\) is the internal output.  The same theorem and
  \Cref{thm:class-euler-product-mertens} give the Mesliza--Noorani row:
  their Frobenius-class Mertens comparison fixes the product scale and
  existence layer, whereas the finite Peter--Weyl Perron-boundary evaluation is
  proved here from the twisted finite matrices.

  For Main Theorem~III, \Cref{thm:class-euler-product-mertens} first converts
  the logarithmic Main Theorem~II estimate and the absolutely convergent product
  tail into \(P_C(N)=K_C N^{-|C|/|G|}(1+O(N^{-1}))\).  The sentence in that
  theorem separates the classical exponent and existence layer from the local
  constant calculation.  The constant is then rewritten through
  \(F_\varrho(\lambda^{-1})\), and
  \Cref{prop:fixed-label-determinant-perron,thm:class-product-determinant-rigidity,prop:class-product-comparison,thm:adams-obstruction-class-product}
  identify its determinant value and the Artin-replacement obstruction.  Hence
  the third row follows.

  Finally, the modern comparison papers cited in the last four rows are
  separated source by source.  Coles concerns geometric periodic-orbit
  observables, while \Cref{thm:quotient-functoriality,thm:abelian-shadow-defect}
  use only finite quotient and character maps.  Humbert concerns
  Ruelle-resonance behaviour, while the eta-gap package uses only spectra of
  finite twisted adjacency matrices.  O'Hare concerns finite-data rigidity,
  while the cyclic-detection and certificate statements here give explicit
  readouts and boundaries rather than a rigidity theorem.  Crismale--Del
  Vecchio--Griseta--Rossi concern non-commutative skew-product structure, while
  the \(S_3\) package uses finite character tables and determinant coordinates.
  None of these 2025 results appears in the hypotheses or proofs of the theorem
  packages just listed.  This proves both the venue calibration and the asserted
  non-dependence of the Main Theorems on those recent comparison sources.
\end{proof}

% END INPUT sec_chebotarev_class_mertens_rigidity_consequences.tex
% BEGIN INPUT sec_chebotarev_class_mertens_equivalence.tex
\subsubsection{Constant equivalence lemmas}

This short unit reconciles the two displayed normalisations of the same
positive product constant.  It is placed after Main Theorem~III so that the
reader can see where the additive constant formula, the fixed-label Euler
formula, and the positive branch convention are proved to agree before the
determinant-only recovery begins.

\begin{lemma}[Compatibility of class Euler-product constants]
  \label{lem:class-product-constant-equivalence}
  Under the hypotheses of \Cref{thm:class-euler-product-mertens}, every
  conjugacy class \(C\subset G\) satisfies
  \begin{equation}\label{eq:class-product-constant-compatibility}
    \mathcal M_C(A,\tau;G)
    +\sum_{r\ge2}\frac{\Pi_{\mathbf1_C}(\lambda^{-r})}{r}
    =
    \frac{|C|}{|G|}
    \left(
      \gamma+\log C(A)
      +\sum_{\varrho\neq\mathbf1}
        \chi_\varrho^\ast(C)F_\varrho(\lambda^{-1})
    \right).
  \end{equation}
  All sums in \eqref{eq:class-product-constant-compatibility} are
  absolutely convergent after the indicated separation of the scalar
  harmonic term. Consequently the constants in
  \eqref{eq:class-euler-product-k-mc} and
  \eqref{eq:class-euler-product-k-f} are exactly the same positive real
  number.
\end{lemma}

\begin{proof}
  The class-indicator expansion gives the identity of class functions
  \[
    \mathbf1_C
    =
    \frac{|C|}{|G|}\mathbf1
    +
    \frac{|C|}{|G|}
      \sum_{\varrho\neq\mathbf1}\chi_\varrho^\ast(C)\chi_\varrho .
  \]
  Since \(r\ge2\), the primitive series are evaluated strictly inside
  their common disc of absolute convergence. Hence linearity may be used
  term by term, and
  \begin{equation}\label{eq:class-product-tail-split}
    \sum_{r\ge2}\frac{\Pi_{\mathbf1_C}(\lambda^{-r})}{r}
    =
    \frac{|C|}{|G|}\sum_{r\ge2}\frac{\Pi_{\mathbf1}(\lambda^{-r})}{r}
    +
    \frac{|C|}{|G|}
      \sum_{\varrho\neq\mathbf1}\chi_\varrho^\ast(C)
      \sum_{r\ge2}\frac{\Pi_\varrho(\lambda^{-r})}{r}.
  \end{equation}
  The convergence used here follows from
  \(\Pi_\phi(\lambda^{-r})\ll_\phi\lambda^{1-r}\) for bounded class
  functions and \(r\ge2\): indeed \(p_n(A)\ll\lambda^n/n\) gives
  \[
    \abs{\Pi_\phi(\lambda^{-r})}
    \le
    \norm{\phi}_\infty\sum_{n\ge1}p_n(A)\lambda^{-rn}
    \ll_\phi
    \sum_{n\ge1}\frac{\lambda^{(1-r)n}}{n}
    \ll_\phi \lambda^{1-r}.
  \]

  Next insert the additive class Mertens constant formula
  \eqref{eq:class-mertens-constant}:
  \[
    \mathcal M_C(A,\tau;G)
    =
    \frac{|C|}{|G|}\mathcal M(A)
    +
    \frac{|C|}{|G|}
      \sum_{\varrho\neq\mathbf1}
        \chi_\varrho^\ast(C)\Pi_\varrho(\lambda^{-1}).
  \]
  Adding this display to \eqref{eq:class-product-tail-split} and
  grouping the non-trivial channels gives
  \begin{multline}\label{eq:class-product-before-scalar-normalisation}
    \mathcal M_C(A,\tau;G)
    +\sum_{r\ge2}\frac{\Pi_{\mathbf1_C}(\lambda^{-r})}{r}
    =
    \frac{|C|}{|G|}
    \left(
      \mathcal M(A)+\sum_{r\ge2}\frac{\Pi_{\mathbf1}(\lambda^{-r})}{r}
    \right) \\
    +
    \frac{|C|}{|G|}
      \sum_{\varrho\neq\mathbf1}\chi_\varrho^\ast(C)
      \left(
        \Pi_\varrho(\lambda^{-1})
        +\sum_{r\ge2}\frac{\Pi_\varrho(\lambda^{-r})}{r}
      \right).
  \end{multline}
  By the definition of the fixed-label Euler transform, together with the
  non-trivial Perron-point convergence supplied by the twisted gap,
  \[
    F_\varrho(\lambda^{-1})
    =
    \Pi_\varrho(\lambda^{-1})
    +\sum_{r\ge2}\frac{\Pi_\varrho(\lambda^{-r})}{r}
    \qquad(\varrho\neq\mathbf1).
  \]
  Thus only the scalar parenthesis in
  \eqref{eq:class-product-before-scalar-normalisation} remains to be
  normalized.

  This scalar normalization is made before any class Fourier inversion is
  used: it is the scalar finite-part identity from \Cref{sec:finite-part},
  not an evaluation of the singular trivial determinant at the Perron
  point as an ordinary logarithm.
  The scalar Mertens convention in this paper is
  \(\mathcal M(A)=\gamma+\FP(A)\). Hence
  \begin{equation}\label{eq:scalar-product-normalisation}
    \mathcal M(A)+\sum_{r\ge2}\frac{\Pi_{\mathbf1}(\lambda^{-r})}{r}
    =
    \gamma+\FP(A)+
    \sum_{\gamma_0\in\cP}\sum_{r\ge2}
      \frac{\lambda^{-r|\gamma_0|}}{r}.
  \end{equation}
  The double series is absolutely convergent because
  \(p_n(A)\ll\lambda^n/n\). From
  \eqref{eq:finite-part-primitive-closed},
  \[
    \FP(A)
    =
    \log C(A)+
    \sum_{\gamma_0\in\cP}
      \left(
        \lambda^{-|\gamma_0|}
        +\log(1-\lambda^{-|\gamma_0|})
      \right).
  \]
  For each primitive orbit \(\gamma_0\), with
  \(x=\lambda^{-|\gamma_0|}\in(0,1)\),
  \[
    x+\log(1-x)+\sum_{r\ge2}\frac{x^r}{r}
    =x+\log(1-x)-\log(1-x)-x=0.
  \]
  Substituting this cancellation into
  \eqref{eq:scalar-product-normalisation} gives
  \[
    \mathcal M(A)+\sum_{r\ge2}\frac{\Pi_{\mathbf1}(\lambda^{-r})}{r}
    =\gamma+\log C(A).
  \]
  Combining this scalar identity with
  \eqref{eq:class-product-before-scalar-normalisation} proves
  \eqref{eq:class-product-constant-compatibility}.

  Finally, \eqref{eq:class-euler-product-k-mc} is the exponential of the
  negative left-hand side of
  \eqref{eq:class-product-constant-compatibility}, and
  \eqref{eq:class-euler-product-k-f} is the exponential of the negative
  right-hand side. The equality of the exponents proves equality of the
  constants. Positivity follows from
  \eqref{eq:class-euler-product-k-mc}: its exponent is real, and the
  exponential is the ordinary real exponential.
\end{proof}

\begin{lemma}[Hermitian normalisation of the class-product constant]
  \label{lem:class-product-hermitian-normalisation}
  Under the hypotheses of \Cref{thm:class-euler-product-mertens}, the
  \(|C|/|G|\) factor in \eqref{eq:class-euler-product-k-f} multiplies both
  the scalar term and every non-trivial fixed-label Fourier contribution.
  Equivalently,
  \begin{equation}\label{eq:class-product-normalisation-identity}
    -\log K_C
    =
    \frac{|C|}{|G|}
    \left(
      \gamma+\log C(A)
      +\sum_{\varrho\neq\mathbf1}
        \chi_\varrho^\ast(C)F_\varrho(\lambda^{-1})
    \right),
  \end{equation}
  and the expression obtained by moving \(|C|/|G|\) off the non-trivial
  summation is not the Hermitian class-indicator expansion except in
  accidental degenerate cases.
\end{lemma}

\begin{proof}
  The Fourier convention fixed in \Cref{sec:chebotarev} gives
  \[
    \mathbf1_C
    =
    \frac{|C|}{|G|}\mathbf1
    +
    \frac{|C|}{|G|}
      \sum_{\varrho\neq\mathbf1}\chi_\varrho^\ast(C)\chi_\varrho .
  \]
  Thus the coefficient of the non-trivial channel \(\chi_\varrho\) in the
  class indicator is \((|C|/|G|)\chi_\varrho^\ast(C)\), not merely
  \(
    \chi_\varrho^\ast(C)
  \).  Applying this identity first to the additive class-Mertens constant
  gives \eqref{eq:class-mertens-constant}, and applying it termwise to the
  absolutely convergent product tail gives
  \eqref{eq:class-product-tail-split}.  After adding those two formulas,
  the scalar channel is normalised by
  \[
    \mathcal M(A)+\sum_{r\ge2}\frac{\Pi_{\mathbf1}(\lambda^{-r})}{r}
    =\gamma+
    \log C(A),
  \]
  while each non-trivial channel is
  \[
    \Pi_\varrho(\lambda^{-1})
    +
    \sum_{r\ge2}\frac{\Pi_\varrho(\lambda^{-r})}{r}
    =F_\varrho(\lambda^{-1}).
  \]
  Therefore the whole bracket in
  \eqref{eq:class-product-normalisation-identity} is multiplied by
  \(|C|/|G|\).  Since \(K_C\) is the exponential of the negative of the
  resulting logarithmic constant, this is exactly
  \eqref{eq:class-euler-product-k-f}.  Removing the factor from the
  non-trivial summation would replace the true class-indicator coefficient
  \((|C|/|G|)\chi_\varrho^\ast(C)\) by
  \(\chi_\varrho^\ast(C)\), so it is not compatible with the stated
  Hermitian convention unless the affected non-trivial fixed-label linear
  combination vanishes accidentally.
\end{proof}

\begin{lemma}[Positive real branch for the character formula]
  \label{lem:class-product-positive-real-branch}
  \label{lem:class-euler-product-real-branch}
  Under the hypotheses of \Cref{thm:class-euler-product-mertens}, the
  character sum
  \begin{equation}\label{eq:class-euler-product-character-exponent}
    \sum_{\varrho\neq\mathbf1}
      \chi_\varrho^\ast(C)F_\varrho(\lambda^{-1})
  \end{equation}
  is real. Consequently the right-hand side of
  \eqref{eq:class-euler-product-k-f} is the ordinary positive real
  exponential and is equal to the positive Euler-product constant defined
  by \eqref{eq:class-euler-product-k-mc} and
  \eqref{eq:class-euler-product-def}.
\end{lemma}

\begin{proof}
  The fixed-label transform is linear in the class function on the open
  disc. Since \(\mathbf1_C\) and the trivial character are real-valued
  class functions, the class-indicator expansion from
  \Cref{prop:class-indicator-expansion} gives, for every \(0<t<1\),
  \[
    F_{\mathbf1_C}(t\lambda^{-1})
    -
    \frac{|C|}{|G|}F_{\mathbf1}(t\lambda^{-1})
    =
    \frac{|C|}{|G|}
      \sum_{\varrho\neq\mathbf1}
        \chi_\varrho^\ast(C)F_\varrho(t\lambda^{-1}).
  \]
  The two transforms on the left have real coefficients in their orbit
  sums, so their difference is real for real \(t\in(0,1)\), even though
  the two summands separately have the same trivial logarithmic
  divergence at the Perron point. Hence the non-trivial character sum is
  real for every \(0<t<1\). The non-trivial boundary values at \(t=1\)
  exist by the twisted gap and the branch convention in
  \Cref{thm:fixed-label-euler-transform}; taking the limit
  \(t\uparrow1\) proves that
  \eqref{eq:class-euler-product-character-exponent} is real. Equivalently,
  the terms may be grouped in contragredient pairs: the transform of the
  contragredient character is the complex conjugate transform on the real
  Abel path, and
  \(\chi_{\varrho^\vee}^\ast(C)=\overline{\chi_\varrho^\ast(C)}\).

  The proof of \Cref{thm:class-euler-product-mertens} identifies
  \eqref{eq:class-euler-product-k-f} with
  \eqref{eq:class-euler-product-k-mc}. In
  \eqref{eq:class-euler-product-k-mc} the exponent is real, since
  \(\mathcal M_C(A,\tau;G)\) and
  \(\Pi_{\mathbf1_C}(\lambda^{-r})\) are real orbit sums, and the
  exponential is positive. This is the same positive constant obtained as
  the limit of
  \(N^{|C|/|G|}P_C(N)\), where every factor
  \(1-\lambda^{-|\gamma|}\) in \eqref{eq:class-euler-product-def} is
  positive because \(\lambda>1\). Therefore the character formula is
  branch-compatible with the Euler product and uses no independent
  complex logarithm choice.
\end{proof}

\begin{lemma}[Sign and normalisation conventions for \texorpdfstring{$K_C$}{KC}]
  \label{lem:class-product-sign-normalisation-conventions}
  Under the hypotheses of \Cref{thm:class-euler-product-mertens}, the two
  displayed formulae for \(K_C\) in
  \eqref{eq:class-euler-product-k-mc} and
  \eqref{eq:class-euler-product-k-f} are identical with the following
  conventions: the Euler--Mascheroni term enters as \(+\gamma\) inside
  the bracket, the scalar reduced-determinant term enters as
  \(+\log C(A)\), the outer sign is the single minus sign coming from
  \(\log(1-x)=-\sum_{r\ge1}x^r/r\), the coefficient of every
  non-trivial character channel is
  \((|C|/|G|)\chi_\varrho^\ast(C)F_\varrho(\lambda^{-1})\), and
  the exponential is the ordinary positive real exponential.
\end{lemma}

\begin{proof}
  Starting from the Euler product, there is only one sign change:
  \[
    \log P_C(N)
    =-
    \sum_{n\le N}\frac{p_{n,C}}{\lambda^n}
    -
    \sum_{r\ge2}\frac1r\sum_{n\le N}p_{n,C}\lambda^{-rn}.
  \]
  The first term is evaluated by
  \eqref{eq:class-mertens-log} as
  \(-(|C|/|G|)\log N-\mathcal M_C(A,\tau;G)+O(N^{-1})\), and the second
  converges to
  \(-\sum_{r\ge2}\Pi_{\mathbf1_C}(\lambda^{-r})/r\).  Hence the constant in
  \(N^{|C|/|G|}P_C(N)\) is precisely the exponential of the negative of the
  left-hand side of
  \eqref{eq:class-product-constant-compatibility}.  No further sign is
  introduced when exponentiating.

  It remains to identify that left-hand side.  The Hermitian class-indicator
  expansion is
  \[
    \mathbf1_C=\frac{|C|}{|G|}\mathbf1+\frac{|C|}{|G|}
      \sum_{\varrho\ne\mathbf1}\chi_\varrho^\ast(C)\chi_\varrho .
  \]
  Therefore the non-trivial coefficient is
  \((|C|/|G|)\chi_\varrho^\ast(C)\), not
  \(\chi_\varrho^\ast(C)\) by itself.  Applying this expansion to the
  additive class-Mertens constant and to the absolutely convergent product
  tail gives
  \[
    \mathcal M_C(A,\tau;G)+
    \sum_{r\ge2}\frac{\Pi_{\mathbf1_C}(\lambda^{-r})}{r}
    =
    \frac{|C|}{|G|}\left(S_0+
      \sum_{\varrho\ne\mathbf1}\chi_\varrho^\ast(C)S_\varrho\right),
  \]
  where
  \[
    S_\varrho=
    \Pi_\varrho(\lambda^{-1})+
    \sum_{r\ge2}\frac{\Pi_\varrho(\lambda^{-r})}{r}
    =F_\varrho(\lambda^{-1})
  \]
  for \(\varrho\ne\mathbf1\).  For the scalar channel,
  \(\mathcal M(A)=\gamma+\FP(A)\) and
  \eqref{eq:finite-part-primitive-closed} give
  \[
  \begin{aligned}
    S_0
    &=
    \mathcal M(A)+\sum_{r\ge2}\frac{\Pi_{\mathbf1}(\lambda^{-r})}{r} \\
    &=\gamma+
      \log C(A)+
      \sum_{\gamma_0\in\cP}
      \left(
        x_{\gamma_0}+\log(1-x_{\gamma_0})+
        \sum_{r\ge2}\frac{x_{\gamma_0}^r}{r}
      \right),
  \end{aligned}
  \]
  with \(x_{\gamma_0}=\lambda^{-|\gamma_0|}\).  Since
  \(x+
  \log(1-x)+\sum_{r\ge2}x^r/r=0\), this scalar term is exactly
  \(\gamma+\log C(A)\).  Substitution gives
  \eqref{eq:class-product-constant-compatibility}, and the preceding sign
  computation turns it into \eqref{eq:class-euler-product-k-f}.  Finally,
  \Cref{lem:class-product-positive-real-branch} proves that the bracket is
  real and that the exponential is taken on the positive real branch, so the
  equality is branch-compatible.
\end{proof}

% END INPUT sec_chebotarev_class_mertens_equivalence.tex
% BEGIN INPUT sec_chebotarev_class_mertens_determinants.tex
\subsubsection{Determinant-only computation of fixed-label Perron values}
\label{subsubsec:determinant-only-fixed-label-values}

This unit is the computational part of Main Theorem~III.  It evaluates the
fixed-label numbers \(F_\varrho(\lambda^{-1})\) from the same finite twisted
determinant data used in Main Theorem~I.  The local branch convention is fixed
before the formula is written: all determinant logarithms are continued along
the Abel path from the open disc, the scalar Perron determinant is not present
in the non-trivial boundary coordinate, and the only omitted summand is the
singular triple \((r,m,\pi)=(1,1,\mathbf1)\).  This is not a replacement of
\(F_\varrho\) by the periodic Artin logarithm; it is the Adams-corrected
determinant expansion of the fixed-label transform.

\begin{lemma}[Singular-triple deletion and branch convention]
  \label{lem:fixed-label-singular-triple-deletion}
  Assume the hypotheses of \Cref{thm:class-mertens-explicit}, and fix a
  non-trivial irreducible representation \(\varrho\in\Irr(G)\).  For
  \(0<t<1\), the Abel-regularized fixed-label determinant expression is
  \begin{equation}\label{eq:fixed-label-singular-deletion-abel}
    F_\varrho(t\lambda^{-1})
    =
    \sum_{r,m\ge1}\frac{\mu(m)}{rm}
      \sum_{\substack{\pi\in\Irr(G)\\(r,m,\pi)\neq(1,1,\mathbf1)}}
        a_{\varrho,\pi}^{(m)}
        \Log\det(I-t^{rm}\lambda^{-rm}A_\pi)^{-1},
  \end{equation}
  where every logarithm is the normalized germ at the origin.  In the limit
  \(t\uparrow1\), the deleted triple \((r,m,\pi)=(1,1,\mathbf1)\) is the
  only scalar Perron singularity. Its coefficient is
  \(a_{\varrho,\mathbf1}^{(1)}=0\), so the singular logarithm is absent
  before the boundary limit is taken.  All remaining \(rm=1\) terms are
  non-trivial twisted branches regular at \(z=\lambda^{-1}\), and all
  \(rm\ge2\) terms, including scalar terms arising from Adams powers, are
  ordinary logarithms inside the Perron disc.  The deleted-series convention
  therefore has a unique Abel limit, and the convergence to that limit is
  justified by an absolute majorant independent of \(t\in(0,1)\).
\end{lemma}

\begin{proof}
  For \(0<t<1\), both the fixed-label Euler transform and every determinant
  logarithm are evaluated inside their discs of convergence. Substituting
  \eqref{eq:class-logarithm-determinant-linearisation} into
  \eqref{eq:fixed-label-euler-transform-adams-mobius} gives the same
  formula as \eqref{eq:fixed-label-singular-deletion-abel} with the scalar
  triple included.  Its coefficient is
  \[
    a_{\varrho,\mathbf1}^{(1)}
    =\langle\psi^1\chi_\varrho,\mathbf1\rangle_G
    =\langle\chi_\varrho,\mathbf1\rangle_G=0,
  \]
  by finite-character orthogonality, because \(\varrho\neq\mathbf1\).
  Thus deleting that summand changes no Abel value for \(0<t<1\), and no
  finite part is assigned to
  \(\Log\det(I-\lambda^{-1}A_{\mathbf1})^{-1}\).

  If \(rm=1\) and \(\pi\neq\mathbf1\), then \(r=m=1\) and the strict twisted
  spectral gap gives \(\rho(A_\pi)<\lambda\); hence
  \(I-	heta\lambda^{-1}A_\pi\) is invertible for \(0\le\theta\le1\), and
  the normalized germ has the stated Abel-path boundary value. If
  \(rm\ge2\), then for every \(\pi\in\Irr(G)\),
  \[
    \rho(\lambda^{-rm}A_\pi)
    \le \lambda^{1-rm}<1,
  \]
  so the normalized logarithm is the ordinary value of the convergent
  determinant power series at \(\lambda^{-rm}\).  The double tail after the
  finitely many \(rm=1\) terms is dominated by
  \Cref{lem:fixed-label-perron-double-majorant}, namely by a constant
  multiple of \(\sum_{rm\ge2}\lambda^{1-rm}/(rm)\).  This majorant is
  summable and independent of \(t\), and dominated convergence gives the
  unique Abel limit with exactly the branch convention stated above.
\end{proof}

\begin{proposition}[Determinant formula for fixed-label Perron values]
  \label{prop:fixed-label-determinant-perron}
  Under the hypotheses of \Cref{thm:class-mertens-explicit}, every
  non-trivial \(\varrho\in\Irr(G)\) satisfies
  \begin{equation}\label{eq:fixed-label-determinant-perron}
    F_\varrho(\lambda^{-1})
    =
    \sum_{r,m\ge1}\frac{\mu(m)}{rm}
      \sum_{\substack{\pi\in\Irr(G)\\(r,m,\pi)\neq(1,1,\mathbf1)}}
        a_{\varrho,\pi}^{(m)}
        \log\det(I-\lambda^{-rm}A_\pi)^{-1}.
  \end{equation}
  The branch convention is local to this proposition. For each displayed
  determinant, start with the analytic germ
  \(\Log\det(I-zA_\pi)^{-1}=0\) at \(z=0\).  If \(rm=1\) and
  \(\pi\neq\mathbf1\), its value is the Abel-path continuation to
  \(z=\lambda^{-1}\) along \(z=t\lambda^{-1}\),
  \(0\le t\uparrow1\), through the zero-free region supplied by the
  twisted spectral gap.  If \(rm\ge2\), the value is the ordinary value of
  the same normalized germ at \(z=\lambda^{-rm}\), where
  \(\rho(\lambda^{-rm}A_\pi)<1\).
  If \(\varrho\neq\mathbf1\), an Adams power \(\psi^m\chi_\varrho\) may
  contain the trivial character for \(m\ge2\).  The corresponding scalar
  determinant terms in \eqref{eq:fixed-label-determinant-perron} are part
  of the formula whenever \(rm\ge2\); they are regular ordinary
  determinant logarithms because \(\lambda^{-rm}A_{\mathbf1}\) is evaluated
  strictly inside the Perron disc.
  For non-real irreducible representations the coefficients
  \(a_{\varrho,\pi}^{(m)}\) and the determinant terms may be complex.  No
  termwise realness is asserted; real class constants result from pairing
  conjugate irreducibles in the Hermitian class expansion.
  The summation sign is a local exclusion: the scalar boundary term
  \((r,m,\pi)=(1,1,\mathbf1)\), and only that term, is not part of the
  expression at \(z=\lambda^{-1}\). It would be
  \(a_{\varrho,\mathbf1}^{(1)}
  \log\det(I-\lambda^{-1}A_{\mathbf1})^{-1}\), but
  \(a_{\varrho,\mathbf1}^{(1)}=0\), so no scalar Perron logarithm occurs
  in the non-trivial fixed-label value.
\end{proposition}

\begin{proof}
  Apply \Cref{lem:fixed-label-singular-triple-deletion} to the Abel identity
  \eqref{eq:fixed-label-singular-deletion-abel}.  That lemma makes the local
  deletion before the boundary limit explicit: the only scalar Perron
  logarithm is \((r,m,\pi)=(1,1,\mathbf1)\), its coefficient is zero by
  non-trivial-channel orthogonality, and every scalar term with \(rm\ge2\)
  lies strictly inside the Perron disc.  The same lemma supplies the branch
  convention for the non-trivial \(rm=1\) terms and the absolute majorant for
  the \(rm\ge2\) double tail.  Dominated convergence as \(t\uparrow1\),
  together with the existence of the non-trivial Perron value in
  \Cref{thm:fixed-label-euler-transform}, proves
  \eqref{eq:fixed-label-determinant-perron} with exactly the displayed
  exclusion and branch choices.
\end{proof}

\begin{corollary}[Convention audit for Main Theorem III]
  \label{cor:main-iii-convention-audit}
  Under the hypotheses of \Cref{thm:class-euler-product-mertens}, the
  determinant form of \eqref{eq:class-euler-product-k-f} has the following
  forced conventions.
  \begin{enumerate}[label=\textup{(\roman*)}]
    \item The sign of \(K_C\) is the single outer minus sign in
    \eqref{eq:class-euler-product-k-f}; equivalently \(K_C\) is the
    positive real exponential of the negative logarithmic constant in
    \eqref{eq:class-product-constant-compatibility}.
    \item The class coefficient is
    \((|C|/|G|)\chi_\varrho^\ast(C)\), where
    \(\chi_\varrho^\ast(C)=\overline{\chi_\varrho(C)}\); inverse recovery of
    Fourier coordinates uses \(\chi_\sigma(C)\) on the summing side, not a
    second conjugated coefficient.
    \item In the determinant expansion of each
    \(F_\varrho(\lambda^{-1})\), the single scalar boundary triple
    \((r,m,\pi)=(1,1,\mathbf1)\) is omitted before evaluation and is not
    replaced by a finite part.
    \item Non-trivial \(rm=1\) determinant logarithms are Abel-path boundary
    values of the normalized germ at the origin, while all \(rm\ge2\) terms
    are ordinary interior-disc values of the same germ.
    \item The determinant double tail is absolutely convergent at the Perron
    boundary after this local omission, so the order of the finite character
    sums and the Abel limit used in \eqref{eq:class-euler-product-k-f} is
    justified by dominated convergence.
  \end{enumerate}
\end{corollary}

\begin{proof}
  Item \textup{(i)} is exactly
  \Cref{lem:class-product-sign-normalisation-conventions}, which derives the
  sign from the identity
  \(\log(1-x)=-\sum_{r\ge1}x^r/r\) and from the scalar normalization
  \(\mathcal M(A)+\sum_{r\ge2}\Pi_{\mathbf1}(\lambda^{-r})/r
  =\gamma+\log C(A)\).  Item \textup{(ii)} is
  \Cref{lem:class-product-hermitian-normalisation} together with the
  Fourier convention in \Cref{sec:chebotarev}: class indicators expand with
  the forward coefficient \(\chi_\varrho^\ast(C)\), and orthogonality recovers
  the \(\sigma\)-coordinate by multiplying by \(\chi_\sigma(C)\) and summing
  over conjugacy classes.

  Items \textup{(iii)} and \textup{(iv)} are the content of
  \Cref{lem:fixed-label-singular-triple-deletion}: the coefficient of the
  omitted scalar boundary triple is
  \(a_{\varrho,\mathbf1}^{(1)}=\langle\chi_\varrho,\mathbf1\rangle_G=0\),
  non-trivial boundary determinants are regular because
  \(\rho(A_\pi)<\lambda\), and terms with \(rm\ge2\) satisfy
  \(\rho(\lambda^{-rm}A_\pi)\le\lambda^{1-rm}<1\).  The same lemma supplies
  the uniform absolute majorant for the double tail, and hence item
  \textup{(v)}.  Substituting the branch-compatible determinant values from
  \Cref{prop:fixed-label-determinant-perron} into
  \eqref{eq:class-euler-product-k-f} gives precisely the determinant form of
  Main Theorem~III, with no remaining sign, conjugation, branch, deletion, or
  convergence ambiguity.
\end{proof}

\begin{proposition}[Finite algorithm for fixed-label Perron constants]
  \label{prop:fixed-label-perron-finite-algorithm}
  Under the hypotheses of \Cref{thm:class-mertens-explicit}, fix a
  non-trivial irreducible representation \(\varrho\in\Irr(G)\) and a tolerance
  \(\epsilon>0\).  There is a deterministic finite algorithm which, from the
  input in \textup{(a)} below and with the branch convention in
  \textup{(c)}, returns a finite determinant sum
  \(F_\varrho^{(M)}\) and a certified radius \(E_M\le\epsilon\) such that
  \[
    F_\varrho(\lambda^{-1})
    \in F_\varrho^{(M)}+\{w\in\CC:|w|\le E_M\}.
  \]
  Consequently every Frobenius-class product constant in Main Theorem~III is
  certified from finitely many twisted determinant polynomials, finite
  character-table operations, the scalar normalisation \(C(A)\), and the stated
  Abel branch; no primitive-orbit enumeration or limiting orbit count is an
  input.  The algorithm is the following.
  \begin{enumerate}[label=\textup{(\alph*)},leftmargin=*]
    \item \emph{Input.}  The determinant inputs are finite: the character table of \(G\),
      the class-power maps \(C\mapsto C^m\) on conjugacy classes, the finite
      determinant polynomials \(D_\pi(z):=\det(I-zA_\pi)\),
      \(\pi\in\Irr(G)\), the Perron root \(\lambda>1\), the twisted
      spectral-gap certificate \(\eta<\lambda\), and an explicit majorant
      constant \(C_\varrho\) furnished by
      \Cref{lem:fixed-label-perron-double-majorant} for
      \(\phi=\chi_\varrho\).
    \item \emph{Finite coefficients.}  For each \(m\ge1\), compute the Adams coefficients
      \begin{equation}\label{eq:fixed-label-algorithm-adams-coefficients}
        a_{\varrho,\pi}^{(m)}
        =\bigl\langle \psi^m\chi_\varrho,\chi_\pi\bigr\rangle_G,
        \qquad
        (\psi^m\chi_\varrho)(g)=\chi_\varrho(g^m),
      \end{equation}
      by the finite character-table inner product.  These coefficients may be
      complex for non-real irreducible pairs; the final class constants are
      real by conjugation symmetry of the full character sum, not by termwise
      realness.
    \item \emph{Branch and deletion rule.}  Omit the single scalar Perron term
      \((r,m,\pi)=(1,1,\mathbf1)\) before any logarithm is evaluated.  Its
      coefficient is
      \(a_{\varrho,\mathbf1}^{(1)}=0\), so the algorithm never assigns a
      finite part or a branch value to
      \(\log\det(I-\lambda^{-1}A_{\mathbf1})^{-1}\).
      Normalize every remaining logarithm from the germ
      \(\Log\det(I-zA_\pi)^{-1}=0\) at \(z=0\).  For \(rm=1\) this means
      Abel continuation along \(z=t\lambda^{-1}\), \(0\le t\uparrow1\), and
      necessarily \(\pi\neq\mathbf1\).  For \(rm\ge2\) it means the ordinary
      value of the same germ at \(z=\lambda^{-rm}\), where the determinant
      is strictly inside its Perron radius.
    \item \emph{Truncation rule.}  Choose \(M\) so that the computable tail
      \begin{equation}\label{eq:fixed-label-perron-finite-algorithm-tail-choice}
        C_\varrho
        \sum_{k>M}\frac{\mathrm d(k)}{k}\lambda^{1-k}
        \le \epsilon/2 .
      \end{equation}
      This finite search terminates because \(\lambda>1\).  Form the finite
      truncation
      \begin{equation}\label{eq:fixed-label-perron-finite-algorithm-truncation}
        F_\varrho^{(M)}
        :=
        \sum_{\substack{r,m\ge1\\rm\le M}}
        \frac{\mu(m)}{rm}
        \sum_{\substack{\pi\in\Irr(G)\\(r,m,\pi)\neq(1,1,\mathbf1)}}
          a_{\varrho,\pi}^{(m)}
          \Log\det(I-\lambda^{-rm}A_\pi)^{-1},
      \end{equation}
      where the scalar Perron term \((1,1,\mathbf1)\) has already been
      omitted, the only possible \(rm=1\) terms use the non-trivial Abel
      branch from \textup{(c)}, and all \(rm\ge2\) terms use the ordinary germ value
      inside the Perron disc.  The truncation error is the effective tail
      estimate
      \begin{equation}\label{eq:fixed-label-perron-finite-algorithm-error}
        \abs{F_\varrho(\lambda^{-1})-F_\varrho^{(M)}}
        \le
        C_\varrho
        \sum_{k>M}\frac{\mathrm d(k)}{k}\lambda^{1-k},
      \end{equation}
      for a constant \(C_\varrho\) depending only on the finite determinant
      collection and \(\varrho\).  For an exact certificate one may replace the
      right-hand side by any rational upper bound.  Concretely, if rational
      numbers \(\theta<1\) and \(Q_\varrho\) satisfy
      \(\lambda^{-1}<\theta\) and \(C_\varrho\lambda<Q_\varrho\), then
      \begin{equation}\label{eq:fixed-label-rational-tail-envelope}
        C_\varrho
        \sum_{k>M}\frac{\mathrm d(k)}{k}\lambda^{1-k}
        \le
        Q_\varrho\frac{\theta^{M+1}}{1-\theta},
      \end{equation}
      because \(\mathrm d(k)\le k\).  Thus the truncations converge absolutely
      and exponentially, and the same finite procedure gives certified
      interval enclosures by truncating after any prescribed rational tail
      bound.
    \item \emph{Output guarantee.}  If
      the finitely many logarithms occurring in
      \eqref{eq:fixed-label-perron-finite-algorithm-truncation} are each
      enclosed on the specified branch with radius at most \(\epsilon_0\), then
      the total output radius is bounded by the tail in
      \eqref{eq:fixed-label-perron-finite-algorithm-error} plus
      \begin{equation}\label{eq:fixed-label-perron-finite-algorithm-roundoff}
        \epsilon_0
        \sum_{\substack{r,m\ge1\\rm\le M}}
          \frac{|\mu(m)|}{rm}
          \sum_{\pi\in\Irr(G)}|a_{\varrho,\pi}^{(m)}| .
      \end{equation}
      Choosing \(\epsilon_0\) so that
      \eqref{eq:fixed-label-perron-finite-algorithm-roundoff} is at most
      \(\epsilon/2\) gives a certified enclosure of radius at most
      \(\epsilon\) by finite interval arithmetic.  This certified disc is the
      algorithmic output.
  \end{enumerate}
\end{proposition}
\begin{proof}
  The proof shows that this finite record is a reproducibility form of the
  Adams-corrected determinant formula, not an additional analytic input.

  Items (a) and (b) are finite because \(G\) has finitely many conjugacy
  classes and irreducible representations, and
  \eqref{eq:fixed-label-algorithm-adams-coefficients} is the standard
  finite character inner product defining the Adams multiplicities.  Item
  (c) has two parts.  The deletion rule follows from
  \(a_{\varrho,\mathbf1}^{(1)}=\langle\chi_\varrho,\mathbf1\rangle_G=0\)
  for \(\varrho\neq\mathbf1\); this is the omitted singular term rule in
  \Cref{prop:fixed-label-determinant-perron}.  The branch rule is exactly the
  branch normalization stated there.  The twisted spectral gap makes the
  non-trivial \(rm=1\) determinants invertible along the Abel path, and
  \(rm\ge2\) gives \(\rho(\lambda^{-rm}A_\pi)\le\lambda^{1-rm}<1\), so the
  ordinary germ value is defined by its absolutely convergent logarithmic
  power series.

  With these conventions, \eqref{eq:fixed-label-determinant-perron} is the
  limit of the finite truncations
  \eqref{eq:fixed-label-perron-finite-algorithm-truncation}.  The tail with
  \(rm>M\) is bounded by
  \Cref{lem:fixed-label-perron-double-majorant}, applied to
  \(\phi=\chi_\varrho\), and regrouping the pairs \((r,m)\) by
  \(k=rm\) gives
  \[
    C_\varrho
    \sum_{\substack{r,m\ge1\\rm>M}}\frac{\lambda^{1-rm}}{rm}
    =
    C_\varrho\sum_{k>M}\frac{\mathrm d(k)}{k}\lambda^{1-k}.
  \]
  This proves \eqref{eq:fixed-label-perron-finite-algorithm-error}.  Since
  the divisor function has subexponential growth and \(\lambda>1\), the
  displayed tail tends to zero exponentially; hence the search in
  \eqref{eq:fixed-label-perron-finite-algorithm-tail-choice} terminates and
  the finite determinant truncations converge to the Perron value.  The
  rational envelope \eqref{eq:fixed-label-rational-tail-envelope} follows by
  choosing rational \(\theta\) and \(Q_\varrho\) with the stated strict
  inequalities, using \(\lambda^{1-k}=\lambda(\lambda^{-1})^k\), and summing
  the resulting geometric series.  The inequalities defining \(\theta\) and
  \(Q_\varrho\) are finite comparisons in the algebraic determinant data.

  It remains only to justify the stated certificate format.  For fixed
  \(M\), the set of triples \((r,m,\pi)\) with \(rm\le M\) is finite, all
  Adams coefficients in it are computed by finite character-table sums, and
  each determinant logarithm is evaluated on the branch fixed in item (c).
  Replacing each such logarithm by an interval of radius \(\epsilon_0\) changes
  the weighted finite sum by at most the triangle-inequality bound in
  \eqref{eq:fixed-label-perron-finite-algorithm-roundoff}.  Choosing
  \(\epsilon_0\) so that this finite bound is at most \(\epsilon/2\), and
  choosing \(M\) so that the analytic tail is at most \(\epsilon/2\), gives a
  rigorous output interval of radius at most \(\epsilon\).  Thus the input,
  branch convention, finite output, and error bound are all finite and
  explicit.
\end{proof}

\medskip
\noindent\textit{Appendix-level reproducibility records.}
  The local rational reproducibility record and five-field \(S_3\) certificate format are
  placed in \Cref{prop:fixed-label-rational-certificate-record,cor:fixed-label-five-field-certificate-schema}.
  They are appendix-level certificate formats for the finite determinant
  computability theorem \Cref{prop:fixed-label-perron-finite-algorithm}; the
  main proof path uses only the determinant formula, the finite truncation
  bound, and finite character-table Fourier inversion.

\medskip
\noindent\textit{Product-level finite consequence.}
  The finite reproducibility record of
  \Cref{prop:fixed-label-perron-finite-algorithm}
  is finite mathematical support for the determinant-coordinate values
  used in Main Theorem~III.  Fix a non-trivial irreducible character
  \(\varrho\), a truncation window \(M\), rational branch rectangles for every
  non-deleted triple \((r,m,\pi)\) with \(rm\le M\), the Adams coefficients
  \(a_{\varrho,\pi}^{(m)}\), and rational radii \(R_M,R_{\mathrm{num}}\) as in
  \Cref{prop:fixed-label-perron-finite-algorithm}.  If these finite records pass
  the exact character-table, branch-tag, interval-propagation, and tail-bound
  checks of that proposition, then they certify the explicit enclosure
  \[
    F_\varrho(\lambda^{-1})
    \in
    I_M+
    \{w\in\CC:|w|\le R_M+R_{\mathrm{num}}\},
  \]
  where \(I_M\) is the rational interval obtained from the finite determinant
  sum.  Applying this finite certificate for all
  \(\varrho\ne\mathbf1\) and then applying the finite Fourier inverse of
  \Cref{cor:class-product-constant-rigidity} certifies rational enclosures for
  the logarithms of all class-product constants \(K_C\).  Thus the displayed
  finite records prove the following four mathematical assertions and no
  broader certified-computation theorem.
  \begin{enumerate}[label=\textup{(\roman*)},leftmargin=*]
    \item The accepted finite record for \(\varrho\) determines a rigorous
      rational enclosure for the single fixed-label Perron value
      \(F_\varrho(\lambda^{-1})\), with the scalar Perron logarithm deleted
      before any boundary value is taken.
    \item The collection of accepted records over
      \(\varrho\ne\mathbf1\) determines, by finite Fourier inversion and the
      fixed scalar normalisation \(C(A)\), rigorous enclosures for the
      class-product constants of \Cref{thm:class-euler-product-mertens}.
    \item The \(S_3\) appendix records instantiate this finite certificate as a
      witness to the non-abelian Artin-substitution obstruction; they add no
      new analytic hypothesis to Main Theorem~III.
    \item Nothing in the finite certificate asserts a general-purpose certified
      computation package for arbitrary finite groups, arbitrary symbolic
      matrices, arbitrary numerical branch tracking, or arbitrary determinant
      identities beyond the determinant formula proved in this article.
  \end{enumerate}

\begin{proof}
  The displayed enclosure is exactly
  \eqref{eq:fixed-label-certificate-interface-inclusion} from
  \Cref{prop:fixed-label-perron-finite-algorithm}.  That proposition proves that
  the four
  finite records enclose the truncated determinant sum and its discarded tail,
  and it also proves that the omitted scalar Perron triple has coefficient
  zero for \(\varrho\ne\mathbf1\).  Thus (i) is a finite consequence of the
  determinant formula in \Cref{prop:fixed-label-determinant-perron}; no
  limiting orbit count, unrecorded branch choice, or numerical assumption is
  introduced.

  Applying the same accepted-record argument to every non-trivial irreducible
  character gives enclosures for the full non-trivial vector
  \((F_\varrho(\lambda^{-1}))_{\varrho\ne\mathbf1}\).  The inverse and forward
  transforms in \Cref{cor:class-product-constant-rigidity} are finite
  character-table operations after \(C(A)\) has fixed the scalar coordinate.
  Therefore those enclosures transport to the logarithms of the class-product
  constants in \Cref{thm:class-euler-product-mertens}, proving (ii).

  Statement (iii) follows from
  \Cref{cor:fixed-label-five-field-certificate-schema}: the appendix files are
  local finite arithmetic transcripts for the same five fields.  The
  concrete \(S_3\) obstruction in
  \Cref{prop:s3-naive-artin-obstruction,lem:s3-example-perron-certificates,lem:appendix-s3-f-level-certificates}
  uses those transcripts only after Main Theorem~III has supplied the analytic
  determinant-coordinate formula.  Hence the witness illustrates and certifies
  the obstruction but does not enlarge the theorem hypotheses.

  For (iv), inspect the inputs permitted in
  \Cref{prop:fixed-label-perron-finite-algorithm}: the group, character table,
  power maps, determinant polynomials, Perron root, strict-gap certificate,
  branch convention, majorant, truncation window, and rational interval bounds
  are all fixed finite records.  The proof of
  \Cref{prop:fixed-label-perron-finite-algorithm} checks only algebraic equalities,
  rational inequalities, interval containments, and the tail bound already
  derived from \Cref{lem:fixed-label-perron-double-majorant}.  These checks are
  sufficient for the finite records used here, but they do not construct a
  general parser, branch tracker, interval-arithmetic library, or universal
  certification theory for all finite groups and matrices.  Thus the result is
  a finite determinant certificate for the constants in this paper, not a
  standalone certified-computation framework.
\end{proof}

\subsubsection{Product-vector rigidity and finite observation budget}
\label{subsubsec:class-product-vector-rigidity}

After the determinant-only evaluation of \(F_\varrho(\lambda^{-1})\), the
remaining statements in this unit concern the finite class vector
\((K_C)_C\).  They are rigidity and observation results for the already
constructed product constants: they record the scalar hyperplane constraint,
Fourier inversion on conjugacy classes, and the exact number of class
coordinates needed to recover the whole positive product vector.  They do not
assert recovery of the cocycle or any branch-free inverse spectral datum.

\begin{corollary}[Rigidity of the Adams-corrected class-product constants]
  \label{cor:class-product-constant-rigidity}
  Assume the hypotheses of \Cref{thm:class-mertens-explicit}.  For each
  conjugacy class \(C\subset G\), let \(K_C\) be the positive constant in
  \Cref{thm:class-euler-product-mertens}, and set
  \begin{equation}\label{eq:class-product-normalised-log-vector}
    \mathfrak K(C)
    :=
    -\frac{|G|}{|C|}\log K_C-\bigl(\gamma+\log C(A)\bigr).
  \end{equation}
  Then, as a structural consequence of the Adams-corrected constant formula
  and finite character Fourier inversion, the class-product constants are
  not a free family of classwise constants.  After the scalar normalisation
  \(C(A)\), they are exactly equivalent to the non-trivial fixed-label
  Perron-boundary vector:
  \begin{equation}\label{eq:class-product-rigidity-forward}
    \mathfrak K(C)
    =
    \sum_{\varrho\neq\mathbf1}
      \chi_\varrho^\ast(C)F_\varrho(\lambda^{-1}),
  \end{equation}
  the normalised logarithmic vector lies in the non-trivial Peter--Weyl
  summand,
  \begin{equation}\label{eq:class-product-rigidity-hyperplane}
    \sum_{C\subset G}\frac{|C|}{|G|}\mathfrak K(C)=0,
    \qquad
    \prod_{C\subset G}K_C=
      \exp\!\bigl(-\gamma-\log C(A)\bigr),
  \end{equation}
  and the Fourier recovery is norm-preserving:
  \begin{equation}\label{eq:class-product-rigidity-parseval}
    \sum_{C\subset G}\frac{|C|}{|G|}\abs{\mathfrak K(C)}^2
    =
    \sum_{\varrho\neq\mathbf1}
      \abs{F_\varrho(\lambda^{-1})}^2.
  \end{equation}
  Moreover, for every non-trivial \(\pi\in\Irr(G)\),
  \begin{equation}\label{eq:class-product-rigidity-inverse}
    F_\pi(\lambda^{-1})
    =
    \sum_{C\subset G}\frac{|C|}{|G|}\mathfrak K(C)\chi_\pi(C).
  \end{equation}
  Consequently the finite determinant data
  \(\{\det(I-zA_\pi):\pi\in\Irr(G)\}\), the finite character table and
  Adams tensors, the normalized Abel branches, and the scalar normalisation
  \(C(A)\) determine the whole Frobenius-class Euler-product vector
  \((K_C)_C\). Conversely, within the same listed finite input data and branch
  convention, the vector \((K_C)_C\) and \(C(A)\) determine all fixed-label
  values \(F_\pi(\lambda^{-1})\), \(\pi\neq\mathbf1\). These are the
  fixed-label Euler-transform values, not the periodic Artin logarithms
  \(\mathcal L_\pi(\lambda^{-1})\). Thus, subject to the single scalar
  constraint in \eqref{eq:class-product-rigidity-hyperplane}, no two
  non-trivial fixed-label Perron-boundary vectors can give the same
  normalised class-product constants.  The converse just described is only
  a coordinate expression in the fixed determinant model; equality of these
  constants is not asserted to imply equality, conjugacy, or cohomology of
  the underlying cocycles or extensions.
\end{corollary}

\begin{proof}
  This is the Fourier-inversion corollary of the constant formula in
  \Cref{thm:class-euler-product-mertens}; it is not an additional
  Chebotarev or Mertens asymptotic.
  Since each \(K_C\) is positive by
  \Cref{lem:class-euler-product-real-branch}, the logarithm in
  \eqref{eq:class-product-normalised-log-vector} is the ordinary real
  logarithm.  Taking the logarithm of
  \eqref{eq:class-euler-product-k-f} gives
  \[
    -\frac{|G|}{|C|}\log K_C-\bigl(\gamma+\log C(A)\bigr)
    =
    \sum_{\varrho\neq\mathbf1}
      \chi_\varrho^\ast(C)F_\varrho(\lambda^{-1}),
  \]
  which is \eqref{eq:class-product-rigidity-forward}.  Since the trivial
  character is orthogonal to every non-trivial irreducible character,
  summing \eqref{eq:class-product-rigidity-forward} with weights
  \(|C|/|G|\) gives the first identity in
  \eqref{eq:class-product-rigidity-hyperplane}.  Substituting the
  definition of \(\mathfrak K(C)\) then gives
  \begin{align*}
    0
    &=
    \sum_{C\subset G}\frac{|C|}{|G|}
      \left(-\frac{|G|}{|C|}\log K_C
        -\bigl(\gamma+\log C(A)\bigr)\right) \\
    &=-\sum_{C\subset G}\log K_C-\gamma-\log C(A),
  \end{align*}
  which is the product constraint in
  \eqref{eq:class-product-rigidity-hyperplane}.  Applying Parseval's
  identity in the orthonormal character basis of class functions gives
  \eqref{eq:class-product-rigidity-parseval}.

  Multiplying \eqref{eq:class-product-rigidity-forward}
  identity by \(|C|\chi_\pi(C)/|G|\), summing over conjugacy classes, and
  using the Hermitian character orthogonality relation
  \[
    \sum_{C\subset G}\frac{|C|}{|G|}
      \chi_\varrho^\ast(C)\chi_\pi(C)=\delta_{\varrho\pi}
  \]
  gives \eqref{eq:class-product-rigidity-inverse} for every
  \(\pi\neq\mathbf1\).

  By \Cref{prop:fixed-label-determinant-perron}, each non-trivial value
  \(F_\pi(\lambda^{-1})\) is obtained from the finite twisted determinant
  family by the branch-compatible Adams--M\"obius determinant formula,
  while the scalar contribution is the single normalisation \(C(A)\).
  Substituting those determinant values into
  \eqref{eq:class-product-rigidity-forward} and then into
  \eqref{eq:class-euler-product-k-f} determines every \(K_C\).  The
  inverse formula \eqref{eq:class-product-rigidity-inverse} proves the
  converse determination from the normalised product constants and shows
  injectivity of the fixed-label vector.  This proves the rigidity
  statement.
\end{proof}

\begin{proposition}[Coordinate observation budget for the product vector]
  \label{prop:class-product-coordinate-observation-budget}
  Assume the hypotheses of \Cref{thm:class-mertens-explicit}, and let
  \(\mathcal C(G)\) be the set of conjugacy classes of \(G\), with
  \(q=|\mathcal C(G)|\).  Put
  \(\alpha_A:=\gamma+\log C(A)\), and let \((K_C)_{C\in\mathcal C(G)}\) be
  the positive product constants of \Cref{thm:class-euler-product-mertens}.
  Inside the scalar-normalised product hyperplane
  \begin{equation}\label{eq:q-minus-one-product-hyperplane}
    \mathcal P_{\alpha_A}:=
    \left\{(L_C)_{C\in\mathcal C(G)}:\ L_C>0,
      \prod_{C\in\mathcal C(G)}L_C=e^{-\alpha_A}\right\},
  \end{equation}
  the following elementary coordinate statement holds.
  \begin{enumerate}[label=\textup{(\roman*)},leftmargin=*]
    \item If \(S\subset\mathcal C(G)\) has complement of size at most one,
      then the observed subvector \((K_C)_{C\in S}\), together with the scalar
      normalisation \(C(A)\), determines the whole product vector
      \((K_C)_C\) and hence determines every non-trivial fixed-label value
      \(F_\pi(\lambda^{-1})\) by
      \eqref{eq:class-product-rigidity-inverse}.
    \item If \(|S|\le q-2\), then the projection
      \(\mathcal P_{\alpha_A}\to(0,\infty)^S\),
      \((L_C)_C\mapsto(L_C)_{C\in S}\), is not injective.  Consequently no
      record consisting only of the scalar normalisation and those
      \(|S|\) class-product constants can determine, even inside the fixed
      Main Theorem~III coordinate chart, all scalar-normalised product
      constants or all non-trivial fixed-label Perron values.
  \end{enumerate}
  Thus the classwise constant vector has the following sharp local observation budget:
  after the scalar coordinate has been fixed, any \(q-1\) class constants are
  enough, while \(q-2\) are never enough.  This is only a statement about the
  branch-compatible product-coordinate vector recovered in
  \Cref{thm:class-euler-product-mertens,cor:class-product-constant-rigidity};
  it strengthens the finite-coordinate consequences of Main Theorem~III, but
  it is not a reconstruction theorem for the underlying cocycle, labelled
  extension, or determinant presentation.
\end{proposition}

\begin{proof}
  The product constraint in
  \eqref{eq:class-product-rigidity-hyperplane} says precisely that the
  dynamical vector \((K_C)_C\) lies in \(\mathcal P_{\alpha_A}\).  If
  \(S=\mathcal C(G)\), there is nothing to recover.  If
  \(\mathcal C(G)\setminus S=\{C_0\}\), then any vector
  \((L_C)_C\in\mathcal P_{\alpha_A}\) satisfies the explicit recovery formula
  \begin{equation}\label{eq:q-minus-one-missing-class-recovery}
    L_{C_0}=e^{-\alpha_A}\prod_{C\in S}L_C^{-1}.
  \end{equation}
  Therefore the observed constants on \(S\), together with \(\alpha_A
  =\gamma+\log C(A)\), determine the whole positive product vector.  Applying
  \Cref{cor:class-product-constant-rigidity}, in particular the inverse
  Fourier formula \eqref{eq:class-product-rigidity-inverse}, then determines
  each non-trivial fixed-label Perron value.

  Conversely suppose \(|S|\le q-2\).  Choose two distinct omitted classes
  \(C_1,C_2\in\mathcal C(G)\setminus S\).  For any
  \((L_C)_C\in\mathcal P_{\alpha_A}\) and any real \(t>0\), define
  \((L_C^{(t)})_C\) by
  \[
    L_C^{(t)}=
    \begin{cases}
      tL_{C_1},& C=C_1,\\
      t^{-1}L_{C_2},& C=C_2,\\
      L_C,& C\notin\{C_1,C_2\}.
    \end{cases}
  \]
  The product of all coordinates is unchanged, so
  \((L_C^{(t)})_C\in\mathcal P_{\alpha_A}\).  The projection to the classes in
  \(S\) is also unchanged.  If \(t\ne1\), the full positive product vector is
  different.  Hence the projection to \(S\) is not injective.

  For the last assertion, write the scalar-normalised logarithmic profile of
  a vector \((L_C)_C\in\mathcal P_{\alpha_A}\) as
  \(H_L(C)=-|G|\log L_C/|C|-\alpha_A\).  The product constraint gives
  \(\sum_C |C|H_L(C)/|G|=0\), so by the character orthogonality used in
  \Cref{cor:class-product-constant-rigidity} the non-trivial Fourier
  coordinates of \(H_L\) are uniquely recovered from the full vector and are
  exactly the fixed-label coordinates for the dynamical vector.  Since the
  construction above gives two different full vectors with the same observed
  coordinates on \(S\), those observations cannot determine the full
  scalar-normalised profile.  If they determined all non-trivial fixed-label
  coordinates, Fourier inversion would determine that profile and hence the
  full vector, contradicting the non-injectivity just proved.
\end{proof}

\begin{proposition}[Secondary-edge refinements do not alter the main constants]
  \label{prop:main-chain-secondary-noninterference}
  Assume the hypotheses of \Cref{thm:class-mertens-explicit}.  The data in
  Main Theorems~II--III are closed under adding any later secondary spectral
  edge analysis in the following precise sense.
  Let an auxiliary refinement consist of finitely many complex numbers of modulus
  strictly smaller than \(\lambda\), arranged in any class-function
  exponential sums \(S_C(n)\), and suppose only that the weighted tails
  \[
    T_C(N):=\sum_{n>N}\frac{S_C(n)}{n\lambda^n}
  \]
  converge and satisfy \(T_C(N)\to0\) for every conjugacy class \(C\).  If a
  refined counting formula has the form
  \begin{equation}\label{eq:secondary-noninterference-refined-count}
    p_{n,C}
    =
    \frac{|C|}{|G|}\frac{\lambda^n}{n}
    +\frac{S_C(n)}{n}+r_C(n)
  \end{equation}
  and its compensated summatory form is
  \begin{equation}\label{eq:secondary-noninterference-refined-sum}
    \sum_{n\le N}\frac{p_{n,C}}{\lambda^n}
    =
    \frac{|C|}{|G|}H_N
    +\mathcal M_C(A,\tau;G)-\frac{|C|}{|G|}\gamma
    -T_C(N)+o(1),
  \end{equation}
  then the refinement contributes no new Perron-boundary coordinate and
  no new Frobenius-class product constant.  More explicitly, the constants
  remain exactly
  \begin{equation}\label{eq:secondary-noninterference-mc}
    \mathcal M_C(A,\tau;G)
    =
    \frac{|C|}{|G|}\mathcal M(A)
    +\frac{|C|}{|G|}
      \sum_{\varrho\neq\mathbf1}
        \chi_\varrho^\ast(C)\Pi_\varrho(\lambda^{-1})
  \end{equation}
  and
  \begin{equation}\label{eq:secondary-noninterference-k}
    K_C=
    \exp\!\left(
      -\frac{|C|}{|G|}
      \left(
        \gamma+\log C(A)
        +\sum_{\varrho\neq\mathbf1}
          \chi_\varrho^\ast(C)F_\varrho(\lambda^{-1})
      \right)
    \right),
  \end{equation}
  with the positive branch of \Cref{thm:class-euler-product-mertens}.  Hence
  the finite determinant data
  \(\{\det(I-zA_\pi):\pi\in\Irr(G)\}\), the scalar normalisation \(C(A)\),
  and the Adams--M\"obius fixed-label transform determine the entire main
  product vector before any secondary-edge hypotheses are introduced.  In
  particular, any two auxiliary refinements satisfying the displayed tail
  condition give the same \((\mathcal M_C)_C\), the same \((K_C)_C\), and the
  same recovered vector \((F_\varrho(\lambda^{-1}))_{\varrho\neq\mathbf1}\).
\end{proposition}

\begin{proof}
  The asserted non-interference is a limit comparison.  The formulas
  \eqref{eq:secondary-noninterference-mc} and
  \eqref{eq:secondary-noninterference-k} are exactly the class-Mertens and
  class-product formulas already proved in
  \Cref{thm:class-mertens-explicit,thm:class-euler-product-mertens}.  In the
  refined summatory identity \eqref{eq:secondary-noninterference-refined-sum}
  the auxiliary contribution is the single compensated tail \(-T_C(N)\).
  Since the hypotheses require \(T_C(N)\to0\), passing to the constant term
  in the normalized partial sums leaves \(\mathcal M_C(A,\tau;G)\) unchanged.

  The Euler-product proof of \Cref{thm:class-euler-product-mertens} depends
  on the counting data only through this constant term and through the
  absolutely convergent \(r\ge2\) primitive tail.  Replacing the pointwise
  count by \eqref{eq:secondary-noninterference-refined-count} changes the
  normalized logarithm only by the same vanishing tail and by a remainder
  which is already absorbed in the displayed \(o(1)\) term.  Therefore the
  limiting positive product constant is still the value in
  \eqref{eq:secondary-noninterference-k}.

  Finally, \Cref{prop:fixed-label-determinant-perron} determines each
  non-trivial value \(F_\varrho(\lambda^{-1})\) from the finite determinant
  data and the fixed Abel branch before any secondary refinement is chosen.
  The inverse Fourier recovery of
  \Cref{thm:class-product-determinant-rigidity} then recovers the same
  fixed-label vector from the normalized product constants.  Thus any two
  auxiliary refinements satisfying the tail condition have identical main
  Mertens constants, product constants, and recovered fixed-label Perron
  values.
\end{proof}

\begin{remark}[Comparison with classical Mertens and Chebotarev
  asymptotics]
  \label{rem:class-mertens-comparison}
  The conclusion of \Cref{thm:class-mertens-explicit} should be read as
  a finite-state classwise constant formula under an already stated
  twisted spectral gap. It does not supply a new Tauberian proof of a
  Mertens theorem of the type proved by Sharp
  \cite{Sharp1991Mertens}, nor a new leading-density Chebotarev theorem
  of the type proved for Axiom~A or homogeneous extensions in
  \cite{ParryPollicott1986Chebotarev,NooraniParry1992ChebotarevShifts}.
  Its extra output is the explicit determinant expression for the
  class-indexed constant term \(\mathcal M_C(A,\tau;G)\) after the leading
  \(\frac{|C|}{|G|}\log N\) behaviour has already been separated. In
  this sense \(\mathcal M_C(A,\tau;G)\) is a classwise finite-state
  analogue of the constant term in classical orbit-counting asymptotics,
  not a new Tauberian theorem of Sharp type.  The theorem-by-theorem
  calibration against the classical and recent comparison results is recorded
  in \Cref{tab:novelty-calibration}.
\end{remark}

% END INPUT sec_chebotarev_class_mertens_determinants.tex
% BEGIN INPUT sec_chebotarev_class_mertens_obstruction.tex
\subsubsection{Artin obstruction and Parseval lower bounds}

This is the obstruction and certification part of the Main Theorem~III block.
All constants in this subsection are the product constants already defined
above; the new question is whether a proposed branch-compatible replacement
for the fixed-label vector can give the same class products.  The finite
\(S_3\) certificates later in the paper are used only for this obstruction,
and their rational interval inequalities are kept as exact certificates rather
than decimal approximations.

\begin{theorem}[Adams obstruction criterion for class-product constants]
  \label{thm:adams-obstruction-class-product}
  Assume the hypotheses of \Cref{thm:class-mertens-explicit}. The true
  class-product constant is the fixed-label Euler-transform constant of
  \Cref{thm:class-euler-product-mertens}. To measure the failure of the
  periodic Artin logarithm to govern primitive-label Euler products,
  define the naive Artin replacement constant
  \begin{equation}\label{eq:naive-artin-class-product-constant}
    K_C^{\mathrm{Art}}
    :=
    \exp\!\left(
      -\frac{|C|}{|G|}
      \left(
        \gamma+\log C(A)
        +\sum_{\varrho\neq\mathbf1}
          \chi_\varrho^\ast(C)\mathcal L_\varrho(\lambda^{-1})
      \right)
    \right),
  \end{equation}
  whenever the displayed non-trivial Artin logarithms are regular at
  \(\lambda^{-1}\). Then the true constants from
  \Cref{thm:class-euler-product-mertens} satisfy
  \begin{equation}\label{eq:adams-obstruction-ratio}
    \log\frac{K_C}{K_C^{\mathrm{Art}}}
    =
    -\frac{|C|}{|G|}
      \sum_{\varrho\neq\mathbf1}
        \chi_\varrho^\ast(C)
        \left(F_\varrho(\lambda^{-1})-
        \mathcal L_\varrho(\lambda^{-1})\right).
  \end{equation}
  In particular, the ordinary Artin substitution gives all class-product
  constants if and only if the class function
  \begin{equation}\label{eq:adams-obstruction-class-function}
    C\longmapsto
    \frac{|C|}{|G|}
      \sum_{\varrho\neq\mathbf1}
        \chi_\varrho^\ast(C)
        \left(F_\varrho(\lambda^{-1})-
        \mathcal L_\varrho(\lambda^{-1})\right)
  \end{equation}
  vanishes identically. Equivalently,
  \(F_\varrho(\lambda^{-1})=\mathcal L_\varrho(\lambda^{-1})\) for every
  non-trivial irreducible character \(\varrho\). More generally, the
  obstruction splits orthogonally along the Peter--Weyl sectors and is
  detected classwise by \eqref{eq:adams-obstruction-ratio}.
\end{theorem}

\begin{proof}
  Formula \eqref{eq:adams-obstruction-ratio} is obtained by subtracting
  the logarithm of \eqref{eq:naive-artin-class-product-constant} from
  the logarithm of \eqref{eq:class-euler-product-k-f}. Hence
  \(K_C=K_C^{\mathrm{Art}}\) for every conjugacy class \(C\) if and only
  if the class function in
  \eqref{eq:adams-obstruction-class-function} is zero.

  It remains only to identify when this class function is zero in terms
  of irreducible channels. Since \(|C|/|G|>0\) for every conjugacy class,
  the function in \eqref{eq:adams-obstruction-class-function} vanishes
  identically if and only if
  \[
    \sum_{\varrho\neq\mathbf1}\chi_\varrho^\ast(C)
      \left(F_\varrho(\lambda^{-1})-
      \mathcal L_\varrho(\lambda^{-1})\right)
    =0
  \]
  for every \(C\). The functions
  \(C\mapsto \chi_\varrho^\ast(C)\), \(\varrho\in\Irr(G)\), form an
  orthonormal basis of class functions for the Hermitian inner product
  \[
    \langle F,H\rangle_G
    :=
    \sum_{C\subset G}\frac{|C|}{|G|}F(C)\overline{H(C)}.
  \]
  Hence all non-trivial coefficients vanish if and only if
  \(F_\varrho(\lambda^{-1})-
  \mathcal L_\varrho(\lambda^{-1})=0\) for every
  \(\varrho\neq\mathbf1\). The same orthogonality gives the final
  Peter--Weyl splitting statement.
\end{proof}

\begin{corollary}[Downstream product-constant Adams obstruction norm]
  \label{thm:adams-obstruction-bound}
  \label{cor:adams-obstruction-bound}
  \label{thm:adams-obstruction-budget}
  \label{cor:adams-obstruction-budget}
  Assume the hypotheses and notation of
  \Cref{thm:adams-obstruction-class-product}.  This is a Parseval statement for
  the product constants of the present finite-state extension: it decomposes
  the classwise error caused by replacing fixed primitive labels with
  periodic Artin logarithms, and it supplies the quantitative interface used
  by the finite \(S_3\) witness in
  \Cref{thm:s3-certified-nonabelian-defect,prop:s3-two-coordinate-product-obstruction-vector,cor:s3-product-constant-obstruction-finite-state}.
  It concerns only the determinant model fixed in this paper.  For
  \(\varrho\neq\mathbf1\), put
  \[
    D_\varrho
    :=F_\varrho(\lambda^{-1})-\mathcal L_\varrho(\lambda^{-1}),
  \]
  and define the normalized class-product error
  \[
    \mathcal E(C)
    :=-
      \frac{|G|}{|C|}
      \log\frac{K_C}{K_C^{\mathrm{Art}}}
      \qquad(C\subset G).
  \]
  Then, inside this determinant/Peter--Weyl model, the possible classwise
  Artin-replacement error is completely confined to the zero-mean
  class-function skeleton:
  \begin{equation}\label{eq:adams-obstruction-bound-zero-mean}
    \sum_{C\subset G}\frac{|C|}{|G|}\mathcal E(C)=0,
    \qquad
    D_\pi=
    \sum_{C\subset G}\frac{|C|}{|G|}\mathcal E(C)\chi_\pi(C)
    \quad(\pi\neq\mathbf1).
  \end{equation}
  Conversely, \(\mathcal E\) is recovered from these non-trivial Fourier
  coordinates by
  \begin{equation}\label{eq:adams-obstruction-bound-forward}
    \mathcal E(C)=
    \sum_{\varrho\neq\mathbf1}\chi_\varrho^\ast(C)D_\varrho .
  \end{equation}
  Hence
  \begin{equation}\label{eq:adams-obstruction-bound-parseval}
    \sum_{C\subset G}\frac{|C|}{|G|}\,|\mathcal E(C)|^2
    =
    \sum_{\varrho\neq\mathbf1}|D_\varrho|^2.
  \end{equation}
  In particular, the naive Artin constants agree with the true constants for
  all conjugacy classes if and only if \(D_\varrho=0\) for every non-trivial
  \(\varrho\).  If they do not agree, the obstruction cannot be hidden in a
  cancellation among conjugacy classes: for
  \(\Delta^2:=\sum_{\varrho\neq\mathbf1}|D_\varrho|^2\), at least one
  conjugacy class satisfies
  \begin{equation}\label{eq:adams-obstruction-bound-sup-lower-bound}
    \left|
      \frac{|G|}{|C|}
      \log\frac{K_C}{K_C^{\mathrm{Art}}}
    \right|
    \ge \Delta .
  \end{equation}
  In particular, if a finite certificate in this model proves
  \(\sum_{\varrho\neq\mathbf1}|D_\varrho|^2\ge \delta^2\), then at least
  one conjugacy class \(C\) satisfies
  \begin{equation}\label{eq:adams-obstruction-bound-lower-bound}
    \left|
      \frac{|G|}{|C|}
      \log\frac{K_C}{K_C^{\mathrm{Art}}}
    \right|
    \ge \delta.
  \end{equation}
  Thus a non-zero Adams discrepancy in any Peter--Weyl sector forces a
  class-level Euler-product error of the same normalized size, and the
  error cannot be hidden by averaging over conjugacy classes.  The conclusion
  is a Parseval refinement of the Adams--M\"obius obstruction in
  \Cref{thm:adams-obstruction-class-product}; it is not an additional main
  theorem determining the constants themselves.
\end{corollary}

\begin{proof}
  By \eqref{eq:adams-obstruction-ratio}, the normalized error is the
  non-trivial Peter--Weyl expansion
  \[
    \mathcal E(C)=
    \sum_{\varrho\neq\mathbf1}\chi_\varrho^\ast(C)D_\varrho.
  \]
  This is \eqref{eq:adams-obstruction-bound-forward}.
  The irreducible characters form an orthonormal basis of the space of
  complex class functions for the inner product
  \[
    \langle f,h\rangle_G
    =
    \sum_{C\subset G}\frac{|C|}{|G|}f(C)\overline{h(C)}.
  \]
  Since the expansion has no trivial-character term, its scalar Fourier
  coefficient is zero, proving the zero-mean identity in
  \eqref{eq:adams-obstruction-bound-zero-mean}.  Multiplying
  \eqref{eq:adams-obstruction-bound-forward} by
  \((|C|/|G|)\chi_\pi(C)\), summing over classes, and using orthogonality
  gives the inverse formula for \(D_\pi\).  Parseval's identity gives
  \eqref{eq:adams-obstruction-bound-parseval}.

  The equivalence between equality of all class constants and vanishing of
  all \(D_\varrho\) follows either from
  \Cref{thm:adams-obstruction-class-product} or directly from Parseval:
  \(K_C=K_C^{\mathrm{Art}}\) for every \(C\) is equivalent to
  \(\mathcal E=0\), and \eqref{eq:adams-obstruction-bound-parseval} is then
  equivalent to \(D_\varrho=0\) for every non-trivial \(\varrho\).  If
  \(\Delta^2=\sum_{\varrho\neq\mathbf1}|D_\varrho|^2\), then
  \eqref{eq:adams-obstruction-bound-parseval} says that the weighted
  average of \(|\mathcal E(C)|^2\) is \(\Delta^2\). Because the weights
  \(|C|/|G|\) are positive and sum to one, some class has
  \(|\mathcal E(C)|^2\ge\Delta^2\), which proves
  \eqref{eq:adams-obstruction-bound-sup-lower-bound}.  The same argument
  with \(\Delta\ge\delta\) proves
  \eqref{eq:adams-obstruction-bound-lower-bound}.
\end{proof}

\begin{statusnote}[Obstruction refinements]
  Under the hypotheses of \Cref{thm:class-euler-product-mertens}, the
  core Main Theorem~III data consist of the Frobenius-class product asymptotic, the
  fixed-label constants \(K_C\), the finite determinant evaluation of
  \(F_\varrho(\lambda^{-1})\), and the Adams--M\"obius obstruction to the
  naive Artin substitution.  The fixed-coordinate rank obstruction of
  \Cref{cor:class-product-minimal-obstruction-skeleton} and the Parseval
  bound of \Cref{cor:adams-obstruction-bound} are determined functorially by
  those data together with the already defined periodic Artin logarithms:
  if two finite-state extensions have the same scalar baseline
  \(\gamma+\log C(A)\), the same non-trivial fixed-label boundary values
  \(F_\varrho(\lambda^{-1})\), and the same periodic Artin logarithms
  \(\mathcal L_\varrho(\lambda^{-1})\), then they have the same
  scalar-normalised fixed-coordinate product vector, the same class constants,
  and the same
  Artin-substitution obstruction bound.  Conversely, equality of the
  scalar-normalised fixed-coordinate product vector forces equality of all
  non-trivial
  fixed-label boundary values.  Hence the fixed-coordinate rank obstruction
  and the Parseval
  obstruction bound are downstream corollary-level refinements of Main
  Theorem~III, not independent main theorems and not additional hypotheses for
  the product asymptotic.
\end{statusnote}

Main Theorem~III, in the form of \Cref{thm:class-euler-product-mertens},
gives the constants
  \[
    K_C=
    \exp\!\left(-\frac{|C|}{|G|}
      \left(\gamma+\log C(A)+
      \sum_{\varrho\ne\mathbf1}\chi_\varrho^\ast(C)
        F_\varrho(\lambda^{-1})\right)\right).
  \]
  Thus the scalar baseline together with the tuple
  \((F_\varrho(\lambda^{-1}))_{\varrho\ne\mathbf1}\) determines every
  \(K_C\), and \Cref{cor:class-product-minimal-obstruction-skeleton} then
  identifies this tuple with the unique zero-mean Peter--Weyl coordinate
  vector of the scalar-normalised product constants.  The same corollary gives
  the converse: equality of scalar-normalised fixed-coordinate product vectors
  implies
  equality of their non-trivial Peter--Weyl coordinates, hence equality of all
  values \(F_\varrho(\lambda^{-1})\).

  The Artin-substitution constants are obtained from the same displayed
  formula after replacing each fixed-label value by
  \(\mathcal L_\varrho(\lambda^{-1})\).  Therefore the difference vector
  \(D_\varrho=F_\varrho(\lambda^{-1})-
  \mathcal L_\varrho(\lambda^{-1})\), the zero-mean classwise error
  \(\mathcal E(C)\), and the norm identity in
  \Cref{cor:adams-obstruction-bound} are all functions of the already fixed
  Main Theorem~III data and of the corresponding periodic Artin logarithms.
If the periodic Artin logarithms also agree, then the difference vector agrees,
and so the Parseval bound agrees.  No additional asymptotic input, determinant
channel, or Perron-point hypothesis is introduced.

\begin{lemma}[Perron-boundary determinant evaluation]
  \label{lem:fixed-label-perron-determinant-evaluation}
  \label{lem:perron-boundary-determinant-evaluation}
  Assume \(\lambda>1\) and \(\eta<\lambda\). For every non-trivial
  \(\varrho\in\Irr(G)\), the boundary determinant series
  \[
    \sum_{m\ge 1}\frac{\mu(m)}{m}
      \sum_{\substack{\pi\in\Irr(G)\\(m,\pi)\neq(1,\mathbf 1)}}
        a_{\varrho,\pi}^{(m)}
        \log\det(I-\lambda^{-m}A_\pi)^{-1}
  \]
  converges absolutely. The displayed summation is a local exclusion of
  the boundary term \((m,\pi)=(1,\mathbf1)\). That term would be
  \(a_{\varrho,\mathbf 1}^{(1)}
  \log\det(I-\lambda^{-1}A_{\mathbf1})^{-1}\) and has coefficient
  \(a_{\varrho,\mathbf 1}^{(1)}=0\); it is therefore absent, not assigned
  a finite part. The remaining \(m=1\) terms are regular, and the
  \(m\ge2\) tail is absolutely summable. The displayed value is
  \(\Pi_\varrho(\lambda^{-1})\).
\end{lemma}

\begin{proof}
  Since \(\psi^1\chi_\varrho=\chi_\varrho\) and
  \(\varrho\neq\mathbf 1\), orthogonality gives
  \(a_{\varrho,\mathbf 1}^{(1)}=0\). Thus the scalar
  \((m,\pi)=(1,\mathbf1)\) term is excluded at the summand level before
  passing to \(z=\lambda^{-1}\); it is never evaluated, and no value of
  \(\log\det(I-\lambda^{-1}A_{\mathbf 1})^{-1}\) occurs. For
  \(m=1\) and \(\pi\neq\mathbf 1\), the twisted gap gives
  \(\rho(A_\pi)<\lambda\), hence \(I-\lambda^{-1}A_\pi\) is
  invertible and the Abel-path continuation along
  \(z=t\lambda^{-1}\), \(0\le t\uparrow1\), is the regular branch used in
  all non-trivial Perron-point channels. For
  \(m\ge2\), one has
  \[
    \lambda^{-m}\rho(A_\pi)\le \lambda^{1-m}<1
    \qquad(\pi\in\Irr(G)).
  \]
  Since the set of matrices \(A_\pi\) is finite, the power-series
  expansion of the determinant logarithm gives a uniform bound
  \[
    \left|\log\det(I-r^m\lambda^{-m}A_\pi)^{-1}\right|
    \le C\lambda^{1-m}
    \qquad(0\le r\le1,\ m\ge2,\ \pi\in\Irr(G))
  \]
  for a constant \(C\) depending only on the finite list of twisted
  matrices. Together with
  \(\abs{a_{\varrho,\pi}^{(m)}}\le\dim\varrho\) from
  \Cref{lem:adams-coefficients-bounded}, this dominates the \(m\ge2\)
  tail by a constant multiple of
  \(\sum_{m\ge2}\lambda^{1-m}/m\). The determinant series is therefore
  absolutely convergent at \(r=1\), apart from the finitely many regular
  \(m=1\), \(\pi\neq\mathbf1\) terms already discussed.

  For \(0\le r<1\), \Cref{cor:primitive-peter-weyl-determinant} applied
  at \(z=r\lambda^{-1}\) gives
  \[
    \Pi_\varrho(r\lambda^{-1})
    =
    \sum_{m\ge 1}\frac{\mu(m)}{m}
      \sum_{\substack{\pi\in\Irr(G)\\(m,\pi)\neq(1,\mathbf 1)}}
        a_{\varrho,\pi}^{(m)}
        \log\det(I-r^m\lambda^{-m}A_\pi)^{-1}.
  \]
  The same majorant permits dominated convergence on the right as
  \(r\uparrow1\). On the left,
  \(\Pi_\varrho(r\lambda^{-1})\to\Pi_\varrho(\lambda^{-1})\) by the
  absolute convergence from
  \Cref{lem:nontrivial-perron-point-convergence}. Passing to the limit
  proves the claimed boundary identity.
\end{proof}

\begin{proposition}[Main Theorem II bridge: class-constant deviations
  and Perron-point determinant evaluation]
  \label{prop:class-mertens-deviations}
  Under the hypotheses of \Cref{thm:class-mertens-explicit}, define
  \[
    \Delta_C(A,\tau;G)
    :=\mathcal M_C(A,\tau;G)-\frac{|C|}{|G|}\mathcal M(A).
  \]
  Then the classwise deviations \(\Delta_C(A,\tau;G)\) are well
  defined, and
  \begin{equation}\label{eq:class-mertens-fourier-expansion}
    \Delta_C(A,\tau;G)
    =
    \frac{|C|}{|G|}
      \sum_{\varrho\neq\mathbf 1}
        \chi_\varrho^\ast(C)\,\Pi_\varrho(\lambda^{-1}).
  \end{equation}
  In the boundary determinant evaluation below, every logarithm
  \(\log\det(I-\lambda^{-m}A_\pi)^{-1}\) is taken in the analytic
  branch fixed by the logarithm convention at the start of
  \Cref{sec:chebotarev}: for \(m=1\) and non-trivial \(\pi\) this
  means the Abel-path boundary value of the normalized determinant germ
  from \(z=0\) to \(z=\lambda^{-1}\) through the zero-free interval given by
  the twisted spectral gap, and for \(m\ge 2\) it is the same germ evaluated
  at \(z=\lambda^{-m}\).
  Moreover, for every non-trivial \(\varrho\),
  \begin{equation}\label{eq:class-mertens-determinant}
    \Pi_\varrho(\lambda^{-1})
    =
    \sum_{m\ge 1}\frac{\mu(m)}{m}
      \sum_{\substack{\pi\in\Irr(G)\\(m,\pi)\neq(1,\mathbf 1)}}
        a_{\varrho,\pi}^{(m)}
        \log\det(I-\lambda^{-m}A_\pi)^{-1}.
  \end{equation}
  The determinant logarithms are taken in the branch fixed above; for
  \(m=1\) this is the Abel regular non-trivial value at \(z=\lambda^{-1}\),
  while possible scalar terms with \(m\ge2\) are ordinary interior-disc
  values.  The formally singular zero-coefficient term
  \((m,\pi)=(1,\mathbf1)\) is omitted before the branch is evaluated.
\end{proposition}

\begin{proof}
  Subtract \((|C|/|G|)\mathcal M(A)\) from
  \eqref{eq:class-mertens-constant} to obtain
  \eqref{eq:class-mertens-fourier-expansion}.
  The boundary determinant identity
  \eqref{eq:class-mertens-determinant} is exactly
  \Cref{lem:perron-boundary-determinant-evaluation}.
\end{proof}


\begin{remark}[On the twisted spectral-gap hypothesis]
  \label{rem:twisted-spectral-gap-role}
  The condition
  \[
    \eta=\max_{\varrho\neq\mathbf 1}\rho(A_\varrho)<\lambda
  \]
  does not enter the determinant identities themselves: the formulas of
  \Cref{thm:class-function-linearisation,cor:class-function-adams-mobius,cor:primitive-peter-weyl-determinant}
  are valid throughout the disc \(|z|<\lambda^{-1}\) without it. The
  role of the hypothesis is exactly to permit evaluation at the Perron
  point \(z=\lambda^{-1}\) in the non-trivial channels and to convert the
  primitive Peter--Weyl series into the class-indexed constants
  \(\mathcal M_C(A,\tau;G)\). From the dynamical point of view this is the
  representation-by-representation analogue of separating the Perron
  eigenvalue from the remaining spectrum.

  The leading Chebotarev and Frobenius-class Mertens phenomena for finite
  homogeneous extensions belong to the existing orbit-counting lineage,
  beginning here with Noorani--Parry
  \cite{NooraniParry1992ChebotarevShifts} and Noorani's Markov-shift
  extension criteria \cite{Noorani1997HomogeneousExtensionsMarkov}, and
  continuing through the Axiom~A, Liv\v{s}ic, skew-product reciprocity, and
  homology-class settings of
  \cite{ParryPollicott1986Chebotarev,Parry1999LivsicNonAbelian,ParryPollicott1997LivsicCompactLie,ParryPollicott2008BauerSkewProducts,PollicottSharp2008ArtinReciprocitySkewProducts,KatsudaSunada1990ClosedOrbitsHomology,Sharp1993ClosedOrbitsHomologyAnosov,PollicottSharp2007Chebotarev,BoyleSchmieding2017FiniteGroupExtensions}.
  The present paper uses the gap as a stated finite-dimensional hypothesis,
  not as a new structural criterion for cocycles.  In the finite-state setting
  it is already checkable:
  \Cref{lem:twisted-gap-skew-product-peripheral} identifies it with the
  absence of non-trivial peripheral spectrum in the skew-product
  adjacency \(\widetilde A\), equivalently with \(\widetilde A\) having
  a simple strictly dominant Perron eigenvalue, and, when
  \(\widetilde A\) is irreducible, with primitivity of
  \(\widetilde A\). What remains outside the present argument is a structural
  criterion on \(\tau\) that would imply this condition without first testing
  the full skew-product matrix.  The new step here is instead the
  finite-state Adams--M\"obius calculus at the level of primitive constants
  and the explicit class-indexed determinant formula
  \(\mathcal M_C(A,\tau;G)\), not a new leading-density theorem of the type
  considered in
  \cite{Lalley1987DistributionPeriodicOrbits,Coles2025VassilievWritheAxiomA}.
\end{remark}
% END INPUT sec_chebotarev_class_mertens_obstruction.tex
% END INPUT sec_chebotarev_class_mertens.tex
% BEGIN INPUT sec_chebotarev_readout.tex
\subsection{Hermitian Fourier readout of class constants}

The class-indexed constants \(\mathcal M_C(A,\tau;G)\) form the readout
interface between Main Theorems~II and~III.  Main Theorem~II first produces
their deviations from the trivial average as Perron-point primitive channels,
using the Adams--M\"obius determinant formula and the strict twisted gap.
The present subsection then applies Hermitian character orthogonality to read
those channels back from the class constants.  This is the final step in Main
Theorem~II and the vector of Perron-point values that Main Theorem~III inserts
into the fixed-label Euler product; it is not a cyclic one-layer sampling
argument.

\begin{theorem}[Main Theorem II(c): class-constant Fourier recovery]
  \label{thm:class-mertens-fourier}
  Under the hypotheses of \Cref{thm:class-mertens-explicit}, let
  \(\Delta_C(A,\tau;G)\) be as in \Cref{prop:class-mertens-deviations}.
  Thus the boundary values \(\Pi_\kappa(\lambda^{-1})\) appearing below
  are the Perron-point channel values supplied by
  \Cref{prop:class-mertens-deviations} with that same fixed analytic
  branch convention.
  Then
  for every non-trivial \(\kappa\in\Irr(G)\),
  \begin{equation}\label{eq:class-mertens-fourier-inversion}
    \Pi_\kappa(\lambda^{-1})
    =
    \sum_{C\subset G}\Delta_C(A,\tau;G)\chi_\kappa(C),
  \end{equation}
  where the sum runs over conjugacy classes. Moreover,
  \begin{equation}\label{eq:class-mertens-plancherel}
    \sum_{C\subset G}\frac{|G|}{|C|}\abs{\Delta_C(A,\tau;G)}^2
    =
    \sum_{\varrho\neq\mathbf 1}
      \abs{\Pi_\varrho(\lambda^{-1})}^2.
  \end{equation}
  Equivalently, this is Fourier inversion and Parseval for the Hermitian
  class-function inner product
  \[
    \langle f,h\rangle_{\mathrm{cl}}
    :=
    \sum_{C\subset G}\frac{|C|}{|G|}f(C)\overline{h(C)},
  \]
  so the forward coefficient is
  \[
    \chi_\varrho^\ast(C).
  \]
  In the inverse recovery, this coefficient is paired with the
  unconjugated class character \(\chi_\kappa(C)\).
\end{theorem}

\begin{proof}
  This is Fourier inversion in the full class-function space of \(G\).
  No cyclic quotient and no root-of-unity sampling condition is involved:
  the inputs are the class-indexed constants
  \(\mathcal M_C(A,\tau;G)\) for all conjugacy classes, and the basis is
  the full irreducible character basis.
  Multiply \eqref{eq:class-mertens-fourier-expansion} by
  \(\chi_\kappa(C)\) and sum over conjugacy classes:
  \[
    \sum_{C\subset G}\Delta_C(A,\tau;G)\chi_\kappa(C)
    =
    \sum_{C\subset G}\frac{|C|}{|G|}
      \sum_{\varrho\neq\mathbf 1}
        \chi_\varrho^\ast(C)\,
        \Pi_\varrho(\lambda^{-1})\chi_\kappa(C).
  \]
  The orthogonality relation
  \[
    \sum_{C\subset G}\frac{|C|}{|G|}
      \chi_\varrho^\ast(C)\,\chi_\kappa(C)
    =\delta_{\varrho,\kappa}
  \]
  yields \eqref{eq:class-mertens-fourier-inversion}.

  For the Parseval identity, multiply
  \eqref{eq:class-mertens-fourier-expansion} by \(|G|/|C|\), take the
  modulus squared, and sum over conjugacy classes. Since
  \[
    |\Delta_C(A,\tau;G)|^2
    =
    \left(\frac{|C|}{|G|}\right)^2
    \left|\sum_{\varrho\neq\mathbf 1}
      \chi_\varrho^\ast(C)\,
      \Pi_\varrho(\lambda^{-1})\right|^2,
  \]
  one obtains
  \[
    \sum_{C\subset G}\frac{|G|}{|C|}|\Delta_C(A,\tau;G)|^2
    =
    \sum_{C\subset G}\frac{|C|}{|G|}
      \left|\sum_{\varrho\neq\mathbf 1}
        \chi_\varrho^\ast(C)\,
        \Pi_\varrho(\lambda^{-1})\right|^2.
  \]
  Expanding the modulus square gives
  \[
    \sum_{\varrho,\kappa\neq\mathbf 1}
      \Pi_\varrho(\lambda^{-1})
      \overline{\Pi_\kappa(\lambda^{-1})}
      \sum_{C\subset G}\frac{|C|}{|G|}
        \chi_\varrho^\ast(C)\,\chi_\kappa(C).
  \]
  The inner sum is exactly
  \[
    \sum_{C\subset G}\frac{|C|}{|G|}
      \chi_\varrho^\ast(C)\,\chi_\kappa(C)
    =\delta_{\varrho,\kappa}.
  \]
  Therefore only the diagonal terms remain, which
  yields \eqref{eq:class-mertens-plancherel}.
\end{proof}

\noindent
Together with \Cref{thm:class-mertens-explicit} and
\Cref{prop:class-mertens-deviations}, this proves Main Theorem~II from
the introduction. The determinant and Adams--M\"obius work has already
produced the Perron-point primitive channels; the final step above is
the standard character-orthogonality recovery on class functions, not the
one-layer DFT calculation used later for the canonical cyclic-lift
comparison.

\begin{remark}[recovery analogy stops at the orthogonality step]
  \label{rem:recovery-analogy-orthogonality-step}
  \label{rem:readout-analogy-orthogonality-step}
  \Cref{thm:class-mertens-fourier} has the same formal recovery shape as
  the root-of-unity reconstruction results from \Cref{sec:finite-part}.
  In the separate canonical cyclic-lift comparison case, one
  root-of-unity layer recovers the reduced spectral data by ordinary
  Fourier inversion.
  Here the full family of
  conjugacy-class Mertens constants \(\mathcal M_C(A,\tau;G)\) recovers
  the non-trivial twisted
  primitive values \(\Pi_\varrho(\lambda^{-1})\) by non-abelian
  character orthogonality. The analogy is only at the final recovery
  step: the class-indexed constants being inverted here have first been
  produced by the
  Adams--M\"obius determinant mechanism and the twisted gap, whereas
  the one-layer statement of \Cref{sec:finite-part} is polynomial
  coefficient recovery.
\end{remark}
% END INPUT sec_chebotarev_readout.tex
% BEGIN INPUT sec_eta_gap.tex
\section{Finite-state criteria for the twisted spectral gap}
\label{sec:eta-gap}

The strict inequality
\(\eta=\max_{\varrho\neq\mathbf 1}\rho(A_\varrho)<\lambda\) is the
finite-state hypothesis that permits the non-trivial Peter--Weyl channels in
Main Theorem~II, and hence the fixed-label Euler coordinates in Main
Theorem~III, to be evaluated at the Perron point; if \(G\) is trivial, this
empty maximum is \(0\).  This section gives finite-state criteria for that hypothesis and describes the equality case
\(\eta=\lambda\).  The equality case is rigid in the precise sense needed
here: a non-trivial channel has Perron spectral radius exactly when the
cocycle admits an equivariant line field.  The resulting loop-character
criteria are used as finite checks for the Perron-point hypotheses in Main
Theorems~II--III.

\begin{theorem}[Line-field rigidity at the Perron boundary]
  \label{thm:eta-gap-rigidity}
  Fix the finite directed graph with primitive adjacency matrix \(A\), the
  fixed finite-group cocycle \(\tau\), and the associated twisted matrices
  \(A_\varrho\) from \eqref{eq:twisted-matrix-def}.  Let
  \(\varrho:G\to U(V_\varrho)\) be an irreducible unitary representation.
  Then
  \[
    \rho(A_\varrho)=\lambda
  \]
  denotes spectral-radius equality for the row-vector twisted matrix from
  \eqref{eq:twisted-matrix-def}. For the line-field formulation define
  the column block operator \(A_\varrho^{\mathrm{col}}\) on
  \(\CC^V\otimes V_\varrho\) by
  \[
    (A_\varrho^{\mathrm{col}} w)_i
    =
    \sum_{e:i\to j}\varrho(\tau(e))w_j .
  \]
  The operator \(A_\varrho^{\mathrm{col}}\) is the transpose-dual block
  realisation of the row operator used in
  \Cref{prop:kernel-peter-weyl-block-diagonalisation,prop:twisted-spectral-radius-bound}; in particular it has the same
  characteristic polynomial and the same spectral radius as
  \(A_\varrho\). The equality above is therefore equivalent to
  \(\rho(A_\varrho^{\mathrm{col}})=\lambda\), and this latter column
  convention is used in the displayed line-field equation below.
  Thus \(\rho(A_\varrho)=\lambda\) if and only if there exist
  \(\omega\in\TT\) and unit vectors
  \(\xi_i\in V_\varrho\), indexed by the vertices \(i\in V\), such that
  \begin{equation}\label{eq:eta-gap-line-field}
    \varrho(\tau(e))\xi_j=\omega\,\xi_i
    \qquad
    \text{for every admissible edge } e:i\to j.
  \end{equation}
  In row form, equality at the Perron boundary is expressed by
  row unit vectors \(\zeta_i\) satisfying
  \[
    \zeta_i\varrho(\tau(e))=\omega\,\zeta_j
    \qquad(e:i\to j),
  \]
  after replacing \(\omega\) by the corresponding unit scalar if one
  passes between row vectors and adjoint column vectors.
  Whenever \eqref{eq:eta-gap-line-field} holds, every closed path
  \(\gamma\) of length \(n\) based at \(i\) satisfies
  \begin{equation}\label{eq:eta-gap-loop-character}
    \varrho(\tau(\gamma))\xi_i=\omega^n\xi_i.
  \end{equation}
  In particular, the subgroup
  \[
    \Gamma_i(\tau)
    :=
    \bigl\langle \tau(\gamma):
      \gamma \text{ closed path based at } i
    \bigr\rangle
  \]
  acts on the line \(\CC\xi_i\) through a one-dimensional character.  This
  theorem is a conditional boundary diagnostic for the fixed finite-state cocycle just
  specified: it identifies exactly when equality \(\rho(A_\varrho)=\lambda\)
  occurs in a non-trivial channel.  Its contrapositive is used below only as a
  finite test for the strict-gap hypothesis of Main Theorems~II--III; it is not
  a classification of cocycles for which \(\eta<\lambda\).
\end{theorem}

\begin{proof}
  We use \(A_\varrho^{\mathrm{col}}\) throughout the proof. The statement
  has already identified its spectral radius with that of the row
  operator \(A_\varrho\), so no spectral assertion changes under this
  convention.

  Assume first that \(\rho(A_\varrho)=\lambda\). Since
  \(A_\varrho^{\mathrm{col}}\) is finite dimensional, some eigenvalue
  \(\alpha\in\spec(A_\varrho^{\mathrm{col}})\) satisfies
  \(|\alpha|=\lambda\). Let
  \(v=(v_i)_{i\in V}\in \CC^V\otimes V_\varrho\) be a corresponding
  eigenvector, so
  \[
    \alpha v_i=\sum_{e:i\to j}\varrho(\tau(e))v_j
    \qquad (i\in V).
  \]
  Set \(u_i:=\norm{v_i}\). Then
  \[
    \lambda u_i
    =
    \norm{\alpha v_i}
    =
    \left\|\sum_{e:i\to j}\varrho(\tau(e))v_j\right\|
    \le
    \sum_{e:i\to j}\norm{v_j}
    =(Au)_i
  \]
  for every \(i\in V\). Let \(\ell>0\) be a left Perron vector of
  \(A\), normalised by \(\ell^{\mathsf T}A=\lambda\ell^{\mathsf T}\).
  Then
  \[
    \lambda\,\ell^{\mathsf T}u
    \le
    \ell^{\mathsf T}Au
    =
    \lambda\,\ell^{\mathsf T}u,
  \]
  so every coordinate inequality is an equality because \(\ell_i>0\)
  for all \(i\). Hence \(Au=\lambda u\). Primitivity of \(A\) implies
  that \(u\) is a positive Perron eigenvector, so \(u_i>0\) for every
  \(i\in V\).

  Equality in the triangle inequality now holds in every vertex equation.
  Fix a vertex \(i\) and list the outgoing edges with multiplicity, including
  parallel edges and edges with the same terminal vertex, as
  \(e_1,
  \ldots,e_q\).  Put
  \(x_a:=\varrho(\tau(e_a))v_{t(e_a)}\).  Each \(x_a\) is non-zero because
  \(u_{t(e_a)}>0\), and
  \[
    \left\|\sum_{a=1}^q x_a\right\|=
    \sum_{a=1}^q\|x_a\| .
  \]
  The Hilbert-space equality case gives the conclusion for every summand in
  this list, not merely for every terminal vertex: equality in Cauchy's
  inequality for
  \(\langle x_a,\sum_bx_b\rangle\) forces each \(x_a\) to be a non-negative
  real multiple of the same vector \(\sum_bx_b\), and the positivity of
  \(\|x_a\|=u_{t(e_a)}\) excludes zero multiples.  Thus there is a unit
  vector \(y_i\) such that
  \(\varrho(\tau(e))v_j=u_jy_i\) for every outgoing edge \(e:i\to j\), with
  repeated terminal vertices treated as separate summands. Since their sum is
  \(\alpha v_i\) and \(Au=\lambda u\),
  \[
    \alpha v_i
    =\sum_{e:i\to j}u_jy_i
    =\lambda u_i y_i,
  \]
  so \(y_i=\alpha v_i/(\lambda u_i)\).  Thus the phase is not chosen
  separately for each edge; it is the single eigenvalue phase
  \(\omega:=\alpha/\lambda\), and
  \[
    \varrho(\tau(e))v_j
    =
    u_j\,\frac{\alpha v_i}{\lambda u_i}
    \qquad
    \text{for every admissible edge } e:i\to j.
  \]
  Write \(\xi_i:=v_i/u_i\in V_\varrho\). Dividing by \(u_j\) yields
  \eqref{eq:eta-gap-line-field}.

  Conversely, assume that \eqref{eq:eta-gap-line-field} holds. Let
  \(h=(h_i)_{i\in V}>0\) be a right Perron vector of \(A\), so
  \(Ah=\lambda h\). Define \(w=(w_i)_{i\in V}\) by
  \(w_i:=h_i\xi_i\). Then
  \begin{align*}
    (A_\varrho^{\mathrm{col}} w)_i
    &=
    \sum_{e:i\to j}\varrho(\tau(e))w_j \\
    &=
    \sum_{e:i\to j}h_j\varrho(\tau(e))\xi_j \\
    &=
    \omega\sum_{e:i\to j}h_j\xi_i \\
    &=
    \omega(Ah)_i\xi_i \\
    &=
    \omega\lambda w_i.
  \end{align*}
  Thus \(\omega\lambda\in\spec(A_\varrho^{\mathrm{col}})\), so
  \(\rho(A_\varrho^{\mathrm{col}})\ge \lambda\), equivalently
  \(\rho(A_\varrho)\ge \lambda\). On the other hand,
  \Cref{prop:twisted-spectral-radius-bound} gives
  \(\rho(A_\varrho)\le \lambda\). Hence
  \(\rho(A_\varrho)=\lambda\).

  Finally, iterating \eqref{eq:eta-gap-line-field} along a closed path
  \(\gamma=e_1\cdots e_n\) based at \(i\) gives
  \eqref{eq:eta-gap-loop-character}. Hence the line \(\CC\xi_i\) is
  \(\Gamma_i(\tau)\)-invariant. The action of \(\Gamma_i(\tau)\) on
  this line defines a one-dimensional character.
\end{proof}

\begin{corollary}[Finite-state criterion for the strict twisted gap]
  \label{cor:eta-gap-boundary-classifier}
This is the operational test for the gap hypothesis.  It contains no
product-constant formula; it only decides whether the finite extension has
a non-trivial Perron-boundary channel.  Thus the diagnostic criterion below is
used only to certify the hypothesis of Main Theorems~II--III before the
Adams--M\"obius determinant recovery is applied.
  Suppose \(G\) is non-trivial. Then the following are equivalent.
  \begin{enumerate}
    \item The Perron-point input used in the class-product theorem fails,
      namely \(\eta=\lambda\).
    \item There exist a non-trivial irreducible representation
      \(\varrho\in\Irr(G)\), a phase \(\omega\in\TT\), and unit vectors
      \(\xi_i\in V_\varrho\), one for each vertex, such that
      \begin{equation}\label{eq:eta-gap-boundary-classifier}
        \varrho(\tau(e))\xi_j=\omega\xi_i
        \qquad
        \text{for every admissible edge } e:i\to j.
      \end{equation}
  \end{enumerate}
  Consequently, for this fixed finite-state cocycle, \(\eta<\lambda\) holds
  exactly when the finite list of edge equations
  \eqref{eq:eta-gap-boundary-classifier} has no solution in any non-trivial
  irreducible representation. In this case the hypotheses of
  \Cref{thm:class-euler-product-mertens} are reduced to the verified absence
  of these boundary line fields.
\end{corollary}

\begin{proof}
  By \Cref{prop:twisted-spectral-radius-bound}, every non-trivial channel
  satisfies \(\rho(A_\varrho)\le\lambda\). Since \(G\) is finite, there
  are only finitely many irreducible representations, so the maximum
  defining \(\eta\) is achieved. Hence \(\eta=\lambda\) if and only if
  \(\rho(A_\varrho)=\lambda\) for at least one non-trivial
  \(\varrho\in\Irr(G)\). Applying \Cref{thm:eta-gap-rigidity} to that
  representation gives exactly the edge equations
  \eqref{eq:eta-gap-boundary-classifier}, and the converse implication in
  the same theorem shows that any solution of those equations forces
  \(\rho(A_\varrho)=\lambda\), hence \(\eta=\lambda\). The final sentence
  is the contrapositive, together with the role of the strict-gap
  hypothesis stated in \Cref{thm:class-euler-product-mertens}.
\end{proof}

\begin{corollary}[Loop-character classification of boundary channels]
  \label{cor:eta-gap-character-classification}
  Fix a non-trivial irreducible unitary representation
  \(\varrho:G\to U(V_\varrho)\). For a phase \(\omega\in\TT\) and a vertex
  \(i\), let
  \[
    W_i(\varrho,\omega)
    :=
    \bigcap_{\gamma:i\to i}
    \ker\bigl(\varrho(\tau(\gamma))-\omega^{|\gamma|}I\bigr),
  \]
  where the intersection is over all closed admissible paths based at
  \(i\). Then \(\rho(A_\varrho)=\lambda\) if and only if there are a phase
  \(\omega\in\TT\) and unit vectors
  \(\xi_i\in W_i(\varrho,\omega)\), one for each vertex, such that
  \[
    \varrho(\tau(e))\xi_j=\omega\xi_i
    \qquad(e:i\to j).
  \]
  For every such solution, the line
  \(L_i:=\CC\xi_i\) is a one-dimensional
  \(\Gamma_i(\tau)\)-subrepresentation of
  \(\varrho|_{\Gamma_i(\tau)}\), with character
  \(\chi_i:\Gamma_i(\tau)\to\TT\) determined by
  \[
    \chi_i(\tau(\gamma))=\omega^{|\gamma|}
    \qquad(\gamma:i\to i).
  \]
  Moreover, if \(p:i\to j\) is any admissible path of length \(|p|\), then
  \(\omega^{-|p|}\varrho(\tau(p))\) gives a unitary similarity
  \(L_j\to L_i\). Thus the Perron-boundary obstruction is precisely a
  compatible system of one-dimensional loop characters and edge
  similarities.
\end{corollary}

\begin{proof}
  If \(\rho(A_\varrho)=\lambda\), then
  \Cref{thm:eta-gap-rigidity} supplies \(\omega\) and unit vectors
  \(\xi_i\) satisfying the displayed edge equations. Iterating those
  equations around a closed path \(\gamma:i\to i\) gives
  \[
    \varrho(\tau(\gamma))\xi_i=\omega^{|\gamma|}\xi_i,
  \]
  so \(\xi_i\in W_i(\varrho,\omega)\) for every \(i\). Conversely, if such
  unit vectors in the spaces \(W_i(\varrho,\omega)\) satisfy the edge
  equations, the converse half of \Cref{thm:eta-gap-rigidity} gives
  \(\rho(A_\varrho)=\lambda\).

  It remains only to identify the character data. The closed-path identity
  just displayed says that every generator \(\tau(\gamma)\) of
  \(\Gamma_i(\tau)\) preserves \(L_i\) and acts on it by the scalar
  \(\omega^{|\gamma|}\). Products and inverses of such generators therefore
  also preserve \(L_i\), and the scalar by which an element of
  \(\Gamma_i(\tau)\) acts on the one-dimensional space \(L_i\) defines a
  well-defined group homomorphism \(\chi_i:\Gamma_i(\tau)\to\TT\). This
  homomorphism has the stated value on every closed-path label. Finally,
  iterating the edge equations along a path \(p:i\to j\) gives
  \(\varrho(\tau(p))\xi_j=\omega^{|p|}\xi_i\); hence
  \(\omega^{-|p|}\varrho(\tau(p))\) maps the unit line \(L_j\) unitarily
  onto \(L_i\). This proves the announced classification.
\end{proof}

\begin{proposition}[Loop-subgroup obstruction to Perron-boundary channels]
  \label{prop:eta-gap-loop-subgroup-obstruction}
  Suppose \(G\) is non-trivial, and fix a vertex \(i_0\).  Put
  \(H:=\Gamma_{i_0}(\tau)\) and \(H':=[H,H]\).  If the Perron-boundary
  obstruction \(\eta=\lambda\) occurs, then there is a non-trivial
  irreducible representation \(\varrho\in\Irr(G)\) such that
  \[
    (V_\varrho)^{H'}
    :=
    \{v\in V_\varrho:\varrho(h)v=v\text{ for every }h\in H'\}
  \]
  is non-zero.  Equivalently, if for some vertex \(i_0\) every
  non-trivial irreducible representation \(\varrho\) of \(G\) satisfies
  \((V_\varrho)^{[\Gamma_{i_0}(\tau),\Gamma_{i_0}(\tau)]}=0\), then
  \(\eta<\lambda\).  In that case the strict twisted spectral-gap
  hypothesis required in Main Theorems~II--III, and hence in
  \Cref{thm:class-mertens-explicit,thm:class-euler-product-mertens}, is
  certified by this single loop-subgroup test.
\end{proposition}

\begin{proof}
  Assume first that \(\eta=\lambda\).  By
  \Cref{cor:eta-gap-boundary-classifier}, there are a non-trivial
  irreducible representation \(\varrho\), a phase \(\omega\in\TT\), and a
  boundary line field \((\xi_i)_i\).  Applying
  \Cref{cor:eta-gap-character-classification} at the fixed vertex
  \(i_0\), the line \(L_{i_0}:=\CC\xi_{i_0}\) is invariant under
  \(H=\Gamma_{i_0}(\tau)\), and the action of \(H\) on that line is by a
  one-dimensional character \(\chi:H\to\TT\).  Since \(\TT\) is abelian,
  \(\chi\) is trivial on the commutator subgroup \(H'=[H,H]\).  Hence
  every element of \(H'\) fixes the non-zero vector \(\xi_{i_0}\), so
  \(\xi_{i_0}\in(V_\varrho)^{H'}\) and the displayed fixed-vector space is
  non-zero.

  The equivalent obstruction statement is the contrapositive of the first
  assertion, together with the finite-state diagnostic
  \Cref{cor:eta-gap-boundary-classifier}: if the fixed-vector space above
  vanishes for every non-trivial irreducible \(\varrho\), then the
  boundary case \(\eta=\lambda\) cannot occur.  Since
  \Cref{prop:twisted-spectral-radius-bound} gives
  \(\rho(A_\varrho)\le\lambda\) in every non-trivial channel and the
  maximum defining \(\eta\) is finite, exclusion of equality is exactly
  \(\eta<\lambda\).  The final sentence records the role of this certified
  strict gap in the two Perron-point product theorems cited in the
  statement.
\end{proof}

\begin{corollary}[Perfect loop groups force the eta gap]
  \label{cor:eta-gap-perfect-group}
  Suppose \(G\) is non-trivial. Assume that for some vertex \(i_0\) the
  loop-label subgroup \(\Gamma_{i_0}(\tau)\) is perfect and that, for every
  non-trivial irreducible representation \(\varrho\) of \(G\), the restricted
  representation \(\varrho|_{\Gamma_{i_0}(\tau)}\) has no non-zero fixed
  vector. Then \(\eta<\lambda\). In particular, if
  \(\Gamma_{i_0}(\tau)=G\) for some vertex \(i_0\) and \(G\) is perfect, then
  the strict eta-gap hypothesis required in Main Theorems~II--III holds.
\end{corollary}

\begin{proof}
  Let \(H=\Gamma_{i_0}(\tau)\).  Since \(H\) is perfect, \([H,H]=H\).  The
  fixed-vector hypothesis is therefore precisely the vanishing condition in
  \Cref{prop:eta-gap-loop-subgroup-obstruction}, and that proposition gives
  \(\eta<\lambda\).

  For the stated special case, let \(\Gamma_{i_0}(\tau)=G\) and assume that
  \(G\) is perfect. If a non-trivial irreducible \(\varrho\) had a non-zero
  \(G\)-fixed vector, the span of that vector would be a
  non-zero invariant subspace on which \(G\) acts trivially; irreducibility
  would force \(\varrho\) itself to be the trivial representation. This is
  impossible for \(\varrho\neq\mathbf1\). Thus the fixed-vector hypothesis
  in the first part is automatic, and the strict gap follows.
\end{proof}

\begin{proposition}[Finite-state certificate for the Perron-point
  hypothesis]
  \label{prop:eta-gap-finite-state-certificate}
  Assume \(\lambda>1\). Let \(\widetilde A\) be the skew-product
  adjacency matrix attached to \(\tau\), and set
  \[
    \chi_{\widetilde A}(t):=\det(tI-\widetilde A).
  \]
  The following finite-state conditions are equivalent.
  \begin{enumerate}[label=\rm(\roman*)]
  \item The strict twisted spectral gap used in the class-Mertens and
    Frobenius-class product theorems holds:
    \[
      \eta=\max_{\varrho\neq\mathbf 1}\rho(A_\varrho)<\lambda .
    \]
  \item The skew-product characteristic polynomial has exactly one root on
    the Perron circle, namely a simple root at \(\lambda\); equivalently
    \(\chi_{\widetilde A}(t)=(t-\lambda)q(t)\) with every root of the
    finite polynomial \(q\) having modulus strictly less than \(\lambda\).
  \item No non-trivial irreducible representation admits a boundary line
    field: for every \(\varrho\neq\mathbf1\), the finite algebraic system
    \[
      \varrho(\tau(e))\xi_j=\omega\xi_i\quad(e:i\to j),
      \qquad \omega\in\TT,
      \qquad \norm{\xi_i}=1\quad(i\in V)
    \]
    has no solution.
  \end{enumerate}
  Consequently the Perron-point assumptions in
  \Cref{thm:class-mertens-explicit,thm:class-euler-product-mertens} are
  exactly \(\lambda>1\) plus any one of the finite certificates above.
  The loop-subgroup obstruction of
  \Cref{prop:eta-gap-loop-subgroup-obstruction} gives a one-vertex
  sufficient certificate for item \rm(i).  In particular, the perfect-loop
  criterion of \Cref{cor:eta-gap-perfect-group} is a structural route into
  the same finite-state certificate, not an additional product-constant
  formula.
  Under these certificates the constants
  \(\mathcal M_C(A,\tau;G)\) and \(K_C\) are those supplied by
  \Cref{thm:class-mertens-explicit,thm:class-euler-product-mertens}; if the
  equivalent conditions fail, this paper asserts no Perron-point
  evaluation for the non-trivial primitive channels.
\end{proposition}

\begin{proof}
  The equivalence of \rm(i) and \rm(ii) is exactly the skew-product
  peripheral-spectrum criterion of
  \Cref{lem:twisted-gap-skew-product-peripheral}: the Peter--Weyl block
  decomposition of \(\widetilde A\) gives the spectrum of
  \(\widetilde A\), with algebraic multiplicities, as the disjoint union
  of the spectra of the blocks \(A_\varrho\), each repeated
  \(\dim\varrho\) times. The trivial block is \(A\), and primitivity of
  \(A\) makes \(\lambda\) the only peripheral eigenvalue of that block and
  makes it simple. Thus all remaining roots of
  \(\chi_{\widetilde A}\) lie inside \(|t|<\lambda\) precisely when every
  non-trivial block has spectral radius strictly smaller than \(\lambda\).

  The equivalence of \rm(i) and \rm(iii) is
  \Cref{cor:eta-gap-boundary-classifier} when non-trivial irreducible
  representations exist. If none exist, then \(G\) is trivial, \(\eta=0\)
  by convention, and the finite line-field condition is vacuous, so both
  \rm(i) and \rm(iii) hold because \(\lambda>1\). In the remaining case,
  \Cref{thm:eta-gap-rigidity} says that equality
  \(\rho(A_\varrho)=\lambda\) for a fixed non-trivial irreducible
  representation is equivalent to the displayed line-field equations for
  that representation. Since \(G\) is finite, only finitely many
  irreducible representations occur, and the maximum defining \(\eta\) is
  attained. Therefore the strict gap fails exactly when one of these
  finite systems has a solution.

  The final assertion is a substitution of this certificate into the
  hypotheses of
  \Cref{thm:class-mertens-explicit,thm:class-euler-product-mertens}. Those
  theorems require \(\lambda>1\) and \(\eta<\lambda\) for the Perron-point
  evaluation of the non-trivial primitive Peter--Weyl channels; the present
  proposition only identifies that hypothesis with finite matrix or
  finite algebraic tests.  The cited loop-subgroup obstruction proves a
  sufficient condition for item \rm(i), and the perfect-loop criterion is
  its immediate structural specialization.  None of these diagnostics
  supplies an additional constant formula when the tests fail.
\end{proof}

\begin{proposition}[Perron-point separation argument for Main Theorems II--III]
  \label{prop:eta-gap-perron-separation}
  Assume \(\lambda>1\).  The strict twisted spectral gap
  \(\eta<\lambda\) is equivalent to each of the following Perron-point
  analytic conditions.
  \begin{enumerate}[label=\textup{(\roman*)},leftmargin=*]
  \item For every non-trivial irreducible representation \(\varrho\), the
    branch
    \[
      \log\det(I-zA_\varrho)^{-1}
      =
      \sum_{n\ge1}\frac{\Tr(A_\varrho^n)z^n}{n}
    \]
    used in Main Theorem~I extends holomorphically to a neighbourhood of
    the closed disc \(|z|\le\lambda^{-1}\).
  \item For every non-trivial class function \(\phi\) with zero average on
    \(G\), all non-trivial Peter--Weyl determinant coordinates entering the
    Adams--M\"obius formula for \(\Pi_\phi\) have branch-compatible values at
    \(z=\lambda^{-1}\), and hence the Abel-path Perron values used in
    \Cref{thm:class-mertens-explicit,thm:class-euler-product-mertens} are
    computed entirely inside the finite determinant data.
  \end{enumerate}
  If the strict gap fails, then there are a non-trivial irreducible
  representation \(\varrho\), a phase \(\omega\in\TT\), and a boundary line
  field satisfying \eqref{eq:eta-gap-line-field}.  Moreover
  \(\omega\lambda\in\spec(A_\varrho)\), so the non-trivial determinant
  factor \(\det(I-zA_\varrho)^{-1}\) has a pole at
  \(z=(\omega\lambda)^{-1}\) on the Perron circle.  Thus the finite-state
  certificate of \Cref{prop:eta-gap-finite-state-certificate} is a necessary
  and sufficient verification gate for the Perron-point constants in Main
  Theorems~II--III; it is not an automatic existence theorem for arbitrary
  cocycles.
\end{proposition}

\begin{proof}
  The equivalence between \(\eta<\lambda\) and item~\textup{(i)} is purely
  finite dimensional.  If \(\eta<\lambda\), then every eigenvalue \(\alpha\)
  of every non-trivial block \(A_\varrho\) satisfies
  \(|\alpha|\le\rho(A_\varrho)<\lambda\).  Hence
  \(1/\alpha\), for \(\alpha\neq0\), lies strictly outside the closed disc
  \(|z|\le\lambda^{-1}\).  The determinant
  \(\det(I-zA_\varrho)
  =\prod_\alpha(1-z\alpha)\) therefore has no zero on a neighbourhood of
  that closed disc, and the logarithm branch normalized at \(0\) extends
  holomorphically there.

  Conversely, if item~\textup{(i)} holds and some non-trivial
  \(A_\varrho\) had \(\rho(A_\varrho)=\lambda\), then an eigenvalue
  \(\alpha\in\spec(A_\varrho)\) would satisfy \(|\alpha|=\lambda\).  The
  determinant \(\det(I-zA_\varrho)\) would vanish at
  \(z=\alpha^{-1}\), a point of the Perron circle, so the
  reciprocal determinant and its normalized logarithm could not be
  holomorphic in any neighbourhood of the closed Perron disc.  This
  contradicts item~\textup{(i)}.  Since
  \Cref{prop:twisted-spectral-radius-bound} gives
  \(\rho(A_\varrho)\le\lambda\) for every \(\varrho\), item~\textup{(i)}
  forces \(\eta<\lambda\).

  Item~\textup{(ii)} is exactly the same condition after finite character
  decomposition.  The zero-average class functions have no trivial
  character coordinate.  Their periodic logarithms are finite linear
  combinations of the non-trivial determinant logarithms by
  \Cref{thm:class-function-linearisation}, and the primitive series are then
  obtained from those finitely many logarithms by the Adams--M\"obius formula
  of \Cref{cor:class-function-adams-mobius}.  Therefore the Perron values
  used by \Cref{thm:class-mertens-explicit,thm:class-euler-product-mertens}
  are available from the determinant data precisely when the
  non-trivial determinant branches in item~\textup{(i)} extend through the
  closed Perron disc.  The trivial channel is deliberately excluded here;
  it is singular at \(z=\lambda^{-1}\) and is handled separately by the
  scalar finite-part normalisation.

  Finally suppose the strict gap fails.  Because only finitely many
  irreducibles occur and each non-trivial block has spectral radius at most
  \(\lambda\), there is a non-trivial \(\varrho\) with
  \(\rho(A_\varrho)=\lambda\).  \Cref{thm:eta-gap-rigidity} supplies a phase
  \(\omega\) and a boundary line field satisfying
  \eqref{eq:eta-gap-line-field}; the converse half of its proof constructs
  an eigenvector of the column realization with eigenvalue
  \(\omega\lambda\).  The row and column realizations have the same spectrum,
  so \(\omega\lambda\in\spec(A_\varrho)\).  Hence
  \(\det(I-zA_\varrho)=0\) at \(z=(\omega\lambda)^{-1}\), and the reciprocal
  determinant has a boundary pole.  This is precisely the obstruction
  detected by \Cref{prop:eta-gap-finite-state-certificate}.
\end{proof}

\begin{corollary}[Line fields are the only Perron-boundary obstruction]
  \label{cor:eta-gap-only-boundary-export}
  Under the standing hypotheses with \(\lambda>1\), the results of this
  section export exactly the following finite-state alternative to the main
  theorem chain:
  \begin{enumerate}[label=\textup{(\roman*)},leftmargin=*]
    \item either no non-trivial boundary line field exists, equivalently
      \(\eta<\lambda\), and the Perron-point hypotheses of
      \Cref{thm:class-mertens-explicit,thm:class-euler-product-mertens} are
      verified;
    \item or a non-trivial boundary line field exists, equivalently
      \(\eta=\lambda\), and some non-trivial determinant factor has a pole on
      the Perron circle, so the Perron-point determinant recovery required by
      the class-Mertens and class-product constants is not supplied by this
      paper.
  \end{enumerate}
  In particular, the loop-character and loop-subgroup statements in this
  section are diagnostic consequences of the boundary alternative.  They are
  not structural sufficient conditions for \(\eta<\lambda\) except where the
  stated hypotheses explicitly eliminate every possible boundary line field.
\end{corollary}

\begin{proof}
  \Cref{cor:eta-gap-boundary-classifier} proves that the absence of
  non-trivial boundary line fields is equivalent to \(\eta<\lambda\), while
  the existence of such a line field is equivalent to \(\eta=\lambda\).  The
  first alternative is exactly the strict-gap hypothesis invoked in
  \Cref{thm:class-mertens-explicit,thm:class-euler-product-mertens}.  For the
  second alternative, \Cref{prop:eta-gap-perron-separation} shows that the same
  line field produces a non-trivial eigenvalue on the Perron circle and hence
  a boundary pole in the corresponding determinant factor.  Therefore the
  non-trivial Perron branches used by the main product theorem are available
  precisely in the first alternative.  The loop-character classifier and the
  loop-subgroup obstruction are proved by starting from such a boundary line
  field and recording necessary consequences, or by imposing explicit
  hypotheses that rule out all such fields.  Thus no broader strict-gap
  criterion is asserted.
\end{proof}

\begin{corollary}[Finite-state criterion for the Perron-point constants]
  \label{thm:eta-gap-diagnostic-interface}
  Assume \(G\) is non-trivial and that the operational test of
  \Cref{prop:eta-gap-finite-state-certificate} has been passed, so no
  non-trivial boundary line field exists.  Then the following finite-state
  implication is the complete eta-gap input to the main theorem chain:
  \begin{enumerate}[label=\textup{(\roman*)},leftmargin=*]
    \item the strict-gap hypothesis \(\eta<\lambda\) required in Main
      Theorems~II and~III is available;
    \item the non-trivial determinant logarithms have exactly the
      Perron-point branches certified by
      \Cref{prop:eta-gap-perron-separation};
    \item the constants \(\mathcal M_C(A,\tau;G)\),
      \(F_\varrho(\lambda^{-1})\), and \(K_C\) are then computed by the
      Adams--M\"obius determinant formulas, the fixed-label Euler transform,
      and class-function Fourier inversion, with no additional datum from the
      eta-gap section.
  \end{enumerate}
  Conversely, if the finite-state certificate fails, the boundary-line-field
  alternative of \Cref{thm:eta-gap-rigidity} occurs and this paper supplies no
  replacement sufficient condition for the missing Perron-point constants.
  Broader structural criteria forcing \(\eta<\lambda\) for general cocycle
  classes are outside the scope of this article, except for the explicit
  line-field exclusions proved in this section.
\end{corollary}

\begin{proof}
  \Cref{prop:eta-gap-finite-state-certificate} proves that passing the
  operational certificate is exactly \(\eta<\lambda\), proving item~(i), while
  \Cref{prop:eta-gap-perron-separation} identifies this same condition with
  the holomorphic determinant branches needed at the Perron point, proving
  item~(ii). The role of this condition in the article is the role stated in
  \Cref{thm:class-mertens-explicit,thm:class-euler-product-mertens}: it
  permits evaluation of the non-trivial primitive Peter--Weyl channels at
  \(z=\lambda^{-1}\). Those evaluated channels are subsequently inserted
  into the Fourier formulas of
  \Cref{prop:class-mertens-deviations,thm:class-mertens-fourier} and the
  fixed-label Euler transform of
  \Cref{thm:class-euler-product-mertens,prop:three-logarithm-product-interface,prop:fixed-label-determinant-perron}.
  This proves item~(iii): the line-field diagnostic is used only to decide
  whether this hypothesis holds, so it contributes no independent constant.
  If the certificate fails,
  \Cref{prop:eta-gap-finite-state-certificate,thm:eta-gap-rigidity} give a
  non-trivial boundary line field, and
  \Cref{prop:eta-gap-perron-separation} shows that the needed Perron branch is
  then obstructed.  The final scope sentence follows because
  \Cref{cor:eta-gap-only-boundary-export,prop:eta-gap-conditional-scope} state
  that this section proves only the finite alternative and its explicit
  line-field exclusions, not a broad sufficient-condition theory.
\end{proof}

\begin{proposition}[Finite-state role of the eta-gap hypothesis]
  \label{prop:eta-gap-conditional-scope}
  The eta-gap results in this section give, for the finite edge cocycle fixed
  in the paper, exactly the following conditional diagnostic and no stronger
  structural criterion.
  \begin{enumerate}[label=\textup{(\roman*)},leftmargin=*]
    \item If the finite-state certificate of
      \Cref{prop:eta-gap-finite-state-certificate} succeeds, then
      \(\eta<\lambda\), and the Perron-point evaluations used in Main
      Theorems~II--III are justified by
      \Cref{prop:eta-gap-perron-separation}.
    \item If the certificate fails, then some non-trivial irreducible channel
      admits the boundary line field of \Cref{thm:eta-gap-rigidity}.  In that
      case this article makes no class-Mertens or class-product assertion for
      the failed cocycle beyond the open-disc determinant identities of Main
      Theorem~I.
    \item The only structural sufficient conditions for the strict inequality
      proved here are those obtained by excluding the classified boundary line
      fields, such as the loop-character obstruction and the perfect-loop-group
      consequence.  A broad sufficient-condition theory forcing
      \(\eta<\lambda\) for general cocycle classes is outside this article.
  \end{enumerate}
\end{proposition}

\begin{proof}
  By \Cref{prop:eta-gap-finite-state-certificate}, success of the operational
  finite-state test is equivalent to the absence of a non-trivial Perron
  boundary channel, hence to \(\eta<\lambda\).  Under that strict inequality,
  \Cref{prop:eta-gap-perron-separation} gives exactly the zero-free Perron
  branches needed to evaluate the non-trivial determinant coordinates at
  \(z=\lambda^{-1}\); this proves (i).

  If the same certificate fails, then
  \Cref{prop:eta-gap-finite-state-certificate} supplies a non-trivial channel
  with spectral radius \(\lambda\).  Applying
  \Cref{thm:eta-gap-rigidity} to that channel gives the boundary line field,
  and the proof of \Cref{prop:eta-gap-perron-separation} shows that the
  Perron-point continuation required by the non-trivial class-Mertens and
  class-product constants is not available from the hypotheses of this paper.
  Main Theorem~I is unaffected because it is proved on the open Perron disc and
  does not require evaluation at \(\lambda^{-1}\).  This proves (ii).

  Finally, \Cref{cor:eta-gap-character-classification,prop:eta-gap-loop-subgroup-obstruction,cor:eta-gap-perfect-group}
  all proceed by starting with the line field classified in
  \Cref{thm:eta-gap-rigidity} and then imposing finite hypotheses that forbid
  that line field.  No argument in this section proves that an arbitrary
  broader cocycle class has no such line field.  Therefore the exported theorem
  chain is precisely the conditional one in (i)--(ii), with only the stated
  finite boundary-obstruction eliminations as sufficient conditions.  This proves
  (iii).
\end{proof}

\begin{remark}[Strict-gap criteria boundary]
  \label{rem:eta-gap-criteria-split}
  The diagnostics above characterize the boundary \(\eta=\lambda\) by a
  finite line-field test and provide only the hypothesis check needed for
  Main Theorems~II--III.  The loop-character classifier
  \Cref{cor:eta-gap-character-classification} and perfect-loop-group
  consequence \Cref{cor:eta-gap-perfect-group} record the structural
  consequences of this boundary test that are used in the present paper.
  We do not attempt here to classify all broader structural hypotheses on
  the cocycle \(\tau\) which force \(\eta<\lambda\); only the finite
  boundary obstruction and the perfect-loop-group elimination of that
  obstruction are part of the finite-determinant product-constant theorem
  chain.  These statements are therefore diagnostics for a fixed finite-state
  input, not a promised classification of cocycles with
  \(\eta<\lambda\).  A general theory of structural criteria for strict twisted spectral
  gaps is therefore outside the scope of the present theorem chain.
\end{remark}

% END INPUT sec_eta_gap.tex
% BEGIN INPUT sec_shadows.tex
% BEGIN INPUT sec_shadows_quotient_boundary.tex
\section{Quotient shadows and cyclic-detection boundary}
\label{sec:shadows}

This section proves structural consequences of the finite-group
theorems after the constants in Main Theorems~II--III have already been
constructed.  Quotient functoriality gives the abelianised primitive
profile; the Peter--Weyl splitting then separates the one-dimensional
shadow from the higher-dimensional defect. The final \(S_3\) model is a
certificate for the Adams obstruction, distinguishing the quadratic
sign-channel effect from the certified standard-channel non-abelian
defect. Cyclic quotients of the fixed cocycle are kept separate from the
canonical cyclic-lift root-of-unity layer used later in
\Cref{sec:finite-part}; neither construction is treated here as data
varying over graph resolutions or folds. None of these quotient-shadow
results supplies a new product constant; they recover and test the
constants already fixed by the Main theorem chain.
Throughout this section, ``boundary'' means the exact linear-algebraic
boundary of what cyclic quotients can detect inside the class-function
space of \(G\).

\subsection{Quotient functoriality and the abelian shadow}

\begin{definition}[Normalised primitive class profile]
  \label{def:normalised-primitive-class-profile}
  Under the hypotheses of \Cref{thm:class-mertens-explicit}, define the
  class function
  \[
    \mathcal B_{A,\tau;G}(C):=\frac{|G|}{|C|}\Delta_C(A,\tau;G),
  \]
  where \(\Delta_C(A,\tau;G)=\mathcal M_C(A,\tau;G)-\frac{|C|}{|G|}\mathcal M(A)\).
  If the target group is a quotient \(Q\) and the cocycle is
  \(\bar\tau:E_A\to Q\), we write
  \(\mathcal B_{A,\bar\tau;Q}\) for the same normalised profile computed
  with \(Q\) and \(\bar\tau\).
\end{definition}

By \Cref{prop:class-mertens-deviations},
\begin{equation}\label{eq:normalised-primitive-class-profile}
  \mathcal B_{A,\tau;G}(C)
  =
  \sum_{\varrho\neq \mathbf 1}
    \chi_\varrho^\ast(C)\,\Pi_\varrho(\lambda^{-1}).
\end{equation}

\begin{theorem}[Finite-state projection boundary for scalar inputs]
  \label{thm:scalar-normalisation-projection-boundary}
  Fix \(A\) with Perron root \(\lambda>1\) and fix \(G\). Let
  \[
    \mathcal T_{\mathrm{gap}}(A,G)
    :=
    \{\tau:E_A\to G:\eta(\tau)<\lambda\}.
  \]
  For \(\tau\in\mathcal T_{\mathrm{gap}}(A,G)\), set
  \[
    \mathcal M^\#_{A,\tau;G}(C)
    :=\frac{|G|}{|C|}\mathcal M_C(A,\tau;G)
  \]
  for conjugacy classes \(C\subset G\), and define the profile map
  \[
    \mathcal B:
    \mathcal T_{\mathrm{gap}}(A,G)\longrightarrow L^2_{\mathrm{cl}}(G),
    \qquad
    \mathcal B(\tau):=\mathcal B_{A,\tau;G}.
  \]
  In the Hermitian class-function space of \(G\),
  \begin{equation}\label{eq:scalar-normalisation-projection}
    \mathcal M^\#_{A,\tau;G}
    =
    \mathcal M(A)\mathbf 1_G+\mathcal B_{A,\tau;G},
    \qquad
    \langle \mathcal B_{A,\tau;G},\mathbf 1_G\rangle_G=0.
  \end{equation}
  Thus the orthogonal projection of
  \(\mathcal M^\#_{A,\tau;G}\) onto the scalar class functions is
  \(\mathcal M(A)\mathbf 1_G\), while the projection onto the
  non-trivial Peter--Weyl sector is exactly
  \(\mathcal B_{A,\tau;G}\).

  Consequently, if \(\tau'\!:\!E_A\to G\) is another cocycle over the
  same matrix \(A\), also satisfying the twisted spectral gap, then
  \begin{equation}\label{eq:scalar-normalisation-two-cocycles}
    \mathcal M^\#_{A,\tau;G}-\mathcal M^\#_{A,\tau';G}
    =
    \sum_{\varrho\neq\mathbf 1}
      \chi_\varrho^\ast
      \bigl(\Pi_\varrho^{\tau}(\lambda^{-1})
            -\Pi_\varrho^{\tau'}(\lambda^{-1})\bigr).
  \end{equation}
  In particular, the normalised class constants for \(\tau\) and
  \(\tau'\) agree for every conjugacy class if and only if their
  non-trivial Perron-point primitive Peter--Weyl channels agree for
  every \(\varrho\neq\mathbf 1\).

  More precisely, let \(\mathscr T\) be any
  non-empty family of cocycles \(E_A\to G\) satisfying the twisted
  spectral gap. For this family the following are equivalent.
  \begin{enumerate}[label=\textup{(\roman*)}]
    \item The map
      \[
        \tau\longmapsto \mathcal M^\#_{A,\tau;G}
      \]
      is constant on \(\mathscr T\). Equivalently, any scalar record
      depending only on \(A\), such as
      \(\FP(A)\), \(\mathcal M(A)\), \(F_A\), or
      \(\{\Psi_q(A)\}_{q\ge1}\), has a well-defined single compatible
      normalised class profile on \(\mathscr T\).
    \item The non-trivial class-profile vector
      \(\mathcal B_{A,\tau;G}\) is independent of
      \(\tau\in\mathscr T\).
    \item For every non-trivial \(\varrho\in\Irr(G)\), the value
      \(\Pi_\varrho^\tau(\lambda^{-1})\) is independent of
      \(\tau\in\mathscr T\).
  \end{enumerate}
  Hence two gap cocycles over the same \(A\) with different non-trivial
  Perron-point channels form a finite-state obstruction to scalar
  normalisation: their scalar fixed-matrix records agree, but their
  normalised primitive class profiles do not.

  Equivalently, the fibers of the profile map are exactly the common
  fibers of the non-trivial Perron-point channel map:
  for \(\tau,\tau'\in\mathcal T_{\mathrm{gap}}(A,G)\),
  \begin{equation}\label{eq:scalar-profile-fibers}
    \mathcal B(\tau)=\mathcal B(\tau')
    \quad\Longleftrightarrow\quad
    \Pi_\varrho^\tau(\lambda^{-1})
    =
    \Pi_\varrho^{\tau'}(\lambda^{-1})
    \quad\text{for all }\varrho\neq\mathbf 1.
  \end{equation}
  If \(\mathcal S(A)\) denotes any collection of scalar invariants
  depending only on \(A\), then \(\mathcal S(A)\) is constant on
  \(\mathcal T_{\mathrm{gap}}(A,G)\). Hence on a subfamily
  \(\mathscr T\subset\mathcal T_{\mathrm{gap}}(A,G)\), such scalar
  invariants identify a single normalised class profile precisely when
  \(\mathcal B\) is constant on \(\mathscr T\).
\end{theorem}

\begin{proof}
  By the definition of \(\Delta_C(A,\tau;G)\) in
  \Cref{prop:class-mertens-deviations},
  \[
    \mathcal M_C(A,\tau;G)
    =
    \frac{|C|}{|G|}\mathcal M(A)+\Delta_C(A,\tau;G).
  \]
  Multiplying by \(|G|/|C|\) gives
  \[
    \mathcal M^\#_{A,\tau;G}(C)
    =
    \mathcal M(A)
    +
    \frac{|G|}{|C|}\Delta_C(A,\tau;G)
    =
    \mathcal M(A)+\mathcal B_{A,\tau;G}(C),
  \]
  which proves the first identity in
  \eqref{eq:scalar-normalisation-projection}. The expansion
  \eqref{eq:normalised-primitive-class-profile} contains only
  non-trivial irreducible characters. Since the irreducible characters
  form an orthonormal basis of the class-function space and
  \(\mathbf 1_G\) is the trivial basis vector, this expansion is
  orthogonal to \(\mathbf 1_G\). Hence
  \(\langle \mathcal B_{A,\tau;G},\mathbf 1_G\rangle_G=0\), and the two
  asserted projection statements follow.

  For two cocycles over the same \(A\), the scalar term
  \(\mathcal M(A)\mathbf 1_G\) is identical in the two decompositions.
  Subtracting the two copies of
  \eqref{eq:scalar-normalisation-projection} and using
  \eqref{eq:normalised-primitive-class-profile} for each cocycle gives
  \eqref{eq:scalar-normalisation-two-cocycles}. Orthogonality of the
  irreducible character basis then gives the stated if and only if:
  equality of the normalised class constants is equality of the
  class-function expansions, hence equality of every non-trivial Fourier
  coefficient, and the converse is immediate from the same expansion.

  Finally, the quantities
  \(\FP(A)\), \(\mathcal M(A)\), \(F_A\), \(\{\Psi_q(A)\}_{q\ge1}\), and
  the root-of-unity layers in \(\Cref{sec:finite-part}\) are functions of
  the single scalar matrix \(A\). In the Main Theorem~II decomposition
  just proved, the only term depending solely on \(A\) is
  \(\mathcal M(A)\mathbf 1_G\). The remaining term is the non-trivial
  Peter--Weyl expansion
  \eqref{eq:normalised-primitive-class-profile}, whose coefficients are
  the cocycle-dependent boundary values
  \(\Pi_\varrho^\tau(\lambda^{-1})\).

  It remains to prove the family classification. Since \(A\) and \(G\)
  are fixed, every listed scalar record has the same value for every
  \(\tau\in\mathscr T\). Thus the map
  \(\tau\mapsto\mathcal M^\#_{A,\tau;G}\) is constant on
  \(\mathscr T\) if and only if the right-hand side of
  \[
    \mathcal M^\#_{A,\tau;G}
    =
    \mathcal M(A)\mathbf 1_G+\mathcal B_{A,\tau;G}
  \]
  is independent of \(\tau\), which is equivalent to
  \(\mathcal B_{A,\tau;G}\) being independent of \(\tau\). This proves
  the equivalence of \textup{(i)} and \textup{(ii)}.

  The equivalence of \textup{(ii)} and \textup{(iii)} is exactly the
  uniqueness of Fourier coefficients in the irreducible-character basis
  applied to
  \eqref{eq:normalised-primitive-class-profile}: equality of the class
  functions \(\mathcal B_{A,\tau;G}\) is equivalent to equality of every
  coefficient \(\Pi_\varrho^\tau(\lambda^{-1})\) with
  \(\varrho\neq\mathbf 1\). This proves
  \eqref{eq:scalar-profile-fibers}. The final obstruction statement follows by
  taking a two-element family \(\mathscr T=\{\tau,\tau'\}\). The scalar
  records agree because they depend only on \(A\), while the full
  profiles disagree by the just-proved equivalence. This proves the
  claimed boundary for scalar fixed-matrix inputs.
\end{proof}

\begin{corollary}[No scalar cyclic-lift descent without a constant
  class profile]
  \label{cor:fixed-matrix-record-no-descent}
  Fix \(A\), \(G\), and a non-empty family
  \(\mathscr T\subset\mathcal T_{\mathrm{gap}}(A,G)\). Let
  \(\mathcal R(A)\) be any collection of fixed-matrix scalar records formed
  from \(A\), for instance
  \(\FP(A)\), \(\mathcal M(A)\), \(F_A\), the branch-lifted cyclic-lift
  record \(\{\Psi_q(A)\}_{q\ge1}\), the torsion layers
  \(\{F_A(\omega_q^j)\}_{0\le j<q}\), or the fixed-matrix exponential
  record of \Cref{def:spectral-fingerprints}. These records are attached
  to the single matrix \(A\), with no fold map, resolution change, or
  cocycle descent structure included. Then the following are
  equivalent.
  \begin{enumerate}[label=\textup{(\roman*)}]
    \item The fixed data \(\mathcal R(A)\), regarded as a candidate record
      for the family \(\mathscr T\), determine a single normalised
      primitive class profile on \(\mathscr T\).
    \item The class-profile map
      \(\tau\mapsto\mathcal B_{A,\tau;G}\) is constant on
      \(\mathscr T\).
    \item For every non-trivial \(\varrho\in\Irr(G)\), the Perron-point
      primitive channel
      \(\Pi_\varrho^\tau(\lambda^{-1})\) is independent of
      \(\tau\in\mathscr T\).
  \end{enumerate}
  Consequently, if two gap cocycles over the same matrix \(A\) have
  distinct non-trivial Perron-point primitive channels, then no
  fixed-matrix cyclic-lift record of \(A\) descends to a stable
  cross-cocycle class-profile invariant on the two-point family they
  form. Any stronger invariant requires additional descent data beyond
  the fixed-matrix records of \Cref{def:spectral-fingerprints}.
\end{corollary}

\begin{proof}
  Since \(\mathcal R(A)\) is formed only from the scalar matrix \(A\), it
  has the same value for every \(\tau\in\mathscr T\). Thus it can
  determine at most one candidate class profile on \(\mathscr T\). It
  determines the actual normalised primitive class profile on the whole
  family exactly when
  \(\mathcal B_{A,\tau;G}\) is independent of \(\tau\). This proves the
  equivalence of \textup{(i)} and \textup{(ii)}. The equivalence of
  \textup{(ii)} and \textup{(iii)} is the family classification in
  \Cref{thm:scalar-normalisation-projection-boundary}. The final
  assertion follows by applying these equivalences to
  \(\mathscr T=\{\tau,\tau'\}\) and using the hypothesis that at least one
  non-trivial Perron-point primitive channel differs.
\end{proof}

When the target group matters, we write superscripts \(G\) and \(Q\) on
\(\mathcal L_\phi\), \(\Pi_\phi\), and \(p_{n,C}\) to indicate whether
they are computed with respect to the original cocycle \(\tau:E_A\to G\)
or to a quotient cocycle \(\bar\tau:E_A\to Q\).

\begin{theorem}[Structural consequence: quotient functoriality]
  \label{thm:quotient-functoriality}
  Let \(q:G\twoheadrightarrow Q\) be a surjective homomorphism, and let
  \(\bar\tau:=q\circ\tau\). For every class function \(f:Q\to\CC\), let
  \(q^\ast f:=f\circ q\). Then, for \(|z|<\lambda^{-1}\),
  \begin{equation}\label{eq:quotient-functoriality-log}
    \mathcal L_{q^\ast f}^{G}(z)=\mathcal L_f^{Q}(z),
    \qquad
    \Pi_{q^\ast f}^{G}(z)=\Pi_f^{Q}(z).
  \end{equation}
  If \(D\subset Q\) is a conjugacy class and
  \[
    p^q_{n,D}:=\sum_{\substack{C\subset G\\ q(C)=D}}p_{n,C},
  \]
  then
  \begin{equation}\label{eq:quotient-primitive-class-decomposition}
    \sum_{n\ge 1}p^q_{n,D}z^n
    =
    \frac{|D|}{|Q|}
      \sum_{\pi\in\Irr(Q)}
        \chi_\pi^\ast(D)\,\Pi_\pi^{Q}(z).
  \end{equation}
\end{theorem}

\begin{proof}
  For every periodic point \(x\in\operatorname{Fix}(\sigma^n)\),
  \[
    (q^\ast f)(\tau_n(x))=f(\bar\tau_n(x)).
  \]
  Therefore
  \[
    T^{G}_{q^\ast f,n}
    =
    \sum_{x\in\operatorname{Fix}(\sigma^n)}(q^\ast f)(\tau_n(x))
    =
    \sum_{x\in\operatorname{Fix}(\sigma^n)}f(\bar\tau_n(x))
    =T^{Q}_{f,n},
  \]
  and hence \(\mathcal L_{q^\ast f}^{G}(z)=\mathcal L_f^{Q}(z)\).

  The Adams operations commute with pullback:
  \[
    \psi^m(q^\ast f)(g)
    =(q^\ast f)(g^m)
    =f(q(g)^m)
    =q^\ast(\psi^m f)(g).
  \]
  Applying \Cref{cor:class-function-adams-mobius} therefore gives
  \[
    \Pi_{q^\ast f}^{G}(z)
    =
    \sum_{m\ge 1}\frac{\mu(m)}{m}\,
      \mathcal L_{\psi^m(q^\ast f)}^{G}(z^m)
    =
    \sum_{m\ge 1}\frac{\mu(m)}{m}\,
      \mathcal L_{\psi^m f}^{Q}(z^m)
    =\Pi_f^{Q}(z),
  \]
  which proves \eqref{eq:quotient-functoriality-log}.

  Now take \(f=\mathbf 1_D\). Since
  \[
    q^\ast\mathbf 1_D=\sum_{\substack{C\subset G\\ q(C)=D}}\mathbf 1_C,
  \]
  the left-hand side of \eqref{eq:quotient-functoriality-log} becomes
  \[
    \Pi_{q^\ast\mathbf 1_D}^{G}(z)
    =
    \sum_{n\ge 1}p^q_{n,D}z^n.
  \]
  Applying \Cref{prop:class-primitive-decomposition} to the quotient
  cocycle \(\bar\tau\) yields
  \eqref{eq:quotient-primitive-class-decomposition}.
\end{proof}

\noindent
The abelian shadow is the one-dimensional part of the expansion
\eqref{eq:normalised-primitive-class-profile}; the non-abelian defect is
the complementary higher-dimensional part:
\begin{equation}\label{eq:abelian-shadow-defect}
  \begin{aligned}
    \mathcal B_{A,\tau;G}^{\mathrm{ab}}(C)
    &:=
    \sum_{\substack{\chi\in\Irr(G)\\ \dim\chi=1,\ \chi\neq \mathbf 1}}
      \chi^\ast(C)\,\Pi_\chi(\lambda^{-1}), \\
    \mathcal B_{A,\tau;G}^{\mathrm{na}}(C)
    &:=
    \sum_{\substack{\varrho\in\Irr(G)\\ \dim\varrho>1}}
      \chi_\varrho^\ast(C)\,\Pi_\varrho(\lambda^{-1}).
  \end{aligned}
\end{equation}
Thus
\begin{equation}\label{eq:abelian-shadow-defect-sum}
  \mathcal B_{A,\tau;G}=\mathcal B_{A,\tau;G}^{\mathrm{ab}}+\mathcal B_{A,\tau;G}^{\mathrm{na}}.
\end{equation}

\noindent
We now fix the separation between two cyclic operations. First, one may
quotient the fixed cocycle \(\tau:E_A\to G\). If
\(q_{\mathrm{ab}}:G\twoheadrightarrow G_{\mathrm{ab}}\) is the
abelianisation map and
\(r:G_{\mathrm{ab}}\twoheadrightarrow\ZZ/q\ZZ\) is cyclic, the inherited
cocycle \(r\circ q_{\mathrm{ab}}\circ\tau\) keeps the original edge
labels after reduction modulo \(q\). Its twisted matrices depend on the
matrices of reduced edge-label counts. Second, and independently, one
may change the data by discarding the inherited labels and resetting the
target to \(\ZZ/q\ZZ\) with the canonical cyclic-lift cocycle
\(\tau_q(e)=1\) on every edge, as in
\Cref{rem:abelian-toy-model-cyclic}. Only this reset model produces the
root-of-unity determinant channels
\(\log\det(I-\omega_q^j zA)^{-1}\) from the scalar matrix \(A\) alone.

\begin{corollary}[Abelianisation pullback, not cyclic sampling]
  \label{cor:abelian-cyclic-shadow}
  Let \(q_{\mathrm{ab}}:G\twoheadrightarrow G_{\mathrm{ab}}\) be the
  abelianisation map. Then the one-dimensional component of
  \(\mathcal B_{A,\tau;G}\) is the pullback of the primitive profile of the
  abelianised cocycle:
  \begin{equation}\label{eq:abelian-shadow-pullback}
    \mathcal B_{A,\tau;G}^{\mathrm{ab}}(C)
    =
    \mathcal B_{A,q_{\mathrm{ab}}\circ\tau;G_{\mathrm{ab}}}
      \bigl(q_{\mathrm{ab}}(C)\bigr)
    \qquad (C\subset G).
  \end{equation}
  If, in addition, the data are replaced by the canonical cyclic-lift
  cocycle of \Cref{rem:abelian-toy-model-cyclic}, namely
  \[
    \tau_q:E_A\to\ZZ/q\ZZ,
    \qquad \tau_q(e)=1,
  \]
  then the corresponding
  one-dimensional twisted determinant channels are
  \[
    \log\det(I-\omega_q^j zA)^{-1},
    \qquad 0\le j\le q-1,
  \]
  and these are the root-of-unity layers of \Cref{sec:finite-part}.
  This final statement is a separate canonical cyclic-lift model, not a
  consequence of passing from \(q_{\mathrm{ab}}\circ\tau\) to an arbitrary
  cyclic quotient. Such a quotient gives the primitive profile of its
  inherited cyclic edge-labelling.
\end{corollary}

\begin{proof}
  The characters of \(G_{\mathrm{ab}}\), pulled back along
  \(q_{\mathrm{ab}}\), are precisely the one-dimensional characters of
  \(G\). The abelianised cocycle inherits the Perron-point gap from the
  original system:
  if \(\chi\in\Irr(G_{\mathrm{ab}})\) is non-trivial, then
  \(\chi\circ q_{\mathrm{ab}}\) is a non-trivial character of \(G\), and
  the quotient twisted matrix satisfies
  \[
    A_{\chi}^{q_{\mathrm{ab}}\circ\tau}
    =
    A_{\chi\circ q_{\mathrm{ab}}}.
  \]
  Hence
  \[
    \max_{\chi\neq\mathbf 1}
      \rho\!\left(A_{\chi}^{q_{\mathrm{ab}}\circ\tau}\right)
    \le
    \max_{\varrho\neq\mathbf 1}\rho(A_\varrho)
    <\lambda,
  \]
  so \(\mathcal B_{A,q_{\mathrm{ab}}\circ\tau;G_{\mathrm{ab}}}\) is
  defined by \Cref{thm:class-mertens-explicit}. Applying
  \Cref{thm:quotient-functoriality} to \(q_{\mathrm{ab}}\)
  gives
  \[
    \Pi^{G}_{\chi\circ q_{\mathrm{ab}}}(z)
    =
    \Pi^{G_{\mathrm{ab}}}_{\chi}(z)
    \qquad
    \bigl(\chi\in\Irr(G_{\mathrm{ab}})\bigr).
  \]
  Substituting this identity into
  \eqref{eq:normalised-primitive-class-profile} gives the displayed
  pullback formula \eqref{eq:abelian-shadow-pullback}.

  For the final assertion, specialize instead to the canonical cyclic
  lift. If \(\chi_j(h)=\omega_q^{jh}\) and \(\tau_q(e)=1\) for every
  edge, then
  \[
    (A_{\chi_j}^{\tau_q})_{uv}
    =
    \sum_{e:u\to v}\chi_j(\tau_q(e))
    =
    \omega_q^j A_{uv}.
  \]
  Hence \(A_{\chi_j}^{\tau_q}=\omega_q^jA\), which gives the displayed
  determinant channels. For an inherited quotient
  \(r\circ q_{\mathrm{ab}}\circ\tau\), the same computation gives a
  Fourier sum over the inherited edge-label count matrices, so it
  satisfies the displayed canonical identity only in the special case
  where all reduced labels are \(1\).
\end{proof}

\begin{proposition}[Canonical cyclic-lift boundary for root-of-unity layers]
  \label{prop:canonical-cyclic-lift-root-layer}
  Let \(q\ge 1\), let \(G=\ZZ/q\ZZ\), and let
  \(\tau_q:E_A\to\ZZ/q\ZZ\) be the canonical cyclic-lift cocycle
  \(\tau_q(e)=1\) from \Cref{rem:abelian-toy-model-cyclic}. For
  \(\chi_j(h)=\omega_q^{jh}\),
  \[
    A_{\chi_j}^{\tau_q}=\omega_q^jA
    \qquad (0\le j\le q-1).
  \]
  Hence the corresponding determinant channels are exactly
  \[
    \log\det(I-zA_{\chi_j}^{\tau_q})^{-1}
    =
    \log\det(I-\omega_q^jzA)^{-1},
  \]
  which are the root-of-unity layers used in the scalar comparison of
  \Cref{sec:finite-part}.
  This conclusion uses the reset hypothesis \(\tau_q(e)=1\) from
  \Cref{rem:abelian-toy-model-cyclic}. It is not obtained by applying an
  arbitrary cyclic quotient to \(q_{\mathrm{ab}}\circ\tau\). Thus this is a
  fixed-matrix scalar comparison, not a non-abelian root-of-unity sampling
  theorem, not a cyclic-arithmetic refinement, and not a descent invariant
  across presentations or resolutions.
  If \(q>1\), this same canonical model does not satisfy the non-trivial
  twisted spectral gap: its non-trivial twisted spectral radius is
  \[
    \max_{1\le j\le q-1}\rho(A_{\chi_j}^{\tau_q})=\lambda.
  \]
\end{proposition}

\begin{proof}
  This is the special cocycle singled out in
  \Cref{rem:abelian-toy-model-cyclic}. Since every edge has label \(1\),
  for vertices \(u,v\) one has
  \[
    (A_{\chi_j}^{\tau_q})_{uv}
    =
    \sum_{e:u\to v}\chi_j(1)
    =
    \omega_q^j A_{uv}.
  \]
  Therefore \(A_{\chi_j}^{\tau_q}=\omega_q^jA\), which gives the
  determinant identity. For \(q>1\) there are non-trivial characters
  \(\chi_j\), and
  \(\rho(A_{\chi_j}^{\tau_q})=\rho(\omega_q^jA)=\rho(A)=\lambda\).
  Thus the required strict inequality \(\eta<\lambda\) fails in this
  canonical comparison model.
\end{proof}

\begin{lemma}[Inherited cyclic quotients are not root-of-unity samples]
  \label{lem:cyclic-quotient-versus-root-layer}
  Let \(\bar\tau:E_A\to\ZZ/q\ZZ\) be any cyclic quotient cocycle and set
  \[
    (N_h)_{uv}:=\#\{e:u\to v:\bar\tau(e)=h\},
    \qquad h\in\ZZ/q\ZZ.
  \]
  Then \(A=\sum_hN_h\), and for \(\chi_j(h)=\omega_q^{jh}\)
  \[
    A_{\chi_j}^{\bar\tau}
    =
    \sum_{h\in\ZZ/q\ZZ}\omega_q^{jh}N_h.
  \]
  Hence the twisted matrices are controlled by the inherited edge-label
  count matrices \(N_h\), not by \(A\) alone. For \(q>1\), the full
  root-of-unity identity
  \(A_{\chi_j}^{\bar\tau}=\omega_q^jA\) for all \(j\) holds if and only
  if \(N_1=A\) and \(N_h=0\) for every \(h\neq1\), i.e. precisely when
  the cyclic labelling is the canonical cyclic lift at the level of edge
  counts. Thus, when
  \(\bar\tau=r\circ q_{\mathrm{ab}}\circ\tau\) for a cyclic quotient
  \(r:G_{\mathrm{ab}}\twoheadrightarrow\ZZ/q\ZZ\), the quotient supplies
  its own primitive profile for the inherited labels
  \(r(q_{\mathrm{ab}}(\tau(e)))\). Without the additional canonical
  reset \(\tau_q(e)=1\) from \Cref{rem:abelian-toy-model-cyclic}, it is
  not one of the root-of-unity layers of \Cref{sec:finite-part}.
\end{lemma}

\begin{proof}
  The displayed formula is the definition of the twisted matrix attached
  to the character \(\chi_j\), grouped according to the label
  \(h=\bar\tau(e)\). If \(A_{\chi_j}^{\bar\tau}=\omega_q^jA\) for every
  \(j\), then Fourier inversion on the finite cyclic group gives
  \[
    N_h=\frac1q\sum_{j=0}^{q-1}\omega_q^{-jh}A_{\chi_j}^{\bar\tau}
    =
    \frac1q\sum_{j=0}^{q-1}\omega_q^{j(1-h)}A.
  \]
  The last sum is \(A\) for \(h=1\) and \(0\) otherwise. The converse is
  immediate by substituting \(N_1=A\) and \(N_h=0\) for \(h\neq1\) in the
  displayed expression for \(A_{\chi_j}^{\bar\tau}\).
\end{proof}

% END INPUT sec_shadows_quotient_boundary.tex
% BEGIN INPUT sec_shadows_abelian_and_cyclic.tex
\subsection{Abelian shadow and genuine non-abelian defect}

The decomposition
\eqref{eq:abelian-shadow-defect-sum} is canonical: it is precisely the
orthogonal splitting of the primitive class profile into the
one-dimensional and higher-dimensional Peter--Weyl sectors.

\begin{theorem}[Structural consequence: abelian shadow and genuine
  non-abelian defect]
  \label{thm:abelian-shadow-defect}
  Equip the class functions on \(G\) with the Hermitian inner product
  \[
    \langle F,H\rangle_G
    :=
    \sum_{C\subset G}\frac{|C|}{|G|}F(C)\overline{H(C)}.
  \]
  Let \(\norm{\cdot}_G\) denote the induced norm.
  With respect to this inner product,
  \(\mathcal B_{A,\tau;G}^{\mathrm{ab}}\) and \(\mathcal B_{A,\tau;G}^{\mathrm{na}}\)
  are orthogonal, and
  \begin{equation}\label{eq:abelian-shadow-defect-energy}
    \norm{\mathcal B_{A,\tau;G}^{\mathrm{ab}}}_G^2
    =
    \sum_{\substack{\chi\in\Irr(G)\\ \dim\chi=1,\ \chi\neq \mathbf 1}}
      \abs{\Pi_\chi(\lambda^{-1})}^2,
    \qquad
    \norm{\mathcal B_{A,\tau;G}^{\mathrm{na}}}_G^2
    =
    \sum_{\substack{\varrho\in\Irr(G)\\ \dim\varrho>1}}
      \abs{\Pi_\varrho(\lambda^{-1})}^2.
  \end{equation}
  Moreover, the following are equivalent:
  \begin{enumerate}[label=\textup{(\roman*)}]
    \item \(\mathcal B_{A,\tau;G}\) factors through the abelianisation
      \(q_{\mathrm{ab}}:G\twoheadrightarrow G_{\mathrm{ab}}\);
    \item \(\mathcal B_{A,\tau;G}^{\mathrm{na}}\equiv 0\);
    \item \(\Pi_\varrho(\lambda^{-1})=0\) for every
      \(\varrho\in\Irr(G)\) with \(\dim\varrho>1\).
  \end{enumerate}
\end{theorem}

\begin{proof}
  The orthonormal basis vectors here are the forward class coefficients
  \(C\mapsto \chi_\varrho^\ast(C)\), where
  \[
    \chi_\varrho^\ast(C):=\overline{\chi_\varrho(C)}
    =\chi_\varrho(C^{-1}).
  \]
  These basis vectors are orthonormal for this Hermitian inner product,
  since
  \[
    \sum_{C\subset G}\frac{|C|}{|G|}
      \chi_\varrho^\ast(C)\,\overline{\chi_\pi^\ast(C)}
    =
    \sum_{C\subset G}\frac{|C|}{|G|}
      \chi_\varrho^\ast(C)\,\chi_\pi(C)
    =
    \delta_{\varrho,\pi}.
  \]
  Therefore the one-dimensional and higher-dimensional parts of the
  expansion \eqref{eq:abelian-shadow-defect-sum} are orthogonal, and
  the squared norms are exactly the sums of squared moduli of their
  Fourier coefficients. This gives
  \eqref{eq:abelian-shadow-defect-energy}.

  If \(\mathcal B_{A,\tau;G}\) factors through \(G_{\mathrm{ab}}\), then it
  belongs to the subspace spanned by the pullbacks of the irreducible
  characters of \(G_{\mathrm{ab}}\), that is, by the one-dimensional
  characters of \(G\). Hence all higher-dimensional Fourier coefficients
  in \eqref{eq:normalised-primitive-class-profile} vanish, which proves
  \((i)\Rightarrow(iii)\). The implication \((iii)\Rightarrow(ii)\) is
  immediate from \eqref{eq:abelian-shadow-defect}. Finally,
  \((ii)\Rightarrow(i)\) because a class function whose irreducible
  expansion involves only one-dimensional characters factors through the
  abelianisation.
\end{proof}

\subsection{The exact boundary of cyclic detection}

The previous theorem identifies the canonical splitting. The next
statement concerns the span obtained by using all cyclic quotient
observables together. It says that this aggregate cyclic information is
exactly the one-dimensional sector. An individual cyclic quotient sees
only the one-dimensional characters that factor through that quotient;
here we take the span over all cyclic quotients. This concerns quotient
pullbacks on the fixed cocycle, not the canonical cyclic lift of
\Cref{sec:finite-part}; no root-of-unity arithmetic for non-abelian sectors is
being asserted here.

For later use, write \(\mathscr C_{\mathrm{cyc}}(G)\) for the linear
span of the pullbacks \(q^\ast f\), where \(q:G\twoheadrightarrow \ZZ/m\ZZ\)
ranges over all cyclic quotients and \(f\) ranges over the class
functions on \(\ZZ/m\ZZ\).

\begin{lemma}[Cyclic quotients and one-dimensional characters]
  \label{lem:cyclic-quotients-one-dimensional-characters}
  One has
  \[
    \mathscr C_{\mathrm{cyc}}(G)
    =
    \operatorname{span}\{\chi\in\Irr(G):\dim\chi=1\}.
  \]
\end{lemma}

\begin{proof}
  Let \(q:G\twoheadrightarrow \ZZ/m\ZZ\) be a cyclic quotient and let
  \(f\) be a class function on \(\ZZ/m\ZZ\). Since \(\ZZ/m\ZZ\) is
  abelian, its irreducible characters form a basis of its class
  functions, so
  \[
    f=\sum_{\psi\in\widehat{\ZZ/m\ZZ}}\widehat f(\psi)\psi.
  \]
  Pulling back gives
  \[
    q^\ast f
    =
    \sum_{\psi\in\widehat{\ZZ/m\ZZ}}
      \widehat f(\psi)(\psi\circ q),
  \]
  and each \(\psi\circ q\) is a one-dimensional character of \(G\).
  Hence
  \[
    \mathscr C_{\mathrm{cyc}}(G)
    \subset
    \operatorname{span}\{\chi\in\Irr(G):\dim\chi=1\}.
  \]

  Conversely, let \(\chi\in\Irr(G)\) be one-dimensional. Its image
  \(\chi(G)\subset\CC^\times\) is a finite subgroup of the multiplicative
  group of complex numbers, hence cyclic. Let
  \(q_\chi:G\twoheadrightarrow \chi(G)\) be the induced quotient map,
  and let \(\iota_\chi:\chi(G)\hookrightarrow \CC^\times\) be the
  tautological character. Then
  \[
    \chi=q_\chi^\ast\iota_\chi,
  \]
  so \(\chi\in\mathscr C_{\mathrm{cyc}}(G)\). This proves the reverse
  inclusion.
\end{proof}

\begin{theorem}[Structural consequence: exact boundary of cyclic
  detection]
  \label{thm:cyclic-detection-boundary}
  With \(\mathscr C_{\mathrm{cyc}}(G)\) as above,
  \[
    \mathscr C_{\mathrm{cyc}}(G)
    =
    \operatorname{span}\{\chi\in\Irr(G):\dim\chi=1\}.
  \]
  If \(P_{\mathrm{cyc}}\) denotes the orthogonal projection onto
  \(\mathscr C_{\mathrm{cyc}}(G)\), then
  \[
    P_{\mathrm{cyc}}\mathcal B_{A,\tau;G}=\mathcal B_{A,\tau;G}^{\mathrm{ab}},
    \qquad
    (I-P_{\mathrm{cyc}})\mathcal B_{A,\tau;G}=\mathcal B_{A,\tau;G}^{\mathrm{na}}.
  \]
  Equivalently, for every cyclic quotient \(q:G\twoheadrightarrow \ZZ/m\ZZ\)
  and every class function \(f\) on \(\ZZ/m\ZZ\),
  \[
    \langle \mathcal B_{A,\tau;G}^{\mathrm{na}},q^\ast f\rangle_G=0.
  \]
  In particular, under the hypotheses of
  \Cref{thm:class-mertens-explicit}, the full primitive class profile at
  the Perron point is generated by cyclic quotient shadows if and only if
  \(\mathcal B_{A,\tau;G}^{\mathrm{na}}\equiv 0\).
\end{theorem}

\begin{proof}
  Set
  \[
    V_{\mathrm{ab}}(G)
    :=
    \operatorname{span}\{\chi\in\Irr(G):\dim\chi=1\}.
  \]
  \Cref{lem:cyclic-quotients-one-dimensional-characters} gives
  \(\mathscr C_{\mathrm{cyc}}(G)=V_{\mathrm{ab}}(G)\).
  Theorem~\ref{thm:abelian-shadow-defect} identifies
  \(\mathcal B_{A,\tau;G}^{\mathrm{ab}}\) and \(\mathcal B_{A,\tau;G}^{\mathrm{na}}\) as
  the orthogonal projections of \(\mathcal B_{A,\tau;G}\) onto the
  one-dimensional and higher-dimensional Peter--Weyl sectors. Since the
  first sector is exactly \(\mathscr C_{\mathrm{cyc}}(G)\), the
  projection identities and orthogonality follow immediately.
\end{proof}

\subsection{A cyclic group test for complex-character conventions}

Before the \(S_3\) model we record a minimal check of the conjugation
conventions using \(G=C_3=\langle a\rangle\) with \(|a|=3\). This
example has only non-real characters and therefore exposes the
Hermitian inner-product requirement directly.

\begin{proposition}[Class indicator for \(C_3\)]
  \label{prop:c3-class-indicator-check}
  Let \(G=C_3=\{1,a,a^2\}\) with dual group
  \(\widehat{C_3}=\{\psi_0,\psi_1,\psi_2\}\), where
  \(\psi_k(a^j)=\omega_3^{kj}\) and \(\omega_3=e^{2\pi i/3}\).
  Then for any \(g\in G\),
  \[
    \mathbf 1_{\{a\}}(g)
    =\frac{1}{3}\sum_{k=0}^{2}\overline{\psi_k(a)}\,\psi_k(g)
    =\frac{1}{3}\sum_{k=0}^{2}\psi_k(a^{-1})\,\psi_k(g),
  \]
  whereas the formula without conjugation gives the wrong answer:
  \[
    \frac{1}{3}\sum_{k=0}^{2}\psi_k(a)\,\psi_k(g)
    =\mathbf 1_{\{a^{-1}\}}(g)
    \ne \mathbf 1_{\{a\}}(g).
  \]
\end{proposition}

\begin{proof}
  Since \(\psi_k(a)=\omega_3^k\), we have \(\overline{\psi_k(a)}=\omega_3^{-k}=\psi_k(a^{-1})\).
  The expansion with conjugate evaluates as
  \[
    \frac{1}{3}\sum_{k=0}^{2}\omega_3^{-k}\omega_3^{kj}
    =\frac{1}{3}\sum_{k=0}^{2}\omega_3^{k(j-1)}
    =\mathbf 1_{\{j\equiv 1\bmod 3\}}
    =\mathbf 1_{\{a\}}(a^j).
  \]
  Without conjugation, the same calculation with
  \(\psi_k(a)=\omega_3^k\) gives
  \[
    \tfrac{1}{3}\sum_k\omega_3^{k(j+1)}
    =\mathbf 1_{\{j\equiv -1\bmod 3\}}
    =\mathbf 1_{\{a^{-1}\}}(a^j).
  \]
\end{proof}

\begin{remark}[Minimal conjugation-error detector]
  \label{rem:minimal-conjugation-error-detector}
  \Cref{prop:c3-class-indicator-check} is precisely the counterexample
  described in the proof of \Cref{prop:class-indicator-expansion}: for
  \(G=C_3\), the correct class-indicator formula requires the conjugate
  (or equivalently the group inverse), and the failure without it sends
  every class \(\{a\}\) to the wrong class \(\{a^{-1}\}\). In the
  subsequent \(S_3\) example all characters happen to be real-valued,
  so this failure does not manifest there; the \(C_3\) check is therefore
  the minimal test that can detect a conjugation error.
\end{remark}

% END INPUT sec_shadows_abelian_and_cyclic.tex
% END INPUT sec_shadows.tex
% BEGIN INPUT sec_shadows_s3_certificates.tex
\subsection{The \texorpdfstring{$S_3$}{S3} obstruction certificate}
\label{subsec:s3-obstruction-certificate}

The preceding structural theorems leave one finite-state question open:
whether the higher-dimensional defect can actually occur in a small
symbolic skew product and whether it survives at the product-constant
level of Main Theorem~III.  The answer is yes: the decisive certificate is
the pair of fixed-label Euler-transform values
\(F_\varepsilon(1/2)\) and \(F_\rho(1/2)\), with the primitive-channel
intervals supplying the cyclic-detection recovery.  We record the witness
here and place the routine matrix arithmetic, primitive coefficient
extraction, and rational interval checks in
\Cref{app:s3-model-computations,app:s3-log-certificates}.
For readability, the certificate chain is as follows.  The finite model,
twisted gap, and vanishing sign Artin logarithm are checked in
\Cref{prop:s3-example-channels}; the primitive fixed-label formulae are
computed in \Cref{prop:s3-example-primitive-channels}; the Perron-point
primitive windows are certified in \Cref{lem:s3-example-perron-certificates},
using the rational logarithm estimates of
\Cref{lem:appendix-s3-sharp-certificates}; and the product-level
fixed-label windows used in Main Theorem~III are certified in
\Cref{lem:appendix-s3-f-level-certificates}.

\begin{lemma}[\(S_3\) certificate summary used in the main proof]
  \label{lem:s3-main-text-certificate-summary}
  In the certified \(S_3\) full-shift model below, the only numerical
  inputs imported from the appendices are the following rational windows:
  \begin{enumerate}[label=\textup{(\roman*)},leftmargin=*]
    \item
    \(-342/1000<\Pi_\varepsilon(1/2)<-341/1000\) and
    \(-719/1000<\Pi_\rho(1/2)<-718/1000\), as stated in
    \eqref{eq:s3-example-certified-intervals} and proved from the exact
    rational truncation-and-tail bounds
    \eqref{eq:s3-example-epsilon-certificate-sum}--
    \eqref{eq:s3-example-rho-tail-certificate} in
    \Cref{lem:appendix-s3-sharp-certificates};
    \item
    \(-381/1000<F_\varepsilon(1/2)<-380/1000\) and
    \(-923/1000<F_\rho(1/2)<-922/1000\), as stated in
    \eqref{eq:s3-f-epsilon-appendix-window} and
    \eqref{eq:s3-f-rho-appendix-window} and proved in
    \Cref{lem:appendix-s3-f-level-certificates} by rational logarithm
    enclosures and rational tail bounds.
  \end{enumerate}
  These four intervals are exact enough for the main-text obstruction because
  the first pair excludes zero in the primitive sign and standard channels,
  while the second pair excludes both zero and the sign-channel Artin value
  \(\mathcal L_\varepsilon(1/2)=0\).  No decimal approximation is used in
  the proof of the obstruction theorem; every endpoint above is a rational
  number and every strict inequality is inherited from the cited appendix
  lemma.
\end{lemma}

\begin{proof}
  The appendix lemma
  \Cref{lem:appendix-s3-sharp-certificates} proves rational lower and upper
  bounds for the finite logarithmic sums and for the remaining positive
  tails in the primitive sign and standard channels.  The proof of
  \Cref{lem:s3-example-perron-certificates} imports precisely those bounds
  and records the coarser rational windows in
  \eqref{eq:s3-example-certified-intervals}; these are item~\textup{(i)}.

  The appendix lemma
  \Cref{lem:appendix-s3-f-level-certificates} applies the same rational
  logarithm enclosures to the fixed-label Euler-transform sums and records
  the rational windows
  \eqref{eq:s3-f-epsilon-appendix-window} and
  \eqref{eq:s3-f-rho-appendix-window}; these are item~\textup{(ii)}.  Since
  the displayed intervals are open rational intervals lying strictly on the
  negative side of zero, the primitive windows prove non-vanishing of
  \(\Pi_\varepsilon(1/2)\) and \(\Pi_\rho(1/2)\), and the fixed-label windows
  prove non-vanishing of \(F_\varepsilon(1/2)\) and
  \(F_\rho(1/2)\).  Finally \Cref{prop:s3-example-channels} proves
  \(\mathcal L_\varepsilon(z)\equiv0\), so the fixed-label sign window also
  proves \(F_\varepsilon(1/2)\ne\mathcal L_\varepsilon(1/2)\).  These are
  exactly the numerical hypotheses used below.
\end{proof}

\begin{theorem}[Certified \(S_3\) fixed-label product obstruction]
  \label{thm:s3-certified-nonabelian-defect}
  Let
  \[
    A=
    \begin{pmatrix}1&1\\1&1\end{pmatrix},
    \qquad \lambda=2,
  \]
  and label the four allowed edges of the full two-shift by
  \[
    \tau(0\to0)=(23),\quad \tau(0\to1)=e,
    \quad \tau(1\to0)=(12),\quad \tau(1\to1)=(123)
  \]
  in \(S_3\).  Let \(\varepsilon\) be the sign character and \(\rho\) the
  standard two-dimensional representation.  Then:
  \begin{enumerate}[label=\textup{(\roman*)}]
    \item The fixed-label Euler transforms at the Perron point satisfy
      \[
        -\frac{381}{1000}<F_\varepsilon(1/2)<-\frac{380}{1000},
        \qquad
        -\frac{923}{1000}<F_\rho(1/2)<-\frac{922}{1000}.
      \]
      These two certified non-zero values are the product-level
      obstruction data.  Explicitly, for
      \[
        H(C):=-\frac{6}{|C|}\log K_C-\bigl(\gamma+\log C(A)\bigr),
      \]
      \Cref{thm:class-euler-product-mertens} gives the two-coordinate
      profile
      \[
        H(C)=\chi_\varepsilon(C)F_\varepsilon(1/2)
          +\chi_\rho(C)F_\rho(1/2),
      \]
      so the displayed rational windows force non-zero sign and standard
      product coordinates.  Since \(\mathcal L_\varepsilon(1/2)=0\), the
      class-product constants of
      \Cref{thm:class-euler-product-mertens} cannot in general be obtained
      by the naive Artin replacement in
      \Cref{thm:adams-obstruction-class-product}.  Quantitatively, for at
      least one conjugacy class \(C\subset S_3\),
      \[
        \left|
          \frac{6}{|C|}
          \log\frac{K_C}{K_C^{\mathrm{Art}}}
        \right|
        >\frac{380}{1000}.
      \]
    \item The twisted gap holds: \(\eta<2\).  Moreover
      \(\mathcal L_\varepsilon(z)\equiv0\), while the fixed-label
      primitive channel is certified by
      \[
        -\frac{342}{1000}
        <\Pi_\varepsilon(1/2)<-\frac{341}{1000}.
      \]
      Hence the sign line gives a one-dimensional quadratic Adams
      obstruction: periodic Artin data vanish, but the primitive
      fixed-label channel does not.
    \item The standard channel is certified by
      \[
        -\frac{719}{1000}
        <\Pi_\rho(1/2)<-\frac{718}{1000}.
      \]
      Consequently the non-abelian component
      \(\mathcal B_{A,\tau;S_3}^{\mathrm{na}}\) is non-zero, is
      orthogonal to every cyclic quotient observable, and the primitive
      class profile does not factor through the abelianisation
      \(S_3\twoheadrightarrow C_2\).
  \end{enumerate}
\end{theorem}

\begin{proof}
  By \Cref{lem:s3-main-text-certificate-summary}, every numerical interval
  used in this proof is an imported rational window with a cited appendix
  certificate.

  The fixed-label Euler-product certificates in
  \Cref{lem:appendix-s3-f-level-certificates} give the displayed intervals
  for \(F_\varepsilon(1/2)\) and \(F_\rho(1/2)\).  Since
  \(\mathcal L_\varepsilon(1/2)=0\), the criterion
  \Cref{thm:adams-obstruction-class-product} applies and rules out the
  naive Artin substitution for the associated class-product constants.
  Moreover
  \(|F_\varepsilon(1/2)-\mathcal L_\varepsilon(1/2)|>380/1000\), so the
  Parseval lower bound in \Cref{cor:adams-obstruction-bound} gives the
  displayed normalized class-level error for some conjugacy class.

  The matrices and character table in
  \Cref{prop:s3-example-channels} give \(\eta<2\) and
  \(\mathcal L_\varepsilon(z)\equiv0\).  The primitive-channel identities
  in \Cref{prop:s3-example-primitive-channels}, together with the exact
  rational windows in \Cref{lem:s3-example-perron-certificates}, give the
  two displayed intervals for \(\Pi_\varepsilon(1/2)\) and
  \(\Pi_\rho(1/2)\).  Since the first interval excludes zero while
  \(\mathcal L_\varepsilon(1/2)=0\), the sign line is a quadratic Adams
  obstruction.

  The abelianisation of \(S_3\) is \(C_2\), so
  \Cref{lem:cyclic-quotients-one-dimensional-characters} identifies the
  cyclic quotient shadow with
  \(\operatorname{span}\{\mathbf1,\varepsilon\}\).  The only
  higher-dimensional irreducible representation is \(\rho\); hence
  \Cref{thm:abelian-shadow-defect,thm:cyclic-detection-boundary} identify
  the standard-channel value \(\Pi_\rho(1/2)\) with the full
  non-abelian defect and show that it is orthogonal to every cyclic
  quotient observable.  The certified non-zero interval for
  \(\Pi_\rho(1/2)\) proves non-factorisation through \(C_2\).

  Thus the same finite model simultaneously supplies the product-level
  fixed-label obstruction and the primitive cyclic-detection witness.
\end{proof}

\begin{proposition}[Two-coordinate closure of the certified \(S_3\)
  product obstruction]
  \label{prop:s3-two-coordinate-product-obstruction-vector}
  \label{prop:s3-two-coordinate-product-obstruction-spine}
  In the certified \(S_3\) model of
  \Cref{thm:s3-certified-nonabelian-defect}, write
  \(C_e=\{e\}\), let \(C_t\) be the transposition class, and let
  \(C_3\) be the class of three-cycles.  For the positive constants
  \(K_C\) in \Cref{thm:class-euler-product-mertens}, set
  \[
    H(C):=-\frac{6}{|C|}\log K_C-\bigl(\gamma+\log C(A)\bigr)
    \qquad(C\in\{C_e,C_t,C_3\}).
  \]
  Then the whole non-scalar product-constant vector is exactly the
  two-coordinate Peter--Weyl vector
  \begin{equation}\label{eq:s3-product-obstruction-vector}
    H(C)=\chi_\varepsilon(C)F_\varepsilon(1/2)
      +\chi_\rho(C)F_\rho(1/2),
  \end{equation}
  with character values
  \[
    \chi_\varepsilon=(1,-1,1),
    \qquad
    \chi_\rho=(2,0,-1)
    \quad\text{on }(C_e,C_t,C_3).
  \]
  Both coordinates are non-zero.  More precisely,
  \[
    -\frac{381}{1000}<F_\varepsilon(1/2)<-\frac{380}{1000},
    \qquad
    -\frac{923}{1000}<F_\rho(1/2)<-\frac{922}{1000}.
  \]
  Consequently the obstruction has no hidden third component: after the
  scalar normalisation, it is completely classified by the sign and
  standard coordinates.  Its sign coordinate already forces a non-zero
  Artin-substitution error, while its standard coordinate is orthogonal to
  every class observable factoring through the abelianisation
  \(S_3\twoheadrightarrow C_2\).  Thus cyclic quotient data cannot recover
  the product-constant vector.
\end{proposition}

\begin{proof}
  The irreducible characters of \(S_3\) are the trivial character, the sign
  character \(\varepsilon\), and the standard character \(\rho\), with the
  displayed class values.  Applying the rigidity formula
  \eqref{eq:class-product-rigidity-forward} from
  \Cref{cor:class-product-constant-rigidity} to this group gives
  \[
    H(C)=\chi_\varepsilon^\ast(C)F_\varepsilon(1/2)
      +\chi_\rho^\ast(C)F_\rho(1/2).
  \]
  The two characters are real-valued, so complex conjugation has no effect;
  this proves \eqref{eq:s3-product-obstruction-vector}.  The certified
  rational intervals are exactly the product-level certificates supplied by
  \Cref{thm:s3-certified-nonabelian-defect}, whose proof cites
  \Cref{lem:appendix-s3-f-level-certificates}.  Since both intervals avoid
  zero, both non-trivial coordinates are present.

  The class functions on \(S_3\) have the orthonormal character basis
  \(\{\mathbf1,\varepsilon,\rho\}\).  Therefore the mean-zero, non-scalar
  class-product data have precisely the two coordinates in
  \eqref{eq:s3-product-obstruction-vector}; there is no further irreducible
  sector in which an additional obstruction could live.  Since
  \(\mathcal L_\varepsilon(1/2)=0\) by
  \Cref{thm:s3-certified-nonabelian-defect}, the non-zero sign coordinate
  implies
  \(F_\varepsilon(1/2)\neq\mathcal L_\varepsilon(1/2)\), and the obstruction
  criterion \Cref{thm:adams-obstruction-class-product} gives a non-zero
  Artin-substitution error.

  Finally, a class function factors through the abelianisation
  \(S_3\twoheadrightarrow C_2\) if and only if it lies in the span of the
  two one-dimensional characters \(\mathbf1\) and \(\varepsilon\).  Character
  orthogonality puts \(\rho\) in the orthogonal complement of that span.
  The certified inequality \(F_\rho(1/2)\neq0\) therefore gives a genuine
  standard-sector component invisible to all abelian, equivalently cyclic,
  quotient class observables.  Hence the cyclic quotient data do not
  determine the full product-constant vector.
\end{proof}

\begin{corollary}[Product-constant obstruction is already finite-state]
  \label{cor:s3-product-constant-obstruction-finite-state}
  In the certified \(S_3\) full-shift model of
  \Cref{thm:s3-certified-nonabelian-defect}, the obstruction to replacing
  the Adams-corrected constants of Main Theorem~III by periodic Artin
  logarithms is completely detected by the two finite rational windows for
  \(F_\varepsilon(1/2)\) and \(F_\rho(1/2)\).  In particular, no cyclic
  quotient recovery can recover the product vector
  \((K_C)_{C\subset S_3}\): the sign window already forces a non-zero
  normalized Artin-substitution error for some class, and the standard
  window lies in the non-abelian Peter--Weyl sector orthogonal to every
  cyclic quotient observable.
\end{corollary}

\begin{proof}
  The first assertion is the fixed-label part of
  \Cref{thm:s3-certified-nonabelian-defect} together with the determinant
  rigidity and obstruction criterion
  \Cref{thm:class-product-determinant-rigidity,thm:adams-obstruction-class-product}:
  the normalized class-product vector is the Peter--Weyl expansion whose
  non-trivial coordinates are the fixed-label values
  \(F_\varrho(1/2)\), and equality with the naive Artin vector would force
  \(F_\varrho(1/2)=\mathcal L_\varrho(1/2)\) for every non-trivial
  \(\varrho\).  The certified sign interval and
  \(\mathcal L_\varepsilon(1/2)=0\) therefore give a non-zero discrepancy;
  \Cref{cor:adams-obstruction-bound} converts it into the stated classwise
  normalized lower bound.  The second assertion follows from
  \Cref{thm:abelian-shadow-defect,thm:cyclic-detection-boundary}: cyclic
  quotient observables see only the one-dimensional characters of
  \(S_3\), while the certified non-zero \(F_\rho(1/2)\) is in the standard
  two-dimensional Peter--Weyl sector.
\end{proof}
% END INPUT sec_shadows_s3_certificates.tex
% BEGIN INPUT sec_finite_part.tex
\section{Scalar normalisation}
\label{sec:finite-part}

Throughout this section, \(A\) denotes a primitive non-negative real
matrix with Perron root \(\lambda>1\). The symbols \(p_n(A)\) denote the
algebraic M\"obius coefficients of \eqref{eq:primitive-mobius}; they are
primitive orbit counts only in the integer multigraph case fixed in
\Cref{sec:preliminaries}. For the finite-group results, the scalar role
was delimited in \Cref{sec:intro-scope-of-results}: only the
trivial-channel finite part \(\mathcal M(A)\) is used in Main
Theorems~II--III. We therefore keep the finite-part formula in full and
record the lift comparison only in a compressed fixed-matrix form; its
determinant manipulations and supporting proofs are deferred to
\Cref{app:cyclic-lift-comparison} and are not part of the
Adams--M\"obius primitive-inversion proof.

\subsection{Finite parts and reduced determinants}

This subsection is classical scalar infrastructure rather than a new
contribution.  It collects the Perron-pole normalisation of the SFT zeta
function from \Cref{prop:prelim-trace-primitive} and the corresponding
dynamical Mertens finite part in the form needed by the finite-group
sections; compare \cite{BowenLanford1970Zeta,Parry1983PrimeOrbit,Sharp1991Mertens,LindMarcus1995,ParryPollicott1990Zeta}.  We keep the
short proofs to make the constants \(C(A)\), \(\FP(A)\), and
\(\mathcal M(A)\) unambiguous.

\begin{lemma}[Classical scalar normalisation at the Perron pole]
  \label{lem:finite-part-primitive-closed-form}
  Let \(A\) be a primitive non-negative matrix with Perron root
  \(\lambda>1\). Then
  \begin{equation}\label{eq:finite-part-determinant-mobius}
    \FP(A)
    =
    \log C(A)
    +\sum_{m\ge 2}\frac{\mu(m)}{m}
      \log\det(I-\lambda^{-m}A)^{-1}.
  \end{equation}
  Equivalently,
  \begin{equation}\label{eq:finite-part-primitive-closed}
    \FP(A)
    =
    \log C(A)
    +\sum_{n\ge 1}p_n(A)
      \bigl(\lambda^{-n}+\log(1-\lambda^{-n})\bigr).
  \end{equation}
  Both series converge absolutely. In particular, if
  \[
    \psi(x):=-x-\log(1-x)=\sum_{k\ge 2}\frac{x^k}{k},
    \qquad 0<x<1,
  \]
  then
  \begin{equation}\label{eq:finite-part-gap-series}
    \log C(A)-\FP(A)
    =
    \sum_{n\ge 1}p_n(A)\psi(\lambda^{-n}).
  \end{equation}
\end{lemma}

\begin{proof}
  This is the standard scalar Euler-product calculation, recorded here
  only to fix the normalisation. Put
  \(\zeta_A(z)=\det(I-zA)^{-1}\). For \(|z|<\lambda^{-1}\),
  \[
    \log\zeta_A(z)
    =
    \sum_{n\ge 1}\frac{\Tr(A^n)}{n}z^n
    =
    \sum_{\ell\ge 1}p_\ell(A)\sum_{k\ge 1}\frac{z^{k\ell}}{k}.
  \]
  M\"obius inversion in the repetition variable gives
  \[
    \sum_{m\ge 1}\frac{\mu(m)}{m}\log\zeta_A(z^m)
    =
    \sum_{\ell\ge 1}p_\ell(A)
      \sum_{q\ge 1}
        \frac{z^{q\ell}}{q}\sum_{m\mid q}\mu(m)
    =
    \sum_{\ell\ge 1}p_\ell(A)z^\ell.
  \]
  With \(z=r/\lambda\), this says
  \[
    P_A(r)
    =
    \log\zeta_A(r/\lambda)
    +\sum_{m\ge 2}\frac{\mu(m)}{m}
      \log\zeta_A\bigl((r/\lambda)^m\bigr).
  \]
  Since
  \[
    \log\zeta_A(r/\lambda)+\log(1-r)
    =
    \log\bigl((1-r)\zeta_A(r/\lambda)\bigr)
    \longrightarrow \log C(A)
  \]
  as \(r\uparrow 1\), it remains to pass to the limit in the \(m\ge2\)
  sum. This is justified by
  \[
    \log\zeta_A(\lambda^{-m})
    =
    \sum_{j\ge 1}\frac{\Tr(A^j)}{j}\lambda^{-mj}
    =
    O(\lambda^{1-m}),
    \qquad m\ge 2,
  \]
  so the determinant series is absolutely summable. Letting
  \(r\uparrow1\) gives \eqref{eq:finite-part-determinant-mobius}.

  For the primitive-orbit form, write the same Euler product at
  \(z=r/\lambda\) as
  \[
    \log\zeta_A(r/\lambda)
    =
    P_A(r)
    +\sum_{n\ge 1}p_n(A)\sum_{k\ge 2}
      \frac{r^{kn}\lambda^{-kn}}{k}.
  \]
  By \Cref{lem:primitive-orbit-asymptotic},
  \(\abs{p_n(A)}=O(\lambda^n/n)\), so the \(k\ge2\) tail is absolutely
  summable at \(r=1\). Hence
  \[
    \FP(A)
    =
    \log C(A)
    -\sum_{n\ge 1}p_n(A)\sum_{k\ge 2}\frac{\lambda^{-kn}}{k}.
  \]
  Finally,
  \[
    -\sum_{k\ge 2}\frac{x^k}{k}=x+\log(1-x),
    \qquad |x|<1,
  \]
  with \(x=\lambda^{-n}\), which proves
  \eqref{eq:finite-part-primitive-closed}. The displayed gap formula is
  the same identity rewritten with \(\psi\).
\end{proof}

\begin{proposition}[Classical reduced determinant identity]
  \label{prop:finite-part-residue-reduced-determinant}
  This is the standard Perron-pole reduced-determinant calculation for a
  finite SFT zeta function; see
  \cite{BowenLanford1970Zeta,LindMarcus1995,ParryPollicott1990Zeta}.  It
  is included only to fix the notation \(C(A)\) used in the scalar
  normalisation.
  Let \(\mu_2,\dots,\mu_d\) be the non-Perron eigenvalues of \(A\). Then
  \[
    C(A)
    =\prod_{j=2}^d\left(1-\frac{\mu_j}{\lambda}\right)^{-1}
    =\det\nolimits'(I-A/\lambda)^{-1}.
  \]
\end{proposition}

\begin{proof}
  Factor
  \[
    \det(I-zA)
    =(1-\lambda z)\prod_{j=2}^d(1-\mu_j z).
  \]
  Multiplying by \(1-\lambda z\) and letting \(z\to\lambda^{-1}\) gives
  the claim.
\end{proof}

\begin{corollary}[Classical scalar orbit-Mertens asymptotic]
  \label{cor:finite-part-mertens-asymptotic}
  \leavevmode\par\noindent
  Let \(A\) be a primitive non-negative matrix with Perron root
  \(\lambda>1\). Then
  \[
    \sum_{n\le N}\frac{p_n(A)}{\lambda^n}
    =\log N+\mathcal M(A)+o(1)
    \qquad (N\to\infty),
  \]
  where \(\mathcal M(A)=\gamma+\FP(A)\).
\end{corollary}

\begin{proof}
  This is the standard dynamical Mertens summation for the scalar
  primitive channel, included only as the normalisation used by
  \Cref{thm:class-mertens-explicit,cor:assembled-adams-frobenius-product-formula,cor:class-product-stability-scalar-independence}; see
  \cite{Sharp1991Mertens,ParryPollicott1990Zeta} for the classical orbit
  product context.
  By \Cref{lem:primitive-orbit-asymptotic}, there exists
  \(\vartheta\in(0,1)\) such that
  \[
    \frac{p_n(A)}{\lambda^n}
    =
    \frac{1}{n}+O\!\left(\frac{\vartheta^n}{n}\right).
  \]
  Hence the correction series
  \[
    \sum_{n\ge 1}\left(\frac{p_n(A)}{\lambda^n}-\frac{1}{n}\right)
  \]
  converges absolutely, and \Cref{lem:abel-correction-sum} gives
  \[
    \sum_{n\ge 1}\left(\frac{p_n(A)}{\lambda^n}-\frac{1}{n}\right)
    =
    \FP(A).
  \]
  Therefore
  \[
    \sum_{n\le N}\frac{p_n(A)}{\lambda^n}
    =
    H_N+\FP(A)
    -\sum_{n>N}\left(\frac{p_n(A)}{\lambda^n}-\frac{1}{n}\right).
  \]
  The tail is \(O(\vartheta^N/N)\) by
  \Cref{lem:geometric-harmonic-tail}, while
  \(H_N=\log N+\gamma+o(1)\). Since \(\mathcal M(A)=\gamma+\FP(A)\),
  the claimed asymptotic follows.
\end{proof}

\subsection{Scalar interface for the finite-group results}

The remaining scalar material is used only as an interface for the
finite-group determinant and product results. It is not a source of
non-trivial class-product coordinates: by
\Cref{thm:class-product-determinant-rigidity}, those coordinates are exactly
the fixed-label Peter--Weyl Perron values. The canonical cyclic lift
\(A^{\langle q\rangle}=A\otimes P_q\) from \Cref{eq:cyclic-lift-def}, the
branch-lifted logarithmic record \(\{\Psi_q(A)\}_{q\ge1}\), and the
root-of-unity layers of \(F_A\) are fixed-matrix records. Their
reconstruction theorem is proved in \Cref{app:cyclic-lift-comparison}; the
main text uses them only through the following boundary statement.

\begin{theorem}[Scalar interface for trivial-channel normalisation]
  \label{thm:finite-part-scalar-interface-no-descent}
  Fix a primitive non-negative matrix \(A\) with Perron root
  \(\lambda>1\), a finite group \(G\), and a non-empty family
  \(\mathscr T\subset\mathcal T_{\mathrm{gap}}(A,G)\) of cocycles
  satisfying the twisted spectral gap; here
  \(\mathcal T_{\mathrm{gap}}(A,G)\) is the set of cocycles
  \(E_A\to G\) with \(\eta(\tau)<\lambda\). For
  \(\tau\in\mathscr T\), let
  \[
    \mathcal M^\#_{A,\tau;G}(C):=\frac{|G|}{|C|}\mathcal M_C(A,\tau;G)
  \]
  and let \(\mathcal B_{A,\tau;G}\) be the normalised primitive class
  profile of \Cref{def:normalised-primitive-class-profile}. Then
  \begin{equation}\label{eq:finite-part-scalar-interface-splitting}
    \mathcal M^\#_{A,\tau;G}
    =
    \mathcal M(A)\mathbf 1_G+
    \mathcal B_{A,\tau;G},
  \end{equation}
  and the first summand is the entire contribution determined solely by
  the scalar matrix \(A\).

  More precisely, let \(\mathcal R(A)\) be any fixed-matrix scalar or
  canonical cyclic-lift record named in \Cref{app:cyclic-lift-comparison}.
  The record \(\mathcal R(A)\)
  descends to a single normalised primitive class profile on
  \(\mathscr T\) if and only if the following equivalent conditions hold:
  \begin{enumerate}[label=\textup{(\roman*)}]
    \item \(\mathcal B_{A,\tau;G}\) is independent of
      \(\tau\in\mathscr T\);
    \item for every non-trivial \(\varrho\in\Irr(G)\), the Perron-point
      primitive channel \(\Pi_\varrho^\tau(\lambda^{-1})\) is independent
      of \(\tau\in\mathscr T\).
  \end{enumerate}
  Consequently, if two gap cocycles over the same matrix \(A\) have
  distinct non-trivial Perron-point primitive channels, then all scalar
  finite-part and fixed-matrix cyclic-lift records agree for the two
  cocycles, but their normalised primitive class profiles differ.
\end{theorem}

\begin{proof}
  By \Cref{prop:class-mertens-deviations}, for every conjugacy class
  \(C\subset G\),
  \[
    \mathcal M_C(A,\tau;G)
    =
    \frac{|C|}{|G|}\mathcal M(A)+\Delta_C(A,\tau;G)
  \]
  and
  \[
    \Delta_C(A,\tau;G)
    =
    \frac{|C|}{|G|}
      \sum_{\varrho\neq\mathbf1}
        \chi_\varrho^\ast(C)\,\Pi_\varrho^\tau(\lambda^{-1}).
  \]
  Multiplying the first identity by \(|G|/|C|\) and using the definition
  of \(\mathcal B_{A,\tau;G}\) gives
  \eqref{eq:finite-part-scalar-interface-splitting}. The displayed
  Peter--Weyl expansion contains no trivial character term, so the scalar
  projection of \(\mathcal M^\#_{A,\tau;G}\) is exactly
  \(\mathcal M(A)\mathbf1_G\), while all non-scalar information lies in
  \(\mathcal B_{A,\tau;G}\).

  Every record \(\mathcal R(A)\) named in the statement is a function of
  the single matrix \(A\); it is therefore constant as \(\tau\) varies in
  \(\mathscr T\). Such a fixed record can assign one normalised primitive
  class profile to the whole family exactly when the actual profiles
  \(\mathcal B_{A,\tau;G}\) are all equal. Finally, the irreducible
  characters form a basis of the class-function space, so equality of the
  expansions
  \[
    \mathcal B_{A,\tau;G}(C)
    =
    \sum_{\varrho\neq\mathbf1}
      \chi_\varrho^\ast(C)\,\Pi_\varrho^\tau(\lambda^{-1})
  \]
  for all classes \(C\) is equivalent to equality of every coefficient
  \(\Pi_\varrho^\tau(\lambda^{-1})\) with
  \(\varrho\neq\mathbf1\). Applying this criterion to a two-point family
  gives the final obstruction statement.
\end{proof}


\begin{corollary}[Only the trivial channel descends from fixed-matrix
  scalar data]
  \label{cor:fixed-matrix-scalar-interface-only}
  In the setting of \Cref{thm:finite-part-scalar-interface-no-descent},
  the projection of the normalised class-Mertens profile onto the constant
  class functions is exactly \(\mathcal M(A)\mathbf1_G\).  Hence any
  scalar finite-part formula or canonical cyclic-lift reconstruction that
  depends only on \(A\) can enter the finite-group theorem chain only
  through the trivial-channel normalisation \(\mathcal M(A)\).  It cannot
  determine the non-trivial class profile for a family \(\mathscr T\)
  unless the non-trivial Perron-point channels
  \(\Pi_\varrho^\tau(\lambda^{-1})\), \(\varrho\neq\mathbf1\), are
  already independent of \(\tau\) on that family.
\end{corollary}

\begin{proof}
  The splitting
  \eqref{eq:finite-part-scalar-interface-splitting} writes the normalised
  profile as the sum of the constant class function
  \(\mathcal M(A)\mathbf1_G\) and
  \(\mathcal B_{A,\tau;G}\).  By definition of
  \(\mathcal B_{A,\tau;G}\) in
  \Cref{def:normalised-primitive-class-profile}, the latter has only
  non-trivial Peter--Weyl coefficients, namely the values
  \(\Pi_\varrho^\tau(\lambda^{-1})\) with
  \(\varrho\neq\mathbf1\).  Orthogonality of irreducible characters
  therefore identifies the constant projection with
  \(\mathcal M(A)\mathbf1_G\) and the non-constant projection with the
  displayed non-trivial channels.  A record depending only on \(A\) is
  constant as \(\tau\) varies, so it can determine this non-constant
  projection throughout \(\mathscr T\) exactly in the independence case
  already characterized in
  \Cref{thm:finite-part-scalar-interface-no-descent}.
\end{proof}

\begin{corollary}[Rank-zero scalar obstruction for class profiles]
  \label{cor:scalar-record-rank-zero-bound}
  Keep the notation and hypotheses of
  \Cref{thm:finite-part-scalar-interface-no-descent}.  Let
  \(\mathscr T\subset\mathcal T_{\mathrm{gap}}(A,G)\) be non-empty, and let
  \[
    \mathcal Z(\mathscr T)
    :=
    \operatorname{span}_{\mathbb{R}}
    \{\mathcal B_{A,\tau;G}-\mathcal B_{A,\tau';G}:
        \tau,\tau'\in\mathscr T\}
  \]
  be the real zero-mean class-function space actually swept out by the
  non-trivial primitive profiles in the family.  Then the scalar finite part
  \(\FP(A)\), the scalar Mertens normalisation \(\mathcal M(A)\), the reduced
  determinant \(C(A)\), and every canonical cyclic-lift fixed-matrix record
  named in \Cref{app:cyclic-lift-comparison} have rank zero on
  \(\mathcal Z(\mathscr T)\): they are constant on \(\mathscr T\) and
  annihilate all profile differences.

  Consequently, any finite collection of real linear invariants whose restriction
  to profile differences is linear and which, together with the scalar
  records, recovers the normalised class-Mertens profile of Main Theorem~II on
  \(\mathscr T\) must have real rank at least
  \(\dim\mathcal Z(\mathscr T)\) on the non-trivial Peter--Weyl sector.  The
  same lower bound applies to any such linear recovery procedure for the
  product constants \((K_C)_{C\subset G}\) of Main Theorem~III after the scalar
  product normalisation is fixed.  In particular, when the profiles in
  \(\mathscr T\) affinely span the
  whole zero-mean class-function space, at least
  \(\#\{\textup{conjugacy classes of }G\}-1\) independent non-scalar
  coordinates are necessary, and scalar finite-matrix records supply none of
  them.
\end{corollary}

\begin{proof}
  The records listed in the statement depend only on the matrix \(A\), not on
  the cocycle \(\tau\).  They are therefore constant on the family
  \(\mathscr T\).  By
  \Cref{eq:finite-part-scalar-interface-splitting}, for every
  \(\tau,\tau'\in\mathscr T\),
  \[
    \mathcal M^\#_{A,\tau;G}-\mathcal M^\#_{A,\tau';G}
    =
    \mathcal B_{A,\tau;G}-\mathcal B_{A,\tau';G}.
  \]
  The right-hand side is a zero-mean class function, since
  \(\mathcal B_{A,\tau;G}\) has only non-trivial Peter--Weyl coefficients by
  \Cref{def:normalised-primitive-class-profile} and
  \eqref{eq:normalised-primitive-class-profile}.  Thus scalar records see no
  vector in \(\mathcal Z(\mathscr T)\), which is the asserted rank-zero
  statement.

  Let \(I\) be any additional finite real invariants whose restriction
  to profile differences is linear and which, after adjoining the scalar
  records, recovers the normalised class-Mertens profile on \(\mathscr T\).
  Since the scalar part is constant on \(\mathscr T\), the induced linear
  difference map must distinguish every vector in
  \(\mathcal Z(\mathscr T)\).  Therefore its restriction to the real vector
  space \(\mathcal Z(\mathscr T)\) is injective, and elementary linear algebra
  forces its rank to be at least \(\dim\mathcal Z(\mathscr T)\).

  For product constants, \Cref{thm:class-euler-product-mertens} and
  \Cref{cor:class-product-hyperplane-stability} identify the logarithmic
  product vector, after scalar normalisation, with the same zero-mean
  Peter--Weyl boundary coordinates.  Thus recovering \((K_C)_C\) on
  \(\mathscr T\) is equivalent to recovering the corresponding zero-mean
  profile differences, so the same rank lower bound applies.  If the affine
  span is the whole zero-mean class-function space, its dimension is the
  number of conjugacy classes of \(G\) minus one, because the full real
  class-function space has one coordinate per conjugacy class and the scalar
  direction is one-dimensional.
\end{proof}

\noindent
The fixed-matrix reconstruction and one-layer Fourier inversion
statements are
\Cref{thm:finite-part-cyclic-lift-spectrum-identifiability,thm:finite-part-cyclic-lift-single-q-torsion-reconstruction} in
\Cref{app:cyclic-lift-comparison}. They identify only the scalar comparison
model and, by the fixed-matrix qualification following \Cref{thm:finite-part-cyclic-lift-spectrum-identifiability}, do not descend to the
non-trivial branch-compatible class-product coordinates. They are not used in
the Adams--M\"obius finite-group arguments;
\Cref{cor:class-product-stability-scalar-independence} records the
corresponding finite-state obstruction on the product side.

% END INPUT sec_finite_part.tex
% BEGIN INPUT sec_conclusion.tex
\section{Discussion and open problems}\label{sec:conclusion}

\subsection{Summary of the main contribution}

For a finite-group extension of a mixing subshift of finite type, the
twisted determinant family records periodic labels. The main point of
this paper is that the corresponding primitive labels are recovered by a
finite Adams--M\"obius calculation in Peter--Weyl coordinates. This is
the finite-state correction forced by the repetition rule
\(g\mapsto g^m\) for primitive orbits. Under the twisted gap
\(\eta<\lambda\), the recovered primitive channels have Perron-point
values and hence determine the class-indexed constants
\(\mathcal M_C(A,\tau;G)\) in the Frobenius-class Mertens sums.

The Euler-product constants require the same distinction. Once the
classical leading exponent \(|C|/|G|\) has been separated, the constant is
governed by the fixed-label Euler transform
\(F_\varrho(\lambda^{-1})\), not by the periodic Artin logarithm
\(\mathcal L_\varrho(\lambda^{-1})\). The determinant formula therefore
gives
\[
  K_C=
  \exp\!\left(
    -\frac{|C|}{|G|}
    \left(
      \gamma+\log C(A)
      +\sum_{\varrho\neq\mathbf1}
        \chi_\varrho^\ast(C)F_\varrho(\lambda^{-1})
    \right)
  \right),
\]
with the positive real branch
(\Cref{cor:class-function-adams-mobius,thm:class-mertens-explicit,thm:class-euler-product-mertens,prop:fixed-label-determinant-perron,cor:class-product-constant-rigidity,thm:class-mertens-fourier}). This is a determinant-level refinement of
the classical Chebotarev/Mertens orbit-counting picture, rather than a
new leading-density theorem. Quotient functoriality identifies the
abelianised primitive profile, and cyclic quotient observables span
precisely the one-dimensional Peter--Weyl shadow. The \(S_3\) example
shows both phenomena concretely: the sign line detects an Adams
correction invisible to the Artin logarithm, while the standard
two-dimensional representation gives a genuine non-abelian component.

The scalar finite-part and cyclic-lift calculations supply only the
trivial-channel normalisation and a fixed-matrix comparison model. The
projection theorem
\Cref{thm:scalar-normalisation-projection-boundary} marks the boundary:
after normalising the class constants, scalar data give the constant
class-function component \(\mathcal M(A)\mathbf 1_G\), whereas the
non-trivial primitive profile is the Peter--Weyl vector
\(\mathcal B_{A,\tau;G}\).

\subsection{Further directions}

Several natural extensions remain outside the finite-dimensional theorem chain
proved here and are not used as hidden inputs.  One direction is to replace the
finite SFT determinant model by an infinite-dimensional H\"older/Ruelle
transfer-operator setting with compatible determinant branches.  A second is a
completed-zeta or xi-normalised theory built after the finite determinant data
have been fixed.  A third is to give structural sufficient criteria for the
strict twisted gap beyond the finite line-field tests of
\Cref{prop:eta-gap-finite-state-certificate}.  A fourth is to develop
secondary spectral-edge asymptotics as a separate refinement of the product
expansion rather than the conditional appendix statement
\Cref{prop:main-chain-secondary-noninterference}.  The scalar cyclic-lift
reconstruction results suggest fixed-matrix inverse problems, but additional
descent data would be needed before they could recover non-trivial finite-group
class-product coordinates.  Cyclic arithmetic or non-abelian root-layer
sampling beyond the boundary theorem is likewise left to a separate treatment.
Sibling determinant, synchronisation, and language-obstruction packages, as
well as unrelated universality-hierarchy material from the project ledger, are
not part of this article's theorem chain.  The explicit \(S_3\) calculation in
this paper should be read only as a finite witness for the Adams correction in
the determinant model.
% END INPUT sec_conclusion.tex
\appendix
% BEGIN INPUT sec_appendices.tex
\section{Appendix: scalar cyclic-lift comparison}
\label{app:cyclic-lift-comparison}

This appendix collects the scalar cyclic-lift comparison arguments behind the
normalisation and comparison statements in the main text.  The reconstruction
statements are fixed-matrix scalar facts: they support the trivial-channel
finite part \(\mathcal M(A)\), the canonical cyclic-lift multiple-sum identity,
and the fixed-matrix branch comparison used to separate scalar data from
non-trivial Peter--Weyl class profiles.  Main Theorems~I--III use the
finite-group determinant and Adams--M\"obius machinery in the main text for the
non-trivial class information; scalar cyclic-lift data alone do not descend to
non-abelian primitive class coordinates.

The appendices collect the scalar normalisation, finite-state criteria, and
explicit \(S_3\) rational witness used at the indicated points in the main
argument.
\begin{center}
\begin{tabular}{@{}L{0.24\textwidth}L{0.21\textwidth}L{0.19\textwidth}L{0.26\textwidth}@{}}
\toprule
Appendix material & Support role & Feeds & Limited use in the paper \\
\midrule
\Cref{app:cyclic-lift-comparison}
& Scalar normalization checks
& \(C(A)\), \(\mathcal M(A)\), \(\Psi_q(A)\)
& Supplies reduced-determinant and cyclic-lift calculations for the scalar
normalisation only. \\
\Cref{app:scalar-boundaries}
& Scalar descent limitations
& Main Theorems~II--III
& Records why scalar, parameter-varying, arithmetic, and growing-family
records do not by themselves recover non-trivial class profiles. \\
\Cref{app:s3-model-computations}
& Finite \(S_3\) model support
& \Cref{thm:adams-obstruction-class-product}
& Gives the explicit finite \(S_3\) matrices behind the sign-channel
Adams obstruction and the standard-channel non-abelian defect. \\
\Cref{app:s3-log-certificates}
& Rational \(S_3\) inequality support
& \(S_3\) numerical windows
& Provides rational finite-sum and tail bounds used to certify the
displayed \(S_3\) inequalities. \\
\Cref{app:fixed-label-local-certificates}
& Local fixed-label reproducibility
& finite algorithm
& Records finite input--output data for the same \(S_3\) witness; it is not a
general certificate theory. \\
\bottomrule
\end{tabular}
\end{center}

\subsection{Logarithm convention for cyclic lifts}

Throughout this appendix, \(\Log\) denotes the principal branch on
\(\CC\setminus(-\infty,0]\). For every \(j\ge 2\) and every \(|t|\le 1\),
\[
  \Re\!\left(1-\frac{\mu_j}{\lambda}t\right)
  \ge 1-\left|\frac{\mu_j}{\lambda}t\right|
  \ge 1-\frac{|\mu_j|}{\lambda}
  >0.
\]
Hence each factor \(1-(\mu_j/\lambda)t\) stays in the right half-plane,
so its principal logarithm is defined on the whole closed unit disc. We
therefore set, for \(|t|\le 1\),
\[
  \log F_A(t)
  :=
  \sum_{j=2}^d \Log\!\left(1-\frac{\mu_j}{\lambda}t\right).
\]
This is a continuous branch of the logarithm of \(F_A\) on \(|t|\le 1\),
normalized by \(\log F_A(0)=0\). If \(|z|<1\), then
\(\Re(1-z)\ge 1-|z|>0\), so
\[
  -\Log(1-z)=\sum_{n\ge 1}\frac{z^n}{n}.
\]
Applying this with \(z=(\mu_j/\lambda)e^{i\theta}\) and summing over the
finitely many non-Perron eigenvalues gives
\[
  -\log F_A(e^{i\theta})
  =
  \sum_{n\ge 1}\frac{u_n(A)e^{in\theta}}{n}.
\]
This branch is defined factorwise; we do not invoke the principal
logarithm of the product \(F_A(t)\). Likewise, for every \(q\ge 1\) we
set
\[
  \Psi_q(A)
  :=
  -\sum_{j=2}^d
    \Log\!\left(1-\left(\frac{\mu_j}{\lambda}\right)^q\right).
\]
All determinant logarithms in this appendix use the indicated branches.

We also write
\[
  \Lambda_{\mathrm{rel}}(A):=\Lambda(A)/\lambda\in[0,1)
\]
and
\[
  u_n(A):=\sum_{j=2}^d\left(\frac{\mu_j}{\lambda}\right)^n
  \qquad (n\ge 1).
\]

\subsection{Reduced determinants and multiple sums}

\begin{proposition}[Fixed-matrix formula for the reduced
  determinants of cyclic lifts]
  \label{prop:finite-part-cyclic-lift-reduced-constant-closed}
  Fix one primitive non-negative matrix \(A\) with Perron root \(\lambda\).
  For its canonical cyclic lift, and only for this fixed-matrix scalar
  comparison model, the reduced constants satisfy, for every \(q\ge 1\),
  \[
    C_\lambda(A^{\langle q\rangle})
    =\frac{1}{q}\prod_{j=2}^d
      \left(1-\left(\frac{\mu_j}{\lambda}\right)^q\right)^{-1}.
  \]
  Equivalently,
  \[
    \exp(\Psi_q(A))
    =
    q\,C_\lambda(A^{\langle q\rangle}).
  \]
\end{proposition}

\begin{proof}
  The spectrum of \(A^{\langle q\rangle}=A\otimes P_q\) is
  \[
    \{\mu\,\omega_q^k:\mu\in\spec(A),\ 0\le k\le q-1\}.
  \]
  The eigenvalue \(\lambda\) of \(A^{\langle q\rangle}\) is simple.
  Applying \Cref{lem:simple-eigenvalue-reduced-constant} to this
  eigenvalue yields
  \[
    C_\lambda(A^{\langle q\rangle})
    =\prod_{k=1}^{q-1}(1-\omega_q^k)^{-1}
      \prod_{j=2}^d\prod_{k=0}^{q-1}
        \left(1-\frac{\mu_j}{\lambda}\omega_q^k\right)^{-1}.
  \]
  The first product equals \(1/q\), while
  \[
    \prod_{k=0}^{q-1}(1-x\omega_q^k)=1-x^q.
  \]
  Substituting \(x=\mu_j/\lambda\) proves the claim.
\end{proof}

\begin{proposition}[Fixed-matrix multiple-sum expansion for
  cyclic lifts with explicit exponential remainder]
  \label{prop:finite-part-cyclic-lift-dirichlet-multiple-sum}
  For every \(q\ge 1\),
  \[
    \Psi_q(A)=\sum_{m\ge 1}\frac{u_{qm}(A)}{m},
  \]
  and therefore
  \[
    \Psi_q(A)=u_q(A)+O\!\left(\Lambda_{\mathrm{rel}}(A)^{2q}\right)
    \qquad (q\to\infty).
  \]
  More precisely,
  \[
    \abs{\Psi_q(A)-u_q(A)}
    \le
    \frac{d-1}{2(1-\Lambda_{\mathrm{rel}}(A))}
    \Lambda_{\mathrm{rel}}(A)^{2q}.
  \]
\end{proposition}

\begin{proof}
  By \Cref{prop:finite-part-cyclic-lift-reduced-constant-closed},
  \[
    \Psi_q(A)
    =-\sum_{j=2}^d\Log\!\left(1-\left(\frac{\mu_j}{\lambda}\right)^q\right).
  \]
  Expanding \(-\Log(1-x)=\sum_{m\ge 1}x^m/m\) for \(|x|<1\) and
  exchanging the finite sum over \(j\) with the absolutely convergent
  series over \(m\) gives the first identity. The remainder estimate
  follows from
  \[
    |u_{qm}(A)|\le (d-1)\Lambda_{\mathrm{rel}}(A)^{qm}.
  \]
  Indeed,
  \[
    \begin{aligned}
      \abs{\Psi_q(A)-u_q(A)}
      &\le
      \sum_{m\ge 2}\frac{\abs{u_{qm}(A)}}{m} \\
      &\le
      \frac{d-1}{2}\sum_{m\ge 2}\Lambda_{\mathrm{rel}}(A)^{qm} \\
      &\le
      \frac{d-1}{2(1-\Lambda_{\mathrm{rel}}(A))}
      \Lambda_{\mathrm{rel}}(A)^{2q}.
    \end{aligned}
  \]
\end{proof}

\begin{lemma}[Fixed-matrix absolute convergence for multiples
  inversion]
  \label{lem:finite-part-cyclic-lift-mobius-absolute}
  Let \(b_n:=u_n(A)/n\). Then, for every \(n\ge 1\),
  \[
    \sum_{m\ge 1}\sum_{r\ge 1}\abs{b_{mnr}}<\infty.
  \]
\end{lemma}

\begin{proof}
  For every \(k\ge 1\),
  \[
    \abs{u_k(A)}
    \le (d-1)\Lambda_{\mathrm{rel}}(A)^k,
  \]
  so
  \[
    \abs{b_k}
    \le \frac{d-1}{k}\Lambda_{\mathrm{rel}}(A)^k.
  \]
  Therefore
  \[
    \sum_{m\ge 1}\sum_{r\ge 1}\abs{b_{mnr}}
    =
    \sum_{s\ge 1}
      \#\{(m,r)\in \NN^2:mr=s\}\,\abs{b_{ns}}.
  \]
  Since every factorization \(mr=s\) is determined by the divisor \(m\),
  the multiplicity \(\#\{(m,r)\in \NN^2:mr=s\}\) is at most \(s\). Hence
  \[
    \sum_{m\ge 1}\sum_{r\ge 1}\abs{b_{mnr}}
    \le
    \sum_{s\ge 1}s\,\abs{b_{ns}}
    \le
    \frac{d-1}{n}\sum_{s\ge 1}\Lambda_{\mathrm{rel}}(A)^{ns}
    <\infty,
  \]
  because \(0\le \Lambda_{\mathrm{rel}}(A)<1\).
\end{proof}

\begin{proposition}[Fixed-matrix M\"obius inversion on multiples]
  \label{prop:finite-part-cyclic-lift-mobius-inversion}
  For every \(n\ge 1\),
  \[
    \frac{u_n(A)}{n}
    =\sum_{m\ge 1}\frac{\mu(m)}{mn}\,\Psi_{mn}(A).
  \]
\end{proposition}

\begin{proof}
  Set \(b_n:=u_n(A)/n\). Then
  \[
    \frac{\Psi_q(A)}{q}
    =\sum_{m\ge 1}b_{qm}
  \]
  by \Cref{prop:finite-part-cyclic-lift-dirichlet-multiple-sum}. By
  \Cref{lem:finite-part-cyclic-lift-mobius-absolute},
  \[
    \sum_{m\ge 1}\sum_{r\ge 1}\abs{\mu(m)b_{mnr}}
    \le
    \sum_{m\ge 1}\sum_{r\ge 1}\abs{b_{mnr}}
    <\infty,
  \]
  so Tonelli--Fubini applies to the double series below. Hence
  \begin{align*}
    \sum_{m\ge 1}\frac{\mu(m)}{mn}\Psi_{mn}(A)
    &=
    \sum_{m\ge 1}\mu(m)\sum_{r\ge 1}b_{mnr} \\
    &=
    \sum_{s\ge 1}b_{ns}\sum_{m\mid s}\mu(m)
    =b_n,
  \end{align*}
  since \(\sum_{m\mid s}\mu(m)=0\) unless \(s=1\).
\end{proof}

\begin{lemma}[Fixed-matrix branch ambiguity in multiples
  inversion]
  \label{lem:cyclic-lift-branch-ambiguity}
  Suppose a different logarithmic lift of the same exponentiated
  cyclic-lift determinant constants is chosen, so that
  \[
    \Psi'_q(A)=\Psi_q(A)+2\pi i\,k_q,
    \qquad k_q\in\ZZ.
  \]
  Let \(u'_n(A)\) be the power sums recovered from \(\{\Psi'_q(A)\}\) by
  the multiples-M\"obius formula
  \[
    \frac{u'_n(A)}{n}
    =
    \sum_{m\ge1}\frac{\mu(m)}{mn}\,\Psi'_{mn}(A),
  \]
  whenever the displayed series is interpreted with the same convergence
  convention as the chosen lift. Then the branch change modifies the
  recovered power sums by
  \[
    u'_n(A)-u_n(A)
    =
    2\pi i\sum_{m\ge1}\frac{\mu(m)}{m}\,k_{mn}.
  \]
  Thus the reconstruction theorem uses supplied logarithmic branches,
  not only the exponentiated constants
  \(\exp(\Psi_q(A))=qC_\lambda(A^{\langle q\rangle})\).
\end{lemma}

\begin{proof}
  Subtract the multiples-M\"obius formula for the two lifts:
  \[
    \frac{u'_n(A)-u_n(A)}{n}
    =
    \sum_{m\ge1}\frac{\mu(m)}{mn}
      \bigl(\Psi'_{mn}(A)-\Psi_{mn}(A)\bigr).
  \]
  Since \(\Psi'_{mn}(A)-\Psi_{mn}(A)=2\pi i\,k_{mn}\), multiplying by
  \(n\) gives the displayed formula. In particular, arbitrary branch
  choices can change the recovered power sums unless the integer
  sequence \(\{k_q\}\) satisfies the corresponding cancellation
  relations for every \(n\).
\end{proof}

\begin{lemma}[Fixed-matrix limsup of a finite exponential sum]
  \label{lem:finite-part-exponential-sum-limsup}
  Let
  \[
    f_q=\sum_{\beta\in B} d_\beta \beta^q
    \qquad (q\ge 1),
  \]
  where \(B\subset\CC\) is finite and nonempty, the numbers \(\beta\) are
  pairwise distinct, and \(d_\beta\in\CC\setminus\{0\}\) for every
  \(\beta\in B\).
  Then
  \[
    \limsup_{q\to\infty}\abs{f_q}^{1/q}
    =\max_{\beta\in B}\abs{\beta}.
  \]
\end{lemma}

\begin{proof}
  Set
  \[
    R:=\max_{\beta\in B}\abs{\beta}.
  \]
  The upper bound
  \(\limsup_{q\to\infty}\abs{f_q}^{1/q}\le R\) is immediate from
  \[
    \abs{f_q}
    \le \sum_{\beta\in B}\abs{d_\beta}\,\abs{\beta}^q
    \le \left(\sum_{\beta\in B}\abs{d_\beta}\right)R^q.
  \]
  If \(R=0\), then \(B=\{0\}\), so \(f_q=0\) for all \(q\ge 1\), and the
  asserted equality is immediate.

  Let
  \[
    B_R:=\{\beta\in B:\abs{\beta}=R\},
    \qquad
    g_q:=\sum_{\beta\in B_R} d_\beta (\beta/R)^q.
  \]
  Since \(B\) is finite, there exists \(\vartheta\in[0,1)\) such that
  \(\abs{\beta}\le \vartheta R\) for every \(\beta\in B\setminus B_R\).
  Hence
  \[
    f_q
    =
    R^q\left(g_q+\sum_{\beta\in B\setminus B_R}d_\beta(\beta/R)^q\right)
    =
    R^q\bigl(g_q+O(\vartheta^q)\bigr).
  \]
  The numbers \(\beta/R\) with \(\beta\in B_R\) lie on \(\TT\) and are
  pairwise distinct. Therefore
  \[
    \frac{1}{N}\sum_{q=1}^N\abs{g_q}^2
    =
    \sum_{\beta,\gamma\in B_R}
      d_\beta\overline{d_\gamma}\,
      \frac{1}{N}\sum_{q=1}^N
        \left(\frac{\beta\overline{\gamma}}{R^2}\right)^q.
  \]
  As \(N\to\infty\), the inner Ces\`aro average tends to \(1\) when
  \(\beta=\gamma\) and to \(0\) otherwise. Hence
  \[
    \frac{1}{N}\sum_{q=1}^N\abs{g_q}^2
    \longrightarrow
    \sum_{\beta\in B_R}\abs{d_\beta}^2>0.
  \]
  Therefore \(\limsup_{q\to\infty}\abs{g_q}>0\), so there exist
  \(c>0\) and infinitely many \(q\) such that \(\abs{g_q}\ge 2c\). Along
  that subsequence,
  \[
    \abs{f_q}\ge cR^q
  \]
  for all sufficiently large \(q\), because \(\vartheta^q\to 0\). This
  yields
  \(\limsup_{q\to\infty}\abs{f_q}^{1/q}\ge R\), completing the proof.
\end{proof}

\subsection{Fixed-matrix reconstruction outputs}

The two reconstruction statements in this subsection support the scalar
comparison model.  They are fixed-matrix
statements about one primitive matrix and its canonical cyclic lift; they are
not promoted to main consequences of the finite-group Adams--M\"obius
theorem chain and do not assert cross-resolution invariance for cocycles,
extensions, or non-trivial Peter--Weyl class profiles.

\begin{theorem}[Fixed-matrix cyclic-lift recovery of the reduced spectrum]
  \label{thm:finite-part-cyclic-lift-spectrum-identifiability}
  \label{thm:intro-cyclic-reconstruction}
  For one fixed primitive non-negative matrix \(A\in M_d(\RR)\), and only
  for the scalar cyclic-lift comparison record attached to that same matrix,
  the input is the full branch-lifted logarithmic record
  \[
    \Psi_q(A)\in \log\!\bigl(qC_\lambda(A^{\langle q\rangle})\bigr),
    \qquad q\ge1,
  \]
  with a specified branch for every \(q\). This fixed-matrix record
  determines the reduced spectral polynomial \(F_A\) of that same matrix and hence the
  relative non-Perron spectrum \(\{\mu_j/\lambda\}_{j=2}^d\) with
  algebraic multiplicity. If \(\lambda\) is supplied as additional data,
  then the unnormalised non-Perron spectrum \(\{\mu_j\}_{j=2}^d\) is also
  recovered.  This is an appendix support statement for the scalar
  normalisation model; it is not a coequal finite-group reconstruction theorem
  and is not used to recover non-trivial Peter--Weyl product constants.
  Moreover,
  \[
    \Psi_q(A)=u_q(A)+O\!\left(\Lambda_{\mathrm{rel}}(A)^{2q}\right),
  \]
  with the explicit bound from
  \Cref{prop:finite-part-cyclic-lift-dirichlet-multiple-sum}, and
  \[
    \limsup_{q\to\infty}\abs{\Psi_q(A)}^{1/q}
    =\Lambda_{\mathrm{rel}}(A).
  \]
\end{theorem}

\begin{proof}
  By \Cref{prop:finite-part-cyclic-lift-mobius-inversion}, the cyclic
  data determine every power sum \(u_n(A)\). Write
  \[
    \alpha_j:=\frac{\mu_j}{\lambda},
    \qquad
    P_A(x):=\prod_{j=2}^d(x-\alpha_j)
    =x^{d-1}F_A(x^{-1}).
  \]
  Then \(P_A\) is monic of degree \(d-1\), even when some
  \(\alpha_j=0\). Newton's identities recover the coefficients of this
  monic degree-\((d-1)\) polynomial from the first \(d-1\) power sums
  \(u_n(A)=\sum_{j=2}^d\alpha_j^n\), and hence from the full sequence
  \(\{u_n(A)\}_{n\ge 1}\). Therefore the multiset
  \(\{\alpha_2,\dots,\alpha_d\}\), including \(0\) with algebraic
  multiplicity, is determined. Equivalently,
  \[
    F_A(t)=\prod_{j=2}^d(1-\alpha_j t)
  \]
  is determined. Its ordinary degree counts the nonzero \(\alpha_j\), so
  the multiplicity of the zero non-Perron eigenvalue is exactly
  \(d-1-\deg F_A\). This reconstructs the relative non-Perron spectrum
  with algebraic multiplicity. Multiplying by \(\lambda\) would recover
  the unnormalised eigenvalues \(\mu_j\), but that requires \(\lambda\)
  as additional data beyond the cyclic-lift record itself. The argument
  is purely fixed-record: it recovers one polynomial \(F_A\) from the
  one sequence attached to the same matrix \(A\).

  For the exponential growth rate, let
  \[
    B:=\{\alpha_j:2\le j\le d,\ \alpha_j\neq 0\}
  \]
  be the set of distinct nonzero values among the relative eigenvalues,
  and for each \(\beta\in B\) let \(m_\beta\ge 1\) be its algebraic
  multiplicity among \(\alpha_2,\dots,\alpha_d\). Then
  \[
    u_q(A)=\sum_{\beta\in B}m_\beta\beta^q.
  \]
  If \(B=\varnothing\), then \(u_q(A)=0\) for every \(q\ge 1\), so
  \(\Lambda_{\mathrm{rel}}(A)=0\) and the asserted limsup identity is
  immediate. Assume henceforth that \(B\neq\varnothing\). Since the
  frequencies in \(B\) are pairwise distinct and each coefficient
  \(m_\beta\) is a positive integer, no cancellation can occur after
  grouping equal frequencies, and
  \Cref{lem:finite-part-exponential-sum-limsup} gives
  \[
    \limsup_{q\to\infty}\abs{u_q(A)}^{1/q}
    =\Lambda_{\mathrm{rel}}(A).
  \]
  The explicit remainder bound from
  \Cref{prop:finite-part-cyclic-lift-dirichlet-multiple-sum} shows that
  \[
    \Psi_q(A)-u_q(A)
    =O\!\left(\Lambda_{\mathrm{rel}}(A)^{2q}\right).
  \]
  If \(\Lambda_{\mathrm{rel}}(A)=0\), then both sequences are
  eventually zero. Otherwise the error term has strictly smaller
  exponential rate than \(u_q(A)\), so
  \[
    \limsup_{q\to\infty}\abs{\Psi_q(A)}^{1/q}
    =\Lambda_{\mathrm{rel}}(A).
  \]
\end{proof}

\begin{proof}[Fixed-matrix reconstruction does not descend to class-product coordinates]
  The reconstruction in
  \Cref{thm:finite-part-cyclic-lift-spectrum-identifiability} is a
  fixed-matrix scalar reconstruction.  More precisely, fix the same primitive
  matrix \(A\) and compare any two finite-group extensions over that base,
  possibly with different cocycles and different finite groups.  Their scalar
  cyclic-lift records \((\Psi_q(A))_{q\ge1}\) in the theorem are identical
  whenever the same branch-lifted logarithms are supplied, because the record
  is attached only to \(A\).  Consequently no function of this fixed-matrix
  cyclic-lift record alone can distinguish two such extensions whose
  non-trivial branch-compatible fixed-label Perron coordinate vectors
  \((F_\varrho(\lambda^{-1}))_{\varrho\ne\mathbf1}\) differ.  In particular,
  the theorem is not a finite-group class-product reconstruction theorem and
  supplies no replacement for the determinant coordinates of
  \Cref{prop:fixed-label-determinant-perron}.

  The data in \Cref{thm:finite-part-cyclic-lift-spectrum-identifiability} are
  exactly the branch-lifted logarithms
  \(\Psi_q(A)\in\log(qC_\lambda(A^{\langle q\rangle}))\) for the canonical
  cyclic lifts of the one matrix \(A\).  They contain no cocycle, no finite
  group, no character table, no Adams tensor, and no twisted determinant
  \(\det(I-zA_\pi)\) for a non-trivial representation.  Hence two extensions
  over the same base matrix, when supplied with the same logarithmic branches
  for this scalar record, give the same input to the appendix theorem.  If
  their non-trivial fixed-label Perron vectors are different, any function of
  the common scalar input takes the same value on both extensions and therefore
  cannot recover or distinguish those vectors.  The determinant-coordinate
  recovery in the main theorem chain uses the additional finite-group data
  listed in \Cref{prop:fixed-label-determinant-perron}, so it is not a
  consequence of this appendix theorem alone.
\end{proof}

\begin{theorem}[Fixed-layer Fourier recovery]
  \label{thm:finite-part-cyclic-lift-single-q-torsion-reconstruction}
  \label{cor:intro-one-layer-reconstruction}
  Write
  \[
    F_A(t)=\sum_{n=0}^{d-1}a_n t^n
    \qquad\text{with } a_0=1.
  \]
  For the same fixed-matrix scalar comparison model, if \(q\ge d\), then
  for this fixed \(q\) the values
  \(\{F_A(\omega_q^j)\}_{0\le j\le q-1}\) determine every coefficient
  \(a_n\). More precisely,
  \[
    a_n=\frac{1}{q}\sum_{j=0}^{q-1}
      F_A(\omega_q^j)\omega_q^{-jn}
    \qquad (0\le n\le d-1).
  \]
  This one-layer inversion is only the finite Fourier step inside the
  fixed-matrix scalar comparison model; it is not a replacement for the
  finite-group Adams--M\"obius determinant recovery of Main Theorem~I or the
  product-constant rigidity of Main Theorem~III.
\end{theorem}

\begin{proof}
  The polynomial identity
  \[
    F_A(\omega_q^j)=\sum_{n=0}^{d-1}a_n\omega_q^{jn}
  \]
  holds for \(0\le j\le q-1\). Averaging against
  \(\omega_q^{-jm}\) and using discrete orthogonality gives
  \[
    \frac{1}{q}\sum_{j=0}^{q-1}
      F_A(\omega_q^j)\omega_q^{-jm}
    =
    \sum_{n=0}^{d-1}a_n
      \frac{1}{q}\sum_{j=0}^{q-1}\omega_q^{j(n-m)}.
  \]
  Since \(q\ge d\), the inner sum equals \(1\) when \(n=m\) and \(0\)
  otherwise. This leaves exactly \(a_m\), proving the displayed recovery
  formula for every \(0\le m\le d-1\).
\end{proof}

% END INPUT sec_appendices.tex
% BEGIN INPUT sec_appendices_scalar_boundaries.tex
\section{Scalar finite-part normalisation}
\label{app:scalar-boundaries}

This appendix supports the scalar normalisation used in Main Theorems~II--III
and records no-descent boundaries for scalar records.  Its role is to keep the
finite part \(\mathcal M(A)=\gamma+\FP(A)\), scalar boundary variation, and
canonical cyclic-lift comparison separate from the non-trivial fixed-label
coordinates \((F_\varrho(\lambda^{-1}))_{\varrho\ne\mathbf1}\).  The results
below are therefore scalar support tools: they do not recover non-abelian class
products, do not assert root-of-unity sampling for an arbitrary fixed cocycle,
and do not supply cross-resolution descent invariants.

\subsection{Parameter variation for scalar records}
\label{app:record-boundary}

The reconstruction statements of \Cref{sec:finite-part} are fixed-matrix
statements. They recover data either from the full record
\[
  \{\Psi_q(A)\}_{q\ge1}
\]
of one primitive matrix \(A\), or from one fixed layer
\[
  \{F_A(\omega_q^j)\}_{0\le j\le q-1},
  \qquad q\ge d .
\]
The following scalar facts locate the normalisation used by the finite-group theorem chain.  They are not extra class-product hypotheses.  Holomorphic variation and effective determinant
truncation provide auxiliary control of the trivial-channel finite part,
while the Adams--M\"obius and class-product theorems use from this appendix
only the fixed-matrix normalisation \(\mathcal M(A)=\gamma+\FP(A)\) and the
canonical cyclic-lift comparison from \Cref{sec:finite-part}.  The
non-scalar coordinates of Main Theorem~III are instead the determinant
coordinate formula isolated in
\Cref{cor:assembled-adams-frobenius-product-formula}.

\begin{theorem}[Holomorphic variation of the scalar finite part]
  \label{thm:finite-part-holomorphic-variation}
  Let \(U\subset\CC\) be a domain and let
  \(\theta\mapsto A(\theta)\in M_d(\CC)\) be holomorphic. Suppose that
  \(\lambda(\theta)\) is a holomorphic simple eigenvalue of
  \(A(\theta)\) on \(U\), that \(|\lambda(\theta)|>1\), and that every
  other eigenvalue \(\nu\) of \(A(\theta)\) satisfies
  \(|\nu|<|\lambda(\theta)|\). Fix \(\theta_0\in U\). After shrinking to
  a simply connected neighbourhood \(V\Subset U\) of \(\theta_0\), the
  scalar finite part admits a holomorphic branch
  \[
    \mathfrak{FP}_{V}(\theta)
    =
    \log C(\theta)
    +\sum_{m\ge 2}\frac{\mu(m)}{m}
      \log\det(I-\lambda(\theta)^{-m}A(\theta))^{-1},
  \]
  where
  \[
    C(\theta):=
    \lim_{z\to\lambda(\theta)^{-1}}
      (1-\lambda(\theta)z)\det(I-zA(\theta))^{-1}.
  \]
  The series converges normally on \(V\). If \(A(\theta)\) is
  primitive non-negative for real \(\theta\in V\cap\RR\), with Perron
  root \(\lambda(\theta)>1\), then
  \(\mathfrak{FP}_{V}(\theta)=\FP(A(\theta))\) for the compatible real
  logarithm branch. Moreover, on \(V\),
  \[
  \begin{aligned}
    \frac{d}{d\theta}\mathfrak{FP}_{V}(\theta)
    &=
    \frac{d}{d\theta}\log C(\theta)  \\
    &\quad+
    \sum_{m\ge 2}\frac{\mu(m)}{m}
      \Tr\!\left(
        (I-\lambda(\theta)^{-m}A(\theta))^{-1}
        \frac{d}{d\theta}
          \bigl(\lambda(\theta)^{-m}A(\theta)\bigr)
      \right),
  \end{aligned}
  \]
  and the differentiated series also converges normally on compact
  subsets of \(V\).
\end{theorem}

\begin{proof}
  Since \(\lambda(\theta_0)\) is simple and spectrally separated from
  the remaining eigenvalues, analytic perturbation theory gives, after
  shrinking \(U\) if necessary, a holomorphic spectral projection for
  this eigenvalue and a uniform number \(\kappa\in(0,1)\) such that the
  non-Perron spectrum of \(A(\theta)\) lies in
  \(|\nu|\le \kappa |\lambda(\theta)|\) for \(\theta\in V\). We also
  shrink \(V\) so that \(|\lambda(\theta)|\ge\lambda_0>1\).

  Factor
  \[
    \det(I-zA(\theta))
    =(1-\lambda(\theta)z)D(\theta,z),
  \]
  where
  \[
    D(\theta,z):=
    \det(I-zA(\theta))/(1-\lambda(\theta)z)
  \]
  is holomorphic near
  \(\{(\theta,\lambda(\theta)^{-1}):\theta\in V\}\). The spectral
  separation gives \(D(\theta,\lambda(\theta)^{-1})\neq0\). Thus
  \(C(\theta)=D(\theta,\lambda(\theta)^{-1})^{-1}\) is holomorphic and
  non-vanishing. Since \(V\) is simply connected, \(\log C(\theta)\)
  has a holomorphic branch.

  For \(m\ge2\), the matrix
  \(I-\lambda(\theta)^{-m}A(\theta)\) is invertible. Indeed the
  eigenvalue \(\lambda(\theta)\) contributes the factor
  \(1-\lambda(\theta)^{1-m}\), which is non-zero because
  \(|\lambda(\theta)|>1\), and every other eigenvalue \(\nu\) contributes
  \(1-\nu\lambda(\theta)^{-m}\), with
  \[
    |\nu\lambda(\theta)^{-m}|
    \le \kappa|\lambda(\theta)|^{1-m}
    \le \kappa\lambda_0^{1-m}<1.
  \]
  Thus the determinants in the displayed series are holomorphic and
  non-vanishing on \(V\). Again \(V\) is simply connected, so their
  logarithms have branches, fixed by the chosen branch at
  \(\theta_0\).

  It remains to prove normal convergence. Let
  \(K\Subset V\). On \(K\), set
  \[
    L:=\inf_{\theta\in K}|\lambda(\theta)|>1,
    \qquad
    B:=\sup_{\theta\in K}\|A(\theta)\|.
  \]
  For \(m\) large enough,
  \(\|\lambda(\theta)^{-m}A(\theta)\|\le BL^{-m}\le1/2\) uniformly on
  \(K\). Using
  \[
    \log\det(I-X)^{-1}
    =
    \sum_{j\ge1}\frac{\Tr(X^j)}{j}
  \]
  for \(\|X\|<1\), we get
  \[
    \left|
      \log\det(I-\lambda(\theta)^{-m}A(\theta))^{-1}
    \right|
    \le
    d\sum_{j\ge1}\frac{(BL^{-m})^j}{j}
    \le 2dB L^{-m},
  \]
  uniformly on \(K\). The Weierstrass test proves normal convergence
  of the logarithmic series. The same estimate applied on a slightly
  larger compact set, followed by Cauchy's integral formula, gives
  normal convergence of the differentiated series on \(K\). Termwise
  differentiation is therefore valid, and the derivative of each summand
  is the standard identity
  \[
    \frac{d}{d\theta}\log\det B(\theta)^{-1}
    =
    -\Tr(B(\theta)^{-1}B'(\theta))
  \]
  with \(B(\theta)=I-\lambda(\theta)^{-m}A(\theta)\), giving the
  displayed formula.

  Finally, for real primitive non-negative matrices the branch can be
  chosen at \(\theta_0\) to agree with the positive real reduced
  determinant branch. The finite-part closed form
  \Cref{lem:finite-part-primitive-closed-form} then identifies the
  holomorphic expression with \(\FP(A(\theta))\) along the real
  parameter interval.
\end{proof}

\begin{proposition}[Scalar perturbations do not create class corrections]
  \label{prop:scalar-boundary-no-class-correction}
  Assume the hypotheses of \Cref{thm:class-mertens-explicit}.  Let
  \(V\) be a holomorphic perturbation neighbourhood for the scalar
  finite part of \(A\) as in
  \Cref{thm:finite-part-holomorphic-variation}.  Suppose one tries to
  replace all class Mertens constants by scalar perturbations, in the
  sense that there are complex numbers \(b_C\), depending only on the
  conjugacy class \(C\), for which
  \[
    \mathcal M_C(A,\tau;G)=b_C+\frac{|C|}{|G|}\mathfrak{FP}_{V}(\theta_0)
  \]
  for the base point \(A=A(\theta_0)\).  Then necessarily
  \[
    b_C=\frac{|C|}{|G|}\gamma
      +\frac{|C|}{|G|}
        \sum_{\varrho\neq\mathbf1}
          \chi_\varrho^\ast(C)\Pi_\varrho(\lambda^{-1}).
  \]
  In particular, the scalar branch can account only for the trivial
  Perron-channel term.  The Peter--Weyl part of this correction vanishes
  for all classes if and only if
  \(\Pi_\varrho(\lambda^{-1})=0\) for every non-trivial
  \(\varrho\in\Irr(G)\).
\end{proposition}

\begin{proof}
  Along the real primitive branch through the base point,
  \Cref{thm:finite-part-holomorphic-variation} gives
  \(\mathfrak{FP}_{V}(\theta_0)=\FP(A)\).  The scalar normalisation in
  \Cref{lem:finite-part-primitive-closed-form} gives
  \(\mathcal M(A)=\gamma+\FP(A)\).  Subtracting
  \((|C|/|G|)\FP(A)\) from the class constant formula
  \eqref{eq:class-mertens-constant} in
  \Cref{thm:class-mertens-explicit} yields the displayed expression for
  \(b_C\).

  Hence the part supplied by the scalar branch is exactly the universal
  factor \((|C|/|G|)\FP(A)\); all non-scalar class information is the
  finite Peter--Weyl sum over non-trivial irreducible representations.
  If all non-trivial Perron-point values vanish, the displayed
  correction is zero for every class after subtracting the universal
  Euler constant contribution. Conversely, if the class-dependent
  correction is zero for every conjugacy class, then the class function
  \[
    C\longmapsto
    \sum_{\varrho\neq\mathbf1}
      \chi_\varrho^\ast(C)\Pi_\varrho(\lambda^{-1})
  \]
  is identically zero.  Irreducible characters form an orthonormal basis
  for class functions on \(G\), so every non-trivial Fourier coefficient
  \(\Pi_\varrho(\lambda^{-1})\) is zero.
\end{proof}

\begin{theorem}[No descent from scalar records to class products]
  \label{thm:scalar-class-product-no-descent}
  Fix a primitive base matrix \(A\) with Perron root \(\lambda>1\), fix the
  finite group \(G\), and let
  \[
    \mathcal T_{\mathrm{gap}}(A,G)
    :=\{\tau:E_A\to G:\eta(\tau)<\lambda\}.
  \]
  Fix a non-empty family
  \(\mathscr T\subset\mathcal T_{\mathrm{gap}}(A,G)\) of gap cocycles over
  the same edge shift, so that \Cref{thm:class-euler-product-mertens}
  applies to every \(\tau\in\mathscr T\).  Let \(\mathcal R(A)\) be any
  fixed-matrix scalar
  record formed only from \(A\), for instance
  \[
    \begin{gathered}
      \FP(A),\qquad \mathcal M(A),\qquad C(A),\qquad F_A,\\
      \{\Psi_q(A)\}_{q\ge1},
      \qquad \{F_A(\omega_q^j)\}_{0\le j<q}.
    \end{gathered}
  \]
  No fold map, resolution inverse system, cocycle quotient, or other descent
  datum is included in \(\mathcal R(A)\).  For each
  \(\tau\in\mathscr T\), let \((K_C^\tau)_C\) be the
  Adams-corrected Frobenius-class product constants of
  \Cref{thm:class-euler-product-mertens}, and put
  \[
    \mathfrak K_\tau(C):=
    -\frac{|G|}{|C|}\log K_C^\tau-\bigl(\gamma+\log C(A)\bigr).
  \]
  Then the following are equivalent.
  \begin{enumerate}[label=\textup{(\roman*)},leftmargin=*]
    \item The scalar record \(\mathcal R(A)\), regarded as a candidate
      product record on \(\mathscr T\), determines a single
      Frobenius-class product vector \((K_C)_C\) on \(\mathscr T\).
    \item The normalised product logarithm
      \(\mathfrak K_\tau\) is independent of
      \(\tau\in\mathscr T\).
    \item The non-trivial fixed-label Euler-transform vector
      \((F_\varrho^\tau(\lambda^{-1}))_{\varrho\ne\mathbf1}\) is
      independent of \(\tau\in\mathscr T\).
  \end{enumerate}
  More precisely, for any \(\tau,\tau'\in\mathscr T\),
  \begin{equation}\label{eq:scalar-product-no-descent-distance}
    \sum_{C\subset G}\frac{|C|}{|G|}
      \abs{\mathfrak K_\tau(C)-\mathfrak K_{\tau'}(C)}^2
    =
    \sum_{\varrho\ne\mathbf1}
      \abs{F_\varrho^\tau(\lambda^{-1})
            -F_\varrho^{\tau'}(\lambda^{-1})}^2 .
  \end{equation}
  Hence, if two cocycles over the same \(A\) have different non-trivial
  fixed-label Euler-transform vectors, every scalar cyclic-lift or finite-part
  record attached only to \(A\) gives the same scalar input for the two systems
  but cannot determine their product constants.  The obstruction is exactly
  the non-trivial Peter--Weyl coordinate vector of Main Theorem~III.
\end{theorem}

\begin{proof}
  The listed scalar records are functions of the single matrix \(A\), so they
  are constant on \(\mathscr T\).  For a fixed cocycle \(\tau\),
  \Cref{cor:class-product-constant-rigidity} gives the normalised logarithmic
  Fourier expansion
  \[
    \mathfrak K_\tau(C)
    =
    \sum_{\varrho\neq\mathbf1}
      \chi_\varrho^\ast(C)F_\varrho^\tau(\lambda^{-1}).
  \]
  Thus a common product vector on \(\mathscr T\) is the same thing as a
  common normalised logarithmic vector \(\mathfrak K_\tau\), because
  \(C(A)\) and the positive logarithm branch are fixed.  This proves the
  equivalence of \textup{(i)} and \textup{(ii)}.

  Subtracting the displayed expansion for two cocycles gives
  \[
    \mathfrak K_\tau(C)-\mathfrak K_{\tau'}(C)
    =
    \sum_{\varrho\neq\mathbf1}
      \chi_\varrho^\ast(C)
      \bigl(F_\varrho^\tau(\lambda^{-1})
            -F_\varrho^{\tau'}(\lambda^{-1})\bigr).
  \]
  Irreducible characters form an orthonormal basis of the Hermitian
  class-function space with weights \(|C|/|G|\), and the trivial character is
  absent.  Uniqueness of Fourier coefficients therefore gives
  \textup{(ii)} if and only if \textup{(iii)}, and Parseval gives
  \eqref{eq:scalar-product-no-descent-distance}.

  If two cocycles have different non-trivial fixed-label Euler-transform
  vectors, the right-hand side of
  \eqref{eq:scalar-product-no-descent-distance} is non-zero, so the product
  vectors are different.  Since \(\mathcal R(A)\) is nevertheless identical
  for the two systems, no scalar fixed-matrix record without an additional
  descent datum can determine the product constants across that family.  The
  last sentence is exactly the coordinate description supplied by
  \Cref{cor:class-product-constant-rigidity}, which is the product-constant
  rigidity part of Main Theorem~III.
\end{proof}

\begin{theorem}[Effective truncation of the scalar finite part]
  \label{thm:logM-truncation-bound}
  Let \(A\) be a primitive non-negative \(d\times d\) matrix with
  Perron root \(\lambda>1\), and put
  \[
    K_{\mathrm{tr}}(A):=
    \sup_{j\ge1}\frac{|\Tr(A^j)|}{\lambda^j}.
  \]
  For \(M\ge2\), define the \(M\)-term determinant truncation
  \[
    \FP_M(A):=
    \log C(A)+\sum_{2\le m\le M}\frac{\mu(m)}{m}
      \log\det(I-\lambda^{-m}A)^{-1}.
  \]
  Then
  \[
    |\FP(A)-\FP_M(A)|
    \le
    \frac{K_{\mathrm{tr}}(A)\lambda^{-M}}
      {(M+1)(1-\lambda^{-M})(1-\lambda^{-1})}.
  \]
  Consequently the same bound controls the scalar orbit-Mertens
  normalisation:
  \[
    |\mathcal M(A)-(\gamma+\FP_M(A))|
    \le
    \frac{K_{\mathrm{tr}}(A)\lambda^{-M}}
      {(M+1)(1-\lambda^{-M})(1-\lambda^{-1})}.
  \]
\end{theorem}

\begin{proof}
  By \Cref{lem:finite-part-primitive-closed-form},
  \[
    \FP(A)-\FP_M(A)
    =
    \sum_{m>M}\frac{\mu(m)}{m}
      \log\det(I-\lambda^{-m}A)^{-1}.
  \]
  For \(m\ge2\), the spectral radius of \(\lambda^{-m}A\) is
  \(\lambda^{1-m}<1\), and hence
  \[
    \log\det(I-\lambda^{-m}A)^{-1}
    =
    \sum_{j\ge1}\frac{\Tr(A^j)}{j}\lambda^{-mj}.
  \]
  Therefore
  \[
  \begin{aligned}
    \left|
      \log\det(I-\lambda^{-m}A)^{-1}
    \right|
    &\le
    K_{\mathrm{tr}}(A)
      \sum_{j\ge1}\frac{\lambda^j\lambda^{-mj}}{j}  \\
    &=
    K_{\mathrm{tr}}(A)
      \sum_{j\ge1}\frac{\lambda^{-(m-1)j}}{j}
    =
    K_{\mathrm{tr}}(A)\bigl[-\log(1-\lambda^{1-m})\bigr].
  \end{aligned}
  \]
  Since \(m>M\), one has \(\lambda^{1-m}\le\lambda^{-M}\). The
  elementary inequality
  \[
    -\log(1-x)\le \frac{x}{1-x}\qquad(0\le x<1)
  \]
  gives
  \[
    -\log(1-\lambda^{1-m})
    \le
    \frac{\lambda^{1-m}}{1-\lambda^{-M}}.
  \]
  Hence
  \[
  \begin{aligned}
    |\FP(A)-\FP_M(A)|
    &\le
    \frac{K_{\mathrm{tr}}(A)}{1-\lambda^{-M}}
      \sum_{m>M}\frac{\lambda^{1-m}}{m} \\
    &\le
    \frac{K_{\mathrm{tr}}(A)}{(M+1)(1-\lambda^{-M})}
      \sum_{m=M+1}^{\infty}\lambda^{1-m}  \\
    &=
    \frac{K_{\mathrm{tr}}(A)\lambda^{-M}}
      {(M+1)(1-\lambda^{-M})(1-\lambda^{-1})}.
  \end{aligned}
  \]
  The assertion for \(\mathcal M(A)=\gamma+\FP(A)\) follows by adding
  the same constant \(\gamma\) to both finite-part quantities.
\end{proof}

\begin{corollary}[Finite input--output certificate for the scalar finite part]
  \label{cor:scalar-finite-part-finite-io-certificate}
  Under the hypotheses of \Cref{thm:logM-truncation-bound}, fix a tolerance
  \(\epsilon>0\).  The finite certificate inputs are the matrix \(A\), its
  Perron root \(\lambda>1\), the scalar normalising constant \(\log C(A)\),
  the determinant polynomial \(\det(I-zA)\), and a trace majorant
  \(K_{\mathrm{tr}}(A)\).  Choose a truncation level \(M\ge2\) for which
  \begin{equation}\label{eq:scalar-finite-io-tail-choice}
    B_M(A):=
    \frac{K_{\mathrm{tr}}(A)\lambda^{-M}}
      {(M+1)(1-\lambda^{-M})(1-\lambda^{-1})}
    \le \epsilon/2 .
  \end{equation}
  If \(I_M\subset\RR\) is any rational interval enclosing the finite
  determinant sum \(\FP_M(A)\), with numerical radius
  \(R_{\mathrm{num}}\le\epsilon/2-B_M(A)\), then
  \begin{equation}\label{eq:scalar-finite-io-output}
    \FP(A)\in I_M+[-B_M(A)-R_{\mathrm{num}},
                       B_M(A)+R_{\mathrm{num}}].
  \end{equation}
  After adding \(\gamma\), the same interval encloses
  \(\mathcal M(A)\).  This scalar certificate depends on \(\lambda\) only
  through the tail bound and on the finite determinant polynomial through the
  finitely many logarithms \(2\le m\le M\); unlike the finite-group
  certificate in \Cref{prop:fixed-label-perron-finite-algorithm}, it has no
  dependence on a twisted gap \(\eta\) or on a finite group \(G\).
\end{corollary}

\begin{proof}
  The search for \(M\) terminates because \(\lambda>1\), so the bound
  \(B_M(A)\) tends to zero.  \Cref{thm:logM-truncation-bound} gives
  \(\abs{\FP(A)-\FP_M(A)}\le B_M(A)\).  Since \(I_M\) encloses
  \(\FP_M(A)\) with numerical radius \(R_{\mathrm{num}}\), adding the two
  independent radii gives \eqref{eq:scalar-finite-io-output}.  The statement
  for \(\mathcal M(A)\) follows by translating the interval by \(\gamma\).
  All logarithms in \(\FP_M(A)\) are ordinary interior-disc determinant
  logarithms, because \(m\ge2\) gives
  \(\rho(\lambda^{-m}A)<1\).  Hence no finite-group character table,
  class-power map, non-trivial Abel branch, or twisted spectral-gap parameter
  enters this scalar certificate.
\end{proof}

\subsection{Arithmetic limitations}
\label{app:arithmetic}

The cyclic-lift formulas in \Cref{sec:finite-part} recover \(F_A\) and
the relative non-Perron spectrum from scalar root-of-unity data.  In this
article those formulas are retained only as appendix-level scalar
normalisation and comparison tools.  Arithmetic refinements of the same data
remain statements about the canonical cyclic lift of the base matrix; cyclic
quotients of a fixed cocycle, treated in \Cref{sec:shadows}, are different
functorial pullbacks and detect only the one-dimensional Peter--Weyl sector.

\subsection{Growing families}
\label{app:rigidity}

The finite-part identities also yield estimates for growing families of
primitive matrices under additional hypotheses on
\(\Lambda(A)/\lambda\), \(d\), and \(\lambda\). In the present fixed-matrix
class-product theorem these estimates document the scalar normalisation
regime; the finite-group Adams--Frobenius mechanism uses only the fixed
matrix and the stated twisted spectral gap.

% END INPUT sec_appendices_scalar_boundaries.tex
% BEGIN INPUT sec_appendices_s3_model.tex
\section[Detailed computations for the S3 obstruction]{Detailed computations for the \texorpdfstring{$S_3$}{S3} obstruction}
\label{app:s3-model-computations}

This appendix contains the matrix calculations and coefficient extractions
used in \Cref{sec:shadows}. The main text keeps only the structural
obstruction theorem; the computational certificates are recorded here to
keep the quotient-shadow calculations from obscuring the Adams--M\"obius
class-Mertens argument.  The appendix verifies one finite obstruction witness,
not a catalogue or classification of small non-abelian examples.
The opening quadratic Adams obstruction is intentionally a proposition,
not a theorem-level main result; all references use the proposition label
accordingly.


\begin{proposition}[Quadratic Adams obstruction]
  \label{prop:quadratic-adams-obstruction}
  Let \(\chi\) be a non-trivial quadratic character of \(G\). Then, for
  \(|z|<\lambda^{-1}\),
  \[
    \Pi_\chi(z)
    =
    \sum_{\substack{m\ge 1\\ m\ \mathrm{odd}}}
      \frac{\mu(m)}{m}\,\mathcal L_\chi(z^m)
    +
    \sum_{\substack{m\ge 1\\ m\ \mathrm{even}}}
      \frac{\mu(m)}{m}\,\mathcal L_{\mathbf 1}(z^m).
  \]
  In particular, if \(\mathcal L_\chi\equiv0\), then
  \[
    \Pi_\chi(z)
    =
    -\frac{1}{2}
      \sum_{\substack{r\ge 1\\ r\ \mathrm{odd}}}
        \frac{\mu(r)}{r}\,
        \mathcal L_{\mathbf 1}(z^{2r}),
  \]
  and hence \(\Pi_\chi\) need not vanish.
\end{proposition}

\begin{proof}
  Since \(\chi\) is quadratic, \(\chi^m=\chi\) for \(m\) odd and
  \(\chi^m=\mathbf1\) for \(m\) even. The Adams--M\"obius inversion
  formula \eqref{eq:class-adams-mobius-inversion} gives the first
  identity. If \(\mathcal L_\chi\equiv0\), only even indices remain;
  writing \(m=2r\) and using \(\mu(2r)=-\mu(r)\) for odd squarefree
  \(r\) gives the second identity.
\end{proof}

Consider the full two-shift
\[
  A=
  \begin{pmatrix}
    1 & 1 \\
    1 & 1
  \end{pmatrix},
  \qquad
  \lambda=\rho(A)=2,
\]
and label the four allowed edges by
\begin{equation}\label{eq:s3-example-edge-labels}
  \begin{aligned}
    \tau(0\to 0)&=(23), & \tau(0\to 1)&=e, \\
    \tau(1\to 0)&=(12), & \tau(1\to 1)&=(123).
  \end{aligned}
\end{equation}
Let \(T\) be the transposition class and \(K\) the \(3\)-cycle class in
\(S_3\). The character table is
\begin{equation}\label{eq:s3-character-table}
  \begin{array}{c|ccc}
    & e & T & K \\ \hline
    \mathbf 1 & 1 & 1 & 1 \\
    \varepsilon & 1 & -1 & 1 \\
    \rho & 2 & 0 & -1
  \end{array}
\end{equation}
where \(\varepsilon\) is the sign character and \(\rho\) the standard
two-dimensional representation. In this \(S_3\) example only, \(\rho\)
also denotes the standard representation; expressions such as
\(\rho(M)\) still denote spectral radius.

In the basis
\[
  v_1=e_1-e_2,\qquad v_2=e_1+e_2-2e_3,
\]
the standard representation matrices of the generators appearing in
\eqref{eq:s3-example-edge-labels} are
\begin{equation}\label{eq:s3-standard-representation-matrices}
\begin{aligned}
  \rho((12))
    &=
    \begin{pmatrix}
      -1 & 0 \\
      0 & 1
    \end{pmatrix}, \\
  \rho((23))
    &=
    \begin{pmatrix}
      \frac12 & \frac32 \\
      \frac12 & -\frac12
    \end{pmatrix}, \\
  \rho((123))
    &=
    \begin{pmatrix}
      -\frac12 & -\frac32 \\
      \frac12 & -\frac12
    \end{pmatrix}.
\end{aligned}
\end{equation}
These matrices are written in a convenient equivalent real basis, not
necessarily a unitary one; the character identities, conjugation
relations, and spectral-radius statements used below are invariant under
equivalence of representations.

\begin{proposition}[Twisted channels for the \(S_3\)-example]
  \label{prop:s3-example-channels}
  One has
  \begin{equation}\label{eq:s3-example-aeps}
    A_\varepsilon=
    \begin{pmatrix}
      -1 & 1 \\
      -1 & 1
    \end{pmatrix},
    \qquad
    A_\varepsilon^2=0,
    \qquad
    \det(I-zA_\varepsilon)=1.
  \end{equation}
  Moreover,
  \begin{equation}\label{eq:s3-example-arho}
    A_\rho=
    \begin{pmatrix}
      \frac12 & \frac32 & 1 & 0 \\
      \frac12 & -\frac12 & 0 & 1 \\
      -1 & 0 & -\frac12 & -\frac32 \\
      0 & 1 & \frac12 & -\frac12
    \end{pmatrix},
  \end{equation}
  and
  \begin{equation}\label{eq:s3-example-arho-charpoly}
    \chi_{A_\rho}(t)=t^3(t+1),
    \qquad
    \det(I-zA_\rho)=1+z.
  \end{equation}
  Consequently,
  \begin{equation}\label{eq:s3-example-logarithms}
    \mathcal L_\varepsilon(z)=0,
    \qquad
    \mathcal L_\rho(z)=\log\det(I-zA_\rho)^{-1}=-\log(1+z).
  \end{equation}
  If \(a_{n,C}\) denotes the number of period-\(n\) points whose
  \(S_3\)-label lies in \(C\), then
  \begin{equation}\label{eq:s3-example-periodic-classes}
    a_{n,e}=\frac{2^n+2(-1)^n}{6},
    \qquad
    a_{n,T}=2^{n-1},
    \qquad
    a_{n,K}=\frac{2^n-(-1)^n}{3}.
  \end{equation}
  In particular, the transposition class is exactly equidistributed at
  every period length:
  \begin{equation}\label{eq:s3-example-periodic-equidistribution}
    a_{n,T}=\frac{|T|}{|S_3|}\,2^n
    \qquad \text{for every } n\ge 1.
  \end{equation}
  Moreover,
  \begin{equation}\label{eq:s3-example-eta-check}
    \rho(A_\varepsilon)=0,
    \qquad
    \rho(A_\rho)=1,
    \qquad
    \eta=1<2=\lambda.
  \end{equation}
  In particular, the hypothesis of
  \Cref{thm:class-mertens-explicit} holds for this example.
\end{proposition}

\begin{proof}
  The sign matrix is obtained by substituting
  \(\varepsilon((23))=-1\), \(\varepsilon(e)=1\),
  \(\varepsilon((12))=-1\), and \(\varepsilon((123))=1\) into the
  \(2\times2\) block matrix of edge labels. This gives
  \[
    A_\varepsilon
    =
    \begin{pmatrix}
      -1 & 1 \\
      -1 & 1
    \end{pmatrix},
    \qquad
    A_\varepsilon^2=0,
  \]
  and therefore \(\det(I-zA_\varepsilon)=1\).

  For the standard channel, \eqref{eq:s3-standard-representation-matrices}
  yields the block form
  \[
    A_\rho
    =
    \begin{pmatrix}
      \rho((23)) & I_2 \\
      \rho((12)) & \rho((123))
    \end{pmatrix},
  \]
  which is exactly \eqref{eq:s3-example-arho}. To compute its
  characteristic polynomial, write
  \[
    tI_4-A_\rho
    =
    \begin{pmatrix}
      tI_2-\rho((23)) & -I_2 \\
      -\rho((12)) & tI_2-\rho((123))
    \end{pmatrix}.
  \]
  Since the upper-right block \(-I_2\) commutes with the lower-right
  block, the block determinant identity gives
  \[
    \det(tI_4-A_\rho)
    =
    \det\!\bigl((tI_2-\rho((123)))(tI_2-\rho((23)))-\rho((12))\bigr).
  \]
  A direct multiplication using
  \eqref{eq:s3-standard-representation-matrices} shows that
  \[
    (tI_2-\rho((123)))(tI_2-\rho((23)))-\rho((12))
    =
    \begin{pmatrix}
      t^2 & 0 \\
      -t & t(t+1)
    \end{pmatrix}.
  \]
  Hence
  \[
    \chi_{A_\rho}(t)=\det(tI_4-A_\rho)=t^3(t+1),
  \]
  so the eigenvalues of \(A_\rho\) are \(0,0,0,-1\). This gives
  \(\rho(A_\rho)=1\), \(\det(I-zA_\rho)=1+z\), and
  \(\Tr(A_\rho^n)=(-1)^n\) for every \(n\ge 1\).

  Since \(\Tr(A^n)=2^n\), \(\Tr(A_\varepsilon^n)=0\), and
  \(\Tr(A_\rho^n)=(-1)^n\), the logarithmic channels are
  \[
  \begin{aligned}
    \mathcal L_\varepsilon(z)
      &=\log\det(I-zA_\varepsilon)^{-1}=0,\\
    \mathcal L_\rho(z)
      &=\log\det(I-zA_\rho)^{-1}=-\log(1+z).
  \end{aligned}
  \]
  The class-counting formula from
  \Cref{prop:class-indicator-expansion} gives
  \[
    \begin{aligned}
      a_{n,C}
      &=
      \frac{|C|}{|S_3|}
        \Bigl(
          \Tr(A^n)
          +\chi_\varepsilon^\ast(C)\Tr(A_\varepsilon^n)
          +\chi_\rho^\ast(C)\Tr(A_\rho^n)
        \Bigr) \\
      &=
      \frac{|C|}{|S_3|}
        \Bigl(2^n+\chi_\rho^\ast(C)(-1)^n\Bigr),
    \end{aligned}
  \]
  from which \eqref{eq:s3-example-periodic-classes} follows.
\end{proof}

\begin{proposition}[Primitive channels for the \(S_3\)-example]
  \label{prop:s3-example-primitive-channels}
  The sign channel satisfies
  \begin{equation}\label{eq:s3-example-pi-eps}
    \Pi_\varepsilon(z)
    =
    -\frac12
      \sum_{\substack{r\ge 1\\ r\ \mathrm{odd}}}
        \frac{\mu(r)}{r}\log(1-2z^{2r})^{-1},
  \end{equation}
  with the Perron-point certificate recorded in
  \Cref{lem:s3-example-perron-certificates}.

  For the standard representation,
  \begin{equation}\label{eq:s3-example-adams-rho}
    \psi^m\rho=
    \begin{cases}
      \rho, & m\equiv \pm 1 \pmod 6, \\
      \mathbf 1-\varepsilon+\rho, & m\equiv \pm 2 \pmod 6, \\
      \mathbf 1+\varepsilon, & m\equiv 3 \pmod 6, \\
      2\mathbf 1, & 6\mid m,
    \end{cases}
  \end{equation}
  and therefore
  \begin{equation}\label{eq:s3-example-pi-rho}
    \Pi_\rho(z)
    =
    \sum_{m\ge 1}\frac{\mu(m)}{m}
      \Bigl(
        \alpha(m)\log(1-2z^m)^{-1}
        -\beta(m)\log(1+z^m)
      \Bigr),
  \end{equation}
  where
  \begin{equation}\label{eq:s3-example-alpha-beta}
    (\alpha(m),\beta(m))=
    \begin{cases}
      (0,1), & m\equiv \pm 1 \pmod 6, \\
      (1,1), & m\equiv \pm 2 \pmod 6, \\
      (1,0), & m\equiv 3 \pmod 6, \\
      (2,0), & 6\mid m,
    \end{cases}
  \end{equation}
  with the Perron-point certificate again recorded in
  \Cref{lem:s3-example-perron-certificates}.
\end{proposition}

\begin{proof}
  Since \(\varepsilon\) is quadratic and
  \(\mathcal L_\varepsilon(z)=0\), \Cref{prop:quadratic-adams-obstruction}
  gives \eqref{eq:s3-example-pi-eps}. The series converges absolutely
  at \(z=1/2\), and the certified non-zero interval is established in
  \Cref{lem:s3-example-perron-certificates}.

  For \(\rho\), the Adams powers are computed from the character table
  \eqref{eq:s3-character-table} together with the power maps on the
  three conjugacy classes. In this model those maps are immediate: a
  transposition has \(m\)-th power in \(T\) for \(m\) odd and in
  \(\{e\}\) for \(m\) even, while a \(3\)-cycle has \(m\)-th power in
  \(K\) when \(3\nmid m\) and in \(\{e\}\) when \(3\mid m\). Therefore
  the class values of \(\psi^m\rho\) are
  \[
    \begin{array}{c|ccc|c}
      m \bmod 6
      & (\psi^m\chi_\rho)(e)
      & (\psi^m\chi_\rho)(T)
      & (\psi^m\chi_\rho)(K)
      & \psi^m\rho \\ \hline
      \pm 1 & 2 & 0 & -1 & \rho \\
      \pm 2 & 2 & 2 & -1 & \mathbf 1-\varepsilon+\rho \\
      3 & 2 & 0 & 2 & \mathbf 1+\varepsilon \\
      0 & 2 & 2 & 2 & 2\mathbf 1
    \end{array}
  \]
  and decomposing these four class functions in the basis
  \(\{\mathbf 1,\varepsilon,\rho\}\) gives
  \eqref{eq:s3-example-adams-rho}. Substituting the decomposition into
  \Cref{cor:primitive-peter-weyl-determinant}, and using
  \eqref{eq:s3-example-logarithms}, yields
  \eqref{eq:s3-example-pi-rho} and \eqref{eq:s3-example-alpha-beta}.
  The resulting series again converges absolutely at \(z=1/2\), and the
  certified non-zero interval is established in
  \Cref{lem:s3-example-perron-certificates}.
\end{proof}

\begin{proposition}[Naive Artin formula obstruction at the Euler-product
  level]
  \label{prop:s3-naive-artin-obstruction}
  In the \(S_3\)-example,
  \begin{equation}\label{eq:s3-f-level-windows}
    -\frac{381}{1000}<F_\varepsilon(1/2)<-\frac{380}{1000},
    \qquad
    -\frac{923}{1000}<F_\rho(1/2)<-\frac{922}{1000}.
  \end{equation}
  At the same point the ordinary Artin logarithm in the sign channel is
  identically zero:
  \begin{equation}\label{eq:s3-artin-zero-f-nonzero}
    \mathcal L_\varepsilon(1/2)=0,
    \qquad
    F_\varepsilon(1/2)\neq0.
  \end{equation}
  Hence Frobenius-class Euler-product constants cannot in general be
  obtained by replacing the fixed-label transforms in
  \Cref{thm:class-euler-product-mertens} with ordinary Artin logarithms.
\end{proposition}

\begin{proof}
  The equality \(\mathcal L_\varepsilon(z)\equiv0\) is
  \eqref{eq:s3-example-logarithms}. The certified intervals for
  \(F_\varepsilon(1/2)\) and \(F_\rho(1/2)\) are proved by exact rational
  interval arithmetic in \Cref{lem:appendix-s3-f-level-certificates}; in
  particular both intervals exclude zero. Since
  \Cref{thm:class-euler-product-mertens} uses
  \(F_\varrho(\lambda^{-1})\), whereas
  \Cref{thm:fixed-label-euler-transform} identifies
  \(\mathcal L_\varrho(z)=\sum_{r\ge1}\Pi_{\psi^r\varrho}(z^r)/r\), the
  certified sign-channel inequality gives a concrete obstruction to the
  naive Artin substitution.
\end{proof}

\begin{lemma}[Certified \(S_3\) non-vanishing]
  \label{lem:s3-example-perron-certificates}
  For the \(S_3\)-extension above, the Perron-point primitive channels
  satisfy the certified inequalities
  \begin{equation}\label{eq:s3-example-certified-intervals}
    -\frac{342}{1000}
    <\Pi_\varepsilon(1/2)<-\frac{341}{1000},
    \qquad
    -\frac{719}{1000}
    <\Pi_\rho(1/2)<-\frac{718}{1000}.
  \end{equation}
  Hence both primitive channels are bounded away from zero, and the
  higher-dimensional term in the non-abelian defect from
  \Cref{thm:abelian-shadow-defect} is genuinely present in this
  example.
\end{lemma}

\begin{proof}
  Appendix~\ref{app:s3-log-certificates} gives a self-contained
  rational interval calculation for the finite truncations and tails of
  the two logarithmic series. Its enclosures imply exactly the coarse
  windows in \eqref{eq:s3-example-certified-intervals}. Since neither
  interval contains \(0\), both primitive channels are non-zero, and the
  concluding statement follows from
  \Cref{thm:abelian-shadow-defect}.
\end{proof}

\begin{corollary}[Primitive profile and deviations for the
  \(S_3\)-example]
  \label{cor:s3-example-profile}
  For the \(S_3\)-extension defined above, one has
  \begin{equation}\label{eq:s3-example-primitive-sign-coefficients}
    \pi_{n,\varepsilon}(A,\tau)
    =
    \frac{1}{n}
      \sum_{\substack{m\mid n\\ 2\mid m}}
        \mu(m)\,2^{n/m},
  \end{equation}
  and
  \begin{equation}\label{eq:s3-example-primitive-rho-coefficients}
    \pi_{n,\rho}(A,\tau)
    =
    \frac{1}{n}
      \sum_{m\mid n}
        \mu(m)\Bigl(
          \alpha(m)2^{n/m}+\beta(m)(-1)^{n/m}
        \Bigr).
  \end{equation}
  Consequently,
  \begin{equation}\label{eq:s3-example-primitive-classes}
    \begin{aligned}
      p_{n,e}
      &=\frac{p_n(A)+\pi_{n,\varepsilon}(A,\tau)+2\pi_{n,\rho}(A,\tau)}{6}, \\
      p_{n,T}
      &=\frac{p_n(A)-\pi_{n,\varepsilon}(A,\tau)}{2}, \\
      p_{n,K}
      &=\frac{p_n(A)+\pi_{n,\varepsilon}(A,\tau)-\pi_{n,\rho}(A,\tau)}{3}.
    \end{aligned}
  \end{equation}
  where
  \begin{equation}\label{eq:s3-example-full-two-shift-primitive}
    p_n(A)=\frac{1}{n}\sum_{d\mid n}\mu(d)\,2^{n/d}.
  \end{equation}
  The first primitive class triples are
  \begin{equation}\label{eq:s3-example-primitive-first-terms}
    (p_{n,e},p_{n,T},p_{n,K})
    =(0,1,1),(0,1,0),(0,1,1),(0,2,1),(1,3,2),(1,5,3),\ldots
  \end{equation}
  Finally,
  \begin{equation}\label{eq:s3-example-delta-values}
    \begin{aligned}
      \Delta_e(A,\tau;S_3)
      &=\frac16\Pi_\varepsilon(1/2)+\frac13\Pi_\rho(1/2), \\
      \Delta_T(A,\tau;S_3)
      &=-\frac12\Pi_\varepsilon(1/2), \\
      \Delta_K(A,\tau;S_3)
      &=\frac13\Pi_\varepsilon(1/2)-\frac13\Pi_\rho(1/2).
    \end{aligned}
  \end{equation}
  and the Perron-point windows imply
  \begin{equation}\label{eq:s3-example-delta-certified-windows}
    -\frac{89}{300}
    <\Delta_e(A,\tau;S_3)<-\frac{1777}{6000},
    \qquad
    \frac{341}{2000}
    <\Delta_T(A,\tau;S_3)<\frac{171}{1000},
  \end{equation}
  \begin{equation}\label{eq:s3-example-delta-k-certified-window}
    \frac{47}{375}
    <\Delta_K(A,\tau;S_3)<\frac{63}{500}.
  \end{equation}
  Equivalently, the class-constant deviations split as
  \begin{equation}\label{eq:s3-example-abelian-shadow}
    \Delta^{\mathrm{ab}}
    =
    \Bigl(
      \frac16\Pi_\varepsilon(1/2),
      -\frac12\Pi_\varepsilon(1/2),
      \frac13\Pi_\varepsilon(1/2)
    \Bigr),
  \end{equation}
  \begin{equation}\label{eq:s3-example-nonabelian-defect}
    \Delta^{\mathrm{na}}
    =
    \Bigl(
      \frac13\Pi_\rho(1/2),
      0,
      -\frac13\Pi_\rho(1/2)
    \Bigr).
  \end{equation}
  In particular,
  \begin{equation}\label{eq:s3-example-energy-ratio}
    \frac{\norm{\mathcal B_{A,\tau;G}^{\mathrm{na}}}_G^2}
         {\norm{\mathcal B_{A,\tau;G}^{\mathrm{ab}}}_G^2}
    =
    \frac{\abs{\Pi_\rho(1/2)}^2}
         {\abs{\Pi_\varepsilon(1/2)}^2}
    >4,
  \end{equation}
  so that, in the normalized \(B\)-profile \(L^2\)-norm of
  \Cref{thm:abelian-shadow-defect}, the genuine non-abelian primitive
  bias is certified to be more than four times stronger than the
  abelian shadow.
\end{corollary}

\begin{proof}
  Expanding \(\log(1-2z^{2r})^{-1}=\sum_{k\ge 1}2^k z^{2rk}/k\) in
  \eqref{eq:s3-example-pi-eps} and reading off the coefficient of
  \(z^n\) gives \eqref{eq:s3-example-primitive-sign-coefficients}.
  Likewise, coefficient extraction in \eqref{eq:s3-example-pi-rho}
  yields \eqref{eq:s3-example-primitive-rho-coefficients}. To obtain the
  three class formulas explicitly, specialise
  \Cref{prop:class-primitive-decomposition} to the character table
  \eqref{eq:s3-character-table}. Since the class sizes are
  \(|e|=1\), \(|T|=3\), and \(|K|=2\), and since
  \(\pi_{n,\mathbf1}(A,\tau)=p_n(A)\), character orthogonality gives
  \[
    \begin{aligned}
    p_{n,e}
      &=\frac{1}{6}\bigl(p_n(A)+\pi_{n,\varepsilon}(A,\tau)
        +2\pi_{n,\rho}(A,\tau)\bigr),\\
    p_{n,T}
      &=\frac{3}{6}\bigl(p_n(A)-\pi_{n,\varepsilon}(A,\tau)
        +0\cdot\pi_{n,\rho}(A,\tau)\bigr),\\
    p_{n,K}
      &=\frac{2}{6}\bigl(p_n(A)+\pi_{n,\varepsilon}(A,\tau)
        -\pi_{n,\rho}(A,\tau)\bigr),
    \end{aligned}
  \]
  which is \eqref{eq:s3-example-primitive-classes}. The scalar primitive
  count \eqref{eq:s3-example-full-two-shift-primitive} is the usual
  M\"obius formula for the full two-shift. Evaluating these formulas for
  \(n=1,\dots,6\) yields
  \eqref{eq:s3-example-primitive-first-terms}.

  The constant formulas are obtained from
  \Cref{prop:class-mertens-deviations} and the character table
  \eqref{eq:s3-character-table}, namely
  \[
    \begin{aligned}
      \Delta_e(A,\tau;S_3)
      &=\frac{1}{6}\Pi_\varepsilon(1/2)+\frac{1}{3}\Pi_\rho(1/2), \\
      \Delta_T(A,\tau;S_3)
      &=-\frac{1}{2}\Pi_\varepsilon(1/2), \\
      \Delta_K(A,\tau;S_3)
      &=\frac{1}{3}\Pi_\varepsilon(1/2)-\frac{1}{3}\Pi_\rho(1/2).
    \end{aligned}
  \]
  Combining the certified intervals in
  \Cref{lem:s3-example-perron-certificates} with these formulas gives
  \eqref{eq:s3-example-delta-certified-windows} and
  \eqref{eq:s3-example-delta-k-certified-window}. The signs and
  non-vanishing of these three deviations follow immediately.
  The decomposition into
  \eqref{eq:s3-example-abelian-shadow} and
  \eqref{eq:s3-example-nonabelian-defect} is exactly the splitting from
  \Cref{thm:abelian-shadow-defect}. Finally,
  \eqref{eq:s3-example-energy-ratio} is the corresponding norm identity,
  and the bound \(>4\) follows from
  \eqref{eq:s3-example-certified-intervals}, since
  \((718/342)^2>4\). The following table is therefore only a compact
  summary of the identities just derived and of the certified windows used
  in the last step.
\end{proof}

\begin{center}
\footnotesize
\begin{tabular}{@{}p{0.16\textwidth}p{0.32\textwidth}p{0.36\textwidth}@{}}
\toprule
Channel / observable & Periodic-level behaviour & Certified primitive information \\
\midrule
\(\varepsilon\) &
\(\mathcal L_\varepsilon(z)=0\) &
\(-\frac{342}{1000}<\Pi_\varepsilon(1/2)<-\frac{341}{1000}\) \\
\(\rho\) &
\(\mathcal L_\rho(z)=-\log(1+z)\) &
\(-\frac{719}{1000}<\Pi_\rho(1/2)<-\frac{718}{1000}\) \\
\(T\) &
\(a_{n,T}=2^{n-1}\) for all \(n\) &
\(\frac{341}{2000}<\Delta_T(A,\tau;S_3)<\frac{171}{1000}\) \\
\bottomrule
\end{tabular}
\end{center}
% END INPUT sec_appendices_s3_model.tex
% BEGIN INPUT sec_appendices_s3_certificates.tex
\section[Exact rational verification for the S3 certificates]{Exact
  rational verification for the \texorpdfstring{$S_3$}{S3}
  certificates}
\label{app:s3-log-certificates}

This appendix contains the exact finite-sum and tail bounds underlying
\Cref{lem:s3-example-perron-certificates}. The main text keeps only the
coarse certified non-vanishing windows; here we record the truncations,
tail estimates, and the final rational enclosures needed for that single
certificate, not a catalogue of finite-group obstruction examples and not a
general finite-certificate theory for arbitrary finite groups. The proof below
is self-contained: every rational truncation, tail estimate, and interval
enclosure used by the example is displayed in the appendix.

\begin{lemma}[Exact rational certificate for the \(S_3\)-channels]
  \label{lem:appendix-s3-sharp-certificates}
  Define
  \begin{equation}\label{eq:s3-example-epsilon-certificate-sum}
    \Pi_\varepsilon(1/2)=E_7+\Theta_\varepsilon,
  \end{equation}
  where
  \[
    E_7
    =
    -\frac12\log 2
    +\frac16\log\frac{32}{31}
    +\frac1{10}\log\frac{512}{511}
    +\frac1{14}\log\frac{8192}{8191},
  \]
  and
  \begin{equation}\label{eq:s3-example-epsilon-tail-certificate}
    \abs{\Theta_\varepsilon}\le \frac{1}{1105920}.
  \end{equation}
  Also define
  \begin{equation}\label{eq:s3-example-rho-certificate-sum}
    \Pi_\rho(1/2)=R_{15}+\Theta_\rho,
  \end{equation}
  where
  \[
  \begin{aligned}
    R_{15}
    ={}&
    -\log\frac32
    -\frac12\log\frac85
    -\frac13\log\frac43
    +\frac15\log\frac{33}{32}
    +\frac13\log\frac{32}{31}
    +\frac17\log\frac{129}{128} \\
    &+\frac1{10}\log\frac{512\cdot 1024}{511\cdot 1025}
    +\frac1{11}\log\frac{2049}{2048}
    +\frac1{13}\log\frac{8193}{8192} \\
    &+\frac1{14}\log\frac{8192\cdot 16384}{8191\cdot 16385}
    +\frac1{15}\log\frac{16384}{16383},
  \end{aligned}
  \]
  and
  \begin{equation}\label{eq:s3-example-rho-tail-certificate}
    \abs{\Theta_\rho}\le \frac{9}{524288}.
  \end{equation}
  Consequently,
  \begin{equation}\label{eq:s3-example-certified-intervals-sharp}
    -\frac{341079}{10^6}-\frac{1}{1105920}
    <\Pi_\varepsilon(1/2)<
    -\frac{341071}{10^6}+\frac{1}{1105920},
  \end{equation}
  and
  \begin{equation}\label{eq:s3-example-certified-intervals-sharp-rho}
    -\frac{718352}{10^6}-\frac{9}{524288}
    <\Pi_\rho(1/2)<
    -\frac{718351}{10^6}+\frac{9}{524288}.
  \end{equation}
  Hence the coarse windows of
  \Cref{lem:s3-example-perron-certificates} follow.
\end{lemma}

\begin{proof}
  The expression for \(E_7\) is obtained from
  \eqref{eq:s3-example-pi-eps} by retaining the odd indices
  \(r=1,3,5,7\). For the tail, put
  \(x_r=2^{1-2r}\). Since \(r\ge 9\) implies \(x_r\le 2^{-17}\),
  \[
    -\log(1-x_r)\le \frac{x_r}{1-x_r}\le 2x_r.
  \]
  Therefore
  \[
    \abs{\Theta_\varepsilon}
    \le
    \sum_{\substack{r\ge 9\\ r\ \mathrm{odd}}}\frac{2^{1-2r}}{r}
    \le
    \frac19\sum_{j\ge 0}2^{1-2(9+2j)}
    =
    \frac{1}{1105920}.
  \]

  The expression for \(R_{15}\) is obtained from
  \eqref{eq:s3-example-pi-rho} by retaining the squarefree indices
  \(m\le 15\). For \(m\ge 16\), set \(y_m=2^{1-m}\). Then
  \(y_m\le 2^{-15}\), and
  \[
    \log(1-y_m)^{-1}\le 2y_m,
    \qquad
    \log(1+2^{-m})\le 2^{-m}.
  \]
  Since \(0\le \alpha(m)\le 2\) and \(0\le \beta(m)\le 1\), this gives
  \[
    \abs{\Theta_\rho}
    \le
    \sum_{m\ge 16}\frac{9\,2^{-m}}{m}
    \le
    \frac{9}{16}\sum_{m\ge 16}2^{-m}
    =
    \frac{9}{524288}.
  \]

  For each rational \(q>1\) appearing in \(E_7\) or \(R_{15}\), put
  \[
    u_q:=\frac{q-1}{q+1},
    \qquad
    L_3(q):=2\sum_{j=0}^{3}\frac{u_q^{2j+1}}{2j+1},
  \]
  and
  \[
    U_3(q):=L_3(q)+\frac{2u_q^9}{9(1-u_q^2)}.
  \]
  Since
  \[
    \log q
    =
    2\sum_{j\ge 0}\frac{u_q^{2j+1}}{2j+1}
    =
    L_3(q)+2\sum_{j\ge 4}\frac{u_q^{2j+1}}{2j+1},
  \]
  the tail is positive and satisfies
  \[
    0<\log q-L_3(q)
    <
    \frac{2u_q^9}{9}\sum_{j\ge 0}u_q^{2j}
    =
    \frac{2u_q^9}{9(1-u_q^2)}.
  \]
  Therefore
  \begin{equation}\label{eq:s3-log-lower-upper-bounds}
    L_3(q)<\log q<U_3(q).
  \end{equation}
  Equivalently, for every real coefficient \(c\),
  \[
    c>0:\qquad cL_3(q)<c\log q<cU_3(q),
  \]
  and
  \[
    c<0:\qquad cU_3(q)<c\log q<cL_3(q).
  \]
  Applying these inequalities termwise gives exact rational endpoints
  \begin{align*}
    E_7^-
    :={}&
    -\frac12\,U_3(2)
    +\frac16\,L_3(32/31)
    +\frac1{10}\,L_3(512/511) \\
    &+\frac1{14}\,L_3(8192/8191), \\
    E_7^+
    :={}&
    -\frac12\,L_3(2)
    +\frac16\,U_3(32/31)
    +\frac1{10}\,U_3(512/511) \\
    &+\frac1{14}\,U_3(8192/8191),
  \end{align*}
  and
  \begin{align*}
    R_{15}^-
    :={}&
    -U_3(3/2)
    -\frac12\,U_3(8/5)
    -\frac13\,U_3(4/3)
    +\frac15\,L_3(33/32)
    +\frac13\,L_3(32/31) \\
    &+\frac17\,L_3(129/128)
    +\frac1{10}\,L_3(524288/523775)
    +\frac1{11}\,L_3(2049/2048) \\
    &+\frac1{13}\,L_3(8193/8192)
    +\frac1{14}\,L_3(134217728/134209535) \\
    &+\frac1{15}\,L_3(16384/16383), \\
    R_{15}^+
    :={}&
    -L_3(3/2)
    -\frac12\,L_3(8/5)
    -\frac13\,L_3(4/3)
    +\frac15\,U_3(33/32)
    +\frac13\,U_3(32/31) \\
    &+\frac17\,U_3(129/128)
    +\frac1{10}\,U_3(524288/523775)
    +\frac1{11}\,U_3(2049/2048) \\
    &+\frac1{13}\,U_3(8193/8192)
    +\frac1{14}\,U_3(134217728/134209535) \\
    &+\frac1{15}\,U_3(16384/16383).
  \end{align*}
  Then \(E_7^-<E_7<E_7^+\) and \(R_{15}^-<R_{15}<R_{15}^+\). Since
  every \(L_3(q)\) and \(U_3(q)\) is rational, the displayed formulas
  are an exact rational interval certificate. Expanding them and
  subtracting the claimed decimal rational endpoints gives positive
  rational margins. Clearing the positive denominators in the rational
  formulas for \(E_7^\pm\) and \(R_{15}^\pm\) gives the following
  exact lower bounds for those margins:
  \[
  \begin{array}{c|c|c}
    \text{comparison} & \text{exact positive margin} & \text{rational lower bound} \\
    \hline
    E_7^-+\frac{341079}{10^6} & >0 & >\frac{94363}{10^{11}}\\
    -\frac{341071}{10^6}-E_7^+ & >0 & >\frac{35281}{5\cdot10^{10}}\\
    R_{15}^-+\frac{718352}{10^6} & >0 & >\frac{2693}{25\cdot10^{9}}\\
    -\frac{718351}{10^6}-R_{15}^+ & >0 & >\frac{1}{2\cdot10^{12}}.
  \end{array}
  \]
  The rightmost column is only a compact way to print the exact
  integer comparisons after denominators have been cleared; no decimal
  approximation is used in the proof of the signs.
  Therefore
  \begin{equation}\label{eq:s3-example-finite-sum-windows}
    \begin{aligned}
      -\frac{341079}{10^6}
      &<E_7^-<E_7<E_7^+<-\frac{341071}{10^6}, \\
      -\frac{718352}{10^6}
      &<R_{15}^-<R_{15}<R_{15}^+<-\frac{718351}{10^6}.
    \end{aligned}
  \end{equation}
  In particular, still by exact rational arithmetic,
  \[
    E_7^+ +\frac1{1105920}<-\frac{341}{1000},
    \qquad
    R_{15}^+ +\frac9{524288}<-\frac{718}{1000},
  \]
  while
  \[
    -\frac{342}{1000}<E_7^- -\frac1{1105920},
    \qquad
    -\frac{719}{1000}<R_{15}^- -\frac9{524288}.
  \]
  Combining these finite-sum bounds with
  \eqref{eq:s3-example-epsilon-tail-certificate} and
  \eqref{eq:s3-example-rho-tail-certificate} gives
  \eqref{eq:s3-example-certified-intervals-sharp} and
  \eqref{eq:s3-example-certified-intervals-sharp-rho}. These sharp
  intervals lie respectively in
  \[
    \left(-\frac{342}{1000},-\frac{341}{1000}\right)
    \qquad\text{and}\qquad
    \left(-\frac{719}{1000},-\frac{718}{1000}\right),
  \]
  which is exactly the certification quoted in
  \Cref{lem:s3-example-perron-certificates}.
\end{proof}

\begin{lemma}[Exact rational certificate for the \(S_3\) product-level
  fixed-label values]
  \label{lem:appendix-s3-f-level-certificates}
  In the \(S_3\)-example of \Cref{sec:shadows}, the product-level
  obstruction coordinates in \Cref{thm:class-euler-product-mertens} are
  certified by the printed rational inequalities below.  The same finite
  rational record is serialized in
  \texttt{certificates/s3\_log\_certificates.cert}, generated by
  \texttt{certificates/s3\_log\_certificates.py}; the SHA--256 digest of the
  script is
  \[
  \begin{array}{l}
    \texttt{301737D892108D115F6F70128D03EBF6291F9E91D5811306}\\
    \texttt{E7D8A49F915D510C},
  \end{array}
  \]
  and the SHA--256 digest of the certificate record is
  \[
  \begin{array}{l}
    \texttt{22C4A0F015F8C579680B87B378BFAC19351CA66421C3988A}\\
    \texttt{025FEA9B433CFF62}.
  \end{array}
  \]
  In particular, independently of the auxiliary serialization, they are
  certified by
  \begin{equation}\label{eq:s3-f-epsilon-appendix-window}
    -\frac{381}{1000}<F_\varepsilon(1/2)<-\frac{380}{1000},
  \end{equation}
  and
  \begin{equation}\label{eq:s3-f-rho-appendix-window}
    -\frac{923}{1000}<F_\rho(1/2)<-\frac{922}{1000}.
  \end{equation}
\end{lemma}

\begin{proof}
  For rational \(q>1\), we use the same rational logarithm bounds as in
  \Cref{lem:appendix-s3-sharp-certificates}:
  \[
    L_3(q)<\log q<U_3(q),
  \]
  where
  \[
    L_3(q)=2\sum_{j=0}^{3}\frac{u_q^{2j+1}}{2j+1},
    \qquad
    U_3(q)=L_3(q)+\frac{2u_q^9}{9(1-u_q^2)},
    \qquad
    u_q=\frac{q-1}{q+1}.
  \]
  Thus all endpoints below are rational numbers.

  From \eqref{eq:s3-example-pi-eps},
  \[
    F_\varepsilon(1/2)
    =-\frac12\sum_{j\ge0}\frac{1}{2^j}
      \log\left(1-2^{1-2^{j+1}}\right)^{-1}.
  \]
  Let
  \[
    Q_8:=-\frac12\log2-\frac14\log\frac87
      -\frac18\log\frac{128}{127}
      -\frac1{16}\log\frac{32768}{32767}.
  \]
  The remaining terms are negative, and, since
  \(-\log(1-x)\le2x\) for \(0<x\le2^{-31}\), their absolute value is at
  most
  \[
    B_\varepsilon:=\frac{1}{16}\,2^{-31}\frac{1}{1-2^{-32}}.
  \]
  Termwise substitution of \(L_3\) and \(U_3\), with inequalities reversed
  on negative coefficients, gives rational numbers
  \[
  \begin{aligned}
    Q_8^- &:=-\frac12U_3(2)-\frac14U_3(8/7)
      -\frac18U_3(128/127)-\frac1{16}U_3(32768/32767),\\
    Q_8^+ &:=-\frac12L_3(2)-\frac14L_3(8/7)
      -\frac18L_3(128/127)-\frac1{16}L_3(32768/32767),
  \end{aligned}
  \]
  with
  \[
    Q_8^- - B_\varepsilon < F_\varepsilon(1/2)<Q_8^+.
  \]
  The rational endpoint values used here are
  \[
  \begin{array}{c|c}
    \text{endpoint} & \text{exact rational value}\\
    \hline
    Q_8^- &
    -\dfrac{77669633184570530219883191591202892933297454719}
      {203890011879382796622109825339358911119360000000}\\[1ex]
    Q_8^+ &
    -\dfrac{13843924721083839036686064834355526}
      {36342196007488170547184558429765625}\\[1ex]
    B_\varepsilon & \dfrac{1}{34359738360}.
  \end{array}
  \]
  Clearing denominators in these rational expressions gives the positive
  margins
  \[
  \begin{array}{c|c|c}
    \text{comparison} & \text{exact positive margin} &
      \text{rational lower bound}\\
    \hline
    Q_8^- -B_\varepsilon +\frac{381}{1000} & >0 & >\frac{1}{10^6}\\
    -\frac{380}{1000}-Q_8^+ & >0 & >\frac{9}{10^4}.
  \end{array}
  \]
  Hence
  \[
    -\frac{381}{1000}<Q_8^- - B_\varepsilon,
    \qquad
    Q_8^+<-\frac{380}{1000},
  \]
  proving \eqref{eq:s3-f-epsilon-appendix-window}.

  For the standard representation, combine
  \eqref{eq:s3-example-pi-rho} with the fixed-label transform. Grouping by
  \(k=rm\) gives
  \[
    F_\rho(1/2)=
    \sum_{k\ge1}\frac{1}{k}
      \left(
        A_k\log(1-2^{1-k})^{-1}
        -B_k\log(1+2^{-k})
      \right),
  \]
  where the first logarithm is absent for \(k=1\), and
  \[
    A_k:=\sum_{m\mid k}\mu(m)\alpha(m),
    \qquad
    B_k:=\sum_{m\mid k}\mu(m)\beta(m).
  \]
  The only non-zero pairs with \(k\le32\) are
  \[
  \begin{array}{c|ccccccccc}
    k&1&2&3&4&8&9&16&27&32\\ \hline
    A_k&0&-1&-1&-1&-1&-1&-1&-1&-1\\
    B_k&1&0&1&0&0&1&0&1&0.
  \end{array}
  \]
  Let \(S_{32}\) be the corresponding finite sum. For \(k\ge33\), the
  elementary bounds \(\abs{A_k}\le2\), \(\abs{B_k}\le1\),
  \(\log(1-2^{1-k})^{-1}\le2^{2-k}\), and
  \(\log(1+2^{-k})\le2^{-k}\) give the rational tail bound
  \[
    B_\rho:=\frac{5}{33}\,2^{-32}.
  \]
  Applying the same \(L_3/U_3\) interval substitution to the finitely many
  logarithms in \(S_{32}\) gives the rational endpoints
  \(S_{32}^-<S_{32}<S_{32}^+\)
  and hence
  \[
    S_{32}^- -B_\rho < F_\rho(1/2)<S_{32}^+ +B_\rho.
  \]
  The verification is completely rational at this point.  For reference,
  the finitely many logarithmic arguments entering \(S_{32}\) are
  \[
  \begin{array}{ccc}
    3/2, & 2, & 4/3,\\
    9/8, & 8/7, & 128/127,\\
    256/255, & 513/512, & 32768/32767,\\
    67108864/67108863, & 134217729/134217728, &
    2147483648/2147483647,
  \end{array}
  \]
  with coefficients respectively
  \[
    -1,
    -\frac12,
    -\frac13,
    -\frac13,
    -\frac14,
    -\frac18,
    -\frac19,
    -\frac19,
    -\frac1{16},
    -\frac1{27},
    -\frac1{27},
    -\frac1{32}.
  \]
  These displayed rational bounds are the mathematical certificate: for
  each term one inserts the lower or upper endpoint according to the sign of
  the rational coefficient, sums the resulting rational numbers, and then
  compares the final rational endpoints with the rational thresholds stated
  below.  No decimal approximation is used, and the main text depends only on
  the displayed rational endpoints, coefficients, and inequalities.
  Substituting \(L_3\) or \(U_3\) according to the sign of these
  coefficients and adding the rational tail \(B_\rho\), denominator
  clearing gives the explicit positive checks
  \[
  \begin{aligned}
    10^6\left(S_{32}^- -B_\rho +\frac{923}{1000}\right) &> 789,\\
    10^6\left(-\frac{922}{1000}-(S_{32}^+ +B_\rho)\right) &> 204.
  \end{aligned}
  \]
  Equivalently, the two margins are larger than \(789/10^6\) and
  \(204/10^6\). They are obtained by clearing the positive denominators in
  the displayed rational logarithmic arguments.  Thus the interval check is
  contained in these rational inequalities. The shorter numerical
  consequences are
  \[
  \begin{array}{c|c|c}
    \text{comparison} & \text{exact positive margin} &
      \text{rational lower bound}\\
    \hline
    S_{32}^- -B_\rho +\frac{923}{1000} & >0 & >\frac{7}{10^4}\\
    -\frac{922}{1000}-(S_{32}^+ +B_\rho) & >0 & >\frac{2}{10^4}.
  \end{array}
  \]
  Therefore
  \[
    -\frac{923}{1000}<S_{32}^- -B_\rho,
    \qquad
    S_{32}^+ +B_\rho<-\frac{922}{1000},
  \]
  proving \eqref{eq:s3-f-rho-appendix-window}.
\end{proof}

\begin{remark}[Exact-arithmetic reproducibility]
  The mathematical certificate in this appendix is the displayed rational
  arithmetic in
  \Cref{lem:appendix-s3-sharp-certificates,lem:appendix-s3-f-level-certificates}.
  A reader or auxiliary exact-arithmetic program can reproduce the calculation
  by computing the rational logarithm endpoints
  \(L_3(q)\le\log q\le U_3(q)\), multiplying them by the signed rational
  coefficients in \(E_7\), \(R_{15}\), \(Q_8\), and \(S_{32}\), adding the four
  rational tails, and clearing positive denominators in the four final interval
  comparisons.  The local reproducibility files are part of the source record:
  \begin{itemize}[leftmargin=*]
  \item \texttt{certificates/s3\_log\_certificates.py}. SHA--256:
    \[\begin{array}{l}
      \texttt{301737D892108D115F6F70128D03EBF6291F9E91D5811306}\\
      \texttt{E7D8A49F915D510C}.
    \end{array}\]
  \item \texttt{certificates/s3\_log\_certificates.cert}. SHA--256:
    \[\begin{array}{l}
      \texttt{22C4A0F015F8C579680B87B378BFAC19351CA66421C3988A}\\
      \texttt{025FEA9B433CFF62}.
    \end{array}\]
  \item \texttt{certificates/s3\_log\_certificates.out}. SHA--256:
    \[\begin{array}{l}
      \texttt{ECC2F9B589781195417952602CF782A209DB513B3C3012DF2603}\\
      \texttt{A3C736A5E18B}.
    \end{array}\]
  \end{itemize}
  The proof uses only the printed rational inequalities; the files give a
  machine-checkable reconstruction of the same rational endpoint and margin
  arithmetic.
\end{remark}
% END INPUT sec_appendices_s3_certificates.tex
% BEGIN INPUT sec_appendices_fixed_label_certificates.tex
\section{Appendix: local finite reproducibility records for the \texorpdfstring{\(S_3\)}{S3} witness}
\label{app:fixed-label-local-certificates}

This appendix records the local finite reproducibility formats used by
\Cref{prop:fixed-label-perron-finite-algorithm}.  The main text uses only
the finite determinant computability theorem and the determinant-to-product
Fourier map; the tabular records below are included to make the concrete
\(S_3\) determinant constants checkable.  They are not promoted to a fourth
main theorem chain or to a general certified-computation framework for
arbitrary finite groups.  In particular, the rational interval records here
support the fixed \(S_3\) Adams-obstruction certificate and do not define a
standalone certificate methodology or an inverse spectral invariant.

\begin{proposition}[Finite rational reproducibility record]
  \label{prop:fixed-label-rational-certificate-record}
  Assume the hypotheses of \Cref{thm:class-mertens-explicit}.  Fix a
  non-trivial \(\varrho\in\Irr(G)\), a tolerance \(\epsilon>0\), and a
  truncation level \(M\) satisfying
  \eqref{eq:fixed-label-perron-finite-algorithm-tail-choice}.  The output of
  \Cref{prop:fixed-label-perron-finite-algorithm} is reproduced by the following
  finite data.  The record consists only of finite algebraic model data and
  rational inequalities of the following four kinds:
  \begin{enumerate}[label=\textup{(\roman*)},leftmargin=*]
    \item the list of triples \((r,m,\pi)\) with \(rm\le M\), with the scalar
      Perron triple \((1,1,\mathbf1)\) deleted;
    \item the Adams coefficients \(a_{\varrho,
      \pi}^{(m)}=\langle\psi^m\chi_\varrho,\chi_\pi\rangle_G\), written from
      the finite character table and class-power table;
    \item for every listed triple, a branch tag specifying either the Abel
      non-trivial Perron branch \((rm=1,\pi\ne\mathbf1)\) or the ordinary germ
      value inside the Perron disc \((rm\ge2)\), together with a rational
      rectangular interval enclosure for
      \(\Log\det(I-\lambda^{-rm}A_\pi)^{-1}\);
    \item two rational error bounds \(R_M\) and \(R_{\mathrm{num}}\): the
      tail bound \(R_M\) is certified either directly from
      \eqref{eq:fixed-label-perron-finite-algorithm-error} or from the
      rational envelope \eqref{eq:fixed-label-rational-tail-envelope}, and
      \(R_{\mathrm{num}}\) is certified from the finite interval-propagation
      bound \eqref{eq:fixed-label-perron-finite-algorithm-roundoff}.
  \end{enumerate}
  Any checker which verifies these four records by exact algebraic comparison
  and rational interval arithmetic produces a rational rectangular interval
  \(I_M\subset\CC\) enclosing the finite sum \(F_\varrho^{(M)}\) and certifies
  the inclusion
  \begin{equation}\label{eq:fixed-label-certificate-interface-inclusion}
    F_\varrho(\lambda^{-1})
    \in
    I_M+
    \left\{w\in\CC:
      \abs{w}\le R_M+R_{\mathrm{num}}\right\},
  \end{equation}
  where all sides of \(I_M\) and both radii are rational.  Consequently the
  mathematical reproducibility support is the displayed finite record itself;
  any external
  file or script for a concrete model is only an optional serialization of the
  same rational checks used in
  \Cref{prop:fixed-label-perron-finite-algorithm,cor:class-product-constant-rigidity}.
  No claim is made here that this record is a polished software package or a
  universal certified-computation infrastructure for arbitrary finite groups,
  arbitrary matrix families, or arbitrary branch conventions.
\end{proposition}

\begin{proof}
  The first two records are finite because the set
  \(\{(r,m):rm\le M\}\) and \(\Irr(G)\) are finite, and because the Adams
  coefficients are finite character-table inner products.  The deleted scalar
  triple is exactly the omitted term in
  \Cref{prop:fixed-label-perron-finite-algorithm}; its coefficient is zero for
  \(\varrho\ne\mathbf1\), so no certificate record must assign a finite value
  to the singular scalar logarithm.

  For each remaining triple, the branch tag is one of the two cases already
  proved in item \textup{(d)} of
  \Cref{prop:fixed-label-perron-finite-algorithm}: either the non-trivial
  Abel branch at \(rm=1\), made regular by the strict twisted gap, or the
  ordinary logarithmic germ at \(rm\ge2\), where
  \(\rho(\lambda^{-rm}A_\pi)<1\).  Thus every tagged interval encloses a
  uniquely specified logarithm.  Multiplying the rational rectangular
  intervals by the certified coefficients \(\mu(m)a_{\varrho,\pi}^{(m)}/(rm)\),
  enclosing each algebraic product by exact comparison in a rational rectangle,
  and summing those rectangles gives a rational rectangular interval \(I_M\)
  which contains the stated finite truncation.  Its excess radius is at most
  \(R_{\mathrm{num}}\) by the triangle inequality used in
  \eqref{eq:fixed-label-perron-finite-algorithm-roundoff}.  The discarded
  terms with \(rm>M\) have total absolute value at most the rational number
  \(R_M\), either because the certificate compares \(R_M\) with the algebraic
  tail in \eqref{eq:fixed-label-perron-finite-algorithm-error}, or because it
  records rational \(\theta,Q_\varrho\) satisfying
  \eqref{eq:fixed-label-rational-tail-envelope} and
  \(R_M\ge Q_\varrho\theta^{M+1}/(1-\theta)\).  All these checks are finite
  algebraic or rational inequalities.  Adding the two independent radii proves
  \eqref{eq:fixed-label-certificate-interface-inclusion}.  Finally,
  \Cref{cor:class-product-constant-rigidity} applies finite Fourier inversion
  to the certified fixed-label Perron vector, so the same certificate record
  serializes the product constants rather than introducing a new analytic
  premise.
\end{proof}

\begin{proposition}[Finite input--output certificate summary]
  \label{prop:fixed-label-finite-io-certificate-summary}
  Assume the hypotheses of \Cref{thm:class-mertens-explicit}, and fix a
  non-trivial irreducible character \(\varrho\).  The following finite record is
  sufficient, and is necessary for the certificates used in this paper, to
  certify the fixed-label Perron coordinate
  \(F_\varrho(\lambda^{-1})\) to any prescribed rational accuracy:
  \begin{enumerate}[label=\textup{(\roman*)},leftmargin=*]
    \item the finite determinant data
      \(A,\lambda,G\), the character table, the power maps on conjugacy
      classes, and the determinant polynomials
      \(D_\pi(z)=\det(I-zA_\pi)\);
    \item a finite window \(M\) and the finitely many Adams coefficients
      \(a_{\varrho,
      \pi}^{(m)}\) with \(rm\le M\), with the sole scalar Perron triple
      \((r,m,\pi)=(1,1,\mathbf1)\) deleted before any logarithm is evaluated;
    \item for every remaining triple, the branch tag fixed by the germ at
      \(z=0\): Abel continuation to \(z=\lambda^{-1}\) when
      \(rm=1\) and \(\pi\ne\mathbf1\), and the ordinary interior-disc value at
      \(z=\lambda^{-rm}\) when \(rm\ge2\);
    \item rational rectangular enclosures for these finitely many logarithms
      and a rational number \(R\) dominating both interval-propagation error
      and the determinant tail
      \(C_\varrho\sum_{k>M}\mathrm d(k)\lambda^{1-k}/k\).
  \end{enumerate}
  From this record one obtains a rational rectangle \(J\subset\CC\) with
  \(F_\varrho(\lambda^{-1})\in J\) and diameter bounded by the requested
  tolerance.  Conversely, every certificate accepted by
  \Cref{prop:fixed-label-perron-finite-algorithm} reduces to such a record after
  replacing its two radii \(R_M\) and \(R_{\mathrm{num}}\) by
  \(R_M+R_{\mathrm{num}}\).  Through finite Fourier inversion in
  \Cref{thm:class-product-determinant-rigidity,cor:class-product-constant-rigidity},
  the same finite certificate determines the corresponding certified enclosure
  for every Adams-corrected Frobenius-class product constant.
\end{proposition}

\begin{proof}
  The sufficiency direction is exactly the finite truncation theorem in
  compressed form.  Items \textup{(i)} and
  \textup{(ii)} are finite by the finiteness of \(\Irr(G)\), the finite
  character table, and the finite set \(\{(r,m):rm\le M\}\).  The scalar
  Perron triple is deleted by
  \Cref{prop:fixed-label-determinant-perron,prop:fixed-label-perron-finite-algorithm}
  because its coefficient is zero for \(\varrho\ne\mathbf1\).  Item
  \textup{(iii)} is the two-branch convention proved in
  \Cref{prop:fixed-label-perron-finite-algorithm}: the non-trivial boundary
  terms are regular Abel limits by the twisted spectral gap, while all terms
  with \(rm\ge2\) lie strictly inside the Perron disc.  Summing the certified
  finite rectangles with coefficients
  \(\mu(m)a_{\varrho,
  \pi}^{(m)}/(rm)\) gives an enclosure of
  \(F_\varrho^{(M)}\), and enlarging by the rational radius in item
  \textup{(iv)} contains the discarded tail by
  \eqref{eq:fixed-label-perron-finite-algorithm-error}.  Choosing \(M\) from
  \eqref{eq:fixed-label-perron-finite-algorithm-tail-choice} and refining the
  finite interval arithmetic gives any prescribed rational diameter.

  Conversely, the certificate consists precisely of the determinant
  record, the finite window and Adams coefficients, the branch-tagged finite
  logarithm enclosures, and the two rational radii appearing in
  \eqref{eq:fixed-label-certificate-interface-inclusion}; replacing those
  two radii by their sum produces item \textup{(iv)} without changing any
  mathematical content.  The last assertion follows because the passage from
  the certified non-trivial Perron vector to the class-product constants is a
    finite character-table computation in
  \Cref{thm:class-product-determinant-rigidity,cor:class-product-constant-rigidity}.
\end{proof}

\begin{corollary}[Local five-field record for the \texorpdfstring{\(S_3\)}{S3} certificate]
  \label{cor:fixed-label-five-field-certificate-schema}
  Under the hypotheses and notation of
  \Cref{prop:fixed-label-perron-finite-algorithm}, the preceding
  analytic certificate can be recorded for the finite verification used in
  this paper by the following five fields:
  \begin{center}
  \small
  \begin{tabular}{@{}L{0.19\textwidth}L{0.33\textwidth}L{0.36\textwidth}@{}}
    \toprule
    field & finite content & verifier check\\
    \midrule
    model
      & \(A,\lambda,G\), the class table, the irreducible character table,
        and the power maps \(C\mapsto C^m\)
      & exact Perron equation for \(\lambda\), character orthogonality, and
        the class-power table\\
    window
      & the tolerance \(\epsilon\), truncation level \(M\), and every
        \((r,m,\pi)\) with \(rm\le M\) except \((1,1,\mathbf1)\)
      & \eqref{eq:fixed-label-perron-finite-algorithm-tail-choice} and the
        stated deletion of the singular scalar channel\\
    Adams
      & all integers \(a_{\varrho,
        \pi}^{(m)}=\langle\psi^m\chi_\varrho,\chi_\pi\rangle_G\)
      & recomputation from the finite character and power tables\\
    branch
      & for each listed triple, either the non-trivial Abel tag
        \((rm=1,\pi\ne\mathbf1)\) or the ordinary Perron-disc tag
        \((rm\ge2)\), with an algebraic or rational interval enclosing the
        corresponding determinant logarithm
      & uniqueness of the tagged branch from the gap or from
        \(\rho(\lambda^{-rm}A_\pi)<1\), followed by interval evaluation of
        \(\Log\det(I-\lambda^{-rm}A_\pi)^{-1}\)\\
    bound
      & the interval-propagation radius \(R_{\mathrm{num}}\) and the tail
        radius \(R_M=C_\varrho\sum_{k>M}\mathrm d(k)\lambda^{1-k}/k\)
      & signed interval summation of the finite window and the tail estimate
        \eqref{eq:fixed-label-perron-finite-algorithm-error}\\
    \bottomrule
  \end{tabular}
  \end{center}
  A checker for these fields returns exactly the interval inclusion
  \eqref{eq:fixed-label-certificate-interface-inclusion}.  Conversely, every
  finite certificate record used by
  \Cref{prop:fixed-label-perron-finite-algorithm} determines the five
  fields above after harmless reordering of the finite window.  Thus the
  appendix files, for example
  \texttt{certificates/s3\_log\_certificates.py} and its generated certificate
  text, are reproducible arithmetic transcripts for the local \(S_3\)
  verification, not new analytic assumptions or a general certificate theory.
\end{corollary}

\begin{proof}
  Starting from the four records in
  \Cref{prop:fixed-label-perron-finite-algorithm}, put the base matrix,
  Perron algebraic number, group table, class table, character table, and
  class-power table in the model field.  The window field is the finite triple
  list in record \textup{(i)} together with \(M\) and \(\epsilon\).  The Adams
  field is record \textup{(ii)}.  The branch field is record \textup{(iii)}.
  The bound field is record \textup{(iv)}.  Therefore the schema contains no
  datum beyond the finite determinant reproducibility proposition.

  The verifier checks listed in the last column are precisely the proof steps
  of \Cref{prop:fixed-label-perron-finite-algorithm}: finite character
  orthogonality recomputes the Adams coefficients; the deleted scalar triple is
  the unique singular Perron logarithm and has zero coefficient for
  \(\varrho\ne\mathbf1\); the two branch tags are the two regularity cases in
  \Cref{prop:fixed-label-perron-finite-algorithm}; interval multiplication by
  \(\mu(m)a_{\varrho,
  \pi}^{(m)}/(rm)\) and finite summation give radius
  \(R_{\mathrm{num}}\); and the discarded triples with \(rm>M\) contribute at
  most \(R_M\).  Hence acceptance gives
  \eqref{eq:fixed-label-certificate-interface-inclusion}.

  Conversely, any four-record certificate of the kind used here already
  specifies the same
  model data, finite window, Adams coefficients, branch intervals, and two
  radii; ordering the finite triple list lexicographically is the only
  normalization needed to place it in the displayed fields.  The cited
  \(S_3\) files instantiate these fields with exact rational arithmetic,
  so their role is a serialized transcript of this finite verification rather
  than a template for an arbitrary finite-group certification framework.
\end{proof}

% END INPUT sec_appendices_fixed_label_certificates.tex

\section*{Competing interests}
The author declares none.

\bibliographystyle{amsplain}
\bibliography{references}

\end{document}